% concatenated from firstdraft8.tex
\documentclass[11pt, reqno]{amsart}
%\documentclass[11pt,letterpaper]{amsart}
%\usepackage[margin=1.2in,marginparsep=0.1in,marginparwidth=1in]{geometry}
\usepackage{amssymb,amsmath,amsthm,amstext,amscd,latexsym,graphics,graphicx,bbm,caption,enumerate}


%\usepackage{showkeys}

\usepackage[dvipsnames,svgnames,table]{xcolor}
%\usepackage{centernot}
%\usepackage{fullpage}
\usepackage{mathtools}
%\usepackage{stmaryrd}
%\usepackage[plainpages=false,colorlinks=true, pagebackref]{hyperref}
\usepackage{hyperref}
\usepackage{tikz-cd}
\usetikzlibrary{matrix,arrows,decorations.pathmorphing}
\usetikzlibrary{positioning}
\usepackage{relsize}
\usepackage{enumerate}
\usepackage[shortlabels]{enumitem}
%\usepackage{tikz-cd}
%\usepackage{tikz}
% set arrows as stealth fighter jets
%\tikzset{>=stealth}
\hypersetup{plainpages=false,colorlinks=true, pagebackref}
%\hypersetup{citecolor=Sepia,linkcolor=blue, urlcolor=blue}
%\hypersetup{citecolor=Sepia,linkcolor=blue, urlcolor=blue}
%\usepackage{extpfeil}



%\makeatletter
%\def\@tocline#1#2#3#4#5#6#7{\relax
	%\ifnum #1>\c@tocdepth % then omit
	%\else
	%\par \addpenalty\@secpenalty\addvspace{#2}%
	%\begingroup \hyphenpenalty\@M
	%\@ifempty{#4}{%
		%\@tempdima\csname r@tocindent\number#1\endcsname\relax
		%}{%
		%\@tempdima#4\relax
		%}%
	%\parindent\z@ \leftskip#3\relax \advance\leftskip\@tempdima\relax
	%\rightskip\@pnumwidth plus4em \parfillskip-\@pnumwidth
	%#5\leavevmode\hskip-\@tempdima
	%\ifcase #1
	%\or\or \hskip 2em \or \hskip 2em \else \hskip 3em \fi%
	%#6\nobreak\relax
	%\dotfill\hbox to\@pnumwidth{\@tocpagenum{#7}}\par
	%\nobreak
	%\endgroup
	%\fi}
%\makeatother

%%%%%%%%%%%%%%%%%%%%%%%%%%%%%%%%%%%%%%%%%%%%%%%%%%%%%%%%%%%%%%%%%

% \swapnumbers
%\usepackage{tikz-cd}
\newtheorem{intro-thm}{Theorem}[]
\theoremstyle{plain}
\newtheorem{thm}{Theorem}[section]
\newtheorem{theorem}[thm]{Theorem}
\newtheorem{q}[thm]{Question}
\newtheorem{lemma}[thm]{Lemma}
\newtheorem{corollary}[thm]{Corollary}
\newtheorem{proposition}[thm]{Proposition}
\newtheorem{conj}[thm]{Conjecture}
\theoremstyle{definition}
\newtheorem{remark}[thm]{Remark}
\newtheorem{note}[thm]{Note}
\newtheorem{conjecture}[thm]{Conjecture}
\newtheorem{question}[thm]{Question}
\newtheorem{comments}[thm]{Comments}
\newtheorem{point}[thm]{}
\newtheorem{notation}[thm]{Notations}
\newtheorem{condition}[thm]{Condition}
\newtheorem{rmks}[thm]{Remarks}
\newtheorem{definition}[thm]{Definition}
\newtheorem{defn}[thm]{Definition}
\newtheorem{example}[thm]{Example}
\newtheorem{claim}[thm]{Claim}
\newtheorem{prop}[thm]{Universal Property}
\newtheorem{construction}[thm]{Construction}
%\numberwithin{equation}{section}
\newcounter{elno}                % This to number lists




%%%%%%%%%%%%%%%%%%%%%%%%%%%%%%%%%%%%%%%%%%%%%%%%%%%%%%%%%%%%%%%%%%%%%%%%%%%%%%

\newcommand{\rup}[1]{\lceil{#1}\rceil}
\newcommand{\rdown}[1]{\lfloor{#1}\rfloor}
\newcommand{\ilim}{\mathop{\varprojlim}\limits} % inverse limit
\newcommand{\dlim}{\mathop{\varinjlim}\limits}  % direct limit
\newcommand{\surj}{\twoheadrightarrow}
\newcommand{\inj}{\hookrightarrow}
\newcommand{\tensor}{\otimes}
\newcommand{\ext}{\bigwedge}
\newcommand{\Intersection}{\bigcap}
\newcommand{\Union}{\bigcup}
\newcommand{\intersection}{\cap}
\newcommand{\union}{\cup}

%%%%%%%%%%%%%%%%%%%%%%%%%%%%% new new commands :) %%%%%%%%%%%%%%%%
\newcommand{\supp}{{\rm Supp}}
\newcommand{\Exceptional}{{\rm Ex}}
\newcommand{\del}{\partial}
\newcommand{\delbar}{\overline{\partial}}
\newcommand{\boldphi}{\mbox{\boldmath $\phi$}}

%%%%%%%%%%%%%%%%%%%%%%%%%%%%%%%%%%%new commands by aritra%%%%%%%%%%%%%%%%%%%%%
%\newcommand{\singa}{\underline{Sing}_*^{\mathbb{A}^1}}
%\usetikzlibrary{between}
%%for short exact sequences %%%
%%fornmat is $\ses{1 ={X}, 2 ={Y}, 3={Z}}$ to get a sequence with 1 at the end, to change the %%ends you can define similar commands%%%%
%\ekvcHash\ses
%{
%	long 1 = {}
%	,long 2 = {}
%	, long 3 ={}
%}
%{1\longrightarrow\ekvcValue{1}{#1}\longrightarrow\ekvcValue{2}{#1}\longrightarrow\ekvcValue{3}{#1}\longrightarrow1}

%\ekvcHash\sesz
%{
%	long 1 = {}
%	,long 2 = {}
%	, long 3 ={}
%}
%{0\longrightarrow\ekvcValue{1}{#1}\longrightarrow\ekvcValue{2}{#1}\longrightarrow\ekvcValue{3}{#1}\longrightarrow0}
%%%%%%%%%%%%%%%%%%%%%%%%%%%%%%%%%%%%%%%%%%%%%%%%%%%%%%%%%%%%%%%%%%%%%%%%%%%%%%
\newcommand{\udiv}{\underline{\Div}}

%%%%%%%%%%%%%%%%%

\newcommand{\Proj}{{\P roj}}
\newcommand{\sEnd}{{\sE nd}}
\newcommand{\mc}{\mathcal}
\newcommand{\mb}{\mathbb}
\newcommand{\an}{{\rm an}} 
\newcommand{\red}{{\rm red}}
\newcommand{\codim}{{\rm codim}}
\newcommand{\Dim}{{\rm dim}}
\newcommand{\rank}{{\rm rank}}
\newcommand{\Ker}{{\rm Ker  }}
\newcommand{\Pic}{{\rm Pic}}
\newcommand{\per}{{\rm per}}
\newcommand{\ind}{{\rm ind}}
\newcommand{\Div}{{\rm Div}}
\newcommand{\Hom}{{\rm Hom}}
\newcommand{\Aut}{{\rm Aut}}
\newcommand{\im}{{\rm im}}
\newcommand{\pre}{\rm{pre}}
\newcommand{\pr}{\mathrm{pr}}%projection%
\newcommand{\gives}{\rightsquigarrow}
\newcommand{\Spec}{{\rm Spec \,}}
\newcommand{\Sing}{{\rm Sing}}
\newcommand{\sing}{{\rm sing}}
\newcommand{\reg}{{\rm reg}}
\newcommand{\Char}{{\rm char}}
\newcommand{\Tr}{{\rm Tr}}
\newcommand{\Gal}{{\rm Gal}}
\newcommand{\Min}{{\rm Min \ }}
\newcommand{\Max}{{\rm Max \ }}
\newcommand{\Alb}{{\rm Alb}\,}
\newcommand{\Mat}{{\rm Mat}}
%\newcommand{\GL}{{\rm GL}\,}        % For the general linear group
\newcommand{\GL}{{\G\L}}
\newcommand{\Ho}{{\rm Ho}}
\newcommand{\ie}{{\it i.e.\/},\ }
%\renewcommand{\ if and only if}{\mbox{ $\Longleftrightarrow$ }}
\renewcommand{\tilde}{\widetilde}
% Skriptbuchstaben
\newcommand{\sA}{{\mathcal A}}
\newcommand{\sB}{{\mathcal B}}
\newcommand{\sC}{{\mathcal C}}
\newcommand{\sD}{{\mathcal D}}
\newcommand{\sE}{{\mathcal E}}
\newcommand{\sF}{{\mathcal F}}
\newcommand{\sG}{{\mathcal G}}
\newcommand{\sH}{{\mathcal H}}
\newcommand{\sI}{{\mathcal I}}
\newcommand{\sJ}{{\mathcal J}}
\newcommand{\sK}{{\mathcal K}}
\newcommand{\sL}{{\mathcal L}}
\newcommand{\sM}{{\mathcal M}}
\newcommand{\sN}{{\mathcal N}}
\newcommand{\sO}{{\mathcal O}}
\newcommand{\sP}{{\mathcal P}}
\newcommand{\sQ}{{\mathcal Q}}
\newcommand{\sR}{{\mathcal R}}
\newcommand{\sS}{{\mathcal S}}
\newcommand{\sT}{{\mathcal T}}
\newcommand{\sU}{{\mathcal U}}
\newcommand{\sV}{{\mathcal V}}
\newcommand{\sW}{{\mathcal W}}
\newcommand{\sX}{{\mathcal X}}
\newcommand{\sY}{{\mathcal Y}}
\newcommand{\sZ}{{\mathcal Z}}
% Sonderbuchstaben mit Doppellinie
\newcommand{\A}{{\mathbb A}}
\newcommand{\B}{{\mathbb B}}
\newcommand{\C}{{\mathbb C}}
\newcommand{\D}{{\mathbb D}}
\newcommand{\E}{{\mathbb E}}
\newcommand{\F}{{\mathbb F}}
\newcommand{\G}{{\mathbb G}}
\newcommand{\HH}{{\mathbb H}}
\newcommand{\I}{{\mathbb I}}
\newcommand{\J}{{\mathbb J}}
\newcommand{\M}{{\mathbb M}}
\newcommand{\N}{{\mathbb N}}
%\renewcommand{\O}{{\mathbb O}}
\renewcommand{\P}{{\mathbb P}}
\newcommand{\Q}{{\mathbb Q}}
\newcommand{\R}{{\mathbb R}}
\newcommand{\T}{{\mathbb T}}
\newcommand{\U}{{\mathbb U}}
\newcommand{\V}{{\mathbb V}}
\newcommand{\W}{{\mathbb W}}
\newcommand{\X}{{\mathbb X}}
\newcommand{\Y}{{\mathbb Y}}
\newcommand{\Z}{{\mathbb Z}}
\newcommand{\Sh}{\sS h}
\newcommand{\deltaop}{\Delta^{op}(\sS h(Sm/\mathbf{k}))}
\newcommand{\pdeltaop}{\Delta^{op}(P\sS h(Sm/\mathbf{k}))}
%\newcommand{\psh}{\pi_0^{\text{\tiny pre}}}
\newcommand{\psh}{\pi_0}
\renewcommand{\k}{\mathbf{k}}
\let\syn\mathsf
\newcommand{\Step}[1]{\underline{\syn{Step \ {#1}}}}
\newcommand{\scr}{\scriptscriptstyle}
\newcommand{\e}[1]{\mathfrak{elim}{( #1)}}
\newcommand{\elim}[1]{\mathfrak{elim}_{\scr #1}}
\newcommand{\dgu}[1]{{{d}}^{\scr 1}_{\scriptstyle #1}}
\newcommand{\dgv}[1]{{{d}}^{\scr 2}_{\scriptstyle #1}}
\newcommand{\dg}{\syn{deg}}
\newcommand{\fixme}[1]{\textcolor{red}{$\bullet$}\marginpar{\small \textcolor{red}{#1}}}
\newcommand{\edit}[1]{\marginpar{\textcolor{Green}{#1}}}
\newcommand{\hash}{\raisebox{.3ex}{\footnotesize \# }}
\newcommand{\blue}[1]{\textcolor{Blue}{#1}}
\input{xy}
\xyoption{all}
\title{$\mathbb{A}^1$-fibration in algebraic geometry and $\mathbb{A}^1$-homotopy type.
}

\begin{document}
\subjclass[2000]{14F42}


\author{Utsav Choudhury}



\author{Aritra Mandal}
	
\author{Biman Roy}
\begin{abstract}
	In this article we show that an $\mathbb{A}^1$ bundle map or a vector bundle map $p: X \to Y$ induces trivial local fibration $\underline{Sing}(X) \to \underline{Sing}(Y)$. Using this, we first show that for Korus Russel threefolds of first kind $X$ the space  $\underline{Sing}(X)$ is $\mathbb{A}^1$  local. Then we show that for any smooth affine complex surface $X$, the $\mathbb{A}^1$- connected component sheaf is homotopy invariant. 
\end{abstract}

\maketitle

\section{Introduction}	
The homotopy information of a manifold $X$ fibered over a base space $Y$ with a fixed fiber $F$ can be recovered from the homotopy information of both $Y$ and $F$. This is because a 
fibration of the form $F \to X \to Y$, gives rise to the following exact sequence of homotopy groups (and pointed sets for $\pi_0$) : 
%$$
\begin{align*}
	\dots \to \pi_n(F) \to \pi_n(X) \to \pi_n(Y) \to & \pi_{n-1}(F) \to \dots  \\
	&\to \pi_0(F) \to \pi_0(X) \to \pi_0(Y).
\end{align*}
%$$

\ 

If the fiber $F$ is moreover contractible (i.e. $\pi_n(F) = 0$ for all $n \geq 0$), a direct application of the long exact sequence or alternatively, the Mayer-Vietoris property for local trivializations demonstrates that the homotopy invariants of the base $Y$ are isomorphic to those of the total space $X$. Such fibrations were classically exploited to compute homotopy groups of spheres and classify fiber bundles (see \cite{steenrod} or \cite{serre thesis}).

\

In the setting of algebraic geometry, one encounters a different class of fibrations. Consider a morphism $p : X \to Y$ of algberaic varieties that is generically trivial with a given variety $F$ as its generic fiber. Consequently, the ``bad" (or singular) fibers are confined to a closed locus of codimension one or greater. In many cases where the base and fibers are smooth, the rigidity inherent in algebraic geometry ensures that the topology of these degenerate fibers is heavily constrained by the generic fiber. For instance, when the base $Y$ is a curve and the generic fiber is an affine curve, bad fibers occur only over finitely many closed points. Furthermore, the Euler characteristic of the components of these bad fibers is explicitly determined by the Euler characteristic of the generic fiber. (For $\mathbb{A}^1$-fibrations on surfaces and their degeneration, see \cite{miyanishicrm} or  \cite{kambayashi}. For more general fiber degenerations and Euler characteristics, see \cite{iitaka}.)



\ 

In such a situation, it is desirable that the homotopy invariants of the total space be recoverable from the homotopy information of the generic fiber, the homotopy information of the base, and the geometry/configuration of the degenerate (bad) fibers. This is readily achieved for vector bundles and, more generally, vector bundle torsors (i.e., affine spaces bundles). These represent the precise cases where the generic and degenerate fibers coincide structurally, thereby forming a genuine fiber bundle (see, for example, \cite{brauer group}). Consequently, non-trivial phenomena arise when genuine degenerations and bad fibers occur. 




\

For instance, the complex sphere $S^2_{\mathbb{C}}$ admits an $\mathbb{A}^1$-fibration over $\mathbb{A}^1$, where the singular fiber is two disjoint copies of $\mathbb{A}^1$. Classically, it is well known that its underlying topological homotopy type is that of $\mathbb{CP}^1$,  which is also equivalent to the homotopy type of $\mathbb{A}^1$ with doubled origin. This can also be deduced from the fact that $S^2_{\mathbb{C}}$ carries the structure of an $\mathbb{A}^1$-bundle over $\mathbb{CP}^1$. The topological and birational classification of smooth affine surafces with negative log Kodaira dimension, which naturally exhibit such $\mathbb{A}^1$-fibrations, has been studied extensively in the literature (see, \cite{miyanishicrm}).


\


In this article, we explore the $\mathbb{A}^1$-algebraic topology of varieties admitting an $\mathbb{A}^1$-fibration over curves or surfaces. Any $\mathbb{A}^1$-fibration $X \to C$ over a curve $C$ has $\mathbb{A}^1$ as its generic fiber, while its bad fibers can consist at most of disjoint unions of "thickened" $\mathbb{A}^1$'s ( non-reduced lines). Within the framework of 
$\mathbb{A}^1$-algebraic topology (as developed by Morel and Voevodsky \cite{mv}), the affine line $\mathbb{A}^1$ is contractible to a point. Generically, therefore, the $\mathbb{A}^1$-homotopy type of $X$ is equivalent to the $\mathbb{A}^1$-homotopy type of $C$. 

\

It is natural to expect that the global $\mathbb{A}^1$-homotopy type of $X$ can be fully recovered from the base $C$ alongside the geometry of the bad fibers - specifically, the number of singular fibers, their component count, and their respective multiplicities. In fact, the fibration $X \to C$ can be factored through a smooth algebraic space $\mathcal{C}$ that is birational to $C$, such that the induced map $X \to \mathcal{C}$ is an étale $\mathbb{A}^1$-bundle. The space $\mathcal{C}$ obtained via this factorization precisely encodes the data of the bad points, the number of components of the fibers over them, and their corresponding multiplicities \cite{doub}. 
%(Cite Duoboloz paper here).









\


We first prove that (Theorem \ref{sing-etale-bundle}) if   $\rho: X \to \mathcal{C}$ is an étale locally trivial $\mathbb{A}^1$-bundle, such that $X$ is a smooth variety over $k$ and $\mathcal{C}$ is a smooth algebraic space over $k$, then the induced morphism $\rho_*: Sing_*(X) \to Sing_*(\mathcal{C})$ is a Nisnevich local weak equivalence. Note that, when $\mathcal{C}$ is a smooth scheme and $\rho$ is a Nisnevich $\mathbb{A}^n$-bundle then this result obviously follows from homotopy invariance and Nisnevich Mayer Vietoris \cite{krishnakumar}.
% (Give reference to Duoboloz student paper). 
Our theorem (Theorem \ref{sing-etale-bundle}) has stronger implications.

\

For instance, if $\rho : X \to \mathcal{C}$ as in the Theorem  \ref{sing-etale-bundle}, we get  $\sS(X) \cong \sS(\mathcal{C})$ as Nisnevich sheaves (see \cite[Definition 2.9]{bhs} for the construction of $\sS(X)$). This observation, together with rigidity property of the base curve and techniques of ghost homotopy, is used to prove Theorem \ref{connected-base-new} and Theorem \ref{rigid-base}. These prove that for any $\mathbb{A}^1$-fibration $X \to C$ over a smooth curve $C$ we have $\pi_0^{\mathbb{A}^1}(X)$ is $\mathbb{A}^1$-invariant. A corollary of these theorems is  that for all smooth affine complex surfaces $X$ of negative log Kodaira dimesnion, $\pi_0^{\mathbb{A}^1}(X)$ is $\mathbb{A}^1$-invariant (Corollary \ref{negative kodaira invariant}).
%(Add a corollary in the end of section 4 stating and proving this result and refer it here).  
Using similar techniques as in \cite[Proposition 2.9]{choudhury2}, 
%proper reference to Biman's thesis papers) 
we get that for non-negative log Kodaira dimension smooth affine complex surfaces $X$, $\sS(X) = \sS^2(X)$ and therefore $\pi_0^{\mathbb{A}^1}(X)$ is homotopy invariant for all smooth affine complex surfaces. Homotopy invariance of $\pi_0^{\mathbb{A}^1}(X)$ of smooth proper complex surfaces has been proven in \cite[Corollary 3.15]{bhs} and \cite[Theorem 1.2]{bs}.  


\

Another direct consequence of our Theorem \ref{sing-etale-bundle} is that $Sing_*(X)$ is $\mathbb{A}^1$-local if and only if $Sing_*(\mathcal{C})$ is $\mathbb{A}^1$-local. The property of $Sing_*(S)$ for space $S$ being $\mathbb{A}^1$-local is extremely useful (see \cite{am}, \cite{asokI}, \cite{asokII}, \cite{balwe reductive}).
% (give references including Morel-Asok-Wendt ....). 
For instance if $S = \mathbb{A}^2 \setminus \left\{(0,0)\right\}$ \cite[Theorem 8.9]{mor} (see also \cite[Corollary 4.2.6]{asokII})
%( Give reference from Morel's book)
or $S = \mathbb{A}^n$, then $Sing_*(S)$ is $\mathbb{A}^1$-local. If $Y$ is a retract of a space $S$, such that $Sing_*(S)$ $\mathbb{A}^1$-local, then $Sing_*(Y)$ is $\mathbb{A}^1$-local. Therefore if a variety $Y$ is retract of $\mathbb{A}^n$ then $Sin_*(Y)$ has to be $\mathbb{A}^1$-local.  

\

We show that if $X$ is a Koras-Russell threefold of the first kind then $Sing_*(X)$ is $\mathbb{A}^1$-local (Theorem \ref{application purity koras}). We use exactly the same (as in 
\cite{df}) using weak unstable five lemma applied to the same diagram of cofibration. It is also useful to note that for all $\mathbb{A}^1$-fibration $\rho : X \to C$ with $X$ smooth affine complex surface  such that the algebraic space curve $\mathcal{C}$ is $\mathbb{A}^1$-rigid we get $Sing_*(X)$ is $\mathbb{A}^1$-local (Corollary \ref{naive spaces}).
% (Add a corollary in section 4 and refer it here).


\

The remainder of this article is organized as follows. In Section 2,  we briefly recall the construction of unstable motivic homotopy category and review the core properties which are relevant to our main results. In particular, we prove that the map  $Sing_*(\mathbb{A}^n) \to Spec(k)$ is a sectionwise trivial fibration (Example \ref{kan fibrant An}, Appendix add general horn filling formula and refer here, Lemma \ref{contractible An}). Finally, we establish the main result,  Theorem \ref{sing-etale-bundle}.

\

In Section 3,  we prove Lemma \ref{key lemma} which shows that $Sing_*$ commutes with good quotient.  Using \ref{sing-etale-bundle}, purity, excision and weak five lemma,  we prove
Theorem  \ref{application purity koras}.


\

The first important result in Section 4 is Proposition \ref{ghost homotopy implies dominant 1}. This shows that existence of nontrivial ghost homotopy on a surface $X$ implies that the surface $X$ is log-uniruled. The method is similar to the method used in (\cite{cb} \cite{choudhury2}).
%(cite Biman and my papers). 
This proposition is used to prove that for smooth affine complex surfaces $X$ of log Kodaira dimension non negative, we have $S(X) = S^2(X)$ and therefore $\pi_0^{\mathbb{A}^1}(X)$ is $\mathbb{A}^1$-invariant (Theorem \ref{non-uniruled}, Corollary \ref{invariant non-uniruled}). Finally, we use Theorem \ref{rigid-base} Theorem \ref{connected-base-new} to conclude that for any smooth affine complex surface $X$, $\pi_0^{\mathbb{A}^1}(X)$ is $\mathbb{A}^1$-invariant. 

\

Section 5 is dedicated to computing the $\mathbb{A}^1$-homotopy types of surfaces with negative log Kodaira dimension. In particular, we provide a complete classification of the $\mathbb{A}^1$-homotopy types for $\mathbb{A}^1$-fibrations $X \to C$ whose fiber components contain no non-reduced (thick) $\mathbb{A}^1$'s. Finally, the Appendix contains several explicit computations related to the horn-filling property; since these results lack a standard reference in the literature, we have included them here for completeness.


\

\textbf{Acknowledgements} :

The authors would like to thank Adrien Dubouloz and Anand Sawant for their helpful comments and suggestions. This article constitutes part of the Ph.D. thesis of the second author, who expresses deep gratitude to the Indian Statistical Institute for providing excellent resources and support throughout this work and to Aranya Kumar Bal for stimulating discussion about $\A^1$-fibration to $\A^1$, especially for Example \ref{rigid example}. The first author was supported by the ANRF/ARGM; this work is part of the project ANRF/ARGM/2025/000240/MTR. The third author thanks R. V. Gurjar for pointing out a mistake in an earlier version of the article. The third author thanks Aravind Asok for his encouragement and discussions. The third author also thanks Charanya Ravi, Peter Haine, Haoyang Liu for helpful discussions. A part of this work was done, while the third author was at Tata Institute of Fundamental Research, Mumbai. The third author would like to thank TIFR and the University of Southern California for providing resources during this work.




	\section{Unstable Motivic homotopy Category}
In this section, we first recall the construction of unstable $\mathbb{A}^1$-homotopy category $\mathcal{H}(k)$ from \cite{mv}, introduced by Morel-Voevodsky and we will discuss a few important properties of $\mathcal{H}(k)$. 
%relevant to this article. 
We prove Theorem \ref{sing-etale-bundle}, which says that an \'etale $\A^1$-bundle $\rho: X \to \mathcal{C}$ induces the Nisnevich local weak equivalence $\rho_*: Sing_*(X) \to Sing_*(\mathcal{C})$ (see also Theorem \ref{general bundle}, for similar result if $\rho$ is an ). Theorem \ref{sing-etale-bundle} is an important ingredient to prove the $\A^1$-locality of $Sing_*(X)$, if $X$ is a Koras-Russell threefold of the first kind (Theorem \ref{application purity koras}) and the $\A^1$-invariance of $\pi_0^{\A^1}(X)$, if $X$ is an $\A^1$-ruled surface (Theorem \ref{connected-base-new}).
%an important theorem of this article (Theorem \ref{general bundle}), which we will use se says that if $\rho: X \to Y$ is a vector bundle, then the induced map $\rho_*: Sing_*(X) \to Sing_*(Y)$ is a Nisnevich local weak equivalence. 
We will also recall the fact that for every $k$-scheme $U$, $Sing_*(\mathbb{A}^n_k)(U)$ is a Kan fibrant simplicial set, by computing the horn filling formula (see Appendix \ref{full hornfill An}). \par


		\subsection{Construction and properties}
 Let $k$ be a field and $Sm/k$ be the category of finite type, smooth schemes over $k$. Let $\Delta^{op}PSh(Sm/k)$ be the category of simplicial presheaves on $Sm/k$, which we call the category of spaces over $k$. The category $\Delta^{op}PSh(Sm/k)$ has a model structure, called the projective model structure; where fibrations are sectionwise Kan fibrations, 
%between the simplicial sets, 
weak equivalences are sectionwise weak equivalences between simplicial sets and the cofibrations are defined using the left lifting property with respect to the trivial fibrations.\par
The category $Sm/k$ is equipped with the Nisnevich topology, in which the stalks are the sections over the Henselian  local schemes, which are essentially smooth schemes over $k$. A morphism $\sX \to \sY$ in $\Delta^{op}PSh(Sm/k)$ is called a Nisnevich local weak equivalence, if for every essentially smooth Henselian local scheme $\sO$ over $k$, the morphism
$$\sX(\sO) \to \sY (\sO)$$ 
is a weak equivalence of simplicial sets. Inverting the local weak equivalences, we get the Nisnevich local model structure on $\Delta^{op}PSh(Sm/k)$. The associated homotopy category is called the simplicial homotopy category, denoted by $\sH_s(k)$. Then taking the left Bousfield localisation of the Nisnevich local model structure on $\Delta^{op}PSh(Sm/k)$ with respect to the projection maps $\A^1_k\times_k X \to X$, for $X \in Sm/k$, we get the $\mathbb{A}^1$-model structure and the associated homotopy category is called the $\A^1$-homotopy category, denoted by $\sH(k)$. If we start with the category of pointed spaces over $k$, the same constructions would give the pointed simplicial homotopy category $\sH_{s, \bullet}(k)$ and the pointed $\mathbb{A}^1$-homotopy category $\mathcal{H}_\bullet(k)$ respectively.
			\begin{remark}
    			\begin{enumerate}
        \item A morphism of spaces $f: \sX \to \sY$ is called a Nisnevich local weak equivalence (or simplicial weak equivalence), if it an isomorphism in $\mathcal{H}_s(k)$ and a morphism of spaces which induces an isomorphism in $\mathcal{H}(k)$, is called an $\mathbb{A}^1$-weak equivalence.
        \item For $X \in Sm/k$ and a Nisnevich covering $f: U \to X$ in $Sm/k$, the \v{C}ech hypercover $\text{\v{C}}(f) \to X$ associated to $f$ is a Nisnevich local weak equivalence \cite[Lemma 1.15, Section 2.1]{mv}. The Nisnevich local model structure is the left Bousfield localisation of the projective model structure at the \v{C}ech hypercovers \cite[Proposition 8.1]{dugger}.
        \item For a space $\mathcal{X}$ over $k$, the projection map $\mathcal{X} \times_k \mathbb{A}^1_k \to \mathcal{X}$ is an $\mathbb{A}^1$-weak equivalence.
        \item If $\rho: X \to Y$ is an  morphism in $Sm/k$, then it is an $\mathbb{A}^1$-weak equivalence. More generally, if $\rho: X \to Y$ is an $\mathbb{A}^n$-bundle in Nisnevich topology (Definition \ref{bundle}), then it is an $\mathbb{A}^1$-weak equivalence \cite[Example 2.3, Section 3.1]{mv}. In Theorem \ref{sing-etale-bundle} and Theorem \ref{general bundle}, we will prove that if $\rho: X \to Y$ is an \'etale $\mathbb{A}^1$-bundle or an , then the induced morphism $\rho_*: Sing_*(X) \to Sing_*(Y)$ (Definition \ref{sing functor}) is a Nisnevich local weak equivalence.
    			\end{enumerate}
			\end{remark}
\begin{definition}($\A^1$-local) \cite[Definition 2.1, Section 3.2]{mv} \label{a1 local definition} \hfill
	%\begin{enumerate}
 A simplicial presheaf $\sX$ 
	%	%in $\Delta^{op}PSh(Sm/k)$ 
		is called $\A^1$-local if for any simplicial presheaf $\sY$, 
		%in $\Delta^{op}PSh(Sm/k)$ 
		the map
		$$Hom_{\sH_s(k)}(\sY,\sX)\to Hom_{\sH_s(k)}(\sY\times\A^1,\sX),$$
		induced by projection $\sY\times\A^1 \to \sY$, is a bijection of sets. 
	%	\item	A presheaf of sets $\sF$ is called an $\A^1$-invariant if for any $U\in Sm/k$ the map
	%	$$\sF(U)\to \sF(U\times\A^1_k)$$ 
	%	induced by the projection $U\times\A^1_k \to U$ is a bijection of sets.
	%	\item A scheme $X$ over $k$ is called an $\A^1$-rigid, if $X$ is $\A^1$-invariant, as representable sheaf \cite[Example 2.4]{mv}. 
%	\end{enumerate}
\end{definition}
\begin{remark}
A space $\sX$ over $k$ is fibrant in the $\A^1$-model structure if and only if it is fibrant in the Nisnevich local model structure and it is $\mathbb{A}^1$–local
\cite[Proposition 3.19, Section 2.3]{mv}.
\end{remark}
%\begin{remark} \label{a1 local properties/example}
%A
%\end{remark}

%\begin{remark}
  %  in Theorem \ref{purity sing}, we will prove a stronger version of purity isomorphism in a particular case. If $X$ admits an étale morphism $q: X \to \mathbb{A}^n_k$ such that the smooth closed subscheme $Z$ is isomorphic to $\mathbb{A}^m_k \times \{0, \dots, 0\}$ ($m \leq n$) then there is a Nisnevich local weak equivalence
  %  $$Sing_*(Th(N_{X, Z})) \sim Sing_*(X/(X\setminus Z)).$$
%\end{remark}
%\begin{example}
  %  \begin{enumerate}
   %     \item The Thom space of the trivial bundle $X \times_k \mathbb{A}^n_k \to X$ 
        %over $X \in Sm/k$ 
    %    is isomorphic to 
        %the pointed Nisnevich sheaf 
    %    $X_+ \wedge T^n$, where $T$ is the quotient sheaf $\mathbb{A}^1_k /\mathbb{G}_m$ and $T^n$ is the $n$-fold smash product of $T$.
        
    %    By purity isomorphism, this
  %  \end{enumerate}
       
%\end{example}


\begin{definition} \cite[Section 2.3]{mv} \label{sing functor}
Suppose, $\mathcal{X}$ is a space over $k$.
%For a space $\mathcal{X}$, 
The functor 
$$Sing_*: \Delta^{op}PSh(Sm/k) \to \Delta^{op}PSh(Sm/k)$$ 
is defined as
$$Sing_*(\mathcal{X})(U)_n=\mathcal{X}(\Delta^n_a \times_k U)_n,$$
where $\Delta^n_a$ is the cosimplicial object in $Sm/k$ defined as, 
$$\Delta^n_a = \Spec\ \frac{k[x_0, x_1,\dots, x_n]}{(\sum_{i=0}^n x_i - 1)}.$$ 
So, $\Delta^n_a \cong \mathbb{A}^n_k$. The boundary maps $d_i$ and the degeneracy maps $s_i$ are described in the Appendix \ref{maps of sing}.
\end{definition}

\begin{remark}
\begin{enumerate}
\item The canonical morphism $\mathcal{X} \to Sing_*(\mathcal{X})$ is an $\mathbb{A}^1$-weak equivalence \cite[Corollary 3.8, Section 2.3]{mv}.
\item If $f: \mathcal{X} \to \mathcal{Y}$ is 
%an $\mathbb{A}^1$-fibration (that is, 
a fibration in the $\A^1$-model structure, then the morphism $Sing_*(f): Sing_*(\mathcal{X}) \to Sing_*(\mathcal{Y})$ is also a fibration in the $\A^1$-model structure
%n $\mathbb{A}^1$-fibration 
\cite[Corollary 3.13, Section 2.3]{mv}.
%\item For $X \in Sm/k$, the connected component sheaf of $Sing_*(X)$ i.e. the Nisnevich sheaf associated to the presheaf $\pi_0(Sing_*(X))$ is the $\mathbb{A}^1$-chain connected component sheaf of $X$ (denoted as $\pi_0^{ch}(X)$), defined by Asok-Morel \cite[Definition 2.2.4]{am}.
\end{enumerate}
\end{remark}

\begin{definition} \label{kan fibrant definition} \cite[Definition 1.3]{may}
A simplicial set $X$ is called Kan fibrant if for every $n$ and 
given $l$-th horn in $X$ i.e. a map $\phi: \Lambda^n_l \to X$, $0 \leq l \leq n$, there is a map $\Phi: \Delta^n \to X$ such that $\phi = \Phi \circ \theta$, where $\theta: \Lambda^n_l \to \Delta^n$ is the inclusion of the $l$-th horn in $\Delta^n$. Equivalently, 
given $n$-many $(n-1)$ simplices $x_0,\dots, x_{l-1}, x_{l+1},\dots, x_n$ of $X$ satisfying the compatibility condition 
$$d_i x_j = d_{j-1} x_i, \text{ for }i < j, i, j \neq l, \ d_i: X_n \to X_{n-1} \text{ are the boundary maps },$$
then there is an $n$-simplex $x$ such that $d_i x  = x_i$ for all $i \neq l$. 
\end{definition}
\begin{example} \label{kan fibrant An}
Any simplicial group is Kan fibrant \cite[Theorem 17.1]{may}. Thus for any $k$-scheme $U$, $Sing_*(\mathbb{A}^m_k)(U)$ is Kan fibrant, since 
%for every $U \in Sm/k$, $Sing_*(\mathbb{A}^m_k)(U)$ is a simplicial abelian group (indeed, 
$\mathbb{A}^m_k$ is an affine group scheme (see also Remark \ref{boundary and horn filling}(3)). The retract of a Kan fibrant simplicial set is Kan fibrant.
\end{example}
Using \cite[Lemma 3.1]{curtis-simplicial} we give an explicit formula of $2$-horn filling of the sections of $Sing_*(\mathbb{A}^m_k)$ in the next example, over $U \in Sm/k$.



%In the next subsection 
%%in Lemma \ref{horn filling}, 
%using these formulas we will describe the horn filling of $Sing_*(\mathbb{A}^m_k)(U)$ in case of $2$-horn (that is, $n=2$ in Definition \ref{kan fibrant definition}).
%is Kan fibrant in degree $2$ (Definition \ref{kan fibrant degree definition}).

%\subsection{Kan Fibrant in degree $2$} 
%\subsection{Horn filling formula of $\Lambda^2_l \to Sing_*(\mathbb{A}^m_k)(U)$, $0 \leq l \leq 2$}
%\pi_0^{\mathbb{A}^1}(X)(\Spec\ k)$ is trivial and Kan Fibrant Property of 
%$Sing_*(X)(\Spec\ k)$ is Kan Fibrant}
%\label{sec:kan fibrant degree}

%In this section we define the notion of a simplicial set being Kan fibrant in degree $n$ (Definition \ref{kan fibrant degree definition}). In Lemma \ref{horn filling}, using the formulas obtained in Subsection \ref{Kan fibrant An} we show that the sections of $Sing_*(\mathbb{A}^m_k)$ are Kan fibrant in degree $2$. In Corollary \ref{connected Kodaira dimension 2}, we prove that if $X$ is a smooth affine surface over an algebraically closed field $k$ of characteristic zero along with $\pi_0^{\mathbb{A}^1}(X)(\Spec\ k)$ is trivial and $Sing_*(X)(\Spec\ k)$ is Kan fibrant in degree $2$, then $X$ has negative logarithmic Kodaira dimension. 
%and we use this condition to show that an


%%\begin{definition} \cite[Definition 2.18]{choudhury2} \label{kan fibrant degree definition}
%%A simplicial set $X$ is called Kan fibrant in degree $n$ if for each $l$-th horn $\Lambda^n_l$ where $0 \leq l \leq n$, a map $\phi: \Lambda^n_l \to X$ can be extended to $\Tilde{\phi}: \Delta^n \to X$. 
%%\%%end{definition}
%%\begin{example} 
%%\begin{enumerate}
%%Any Kan fibrant simplicial set is Kan fibrant in degree $n$ for any $n$. 
%%The retract of a Kan fibrant simplicial set is Kan fibrant. 
%%Any simplicial group is Kan fibrant in degree $n$ \cite[Theorem 17.1]{may}. For example, the sections of $Sing_*(\mathbb{A}^m_k)$ are Kan fibrant in degree $n$ for any $n$.
%%Any simplicial group is Kan fibrant. For example, $Sing_*(\mathbb{A}^m)$ is sectionwise Kan fibrant. By the phrase "$Sing_*(X)$ is Kan fibrant" for a $k$-variety $X$, we mean $Sing_*(X)(\Spec\ k)$ is Kan fibrant. If $Sing_*(X)$ and $Sing_*(Y)$ are Kan fibrants for $k$-varieties $X$ and $Y$, $Sing_*(X \times Y)$ is also Kan fibrant. 
%
%%Thus in particular if we are given two morphisms $f, g : \mathbb{A}^1_k \to \mathbb{A}^m_k$ such that $f(1) = g(0)$, then the morphism
%%$h: \mathbb{A}^2_k \to \mathbb{A}^m_k$ defined as
%%$$(t_1, t_2) \mapsto (f(1-t_1)+g(t_2) - f(1))$$
%%satisfies $h(1-x, 0) = f(x)$ and $h(0, x) = g(x)$. So $Sing_*(\mathbb{A}^m_k)(\Spec\ k)$ is Kan fibrant in degree $2$. Therefore, if there are two intersecting $\mathbb{A}^1$'s in $\mathbb{A}^2_k$, then we can extend it to get a morphism from $\mathbb{A}^2_k$ to $\mathbb{A}^2_k$.
%%\end{enumerate}
%%\end{example}
%%In the next lemma (Lemma \ref{horn filling}), using the formulas in Subsection \ref{Kan fibrant An} we show that the sections of $Sing_*(\mathbb{A}^m_k)$ are Kan fibrant in degree $2$.
\begin{example}{[$2$-horn filling formula for $Sing_*(\A^m_k)(U)$, see also Section \ref{full hornfill An}]} \label{2 horn}
	Let $U$ be any $k$-scheme and
	suppose, we are given a morphism 
	$$\phi: \Lambda^2_l \to Sing_*(\mathbb{A}^m_k)(U), \text{ where } U \in Sm/k, l \in \{0, 1, 2\}.$$
	We will extend $\phi$ to a $2$-simplex of $Sing_*(\mathbb{A}^m_k)(U)$.
	We have the following isomorphism:
	\begin{align*}
		Sing_*(\mathbb{A}^m_k)(U)_n & = Hom_{Sm/k}(\Delta^n_a \times_k U, \mathbb{A}^m_k) \\
		& \cong Hom_{k-alg}(k[T_1,\dots, T_m], \frac{A[x_0,\dots, x_n]}{(\sum_{i=0}^n x_i - 1)}),
	\end{align*}
	where $A$ is the ring of regular functions $\mathcal{O}(U)$.
	Here are three cases according to the $l$-th horn, $l = 0,1,2$. \par
	\begin{enumerate}[start=1,label={\bfseries Case \arabic*:}]
		\item $l=0$.

 Suppose we are given
		$$\phi_1, \phi_2: k[T_1,\dots,T_m] \to \frac{A[x_0, x_1]}{(x_0 + x_1 - 1)}$$
		such that $d^1 \circ \phi_2 = d^1 \circ \phi_1$.  Since $T_j$'s are the independent variables, suppose $\phi_1, \phi_2$ are given by
		$$T_j \xmapsto{\phi_1} \overline{f_j}, \ f_j \in A[x_0, x_1] \text{ and}$$
		$$T_j \xmapsto{\phi_2} \overline{g_j}, \ g_j \in A[x_0, x_1]$$
		and $f_j, g_j$ satisfy
		$$\overline{g_j(x_0, 0)} = \overline{f_j(x_0, 0)}.$$
		Define $\Phi: k[T_1,\dots,T_m] \to \frac{A[x_0, x_1, x_2]}{(x_0 + x_1 +x_2-1)}$
		as 
		$$T_j \mapsto \overline{g_j(x_0, x_1+x_2) + f_j(x_0+x_1, x_2) - g_j(x_0 + x_1 , x_2) }.$$
		Then $d^1 \circ \Phi$ is given by 
		$$T_j \mapsto \overline{g_j(x_0, x_1) - g_j(x_0, x_1)+ f_j(x_0, x_1)}$$
		so, $T_j \xmapsto{d^1 \circ \Phi} \overline{f_j(x_0, x_1)}$. Thus, $d^1 \circ \Phi = \phi_1$. \\
		The map $d^2 \circ \Phi$ is given by
		$$T_j \mapsto \overline{g_j(x_0, x_1) - g_j(x_0+x_1, 0) + f_j(x_0+x_1, 0)}.$$
		Since $\overline{g_j(x_0, 0)} = \overline{f_j(x_0, 0)} \in \frac{A[x_0]}{(x_0 -1)}$, so $T_j \xmapsto{d^2 \circ \Phi} \overline{g_j(x_0, x_1)}$. Thus, $d^2 \circ \Phi = \phi_2$.
		\item $l = 1$.

		Suppose, we are given
		$$\phi_0, \phi_2: k[T_1,\dots,T_m] \to \frac{A[x_0, x_1]}{(x_0 + x_1 -1)}$$
		such that $d^0 \circ \phi_2 = d^1 \circ \phi_0$. Suppose $\phi_0, \phi_2$ are given by
		$$T_j \xmapsto{\phi_0} \overline{f_j}, \ f_j \in A[x_0, x_1] \text{ and}$$
		$$T_j \xmapsto{\phi_2} \overline{g_j}, \ g_j \in A[x_0, x_1]$$
		and $f_j, g_j$ satisfy
		$$\overline{g_j(0, x_0)} = \overline{f_j(x_0, 0)}.$$
		Define $\Phi: k[T_1,\dots,T_m] \to \frac{A[x_0, x_1, x_2]}{(x_0 + x_1 +x_2-1)}$
		as 
		$$T_j \mapsto \overline{f_j(x_0+ x_1,x_2) + g_j(x_0, x_1+ x_2) - f_j(x_0 + x_1+ x_2, 0) }, $$
		$$\text{so } T_j \xmapsto{\Phi} \overline{f_j(x_0+x_1, x_2) + g_j(x_0, x_1 + x_2) - f_j(1, 0) }.$$
		Similarly as in Case 1, we can check $d^0 \circ \Phi = \phi_0$ and $d^2 \circ \Phi = \phi_2$. 
		\item $l = 2$.

		Suppose, we are given
		$$\phi_0, \phi_1: k[T_1,\dots,T_m] \to \frac{A[x_0, x_1]}{(x_0 + x_1 -1)}$$
		such that $d^0 \circ \phi_1 = d^0 \circ \phi_0$. Suppose $\phi_0, \phi_1$ are given by
		$$T_j \xmapsto{\phi_0} \overline{f_j}, \ f_j \in A[x_0, x_1] \text{ and}$$
		$$T_j \xmapsto{\phi_1} \overline{g_j}, \ g_j \in A[x_0, x_1]$$
		and $f_j, g_j$ satisfy
		$$\overline{g_j(0, x_0)} = \overline{f_j(0, x_0)}.$$
		Define $\Phi: k[T_1,\dots,T_m] \to \frac{A[x_0, x_1, x_2]}{(x_0 + x_1 +x_2-1)}$
		as 
		$$T_j \mapsto \overline{f_j(x_0+ x_1,x_2) + g_j(x_0, x_1+ x_2) - f_j(x_0, x_1+ x_2) }.$$
		$$T_j \xmapsto{\Phi} \overline{f_j(x_0+x_1, x_2) - f_j(1, 0) + g_j(x_0, x_1 + x_2)}.$$
		Similarly as in Case 1, we can check $d^0 \circ \Phi = \phi_0$ and $d^1 \circ \Phi = \phi_1$. 
		
		Therefore, for any horn $\phi: \Lambda^2_l \to Sing_*(\A^m_k)(U)$, where $0 \leq l \leq 2$, there is a $2$-simplex $\Phi:\Delta^2 \to Sing_*(\A^m_k)(U)$ that extends the horn $\phi$, that is, $\phi= \Phi \circ \theta$, where $\theta: \Lambda^2_l \to \Delta^2$ is the inclusion (Definition \ref{kan fibrant definition})
	\end{enumerate}
\end{example}
In \ref{full hornfill An} using \cite[Lemma 3.1]{curtis-simplicial}, we write $l$-th horn (inside $\Delta^n$) filling formula, for any $l,n$ with $0 \leq l \leq n$.



%	Therefore, $Sing_*(\mathbb{A}^m_k)(U)$ is Kan fibrant. 


%%
%%\
%%

%%
%%\
%%



%The general horn filling formaula can be done similarly.
\begin{remark} \cite[Example 2.19]{choudhury2}
Thus in particular, if we are given two morphisms $f, g : \mathbb{A}^1_k \to \mathbb{A}^m_k$ such that $f(1) = g(0)$, then the morphism
$h: \mathbb{A}^2_k \to \mathbb{A}^m_k$ defined as
$$(t_1, t_2) \mapsto (f(1-t_1)+g(t_2) - f(1))$$
satisfies $h(1-x, 0) = f(x)$ and $h(0, x) = g(x)$. 
%So $Sing_*(\mathbb{A}^m_k)(\Spec\ k)$ is Kan fibrant in degree $2$. 
Therefore, if there are two intersecting $\mathbb{A}^1$'s in $\mathbb{A}^2_k$, then we can extend it to get a morphism from $\mathbb{A}^2_k$ to $\mathbb{A}^2_k$.
%\end{enumerate}
\end{remark}
We end this subsection with the following lemma about section-wise simplicial contractibility of $Sing_*(\A^m_k)$ (see also Remark \ref{boundary and horn filling}(5)). 
%(Corollary \ref{connected Kodaira dimension 2}) to Proposition \ref{connected implies}. The following corollary is taken from \cite[Corollary 2.12]{cb2}.

%\begin{corollary} \label{connected Kodaira dimension 2}
%Let $X \in Sm/k$ be an affine surface, where $k$ is an algebraically closed field of characteristic zero. Suppose that $\pi_0^{\mathbb{A}^1}(X)(\Spec\ k)$ is trivial and $Sing_*(X)(\Spec\ k)$ is Kan fibrant in degree $2$. Then $X$ has negative logarithmic Kodaira dimension.
%\end{corollary}
%\begin{proof}
%Since $\pi_0^{\mathbb{A}^1}(X)(\Spec\ k)$ is trivial, so by the Proposition \ref{connected implies} either there is some $Y \in Sm/k$ irreducible along with a non-constant homotopy $H: \mathbb{A}^1_k \times_k Y \to X$ which is dominant or there are two intersecting $\mathbb{A}^1$'s in $X$. For the first case, we proceed as in Corollary \ref{connected Kodaira dimension 1} to conclude that $X$ has negative logarithmic Kodaira dimension. For the second case, since $Sing_*(X)(\Spec\ k)$ is Kan fibrant in degree $2$, there is a dominant morphism $\phi: \mathbb{A}^2_k \to X$. Therefore $X$ has logarithmic Kodaira dimension $-\infty$ by \cite[Proposition 11.4]{iitaka}.
%\end{proof}

\begin{lemma} \cite[Section 2.3, Corollary 3.5]{mv} \label{contractible An}
    Suppose, $U\in Sm/k$. Then the Kan fibrant simplicial set $Sing_*(\mathbb{A}^n_k)(U)$ is contractible.
\end{lemma}
\begin{proof}
    The multiplication map 
    $$\gamma: \mathbb{A}^n_k \times_k \mathbb{A}^1_k \to \mathbb{A}^n_k \text{ given by } (s,t) \mapsto st$$
    gives naive $\mathbb{A}^1$-homotopy between the identity map and the constant map to a $k$-point. Since $Sing_*$ functor commutes with the product, so $\gamma$ induces a map
    $$Sing_*(\gamma): Sing_*(\mathbb{A}^n_k) \times Sing_*(\mathbb{A}^1_k) \to Sing_*(\mathbb{A}^n_k).$$
    Consider the $1$-simplex in $\theta= \mathrm{Id}_{\A^1}\in Sing_*(\A^1_k)$, given by $\theta(0)=\Spec(k)\xrightarrow{s_0}\A^1_k,\mathrm{\ and \ } \theta(1)=\Spec(k)\xrightarrow{s_1}\A^1_k$.
    Consider the composition of morphisms
    $$\Gamma: Sing_*(\mathbb{A}^n_k) \times \Delta^1 \xrightarrow{(Id, \theta)} Sing_*(\mathbb{A}^n_k) \times Sing_*(\mathbb{A}^1_k) \xrightarrow{Sing_*(\gamma)} Sing_*(\mathbb{A}^n_k).$$ 
    Then we have -
    \begin{align*}
    	\Gamma(0) &= Sing_*(\A^n_k)\xrightarrow{\mathrm{Const}_{(0,\dots,0)}} Sing_*(\A^n_k)\\
    	\Gamma(1) &= Sing_*(\A^n_k) \xrightarrow{\mathrm{Id}} Sing_*(\A^n_k).
    \end{align*}
     This defines a simplicial homotopy between identity map and a constant map. Hence $Sing_*(\A^n_k)(U)$ is contractible for any $U\in Sm/k$.
%    The simplicial homotopy $\Gamma: Sing_*(\mathbb{A}^1_k) \times \Delta^1 \to Sing_*(\mathbb{A}^1_k)$ is a simplicial homotopy equivalence.
%    Suppose, $d_0, d_1: Sing_*(\mathbb{A}^1_k) \to Sing_*(\mathbb{A}^1_k) \times \Delta^1$ are two inclusions, induced by two $0$-simplices $\Delta^0 \to \Delta^1$. The morphism of simplicial sheaves
%    $$\Gamma \circ d_0: Sing_*(\mathbb{A}^1_k) \to Sing_*(\mathbb{A}^1_k)$$
%    is the constant map given by 
%    $$Sing_*(\mathbb{A}^1_k) \to \Spec\ k \xrightarrow{ \text{inclusion } \{0\} \hookrightarrow \mathbb{A}^1_k} Sing_*(\mathbb{A}^1_k)$$
%    and 
%    $$\Gamma \circ d_1: Sing_*(\mathbb{A}^1_k) \to Sing_*(\mathbb{A}^1_k)$$
%    is the identity map. For a $k$-scheme $U$, $\Gamma$ induces simplicial homotopy equivalence
%    $$\Gamma_U: Sing_*(\mathbb{A}^1_k)(U) \times \Delta^1 \to Sing_*(\mathbb{A}^1_k)(U)$$
%    between the identity map and the constant map on $Sing_*(\mathbb{A}^1_k)(U)$. Therefore, the Kan fibrant simplicial set $Sing_*(\mathbb{A}^1_k)(U)$ is contractible. Since $Sing_*$ functor commutes eith limits we have $Sing_*(\A^m_k)(U)$ is also contractible.
\end{proof}


%\begin{lemma} \cite[Corollary 3.5, Section 2.3]{mv} \label{contractible An}
%    Suppose, $U$ is a $k$-scheme. Then the Kan fibrant simplicial set $Sing_*(\mathbb{A}^m_k)(U)$ is contractible.
%\end{lemma}
%\begin{proof}
%    The multiplication map 
%    $$\gamma: \mathbb{A}^1_k \times_k \mathbb{A}^1_k \to \mathbb{A}^1_k \text{ given by } (s,t) \mapsto st$$
%    gives naive $\mathbb{A}^1$-homotopy between the identity map and the constant map to a $k$-point. Since $Sing_*$ functor commutes with the fiber product, so $\gamma$ induces a map
%    $$Sing_*(\gamma): Sing_*(\mathbb{A}^1_k) \times Sing_*(\mathbb{A}^1_k) \to Sing_*(\mathbb{A}^1_k).$$
%    Consider the map,
%    $$\theta: \Delta^1 \to Sing_*(\mathbb{A}^1_k)$$
%    given by the identity map on $\mathbb{A}^1_k$ as an $1$-simplex of $Sing_*(\mathbb{A}^1_k)(Spec \ k)$.
%    Consider the composition of morphisms
%    $$\Gamma: Sing_*(\mathbb{A}^1_k) \times \Delta^1 \xrightarrow{(Id, \theta)} Sing_*(\mathbb{A}^1_k) \times Sing_*(\mathbb{A}^1_k) \xrightarrow{Sing_*(\gamma)} Sing_*(\mathbb{A}^1_k)$$
%    The simplicial homotopy $\Gamma: Sing_*(\mathbb{A}^1_k) \times \Delta^1 \to Sing_*(\mathbb{A}^1_k)$ is a simplicial homotopy equivalence.
%    Suppose, $d_0, d_1: Sing_*(\mathbb{A}^1_k) \to Sing_*(\mathbb{A}^1_k) \times \Delta^1$ are two inclusions, induced by two $0$-simplices $\Delta^0 \to \Delta^1$. The morphism
%    $$\Gamma \circ d_0: Sing_*(\mathbb{A}^1_k) \to Sing_*(\mathbb{A}^1_k)$$
%    is the constant map given by 
%    $$Sing_*(\mathbb{A}^1_k) \to Spec \ k \xrightarrow{ \text{inclusion } \{0\} \hookrightarrow \mathbb{A}^1_k} Sing_*(\mathbb{A}^1_k)$$
%    and 
%    $$\Gamma \circ d_1: Sing_*(\mathbb{A}^1_k) \to Sing_*(\mathbb{A}^1_k)$$
%    is the identity map. For a $k$-scheme $U$, $\Gamma$ induces simplicial homotopy equivalence
%    $$\Gamma_U: Sing_*(\mathbb{A}^1_k)(U) \times \Delta^1 \to Sing_*(\mathbb{A}^1_k)(U)$$
%    between the identity map and the constant map on $Sing_*(\mathbb{A}^1_k)(U)$. Therefore, the Kan fibrant simplicial set $Sing_*(\mathbb{A}^1_k)(U)$ is contractible.
%\end{proof}
\begin{remark} (see also Remark \ref{boundary and horn filling}(5)) \label{boundary lifting question}
     Let $U$ be a smooth $k$-scheme. Then $Sing_*(\mathbb{A}^m_k)(U)$ is Kan fibrant simplicial set, which is also contractible (Example \ref{kan fibrant An} and Lemma \ref{contractible An}). Therefore, the constant map $Sing_*(\mathbb{A}^m_k)(U) \to \bullet$ is a trivial Kan fibration. Therefore, $Sing_*(\mathbb{A}^m_k)(U)$ has boundary lifting property:
    
		\[\xymatrix{\partial \Delta^n \ar@{^{(}->}[d]^i \ar[r] &Sing_*(\mathbb{A}^m_k)(U)\\
		\Delta^n \ar@{.>}[ru]^{\exists}} \]
        This means, given a collection of morphisms $f_0, \dots, f_n: \Delta^{n-1}_a \times U \to \mathbb{A}^m_k$ satisfying the compatibility conditions (i.e., $d^if_j = d^{j-1}f_i$ for $0 \leq i<j \leq n$), there exists an $n$-simplex $f: \Delta^n_a \times U \to \mathbb{A}^m_k$ such that $d_i(f) = f_i$, for every $i$. 
        
%        Therefore, it is natural to ask if we are given such $f_i$'s (that satisfy the compatibility conditions), then to find a formula, that is a morphism $f: \A^n_U \to \A^m_k$ (like the horn-filling formula in Example \ref{2 horn} or Section \ref{full hornfill An}) that would lift the boundary given by $\{f_i\}_i$ to an $n$-simplex $f$.
\end{remark}


\subsection{A local weak equivalence}
	In this subsection we prove Theorem \ref{sing-etale-bundle},
% one of the main theorem of this article 
which will be used substantially in the later sections. We start by recalling few definitions.  

 We recall the definition and few properties of $\mathbb{A}^1$-fibration from \cite[Chapter 2]{masuda} and \cite[Chapter 3]{miyanishicrm}. Let 
%$X$ be a smooth 
%affine surface 
%variety over
$k$ be an algebraically closed field of characteristic zero.
%and $\mathcal{C}$ be a smooth algebraic space over $k$. 
%Suppose, $\pi: X \to \mathcal{C}$ is a morphism.
%over a smooth curve (\cite[Chapter 2]{masuda}, \cite[Chapter 3]{miyanishicrm}).
	\begin{definition} \label{bundle}
    %Suppose, $\pi: X \to \mathcal{C}$ is a morphism, where $X$ is a smooth variety over $k$ and $\mathcal{C}$ is a smooth algebraic space over $k$.
 Let $\mathcal{C}$ be a smooth algebraic space over $k$ and $X$ be a smooth variety over $k$. A morphism $\pi: X \to \sC$ is called an $\A^1$-bundle (with respect to the Zariski, Nisnevich or \'etale topology),
        %locally trivial) 
        if there exists a covering (Zariski, Nisnevich or \'etale covering respectively)
        %Zariski, nisnevich or étale resp.) 
        $U \to \sC$ such that the fiber product $X \times_\sC U$, 
        which is a scheme 
        %by definition of algebraic spaces 
       \cite[Corollary 4.2.3]{alpher},
        is isomorphic to $\A^1_k \times_k U$ over $U$. If $\pi$ is an $\mathbb{A}^1$-bundle with respect to the Zariski topology, then we will simply call $\pi$ to be an $\mathbb{A}^1$-bundle.
	\end{definition}
	\begin{remark} \label{a1 bundle properties}
	    \begin{enumerate}
	    
          \item Every fiber over a point $x$ of an $\mathbb{A}^1$-bundle is isomorphic to $\mathbb{A}^1_{\kappa(x)}$, where $\kappa(x)$ is the residue field of $x$. 
\item An $\A^1$-bundle $\pi: X \to \sC$ in the Nisnevich topology (hence in the Zariski topology also) is an $\A^1$-weak equivalence \cite[Section 3.1, Example 2.3]{mv}. In Theorem \ref{sing-etale-bundle}, we will prove that if $\pi$ is an \'etale $\A^1$-bundle, then $\pi_*: Sing_*(X) \to Sing_*(\sC)$ is a Nisnevich local weak equivalence. Therefore, an $\A^1$-bundle (in any topology Zariski, Nisnevich or \'etale) is an $\A^1$-weak equivalence (Corollary \ref{factorisation A1 equivalence}).
%So an $\mathbb{A}^1$-bundle $\pi: X \to C$ from a smooth affine surface to a smooth curve $C$ is an $\mathbb{A}^1$-fibration. However, the converse is true if every scheme-theoretic fiber of $\pi$ is irreducible and reduced \cite[Proposition 3]{doub}.
          %\item An $\mathbb{A}^1$-bundle is an $\mathbb{A}^1$-weak equivalence \cite[Theorem 5.3.1.3]{as}.
	        \item A line bundle is an $\mathbb{A}^1$-bundle (with respect to the Zariski topology). But the converse is not true. The $\A^1$-bundle in \cite[Example 2]{doub} does not admit a section, therefore it is not a line bundle.
            \item An $\mathbb{A}^1$-bundle corresponds to a line bundle, if the base is an affine scheme. Indeed, a line bundle over a scheme $X$ correspond to a $\mathbb{G}_m$-torsor over $X$; which correspond to an element of $H^1_{\text{Zar}}(X, \mathbb{G}_m)$ \cite[Corollary 4.7]{milne}. On the other hand, an $\mathbb{A}^1$-bundle over $X$ corresponds to an $\Aut(\A^1_k)$-torsor over $X$ (where, $\Aut(\A^1_k)$ is the automorphism group of the affine line $\A^1_k$); which corresponds to an element of $H^1_{\text{Zar}}(X, \Aut(\A^1_k))$. But the Zariski group sheaf $\Aut(\A^1_k)$ is isomorphic to the semidirect product
            %The étale group sheaf $\text{Aff}_1$ is isomorphic to the semidirect product 
            $\mathbb{G}_m \ltimes \mathbb{G}_a$, which gives the short exact sequence of Zariski sheaves 
            %on essentially smooth schemes over $k$
		$$0 \rightarrow \mathbb{G}_a \rightarrow \Aut(\A^1_k) \rightarrow \mathbb{G}_m \rightarrow 0.$$
		This gives a long exact sequence in cohomology of $X$ (in Zariski topology)
		$$\dots \rightarrow H^1_{\text{Zar}}(X, \mathbb{G}_a) \rightarrow H^1_{\text{Zar}}(X, \Aut(\A^1_k)) \rightarrow H^1_{\text{Zar}}(X, \mathbb{G}_m) \rightarrow 
H^2_{\text{Zar}}(X, \mathbb{G}_a) \rightarrow 
\dots$$
		%Now by Hilbert's Theorem $90$, 
		%$$H^1_{\text{ét}}(\mathbb{A}^n_{\mathcal{O}}, \mathbb{G}_m) \cong H^1_{\text{Zar}}(\mathbb{A}^n_{\mathcal{O}}, \mathbb{G}_m) \cong Pic(\mathbb{A}^n_{\mathcal{O}}).$$
		%Since the Picard group functor is $\mathbb{A}^1$-invariant and any projective module over a local ring is free, so $H^1_{\text{ét}}(\mathbb{A}^n_{\mathcal{O}}, \mathbb{G}_m)$ is trivial. 
		%sheaf $\Spec\ k$. 
        Since $X$ is affine, so 
        $H^1_{\text{Zar}}(X, \mathbb{G}_a) \cong H^2_{\text{Zar}}(X, \mathbb{G}_a) \cong \{0\}$. Therefore,
        $$H^1_{\text{Zar}}(X, \Aut(\A^1_k)) \xrightarrow{\cong} H^1_{\text{Zar}}(X, \mathbb{G}_m).$$
		%On the other hand, si étale cohomology vanishes with respect to the structure sheaf of an affine scheme; so
		%$$H^1_{\text{ét}}(\mathbb{A}^n_{\mathcal{O}}, \mathbb{G}_a) \cong H^1_{\text{Zar}}(\mathbb{A}^n_{\mathcal{O}}, \mathbb{G}_a) \cong \{0\}.$$
	    \end{enumerate}
	\end{remark}
\begin{definition} \cite[Definition 2.1.1]{masuda} \label{a1 fibration definition}
Suppose, $F$ is a variety over $k$. A surjective morphism of smooth varieties 
        $$\pi: X \to C$$ 
   over $k$ is called an $F$-fibration, if the generic fiber $X_\eta := X\times_C \Spec \kappa(\eta)$ is isomorphic to $F \times_k \Spec \kappa(\eta)$ over $\kappa(\eta)$,
%i.e., $X_\eta \cong F\times_C \Spec(\kappa(\eta))$. H
where $\eta$ is the generic point of $C$. If $F$ is isomorphic to $\A^1_k$, then we call the morphism $\pi$ to be an $\A^1$-fibration. 
%         there is an open subset $U$ of $C$ such that for every $x \in U(k)$, 
%       % the scheme-theoretic fiber 
%        $\pi^{-1}(x)$ 
%        %(that is, $\pi^{-1}(x):=X \times_{C,x} \Spec\ k$) 
%        is isomorphic to $\mathbb{A}^1_k$. Here for a point $a \in C$ with residue field $F$, the fiber $\pi^{-1}(a)$ denotes the scheme-theoretic fiber of $\pi$ over the point $a$, which is by definition the fiber product $X \times_{C,a} \Spec\ F$.
	\end{definition}
\begin{notation}
For a morphism $\pi: X \to C$ of smooth varieties over $k$ and a point $a \in C$ with the residue field $\kappa(a)$, the fiber $\pi^{-1}(a)$ denotes the scheme-theoretic fiber of $\pi$ over $a$, which is by definition the fiber product $X \times_{\pi, a} \Spec (\kappa(a))$.
\end{notation}	
	\begin{remark} \label{properties a1 fibration}
Suppose, $X$ is a smooth affine surface and the base of the $\mathbb{A}^1$-fibration $\pi: X \to C$ is a smooth curve $C$. The base curve $C$ consists of the closed points (which have residue fields $k$) and the generic point $\eta$. 
    \begin{enumerate}
        \item The generic fiber of $\pi$ is isomorphic to $\mathbb{A}^1_K$, where $K=\kappa(\eta)$ is the function field of $C$ over $k$ \cite[Section 2.2]{kambayashi}.
        \item A closed fiber of $\pi$ is isomorphic to $\mathbb{A}^1_k$ if it is reduced and irreducible. 
        \item A closed fiber $F=\pi^{-1}(x)$ is called degenerate, if it is either reducible or non-reduced \cite[Definition 4]{doub}. Suppose,
        $F = \cup_{i} F_i$; $F_i$'s are the connected components of $F$. Then the local ring $\mathcal{O}_{F_i, \xi_i}$ of $F_i$ at its generic point $\xi_i$ is an Artinian local ring of length $m_i$ (say). The length $m_i =1$ if and only if $F_i$ is reduced. Scheme-theoretically, the fiber $F$ can be expressed as $$F=\sum_{i=0}^n m_i F_i.$$
        The associated reduced structure $(F_i)_{\text{red}}$ of $F_i$ is isomorphic to $\mathbb{A}^1_k$
        %Otherwise, every irreducible component of $\pi^{-1}(x)$ (which is also a connected component) is isomorphic to $\mathbb{A}^1_k$
        \cite[Chapter 3, Lemma 1.4.2]{miyanishicrm}.
        %, with respect to the associated reduced structure. 
    %For a closed point $x \in C$, reduced structure of the fiber $\pi^{-1}(x)$ is either isomorphic to the affine line $\mathbb{A}^1_k$ or every irreducible component of $\pi^{-1}(x)$ ( is isomorphic to $\mathbb{A}^1_k$ \cite[Chapter 3, Lemma 1.4.2]{miyanishicrm}.  
   \item An $\mathbb{A}^1$-fibration is generically trivial i.e. there is an open subset $U$ of $C$ such that the fiber $X \times_C U$ over $U$ is isomorphic to $\mathbb{A}^1_U$ \cite[Lemma 1.1]{threefolds}. So in particular, there are only finitely many points $x \in C(k)$ such that $\pi^{-1}(x)$ is not isomorphic to $\mathbb{A}^1_k$. 
   \end{enumerate}
	\end{remark}
	
	\begin{example} [\textbf{$\mathbb{A}^1$-fibrations}]\label{d-sphere}\hfill
    \begin{enumerate}
        \item The first family of examples of $\mathbb{A}^1$-fibrations are $\mathbb{A}^1$-bundles, since every fiber over a point of an $\mathbb{A}^1$-bundle
% $\pi: X \to C$
%         to a smooth curve $C$ 
         is isomorphic to $\mathbb{A}^1$ (Remark \ref{a1 bundle properties}).
% over a curve. 
The converse is true for an $\A^1$-fibration $\pi: X \to C$ from a smooth affine surface $X$ to a smooth curve $C$, if every scheme-theoretic fiber of $\pi$ is irreducible and reduced \cite[Proposition 3]{doub}.
        \item (\textbf{Danielewski surface} \cite[Chapter 5, Section 5.3.4]{as}) An $\A^1$-fibration is not necessarily an $\A^1$-bundle or an $\A^1$-weak equivalence. %hence is not always an $\mathbb{A}^1$-bundle. 
For example, suppose $X$ is a Danielewski surface given by
$$X \equiv \{x^nz=P(y)\} \subset \mathbb{A}^3_k = \Spec\ k[x,y,z],$$ 
where $P(y) \in k[t]$ is a separable polynomial.
%\cite[Section 5.3.4, Chapter 5]{as}. 
If the degree of $P(y)$ is at least $2$ then $k[x,y,z]/(x^nz-P(y))$ is not a UFD, so the surface $X$ is not isomorphic to affine plane. Indeed, if $deg(P(y))=d \geq 2$, then we can write $P(y) = a \prod_{i=1}^{d}(y - \alpha_i)$, for some $a, \alpha_1, \dots, \alpha_d \in k$. So in $\sO(X)$, there are two different ways to express $x^nz$ as products of irreducibles
$$x^nz = a \prod_{i=1}^{d}(y - \alpha_i) \in \sO(X).$$
Thus $\sO(X)$ is not a UFD.

Consider the morphism 
$$\text{pr}_x: X \to \mathbb{A}^1_k \text{ defined as } (x,y,z) \mapsto x$$
is an $\A^1$-fibration, since 
$$\text{pr}_x^{-1}(\mathbb{G}_m = \{x \neq 0\}) \cong \mathbb{A}^1_k \times_k \mathbb{G}_m.$$ 
In this example, all fibers of $\text{pr}_x$ are reduced and only the fiber over $0$ is reducible, which is the disjoint union of $d$-copies of $\mathbb{A}^1_k$'s, where $d=deg(P(y))$. The surface $X$ is $\A^1$-weak equivalent to the wedge sum of $(d-1)$-copies of $\mathbb{P}^1_k$ \cite[Chapter 5, Proposition 5.3.4.3]{as} (which is again $\A^1$-weakly equivalent to the affine line with $d$-many origins, see Section \ref{section a1 homotopy types}). 
%So if $d \geq 2$, $X$ is not $\A^1$-contractible 
In particular, if we take $d=2$ and $X \equiv \{x^2z=y(y-1)\}$, then $X$ is $\A^1$-weakly equivalent to $\mathbb{P}^1_k$ (which is again $\A^1$-weakly equivalent to the affine line with double origin, see Section \ref{section a1 homotopy types}).
%is isomorphic to $\mathbb{}$$
		%Let 
      %  $S \subset \A^3$ is a Danielewski surface defined by the polynomial $x^2z=y(y-1)$ \cite[Section 5.3.4, Chapter 5]{as}.  Consider the projection $\text{pr}_x: S\to \A^1_k$ be projection to the $x$-axis. The fiber over $\mathbb{G}_m = \A^1_k \backslash \{0\}$, is isomorphic to $\G_m\times \A^1_k$ which is given by 
    %    \begin{align*}
      %      \text{pr}_x^{-1}(\mathbb{G}_m) & \cong \Spec\ (k[x,y,z]/(x^2z-y^2+y))_x \\
    %        & \cong \Spec\ k[x,y,z]/(z=(y^2-y)/x^2))_x \\
    %        & \cong \Spec\ k[x,y]_x 
   %     \end{align*}
        %$$(k[x,y,z]/(x^2z-y^2+y))_x = (k[x,y,z]/(z=(y^2-y)/x^2))_x \to k[x,y]_x $$ 
  %      where the last isomorphism is given by 
   %     $$x\mapsto x,\ y\mapsto y,\ z\mapsto (y^2-y)/x^2.$$ 
     %   Thus it is an $\A^1$-fibration. 

%More generally, the surface $S$ is an example of a Danielewski surface defined by the polynomial 

        %By \cite[Theorem 11]{doub}, the $\mathbb{A}^1$-fibration $\text{pr}_x$ factors through an étale locally trivial $\mathbb{A}^1$-bundle over the affine line with doubled origin. We will see that $S$ is $\mathbb{A}^1$-weakly equivalent to the affine line with doubled origin (Lemma \ref{sing-etale-bundle}).
        \item An $\mathbb{A}^1$-fibration may have a non-reduced fiber. Consider the surface, $X \subseteq \Spec\ \mathbb{C}[x,y,z]$ given by $x^2z = y^2-x$, which admits an $\mathbb{A}^1$-fibration \cite[Example 5]{doub}
        $$\text{pr}_x: X \to \mathbb{A}^1_{\mathbb{C}} \text{ given by } (x,y,z) \mapsto x..$$
        Here the fiber $\text{pr}_x^{-1}(a)$ is isomorphic to $\mathbb{A}^1_\mathbb{C}$, if $a$ is a non-zero closed point of $\mathbb{A}^1_{\mathbb{C}}$. If $a=0$, the fiber $\text{pr}_x^{-1}(0)$ is isomorphic to the scheme $\Spec\ (\frac{\mathbb{C}[y]}{(y^2)}[z])$, therefore it is non-reduced, but the associated reduced structure $\text{pr}^{-1}_x(0)_{\text{red}}$ is isomorphic to $\A^1_k$ (see also Example \ref{multi-section example}).
      %  But in case of a $\A^1$-bundle with the base being a smooth algebraic space of dimension 1, this is indeed the case.
    \end{enumerate}
	\end{example}
\begin{remark}
The $\A^1$-fibrations are a particular class of morphisms, used to study the affine surfaces. An $\A^1$-fibration is not a fibration in the $\A^1$-model structure, in general. For example, suppose $X$ is the affine surface as in Example \ref{d-sphere}(3). The morphism $\text{pr}_x$ is not a fibration in the $\A^1$-model structure. Indeed, $\text{pr}_x$ does not admit a section such that the following diagram has a lift;
% is the $\$ 
\[\xymatrixcolsep{3pc}
\xymatrix{\Spec\ k \ar[d]_{0} \ar[r]^{(0,0,0)} & X \ar[d]^{\text{pr}_x}\\
			\A^1_k\ar@{.>}[ru]^{\nexists} \ar@{=}[r]&\A^1_k} 
\]
that is, there does not exist any morphism $\A^1_k \to X$ such that both the triangles commute (see Example \ref{multi-section example}(3)). So $\text{pr}_x$ does not have right lifting property with respect to the trivial $\mathbb{A}^1$- cofibration $\Spec k \xrightarrow{s_0} \A^1_k$.
\end{remark}
	Let us recall the structure theorem for $\A^1$-fibration:
    %on smooth quasi projective surfaces.
	\begin{theorem} \cite[Theorem 9																																																																																																																																]{doub}\label{factor-A1-fibration}
		Suppose $X$ is a smooth affine surface over $k$ and $\pi: X \to C$ is an $\mathbb{A}^1$-fibration over a smooth curve $C$. Then $\pi$ factors as an \'etale locally trivial $\mathbb{A}^1$-bundle $\theta: X \to \mathcal{C}$ over a smooth algebraic space $\mathcal{C}$ of dimension $1$ (possibly not a scheme) along with the structure morphism $\alpha: \mathcal{C} \to C$, which is surjective, quasi-finite and birational morphism of finite type.
	\end{theorem}
   
    \begin{construction} \label{construction of algebraic space}
Following \cite[Theorem 9]{doub}, the construction of the algebraic space $\mathcal{C}$ in Theorem \ref{factor-A1-fibration} consists of two steps: 
\begin{enumerate}
	\item 
	%If $\pi$ has a reducible fiber then, 
	Suppose, $c_1, \dots, c_r$ are finitely many $k$-points of $C$ such that $\pi^{-1}(c_i)$ is not irreducible and $\pi^{-1}(c_i)$ consists $d_i$-many irreducible components.
	The morphism $\pi: X \to C$ is factored as 
	$$X \xrightarrow{\check{\pi}} \check{C} \xrightarrow{\check{p}} C,$$
	where $\check{C}$ is a non-separated scheme obtained by replacing the points $c_i$ by $d_i$-many points and $\check{p}: \check{C} \to C$ is the canonical morphism (see also \cite[Example 7]{doub}). Every fiber of $\check{\pi}$ is irreducible 
	(fiber of $\check{\pi}$ is non-reduced if and only if the corresponding fiber of $\pi$ is non-reduced) and thus the reduced structure of each fiber is isomorphic to $\A^1$. If every fiber of $\pi$ is irreducible, then $\check{C}$ is just $C$ itself and $\check{p}$ is the identity map.
	\item The algebraic space $\sC$ is obtained in finitely many steps, one for each non-reduced fiber of $\check{\pi}$. At each step, for a non-reduced fiber of $\check{\pi}$ at $\check{c} \in \check{C}$, an intermediate algebraic space is constructed by replacing a Zariski neighbourhood of $\check{c}$ in $\check{C}$, using an \'etale equivalence relation. The morphism $\check{\pi}$ is further factored as 
	$$X \xrightarrow{\theta} \mathcal{C} \xrightarrow{p} \check{C},$$
	where $\theta$ is an \'etale locally trivial $\A^1$-bundle and $\alpha = \check{p} \circ p$. As in (a), if every fiber of $\pi$ is reduced, then the algebraic space $\sC$ is just $\check{C}$ itself (a scheme) and $p$ is the identity map. 
	\item A closed point of $\sC$ corresponds to a closed point of $\check{C}$ and a closed point $\check{c} \in \check{C}$ corresponds to an irreducible component (which is also a connected component) of the fiber $\pi^{-1}(\check{p}(\check{c}))$.
	%$$X \xrightarrow{\theta} \mathcal{C} \xrightarrow{p} \check{C} \xrightarrow{\check{p}}C,$$ 
\end{enumerate}
    \end{construction}    
          

%             \item An $\A^1$-fibration $\pi: X \to C$ is not necessarily an $\A^1$-weak equivalence (Example \ref{d-sphere}(2)). In Lemma \ref{sing-etale-bundle}, we will see that the morphism $\rho: X \to \mathcal{C}$ in Theorem \ref{factor-A1-fibration} induces Nisnevich local weak equivalence $\rho_*: Sing_*(X) \to Sing_*(\mathcal{C})$. So $\rho$ is an $\mathbb{A}^1$-weak equivalence (Corollary \ref{factorisation A1 equivalence}).
          \begin{remark} \label{comments of structure theorem}
          	If every fiber $\pi^{-1}(x)$ is isomorphic to $\mathbb{A}^1$, then $\pi$ is actually an $\mathbb{A}^1$-bundle \cite[Proposition 3]{doub}. An $\mathbb{A}^1$-bundle between schemes is an $\mathbb{A}^1$-weak equivalence \cite[Theorem 5.3.1.3]{as}.
         
    \end{remark}



	%\begin{definition} \cite[Section 1.3, Chapter 3]{miyanishicirm} 
	%	With same notation as above $\pi: X \to \sC$ is called an $\A^1$-fibration if a general fiber of $\pi$ is isomorphic to $\A^1$. 
%	\end{definition}
	%Unlike $\A^1$-bundle, an $\A^1$-fibration can have fibers not isomorphic to $\A^1$. Those are called degenerate fibers.
%	\begin{definition}
	%	For an $\A^1$-fibration $\pi: X \rightarrow \sC$ to a smooth $1$-dimentional algebraic space. Let $F=p^{-1}(c)$ a fiber, it can be scheme theoretically expressed as $$F=\sum_{i=0}^n m_i\cdot(F_i)_{\red}$$ where $(F_i)_{\red}\cong \A^1$ and $m_i = \mathrm{length}(\sO_{F_i,\xi_i})$ and each $F_i$ is the irreducible component of the fiber. m$(F)=\gcd(m_1,\dots,m_n)$ is called total multiplicity of the fiber $F$. A fiber is non-degenerate (i.e., isomorphic to $\A^1$) if $n=1,\ m_1=1$.
		
		
%	\end{definition}
%	\begin{remark}
%	    \begin{enumerate}
  %        \item Every fiber over a point of an $\mathbb{A}^1$-bundle is isomorphic to $\mathbb{A}^1$. So an $\mathbb{A}^1$-bundle $\pi: X \to C$ from a smooth affine surface to a smooth curve $C$ is an $\mathbb{A}^1$-fibration. However, the converse is true if every scheme-theoretic fiber of $\pi$ is irreducible and reduced \cite[Proposition 3]{doub}.
          %\item An $\mathbb{A}^1$-bundle is an $\mathbb{A}^1$-weak equivalence \cite[Theorem 5.3.1.3]{as}.
%	        \item A line bundle (algebraic) is an $\mathbb{A}^1$-bundle (with respect to the Zariski topology). But the converse is not true. An $\mathbb{A}^1$-bundle may not have a section, where a line bundle always admits a section (zero section). The $\mathbb{A}^1$-bundle in \cite[Example 2]{doub} does not admit a section, hence it is not a line bundle.
    %        \item An $\mathbb{A}^1$-bundle corresponds to a line bundle, if the base is an affine scheme. Indeed, a line bundle over a scheme $X$ correspond to a $\mathbb{G}_m$-torsor over $X$; which in turn correspond to an element of $H^1_{\text{Zar}}(X, \mathbb{G}_m)$ \cite[Corollary 4.7]{milne}. On the other hand, an $\mathbb{A}^1$-bundle over $X$ corresponds to an $\text{Aff}_1:=\text{Aut}(\mathbb{A}^1)$-torsor over $X$; which in turn is isomorphic to $H^1_{\text{Zar}}(X, \text{Aff}_1)$. But the Zariski group sheaf $\text{Aff}_1$ is isomorphic to the semidirect product
            %The étale group sheaf $\text{Aff}_1$ is isomorphic to the semidirect product 
    %        $\mathbb{G}_m \ltimes \mathbb{G}_a$, which gives the short exact sequence of Zariski sheaves 
            %on essentially smooth schemes over $k$
%		$$0 \rightarrow \mathbb{G}_a \rightarrow \text{Aff}_1 \rightarrow \mathbb{G}_m \rightarrow 0.$$
%		This gives a long exact sequence in cohomology of $X$ (in Zariski topology)
%		$$\dots \rightarrow H^1_{\text{Zar}}(X, \mathbb{G}_a) \rightarrow H^1_{\text{Zar}}(X, \text{Aff}_1) \rightarrow H^1_{\text{Zar}}(X, \mathbb{G}_m) \rightarrow H^2_{\text{Zar}}(X, \mathbb{G}_a) \rightarrow \dots$$
		%Now by Hilbert's Theorem $90$, 
		%$$H^1_{\text{ét}}(\mathbb{A}^n_{\mathcal{O}}, \mathbb{G}_m) \cong H^1_{\text{Zar}}(\mathbb{A}^n_{\mathcal{O}}, \mathbb{G}_m) \cong Pic(\mathbb{A}^n_{\mathcal{O}}).$$
		%Since the Picard group functor is $\mathbb{A}^1$-invariant and any projective module over a local ring is free, so $H^1_{\text{ét}}(\mathbb{A}^n_{\mathcal{O}}, \mathbb{G}_m)$ is trivial. 
		%sheaf $\Spec\ k$. 
    %    Since $X$ is affine, so 
    %    $H^1_{\text{Zar}}(X, \mathbb{G}_a) \cong H^2_{\text{Zar}}(X, \mathbb{G}_a) \cong \{0\}$. Therefore,
   %     $$H^1_{\text{Zar}}(X, \text{Aff}_1) \xrightarrow{\cong} H^1_{\text{Zar}}(X, \mathbb{G}_m).$$
		%On the other hand, si étale cohomology vanishes with respect to the structure sheaf of an affine scheme; so
		%$$H^1_{\text{ét}}(\mathbb{A}^n_{\mathcal{O}}, \mathbb{G}_a) \cong H^1_{\text{Zar}}(\mathbb{A}^n_{\mathcal{O}}, \mathbb{G}_a) \cong \{0\}.$$
%	    \end{enumerate}
%	\end{remark}
	
Now we will prove the main theorem in this article (Theorem \ref{sing-etale-bundle}), which implies that the morphism $\theta$ in Theorem \ref{factor-A1-fibration} induces
%says that an étale $\mathbb{A}^1$-bundle $\rho: X \to \mathcal{C}$ from a smooth variety to a smooth algebraic space $\mathcal{C}$ induces 
a Nisnevich local weak equivalence $\theta_*: Sing_*(X) \to Sing_*(\mathcal{C})$. So, $\theta$ is an $\A^1$-weak equivalence.
\begin{theorem} \label{sing-etale-bundle}
		Let $X$ be a smooth variety over $k$ and $\mathcal{C}$ be a smooth algebraic space over $k$. Suppose, $\rho: X \to \mathcal{C}$ is an \'etale locally trivial $\mathbb{A}^1$-bundle.
		%, where . 
		Then the induced morphism $\rho_*: Sing_*(X) \to Sing_*(\mathcal{C})$ is a Nisnevich local weak equivalence.
	\end{theorem}
	\begin{proof}
		We will show that for every essentially smooth Henselian local scheme $\mathcal{O}$ over $k$, the morphism between the simplicial sets
		$$\rho_*: Sing_*(X)(\mathcal{O}) \to Sing_*(\mathcal{C})(\mathcal{O})$$
		is a trivial Kan fibration.
		Therefore, we need to find a lift of the following given commutative square
	
		
	
        \begin{equation}
		\label{lift}
%        \begin{tikzcd}[cramped]
%			{\bigcup_j \mathbb{A}^{n-1}_\mathcal{O}} && {X} \\
%			{\mathbb{A}^n_\mathcal{O}} && {\mathcal{C}}
%			\arrow["f", from=1-1, to=1-3]
%			\arrow["i"', hook, from=1-1, to=2-1]
%			\arrow["\rho", from=1-3, to=2-3]
%			\arrow["\exists", dashed, from=2-1, to=1-3]
%			\arrow["g", from=2-1, to=2-3]
%		\end{tikzcd}
%		\xymatrix{\bigcup_j \mathbb{A}^{n-1}_\mathcal{O} \ar@{^{(}->}[d]^i \ar[r]^f & X \ar[d]^{\rho}\\
%			\A^n_{\sO}\ar@{.>}[ru]^{\exists} \ar[r]^g &\sC} 
%\[
	\xymatrix{\partial \Delta^n \ar@{^{(}->}[d]^i \ar[r]^{f \ \ \ \ } &Sing_*(X)(\sO)\ar[d]^{\rho_*}\\
	\Delta^n \ar@{.>}[ru]^{\exists} \ar[r]^-{g} &Sing_*(\sC)(\sO)} 
%\]
       \end{equation}  
The $n$-simplex $g: \Delta^n \to Sing_*(\sC)(\sO)$ is given by a morphism (we denote it by $g$ again) $g: \A^n_\sO \to \sC$.	 The morphism $f: \partial \Delta^n \to Sing_*(X)(\mathcal{O})$ is a collection of $(n+1)$-many morphisms $f_j:\mathbb{A}^{n-1}_{\mathcal{O}} \to X$ ($0 \leq j \leq n$) satisfying the compatibility conditions ($d^if_j = d^{j-1}f_i$ for $0 \leq i<j \leq n$). This is 
equivalent to find the lift of the diagram (see Remark \ref{boundary and horn filling}(5))
 \begin{equation}
  \label{liftequivalent}
\xymatrix{
\partial\Delta^n_\sO \ar[r]^{f} \ar[d]_i & X  \ar[d]^\rho\\
\mathbb{A}^n_{\mathcal{O}} \ar[r]^{g}\ar@{..>}[ur]^{\exists} & \mathcal{C}
}
\end{equation}
where $\partial\Delta^n_\sO:= \partial \Delta^n_a \times_k \sO$ (Remark \ref{boundary and horn filling}(4)).
%is defined as the coequaliser in the category of Nisnevich sheaves  \cite[Section 3, Chapter I]{jardine}
%\[
%\begin{tikzcd}
%\coprod_{\substack{0 \leq i < j \leq n}}
% \A^{n-2}_\sO \ar[r,shift left=.75ex, "d^{j-1}_i"]
%  \ar[r,shift right=.75ex, swap, "d^i_j"]
%&
%\coprod_{i=0}^n \A^{n-1}_\sO \ar[r] 
%&
% \partial(\A^n_\sO)
%\end{tikzcd}
%\]
where $d^p_q$ is the boundary map $d^p:\A^{n-1}_\sO \to \A^n_\sO$ to the $q$-th component in the coproduct (see Section \ref{detailed horn formula}, for the boundary maps $d^i$).
		Consider the pullback diagram in the category of algebraic spaces
	
        
        \begin{equation}
        	\label{pullback}
        \xymatrix{\mathbb{A}^n_\mathcal{O}\times_\mathcal{C}X \ar[d]^{\pr_1} \ar[r]^{\ \ \ \pr_2} &X \ar[d]^{\rho}\\
        	\mathbb{A}^n_\mathcal{O} \ar[r]^g & \sC }
        \end{equation}
        
		Since $\rho$ is an étale locally trivial $\mathbb{A}^1$-bundle, so $\text{pr}_1$ is also an étale locally trivial $\mathbb{A}^1$-bundle.
% and $g$ is induced from the map $g'$. 
The pullback algebraic space $\mathbb{A}^n_{\mathcal{O}} \times_{\mathcal{C}} X$ is actually a scheme \cite[Corollary 4.2.3]{alpher}.
        %, since $\mathcal{C}$ is an algebraic space. 
        We claim that 
		%the étale locally trivial $\mathbb{A}^1$-bundle 
		$\text{pr}_1$
		%: \mathbb{A}^n_{\mathcal{O}} \times_{\mathcal{C}} X \to \mathbb{A}^n_{\mathcal{O}}$ 
		is isomorphic to the trivial line bundle over $\mathbb{A}^n_{\mathcal{O}}$. Indeed, an \'etale $\A^1$-bundle over $\A^n_\sO$ is given by a class in $H^1_{\text{\'et}}(\A^n_\sO,\Aut(\A^1_k)) \cong H^1_{\text{Zar}}(\A^n_\sO,\Aut(\A^1_k))$ and furthermore, a Zariski $\A^1$-bundle is isomorphic to a line bundle, if the base is affine
 %which is $0$ from 
(Remark \ref{a1 bundle properties}(4)).
So,
$$H^1_{\text{\'et}}(\A^n_\sO,\Aut(\A^1_k)) \cong H^1_{\text{Zar}}(\A^n_\sO,\Aut(\A^1_k)) \cong H^1_{\text{Zar}}(\A^n_\sO, \G_m) \cong Pic(\mathbb{A}^n_{\mathcal{O}}) = 0,$$
% and Hilbert's Theorem 90 \cite[Proposition 4.9]{milne} we have,
%		 Indeed, %the isomorphism classes of 
%		an étale locally trivial $\mathbb{A}^1$-bundle over $\mathbb{A}^n_{\mathcal{O}}$ corresponds to an étale $\text{Aff}_1$-torsor over $\mathbb{A}^n_{\mathcal{O}}$, where $\text{Aff}_1$ is the automorphism group of the affine line $\mathbb{A}^1_k$. Thus, the isomorphism classes of étale locally trivial $\mathbb{A}^1$-bundles over $\mathbb{A}^n_{\mathcal{O}}$ is isomorphic to the non-abelian étale cohomology $H^1_{\text{ét}}(\mathbb{A}^n_{\mathcal{O}}, \text{Aff}_1)$. Now the étale group sheaf $\text{Aff}_1$ is isomorphic to the semidirect product $\mathbb{G}_m \ltimes \mathbb{G}_a$, which gives the short exact sequence of étale sheaves on essentially smooth schemes over $k$
%		$$0 \rightarrow \mathbb{G}_a \rightarrow \text{Aff}_1 \rightarrow \mathbb{G}_m \rightarrow 0.$$
%		This gives a long exact sequence in étale cohomology of $\mathbb{A}^n_{\mathcal{O}}$
%		$$\dots \rightarrow H^0_{\text{ét}}(\mathbb{A}^n_{\mathcal{O}}, \mathbb{G}_m) \rightarrow H^1_{\text{ét}}(\mathbb{A}^n_{\mathcal{O}}, \mathbb{G}_a) \rightarrow H^1_{\text{ét}}(\mathbb{A}^n_{\mathcal{O}}, \text{Aff}_1) \rightarrow H^1_{\text{ét}}(\mathbb{A}^n_{\mathcal{O}}, \mathbb{G}_m) \rightarrow \dots$$
%		Now by Hilbert's Theorem $90$ , 
%		$$H^1_{\textit{ét}}(\mathbb{A}^n_{\mathcal{O}}, \mathbb{G}_m) \cong H^1_{Zar}(\mathbb{A}^n_{\mathcal{O}}, \mathbb{G}_m) \cong Pic(\mathbb{A}^n_{\mathcal{O}}) = 0.$$
		since the Picard group functor is $\mathbb{A}^1$-invariant and any projective module over a local ring is free. %and étale cohomology vanishes with respect to the structure sheaf of an affine scheme.
%		; so
%		$$H^1_{\text{ét}}(\mathbb{A}^n_{\mathcal{O}}, \mathbb{G}_a) \cong H^1_{\text{Zar}}(\mathbb{A}^n_{\mathcal{O}}, \mathbb{G}_a) \cong \{0\}.$$
%		Therefore, from the long exact sequence, $H^1_{\text{ét}}(\mathbb{A}^n_{\mathcal{O}}, \text{Aff}_1)$ is also trivial.
		 Hence, the \'etale $\mathbb{A}^1$-bundle $\text{pr}_1: \mathbb{A}^n_{\mathcal{O}} \times_{\mathcal{C}} X \to \mathbb{A}^n_{\mathcal{O}}$ is isomorphic to the projection $p_1: \mathbb{A}^n_{\mathcal{O}} \times_k \mathbb{A}^1_k \to \mathbb{A}^n_{\mathcal{O}}$.
		Thus we have the following commutative diagram
	
		
		\[
		\xymatrix{\mathbb{A}^n_\mathcal{O}\times_k \mathbb{A}^1_k \ar[rd]_{p_1} \ar[r]^{\cong}_\phi &\mathbb{A}^n_\mathcal{O}\times_\mathcal{C}X \ar[d]^{\pr_1} \ar[r]^<<<<{\pr_2} &X \ar[d]^{\rho}\\
		 &\A^n_{\sO} \ar[r]^g &\sC}
		\]
		Since $p_1$ is the trivial bundle, so it admits a global section and a global section of $p_1$ is determined by a morphism $\mathbb{A}^n_{\mathcal{O}} \to \mathbb{A}^1_k$. 
		The commutative diagram \eqref{liftequivalent} and the pullback diagram \eqref{pullback} induces the morphism 
		$$h:  \partial\Delta^n_\sO \to  \mathbb{A}^n_{\mathcal{O}} \times_{\mathcal{C}} X. $$
		Consider the composition of the morphisms 
		$$p_2 \circ \phi^{-1} \circ h:  \partial\Delta^n_\sO \to \mathbb{A}^1_k,$$
		where $p_2: \mathbb{A}^n_{\mathcal{O}} \times_k \mathbb{A}^1_k \to \mathbb{A}^1_k$ is the projection.
		%Com
		Now for any $k$-scheme $U$, the simplicial set $Sing_*(\mathbb{A}^1_k)(U)$ is a Kan fibrant simplicial set, which is also contractible simplicial set (Example \ref{kan fibrant An} and Lemma \ref{contractible An}). So the following commutative diagram has a lift
		
		
		\[\xymatrixcolsep{5pc}
		\xymatrix{ \partial\Delta^n_\sO \ar@{^{(}->}[d]^i \ar[r]^{\ \ \ p_2\circ \phi^{-1}\circ h} &{\mathbb{A}^1_k}\\
		\A^n_\sO \ar@{.>}[ur]^{\exists s}}
		\]
%		A lift to the diagram means 
Let $\Tilde{s}: \mathbb{A}^n_{\mathcal{O}} \to \mathbb{A}^n_{\mathcal{O}} \times_k \mathbb{A}^1_k$ be the section  of $p_1$, which is determined by the lift $s: \mathbb{A}^n_{\mathcal{O}} \to \mathbb{A}^1_k$. Then the desired lift of the diagram \eqref{liftequivalent} is given by 
		$$\text{pr}_2 \circ \phi \circ \Tilde{s}: \mathbb{A}^n_{\mathcal{O}} \to X.$$
		
		Indeed,
		\begin{align*}
			\rho \circ \text{pr}_2 \circ \phi \circ \Tilde{s} &= g \circ \text{pr}_1 \circ \phi \circ \Tilde{s} \\
			& = g \circ p_1 \circ \Tilde{s} \\
			& = g,
		\end{align*}
so the lower triangle in diagram \ref{liftequivalent} commutes and for the commutativity of the upper triangle in diagram \ref{liftequivalent}, we will prove that 
$$\text{pr}_2 \circ \phi \circ \Tilde{s} \circ d^i = f_i, \text{ for every }i, 0 \leq i \leq n$$
where $d^i: \A^{n-1}_\sO \to \A^n_\sO$ is the $i$-th face inclusion (see Section \ref{detailed horn formula}).
Now, for $u \in  \mathbb{A}^{n-1}_{\mathcal{O}}$,
		\begin{align*}
			\text{pr}_2 \circ \phi \circ \Tilde{s} \circ d^i (u) &=  \text{pr}_2 \circ \phi(d^i(u),(s\circ d^i)(u))\\
			&=\text{pr}_2 \circ \phi(d^i(u), p_2\circ \phi^{-1} \circ h(u))\\
			&=\text{pr}_2 \circ \phi(d^i(u), p_2\circ \phi^{-1}(d^i(u),f_i(u))) \\
			&=f_i(u).
		\end{align*}
Hence, diagram \ref{liftequivalent} has a lift.
		Therefore, 
%the morphism 
$\rho_*: Sing_*(X)(\mathcal{O}) \to Sing_*(\mathcal{C})(\mathcal{O})$ is a trivial Kan fibration. Hence the morphism $\rho_*: Sing_*(X) \to Sing_*(\mathcal{C})$ is a Nisnevich local weak equivalence.
		
	\end{proof}
    Since the canonical morphism $\mathcal{X} \to Sing_*(\mathcal{X})$ is an $\mathbb{A}^1$-weak equivalence \cite[Corollary 3.8, Section 2.3]{mv}, therefore we have the following immediate corollary:
	\begin{corollary} \label{factorisation A1 equivalence}
	    Let $X$ be a smooth variety over $k$ and $\mathcal{C}$ be a smooth algebraic space over $k$. Suppose, $\rho: X \to \mathcal{C}$ is an \'etale locally trivial $\mathbb{A}^1$-bundle. Then $\rho$ is an $\mathbb{A}^1$-weak equivalence. So, the morphism $\theta$ which appears in the factorisation of the $\A^1$-fibration $\pi:X \to C$ in Theorem \ref{factor-A1-fibration} is an $\A^1$-weak equivalence.
	\end{corollary}
	\begin{remark}~
				 \begin{enumerate}
			
			\item If every fiber of $\pi$ is reduced and some fiber is reducible, then the algebraic space in Theorem \ref{factor-A1-fibration} is a non-separated scheme. In Example \ref{d-sphere}(2), the $\mathbb{A}^1$-fibration $\text{pr}_x: X \to \A^1_k$ factors through an \'etale locally trivial $\mathbb{A}^1$-bundle over the affine line with $d$-origins, where $d=deg(P(y))$. So, $X$ is $\mathbb{A}^1$-weakly equivalent to the affine line with $d$-origins (Corollary \ref{factorisation A1 equivalence}).
			\item If $\pi$ has a non-reduced fiber, then the algebraic space in Theorem \ref{factor-A1-fibration} is not a scheme. In Example \ref{d-sphere} (3), the $\mathbb{A}^1$-fibration $\text{pr}_x: X \to \A^1_k$ factors through an \'etale locally trivial $\mathbb{A}^1$-bundle over an algebraic space $\mathcal{C}$, which is the quotient of the affine line with double origin with free $\mathbb{Z}/2\mathbb{Z}$-action \cite[Example 12]{doub}. So, $X$ is $\A^1$-weakly equivalent to $\sC$ (Corollary \ref{factorisation A1 equivalence}). Thus in particular, $Pic(X)$ is isomorphic to $\Z/2\Z$.
		\end{enumerate}
	\end{remark}
%Since étale, Zariski and Nisnevich cohomology agree over the structure she The same method
By proceeding the same way as of proof of Theorem \ref{sing-etale-bundle}, the diagram \ref{lift} has a lift, if $\rho: X \to \sC$ is an vector bundle. Indeed, a vector bundle over $\mathbb{A}^n_\mathcal{O}$ corresponds to a $GL_n$-torsor over $\mathbb{A}^n_\mathcal{O}$ and the $\mathbb{A}^1$-invariance of the isomorphism class of vector bundles over affine base \cite[Theorem 1]{asokI} implies that a vector bundle over $\mathbb{A}^n_{\mathcal{O}}$ corresponds to a vector bundle over $\mathcal{O}$.  By Hilbert's Theorem 90, we have the isomorphism
$$H^1_{Zar}(\mathcal{O}, GL_n) \cong H^1_{Nis}(\mathcal{O}, GL_n) \cong H^1_{\text{\'et}}(\mathcal{O}, GL_n).$$
Finally in Zariski topology, a vector bundle (or equivalently, a projective module) over a local ring is trivial (or free) implies that diagram \ref{lift} in the proof of Theorem \ref{sing-etale-bundle} has a lift. Therefore, we have the following result:
%same Nisnevich local weak equivalence for any s: %Proposition \ref{sing-etale-bundle} to étale 
\begin{theorem} \label{general bundle}
		Let $X$ be a smooth variety over $k$ and $\mathcal{C}$ be a smooth algebraic space over $k$. Suppose, $\rho: X \to \mathcal{C}$ is an  in \'etale, Nisnevich or Zariski topology. %locally trivial $\mathbb{A}^1$-bundle.
		%, where . 
		Then the induced morphism $\rho_*: Sing_*(X) \to Sing_*(\mathcal{C})$ is a Nisnevich local weak equivalence.
	\end{theorem} 
%	\begin{corollary}\label{sing-line-bundle}
	%	Let X,C be smooth varieties over k and $p: X\rightarrow C$ is any $\A^1$-bundle then the induced morphism $$p_*:Sing_*(X)\rightarrow Sing_*(C)$$ is a Nisnevich local weak equivalence.
	%\end{corollary}
%	\begin{remark}
%(\textbf{??})      Suppose, $X, Y \in Sm/k$ and $\rho: X \to Y$ is an Nisnevich $\mathbb{A}^n$-bundle. For Nisnevich topology the induced morphism $\rho_*: Sing_*(X) \to Sing_*(Y)$ a Nisnevich local weak equivalence by \cite[Example 2.3, Section 3.2]{mv}.
%
%      And by Theorem \ref{sing-etale-bundle}, applying $Sing_*$ functor to an \'etale $\mathbb{A}^1$-bundle induces a Nisnevich local weak equivalence. However similar result for an étale $\A^n$-bundle (for $n>1$) is not known. 
%      In the proof of Theorem \ref{sing-etale-bundle}, we have used the structure of the automorphism group $\text{Aut}(\mathbb{A}^1_k)$. However the structure of $\text{Aut}(\mathbb{A}^n_k)$ is not known, for $n > 1$.
%  \end{remark} 

  \section{Simplicial Contractibility of Koras-Russell Threefolds}
  The family of varieties known as Koras-Russell threefolds are the first examples of exotic algebraic structure on the 6-dimentional real manifold $\C^3$. They are diffeomorphic to $\C^3$, but not isomorphic to $\C^3$ as varieties. It is also remarkable that the Makar-Limanov invariant, i.e. ML($X$) is non-trivial, therefore $X \ncong \A^3_\C$ but ML($X\times \A^1_\C)=\C$, for the Koras-Russell threefolds $X$. Though it is not known whether $X \times \A^1_\C \cong \A^4_\C$. If it were true, i.e. $X \times \A^1_\C \cong \A^4_\C$, then $X$ has to be $\A^1$-contractible and $Sing_*(X)$ has to be in fact $\A^1$-fibrant as it is a retract of $Sing_*(\A^4_\C)$ which is $\A^1$-fibrant. Koras-Russell threefolds of the first kind $X$ are $\A^1$-contractible has already been proved in \cite[Theorem 3.1]{df}. Here we prove that $Sing_*(X)$ is $\A^1$-local.
 Before proving Theorem \ref{application purity koras}, we need the following technical lemma (Lemma \ref{key lemma}), which says that in case of open or closed immersions, the functor $Sing_*$ commutes with taking quotient of schemes (as Nisnevich sheaves).
\begin{lemma} \label{key lemma}
    Let $i:V \to U$ be a morphism between smooth varieties over $k$, which is either an open immersion or a closed immersion of schemes. Then the canonical morphism
    $$i_*: Sing_*(U)/Sing_*(V) \to Sing_*(U/V)$$
    is an isomorphism of simplicial sheaves.
    %Nisnevich local weak equivalence.
\end{lemma}
\begin{proof}
%    $Sing_*$ is a homotopy colimit of a simplicial diagram.
There is a commutative diagram of simplicial sheaves
  \[
\xymatrix{
Sing_*(V) \ar[r] \ar[d] & Sing_*(U) \ar[d] \\
\text{*} \ar[r] & Sing_*(U/V)
}
\]
This induces a morphism of simplicial sheaves
$$i_*: Sing_*(U)/Sing_*(V) \to Sing_*(U/V).$$
We will show that for every essentially smooth Henselian local scheme $\mathcal{O}$ over $k$, the morphism
$$i_*: (Sing_*(U)/Sing_*(V))(\mathcal{O}) \to Sing_*(U/V)(\mathcal{O})$$
is an isomorphism of simplicial sets, that is for every $n \geq 0$ the morphism between $n$-simplices
$$i_*: (Sing_*(U)/Sing_*(V))(\mathcal{O})_n \to Sing_*(U/V)(\mathcal{O})_n$$
is a bijection.
%n isomorphism.

Since over the Henselian local scheme, the presheaf section and sheaf section coincide, so
\begin{align*}
(Sing_*(U)/Sing_*(V))(\mathcal{O})_n & \cong Sing_*(U)(\sO)_n/Sing_*(V)(\sO)_n
%& \cong (Hom_{Sch/k}(\A^n_\sO, U)\coprod \{\bullet\})/\sim
\end{align*}
For any scheme $T$, the section $Sing_*(U)(T)_n/Sing_*(V)(T)_n$ is the pushout in the category of sets
 \[
\xymatrix{
Hom_{Sch/k}(\A^n_T, V) \ar[r] \ar[d] & Hom_{Sch/k}(\A^n_T, U) \ar[d] \\
\text{*} \ar[r] & Sing_*(U)(T)_n/Sing_*(V)(T)_n
}
\]
where $Sch/k$ is the category of $k$-schemes. So, $Sing_*(U)(T)_n/Sing_*(V)(T)_n$ is a pointed set canonically, denote the base point by $\bullet$. Also for the simplicial presheaf 
$$T \mapsto Sing_*(U)(T)/Sing_*(V)(T),$$
the restriction map
$$Sing_*(U)(T)/Sing_*(V)(T) to Sing_*(U)(T^\prime)/Sing_*(V)(T^\prime)$$
induced by a morphism $T^\prime \to T$ of $k$-schemes, respect the base points.

%$(Sing_*(U)/Sing_*(V))_n(\mathcal{O})$ is a pointed set canonically, denote it by $\bullet$. 
We will prove that the morphism  $i_*$ is a bijection between the pointed sets
$$i_*: ((Sing_*(U)/Sing_*(V))(\mathcal{O})_n, \bullet) \to (Sing_*(U/V)(\mathcal{O})_n, i_*(\bullet)),$$
for every Henselian local scheme $\sO$.

%$$

%Fix a $k$-point $v \in V(k)$. The point $v$ determines an element in the simplicial set $(Sing_*(U)/Sing_*(V))(\mathcal{O})$ (which we will denote by $v$ again),
%and $Sing_*(U/V)(\mathcal{O})$ 
%by taking the composition of the morphisms
%$$\mathbb{A}^n_{\mathcal{O}} \to \Spec\ k \xrightarrow{v} V.$$
%The elements $v \in (Sing_*(U)/Sing_*(V))(\mathcal{O})$ and $i_*(v) \in Sing_*(U/V)(\mathcal{O})$ make both the simplicial sets to be pointed. 
%For each $n$, fix a distinguished element $c: \mathbb{A}^n_{\mathcal{O}} \to V$ given by the constant map to a $k$ -point of $V$; that makes both the simplicial sets to the pointed.

\textbf{Injectivity:} Suppose, $\phi, \psi \in (Sing_*(U)/Sing_*(V))(\mathcal{O})_n$ and 
$$i_*(\phi) = i_*(\psi) \in Sing_*(U/V)(\mathcal{O}).$$ 
For presheaf and the associated sheaf, the sections over $\mathcal{O}$ are same; thus $\phi, \psi \in (Sing_*(U)/Sing_*(V))^{\text{pre}}(\mathcal{O})$. So the sections are given by the morphisms $\phi, \psi: \mathbb{A}^n_{\mathcal{O}} \to U$ (again, denote by $\phi$ and $\psi$). 
%We denote the images of $\phi$ and $\psi$ in the quotient $U(\mathbb{A}^n_{\mathcal{O}})/V(\mathbb{A}^n_{\mathcal{O}})$ by $[\phi]$ and $[\psi]$ respectively.
% such that their images $[\phi], [\psi] \in U(\mathbb{A}^n_{\mathcal{O}})/V(\mathbb{A}^n_{\mathcal{O}})$ are same. 
Since $i_*(\phi) = i_*(\psi) \in (U/V)(\mathbb{A}^n_{\mathcal{O}})$, so there is a Nisnevich covering $f: W \to \mathbb{A}^n_{\mathcal{O}}$ such that 
$$\phi \circ f = \psi \circ f \in (U/V)^{\text{pre}}(W) = U(W)/V(W).$$
This means either $\phi \circ f, \psi \circ f: W \to V \hookrightarrow U$ or $\phi \circ f = \psi \circ f: W \to U$.
Since $Hom(-, U)$ is a Nisnevich sheaf, so for the second case we conclude $\phi = \psi: \mathbb{A}^n_{\mathcal{O}} \to U$. For the first case, since $f$ is a covering, in particular surjective; so the images of $\phi$ and $\psi$ are in $V$, thus both the morphisms $\phi$ and $\psi$ factor as
$$\phi: \mathbb{A}^n_{\mathcal{O}} \to V \hookrightarrow U \text{ and } \psi: \mathbb{A}^n_{\mathcal{O}} \to V \hookrightarrow U.$$
Hence 
$$\phi = \psi \in (Sing_*(U)/Sing_*(V))^{\text{pre}}(\mathcal{O}) = U(\mathbb{A}^n_{\mathcal{O}})/V(\mathbb{A}^n_{\mathcal{O}}).$$
So $i_*$ is injective.

\textbf{Surjectivity: } Suppose, $\alpha \in Sing_*(U/V)(\mathcal{O})_n$. So $\alpha \in (U/V)(\mathbb{A}^n_{\mathcal{O}})$. Thus there is a Nisnevich covering $W \to \mathbb{A}^n_{\mathcal{O}}$ such that $\alpha$ is given by an element $\phi \in (U/V)^{\text{pre}}(W)$ (that is, $\phi: W \to U$ is a morphism) such that there is a Nisnevich covering $\Tilde{W} \to W \times_{\mathbb{A}^n_{\mathcal{O}}} W$ such that 
$$\text{pr}_1^*(\phi)|_{\Tilde{W}} = \text{pr}_2^*(\phi)|_{\Tilde{W}} \in (U/V)^{\text{pre}}(\Tilde{W}) = U(\Tilde{W})/V(\Tilde{W}).$$
This means either 
$$\text{pr}_1^*(\phi)|_{\Tilde{W}}, \text{pr}_2^*(\phi)|_{\Tilde{W}} : \Tilde{W} \to V \hookrightarrow U$$
or
$$\text{pr}_1^*(\phi)|_{\Tilde{W}} = \text{pr}_2^*(\phi)|_{\Tilde{W}} : \Tilde{W} \to U.$$
For the second case, since $Hom(-, U)$ is a Nisnevich sheaf and the morphism $\Tilde{W} \to W \times_{\mathbb{A}^n_{\mathcal{O}}} W$ is a Nisnevich covering; so 
$$\text{pr}_1^*(\phi) = \text{pr}_2^*(\phi) : W \times_{\mathbb{A}^n_{\mathcal{O}}} W \to U$$
and again as $W \to \mathbb{A}^n_{\mathcal{O}}$ is a Nisnevich covering; so $\phi$ can be lifted to a morphism $\bar{\phi}: \mathbb{A}^n_{\mathcal{O}} \to U$, which proves surjectivity in this case. \par
For the first case, if $\text{pr}_1^*(\phi)|_{\Tilde{W}}, \text{pr}_2^*(\phi)|_{\Tilde{W}} : \Tilde{W} \to V$, then we claim that the morphism $\phi$ factors through $V$. Indeed, the composition of morphisms
$$\Tilde{W} \to W \times_{\mathbb{A}^n_{\mathcal{O}}} W \xrightarrow{\text{pr}_1} W$$
and 
$$\Tilde{W} \to W \times_{\mathbb{A}^n_{\mathcal{O}}} W \xrightarrow{\text{pr}_2} W$$
are Nisnevich coverings of $W$ and the images of $\text{pr}_1^*(\phi)|_{\Tilde{W}}$ and $\text{pr}_2^*(\phi)|_{\Tilde{W}}$ are in $V$, so the image of $\phi$ is also in $V$. Now we claim that $\alpha = i_*(\bullet) \in (U/V)(\mathbb{A}^n_{\mathcal{O}})$, where $\bullet$ is the distinguished element in $(Sing_*(U)/Sing_*(V))(\mathcal{O})$ mentioned before. This is equivalent to show that 
$$\alpha|_W = i_*(\bullet)|_W \in (U/V)(W),$$
where $W \to \mathbb{A}^n_{\mathcal{O}}$ is the Nisnevich covering and $-|_W$ denotes the restriction maps, induced by $W \to \A^n_\sO$, for both the presheaf $(U/V)^{\text{pre}}$, given by 
$$(U/V)^{\text{pre}}(T) = U(T)/V(T)$$
and the associated Nisnevich sheaf $U/V$.
% the restriction map of the quotient sheaf $U/V$
%$$-|_W:(U/V)(\mathbb{A}^n_{\mathcal{O}}) \to (U/V)(W).$$
%The quotient sheaf $U/V$ is the Nisnevich sheaf associated to the presheaf $(U/V)^{\text{pre}}$, which is given by 
%$$(U/V)^{\text{pre}}(T) = U(T)/V(T).$$
Now, $\alpha|_W = \eta(\phi)$. Consider the commutative diagram
 \[
\xymatrix{
U(\A^n_\sO)/V(\A^n_\sO) \ar[r]^\eta \ar[d]^{|_W} & (U/V)(\A^n_\sO) \ar[d]^{|_W} \\
U(W)/V(W) \ar[r]_\eta & (U/V)(W)
}
\]
%$i_*(\bullet)|_W = \eta(\bullet|_W)$, 
where 
$$\eta: (U/V)^{\text{pre}} \to U/V$$ 
is the sheafification morphism and in the square, the left vertical morphism respects the base point. So,
$$i_*(\bullet)|_W = \eta(\bullet|_W).$$
Since the image of $\phi$ is contained in $V$, so $\phi = \bullet|_W \in U(W)/V(W)$, so 
$$\alpha|_W =\eta(\bullet|_W)= i_*(\bullet)|_W \in (U/V)(W)$$ 
and consequently, $\alpha = i_*(\bullet)$. This proves the surjectivity. 

This completes the proof.
\end{proof}

 %Before proving the simplicial contractibility of the Koras-Russell threefold of first kind, first we will prove the following version of homotpy purity, in the case when the pair $(X, Z)$ is étale over $(\mathbb{A}^n_k, \mathbb{A}^{n-c}_k \times \{(0, \dots 0)\}))$ (Theorem \ref{purity sing}). This is an application of Proposition \ref{sing-etale-bundle} and the simplicial contractibility of Koras-Russell threefold will be an application of Theorem \ref{purity sing}.
%\begin{theorem} \label{purity sing}
  %  Suppose, $i: Z \to X$ is a closed immersion of smooth varieties over a field $k$. Also assume that $X$ admits an étale morphism $q: X \to \mathbb{A}^n_k$ such that $i(Z) = q^{-1}(\mathbb{A}^{n-c}_k \times \{(0, \dots 0)\})$. Then there is a Nisnevich local weak equivalence
    %$$Sing_*(Th(N_{X, Z})) \sim Sing_*(X/(X\setminus Z)).$$
%\end{theorem}
%\begin{proof}
  %  The proof essentially follows from the combinations of the proof of \cite[Lemma 2.28, Section 3.2]{mv} and Proposition \ref{sing-etale-bundle}. 
    %after applying the $Sing_*$ functor.

   % Suppose, $p: B(X \times_k \mathbb{A}^1_k, Z \times \{0\}) \to X \times_k \mathbb{A}^1_k$ is the blow-up of $X \times_k \mathbb{A}^1_k$ at the center $Z \times \{0\}$. The normal bundle $N_{X \times_k \mathbb{A}^1_k, Z \times \{0\}}$ is isomorphic to $N_{X,Z} \oplus \mathcal{O}$. So the exceptional divisor of $p$ is isomorphic to $\mathbb{P}(N_{X,Z} \oplus \mathcal{O})$. This gives an inclusion $\alpha: N_{X,Z} \to B(X \times_k \mathbb{A}^1_k, Z \times \{0\})$. By the universal property of the blow-up, the natural inclusion $j: Z \times_k \mathbb{A}^1_k \to X \times_k \mathbb{A}^1_k$ induces a morphism $f: Z \times_k \mathbb{A}^1_k \to B(X \times_k \mathbb{A}^1_k, Z \times \{0\})$ such that $p \circ f = j$. Thus in particular,
%$$p^{-1}(Z \times \{0\}) \cap f(Z \times_k \mathbb{A}^1_k) =  f(Z \times \{0\})$$
%and $f|_{Z \times \{0\}}$ is given by the zero section of the normal bundle $N_{X,Z} \to Z$. This gives the following isomorphism
%$$p^{-1}(Z \times \{0\}) \setminus f(Z \times_k \mathbb{A}^1_k) \cong \mathbb{P}(N_{X,Z} \oplus \mathcal{O}) \setminus \mathbb{P}(\mathcal{O}).$$
% following elementary distinguished square (in the Zariski topology)
%\[
%\xymatrix{
%N_{X,Z} \setminus Z  \ar[r] \ar[d] & \mathbb{P}(N_{X,Z} \oplus \mathcal{O}) \setminus \mathbb{P}(\mathcal{O}) \ar[d]\\
%N_{X, Z} \ar[r] & \mathbb{P}(N_{X,Z} \oplus \mathcal{O})
%}
%\]
%gives the isomorphism of the Nisnevich sheaves
%$$Th(N_{X,Z}) \cong \mathbb{P}(N_{X,Z} \oplus \mathcal{O})/(\mathbb{P}(N_{X,Z} \oplus \mathcal{O}) \setminus \mathbb{P}(\mathcal{O})).$$
%is the canonical inclusion $Z \times \mathbb{A}^1_k \hookrightarrow X \times \mathbb{A}^1_k$. 
%Suppose, $g: X \to B(X \times_k \mathbb{A}^1_k, Z \times \{0\})$ is the inclusion such that $p \circ g$ is the inclusion $X \times\{1\} \hookrightarrow X \times \mathbb{A}^1_k$.

   % The morphisms $\alpha$ and $g$ induces the monomorphisms of the pairs
    %$$g: (X, X \setminus Z) \to (B(X \times_k \mathbb{A}^1_k, Z \times \{0\}), B(X \times_k \mathbb{A}^1_k, Z \times \{0\}) \setminus f(Z \times\mathbb{A}^1_k))$$
    %and
    %$$\alpha: (N_{X,Z}, N_{X,Z} \setminus Z) \to (B(X \times_k \mathbb{A}^1_k, Z \times \{0\}), B(X \times_k \mathbb{A}^1_k, Z \times \{0\}) \setminus f(Z \times \mathbb{A}^1_k)).$$
%These induce morphisms between the quotients (from now on, we write $B$ instead of $B(X \times_k \mathbb{A}^1_k, Z \times \{0\})$)
%$$\Tilde{g}: X/(X \setminus Z) \to B/(B \setminus f(Z \times \mathbb{A}^1))$$
%and
%$$\Tilde{\alpha}: Th(N_{X, Z}) \to B/(B \setminus f(Z \times \mathbb{A}^1)).$$
%Applying the $Sing_*$ functor, the morphisms $\Tilde{g}, \Tilde{\alpha}$ induce
%$$g_*: Sing_*(X/(X \setminus Z)) \to Sing_*(B/(B \setminus f(Z \times \mathbb{A}^1)))$$
%and
%$$\alpha_*: Sing_*(Th(N_{X, Z})) \to Sing_*(B/(B \setminus f(Z \times \mathbb{A}^1))).$$
%Theorem \ref{purity sing} would follow from Lemma \ref{purity sing lemma}.
%\end{proof}
%\begin{lemma} \label{purity sing lemma}
   % The morphisms $g_*$ and $\alpha_*$ in Theorem \ref{purity sing} are Nisnevich local weak equivalences.
%\end{lemma}
%\begin{proof}
 %   We will follow the argument exactly as in \cite[Proposition 2.24, Section 3.2]{mv}. The proof of Lemma \ref{purity sing lemma} would follow from the following three steps.

 %   \textbf{Step 1:} Suppose the closed immersion $i$ is given by the zero section $i: X \to \mathbb{A}^n_X$. So in this case, $\alpha_*$ and $g_*$ are same maps; here we will prove that $\alpha_*$ is a the Nisnevich local weak equivalences for the smooth pair $(\mathbb{A}^n_X, X)$.
    
  %  The total space of the blow-up $p: B(\mathbb{A}^{n+1}_k \times_k X, X \times \{0\})\to \mathbb{A}^{n+1}_k \times_k X$ admits the canonical line bundle 
   % $$q: B \to \mathbb{P}^n_k \times X$$ 
%where, we write $B$ instead of $B(\mathbb{A}^{n+1}_k \times_k X, X \times \{0\})$) as before. By the universal property of the blow-up, the inclusion $X \times_k \mathbb{A}^1_k \to \mathbb{A}^{n+1}_X$ given by $(x,t) \mapsto (0,\dots,0,t,x)$ lifts to the morphism $f: \mathbb{A}^1_X \to B$ which makes the following square cartesian:



%\[    
%\xymatrix @R=.4in @C=1.5in{  
%\mathbb{A}^1_{X} \ar[r]_{f} \ar[d] &B \ar[d]_q \\
%X \ar[r]_{x \mapsto ([0:\dots 0:1],x)} &\mathbb{P}^n_X
%}
%\]



   %So the morphism $q$ restricts to the line bundle 
    %$$B \setminus f (\mathbb{A}^1_X) \to (\mathbb{P}^n_X \setminus X).$$
    %Thus by Proposition \ref{sing-etale-bundle}, the induced morphisms
    %$$q_*: Sing_*(B) \to Sing_*(\mathbb{P}^n_X)$$
    %and
   % $$q_*: Sing_*(B \setminus f(\mathbb{A}^1_X)) \to Sing_*(\mathbb{P}^n_X \setminus X)$$
    %are the Nisnevich local weak equivalences.
   % So by \cite[Lemma 2.11, Section 2.2]{mv}, the induced morphism
    %$$Sing_*(B)/Sing_*(B \setminus f(\mathbb{A}^1_X)) \to Sing_*(\mathbb{P}^n_X)/Sing_*(\mathbb{P}^n_X \setminus X)$$
    %is a Nisnevich local weak equivalence.
    %Therefore, by Lemma \ref{key lemma}, the morphism
    %$$q^\prime: Sing_*(B/(B \setminus f(\mathbb{A}^1_X))) \to Sing_*(\mathbb{P}^n_X/(\mathbb{P}^n_X \setminus X))$$
    %is a Nisnevich local weak equivalence.
    
   % Consider the following Cartesian square of Zariski covering of $\mathbb{P}^n_X$:
   % \[
%\xymatrix{
%\mathbb{A}^n_{X} \setminus X \ar[r] \ar[d] & \mathbb{A}^n_X \ar[d]^j \\
%\mathbb{P}^n_{X} \setminus X \ar[r] & \mathbb{P}^n_X
%}
%\]
%where $j$ is the inclusion $(x_1, \dots x_n, y) \mapsto ([x_1:\dots x_n:1],y)$ and the horizontal map $\mathbb{P}^n_X \setminus X \to \mathbb{P}^n_X$ is the complement of the closed immersion $X \to \mathbb{P}^n_X$ as $y \mapsto ([0:\dots:0:1], y)$ This induces an isomorphism of Nisnevich sheaves
%$$\Tilde{j}: \mathbb{A}^n_X / (\mathbb{A}^n_{X} \setminus X) \xrightarrow{\cong} \mathbb{P}^n_X/(\mathbb{P}^n_{X} \setminus X).$$
%There is a the following commutative triangle:
 % \[
%\xymatrix{
%Sing_*(\mathbb{A}^n_{X}/(\mathbb{A}^n_X \setminus X) \ar[r]_(.5){\alpha_*} \ar[d]^{j_*}_{\cong} & Sing_*(B/(B \setminus f(\mathbb{A}^1_X)) \ar[dl]^{q^\prime} \\
%Sing_*(\mathbb{P}^n_{X}/(\mathbb{P}^n_X \setminus X) 
%}
%\]
%Since $j_*$ and $q^\prime$ are Nisnevich local weak equivalences, so $\alpha_*$ is a Nisnevich local weak equivalence.\par
%Similarly, $q^\prime \circ g_*$ is an isomorphism of simplicial sheaves, which is induced by the inclusion $\mathbb{A}^n_X \to \mathbb{P}^n_X$ given by $(a_1, \dots, a_n) \mapsto (a_1: \dots a_n:1)$ and $q^\prime$ is a Nisnevich local weak equivalence, so $g_*$ is also a Nisnevich local weak equivalence. \\
%This completes Step 1.

%\

%Before going to Step 2, here is the setup:
%Suppose, $\phi: U \to X$ is an étale morphism such that $\phi$ induces isomorphism $Z_U:= \phi^{-1}(Z) \to Z$.
%Then there is the following commutative diagram:
%\begin{equation}
 %   \label{purity diagram}
  %  \xymatrix{
%Sing_*(U/(U \setminus Z_U)) \ar[r] \ar[d] & Sing_*(B_U/(B_U \setminus f_U(Z_U \times \mathbb{A}^1))) \ar[d] & Sing_*(Th(N_{U, Z_U})) \ar[l] \ar[d] \\
%Sing_*(X/(X \setminus Z)) \ar[r] & Sing_*(B/(B \setminus f(Z \times \mathbb{A}^1))) & Sing_*(Th(N_{X, Z})) \ar[l]
%\mathbb{A}^n_{\mathcal{O}} \ar@{=}[r] & \mathbb{A}^n_{\mathcal{O}} \ar[r]_{g} & \mathcal{C}
%}
%\end{equation}
%The étale morphism $\phi$ induces elementary distinguished square of Nisnevich covering of $X$ given by
%\[  
%\xymatrix{
%U \setminus Z_U \ar[r] \ar[d] & U \ar[d]^{\phi} \\
%X \setminus Z \ar[r] & X
%}
%\]
%This induces an isomorphism of Nisnevich sheaves
%$$U/(U \setminus Z_U) \xrightarrow{\cong} X/(X \setminus Z).$$
%Similarly, 
%Here all the vertical arrows in the diagram \ref{purity diagram} are isomorphisms of simplicial sheaves. Thus we have the following conclusion:

%\begin{lemma} \label{sub lemma}
  %  Suppose, $\phi: U \to X$ is an étale morphism such that $\phi$ induces isomorphism $Z_U:= \phi^{-1}(Z) \to Z$. Then in Lemma \ref{purity sing lemma}, $\alpha_*$ and $g_*$ are the Nisnevich local weak equivalences for the pair $(U, Z_U)$ if and only if $\alpha_*$ and $g_*$ are the Nisnevich local weak equivalences for the pair $(X, Z)$.
%\end{lemma}

%Here is the final step of the proof.

%\textbf{Step 2:} 
%Suppose $i: Z \to X$ is a closed immersion. Also assume that $X$ admits an étale morphism $q: X \to \mathbb{A}^n$ such that 
%$$i(Z) = q^{-1}(\mathbb{A}^{n-c} \times \{(0, \dots 0)\}).$$
%We will show that for this pair $(X, Z)$, $\alpha_*$ and $g_*$ in Lemma \ref{purity sing lemma} are Nisnevich local weak equivalences.

%Following the proof in \cite[Lemma 2.28, Section 3.2]{mv}, there is a scheme $U$ which is an open subscheme of $X \times_{\mathbb{A}^n} (X \times \mathbb{A}^c)$ along with the 
%. Consider the fiber product
%\[
%\xymatrix{
%X \times_{\mathbb{A}^n} (Z \times \mathbb{A}^c) \ar[r] \ar[d] & X \ar[d]^q \\
%Z\times \mathbb{A}^c \ar[r]^{(q \circ i, Id)} & \mathbb{A}^n
%}
%\]
%Here all the morphisms in the diagram are étale morphisms. The fiber of the morphism $X \times_{\mathbb{A}^n} (Z \times \mathbb{A}^c) \to \mathbb{A}^n$ over the inclusion $\mathbb{A}^{n-c}_k \to \mathbb{A}^n_k$ is $Z \times_{\mathbb{A}^{n-c}} Z$, which is a closed subscheme of $X \times_{\mathbb{A}^n} (Z \times \mathbb{A}^c)$. The morphism $Z \to \mathbb{A}^{n-c}$ is étale, so the diagonal $\Delta: Z \to Z \times_{\mathbb{A}^{n-c}} Z$ is also an étale morphism. Hence, we can write $Z \times_{\mathbb{A}^{n-c}} Z = \Delta(Z) \coprod Y$, for some closed subscheme $Y$ of $Z \times_{\mathbb{A}^{n-c}} Z$. Thus in particular, $Y$ is closed in $X \times_{\mathbb{A}^n}(Z \times \mathbb{A}^c)$. 
%Consider the scheme, $U = X \times_{\mathbb{A}^n}(Z \times \mathbb{A}^c) \setminus Y$; which is an open subscheme of $X \times_{\mathbb{A}^n}(Z \times \mathbb{A}^c)$. There are 
%two projections which are étale morphisms
%$$\text{pr}_1: U \to X$$
%and 
%$$\text{pr}_2: U \to Z \times \mathbb{A}^c$$
%The restrictions 
%$$\text{pr}_1^{-1}(i(Z)) \to i(Z)$$ 
%and 
%$$\text{pr}_2^{-1}(Z \times \{0\}) \to Z \times \{0\}$$ 
%are isomorphisms. Also,
%$$\text{pr}_1^{-1}(i(Z)) = \Delta(Z) = \text{pr}_2^{-1}(Z \times \{0\})\subset X \times_{\mathbb{A}^n}(Z \times \mathbb{A}^c).$$
%For the proof, if $(x, z, t) \in U \subset X \times_{\mathbb{A}^n}(Z \times \mathbb{A}^c)$ such that $x \in Z$, then $q(x) = (s,0)$ with $s =q(z) \in \mathbb{A}^{n-c}$ and $t = 0 \in \mathbb{A}^c$. But $(x,z) \notin Y$, so $x=z$ and $(x,z,t) \in \Delta(Z)$. Conversely, $(z,z) \in \text{pr}_1^{-1}(i(Z))$. Thus, $\text{pr}_1^{-1}(i(Z)) = \Delta(Z)$. Similarly, $ \text{pr}_2^{-1}(Z \times \{0\}) = \Delta(Z)$.

%Observe that, we can see $\Delta(Z)$ as a closed subscheme of $U$ viewing it as the inverse image of $\mathbb{A}^{n-c} \times \{0\}$ of the morphism $U \to \mathbb{A}^n$.)

%%\end{proof}




\begin{theorem} \label{application purity koras}
    Suppose, $X$ is the Koras-Russell threefold of the first kind given by
    $$X \equiv \{x^mz = y^r +t^s+x\} \subset \mathbb{A}^4_k=\Spec\ k[x,y,z,t],$$
    where $m \geq 2$ and $r,s \geq 2$ are coprime integers. Then the structure map $Sing_*(X) \to \Spec\ k$ is a Nisnevich local weak equivalence or $Sing_*(X)$ is simplicially contractible and $Sing_*(X)$ is $\A^1$-local.
\end{theorem}
\begin{proof}
    We will follow the proof in \cite[Section 3.1]{df}. 
    %Consider the hyperplane intersection given by $X \cap \{z=0\}$ which is isomorphic to $\mathbb{A}^2_k$ and this induces a closed immersion
   % $$i: \mathbb{A}^2_k \to X \text{ given by } (y,t) \mapsto (-y^r-t^s, y,0,t).$$
    %The closed subscheme $X \cap \{x=0\} \cap \{y = 0\}$ is isomorphic to $\mathbb{A}^1_k$; there is a closed immersion
    %$$j: \mathbb{A}^1_k \to X \text{ given by } z \mapsto (0, 0, z,0).$$
   % The image of $j$ will be denoted by $L$. Note that, $L$ intersects the image of $i$ at the origin $(0,0,0,0)$. Thus the following square is a pullback
 %   \[
%\xymatrix{
%\mathbb{A}^2 \setminus \{(0,0)\} \ar[r] \ar[d] & \mathbb{A}^2 \ar[d]^i \\
%X \setminus L \ar[r] & X
%}
%\]
%Therefore, we have the following commutative diagram of cofiber sequences
 % \[
%\xymatrix{
%\mathbb{A}^2 \setminus \{(0,0)\} \ar[r] \ar[d] & \mathbb{A}^2 \ar[d]^i \ar[r] & \mathbb{A}^2/(\mathbb{A}^2 \setminus \{(0,0)\}) \ar[d] \\
%X \setminus L \ar[r] & X \ar[r] & X/(X \setminus L)
%}
%\]
%Since $Sing_*$ preserves cofibrations, applying the $Sing_*$ functor and by Lemma \ref{key lemma}, we have the following commutative diagram where each row is a cofiber sequence
 %\[
%\xymatrix{
%Sing_*(\mathbb{A}^2 \setminus \{(0,0)\}) \ar[r] \ar[d] & Sing_*(\mathbb{A}^2) \ar[d]^i %\ar[r] & Sing_*(\mathbb{A}^2/(\mathbb{A}^2 \setminus \{(0,0)\})) \ar[d] \\
%Sing_*(X \setminus L) \ar[r] & Sing_*(X) \ar[r] & Sing_*(X/(X \setminus L))
%}
%\]
Let us write, 
$$X(s) \equiv \{x^mz = y^r +t^s+x\}; \ s \geq 1.$$
The morphism
$$\phi: X(1) \equiv \{x^mz = y^r+t+x\} \to \mathbb{A}^3_k \text{ defined as }(x,y,z,t) \mapsto (x, y ,z)$$
is an isomorphism.
%n $X(1)$ is isomorphic to $\mathbb{A}^3_k = \Spec\ k[x,y,z]$. 
Consider the closed immersion 
$$i: \mathbb{A}^2_k \equiv \Spec\ k[y,t] \to X(s) \text{ defined as } (y,t) \mapsto (-y^r-t^s, y,0, t).$$
Let $L$ be the line contained in $X(s)$:
$$L:=X(s) \cap \{x=y=0\}.$$
Then the closed immersion $i$ restricts to the morphism (we denote the restriction also by $i$)
$$i: \mathbb{A}^2_k \setminus \{(0,0)\} \to X(s) \setminus L.$$
Consider the following commutative diagram where each row is a $\A^1$-cofiber sequence
  \[
\xymatrix{
\mathbb{A}^2_k \setminus \{(0,0)\}  \ar[r] \ar[d]^i & \mathbb{A}^2_k \ar[d]^{i} \ar[r] & \mathbb{A}^2_k/(\mathbb{A}^2_k \setminus \{(0,0)\})  \ar[d]^q \\
X(s) \setminus L \ar[r] & X(s)  \ar[r] & X(s)/(X(s) \setminus L)  
}
\]
here $q$ is the morphism, induced by $i$. Now, applying the $Sing_*$ functor and by Lemma \ref{key lemma}, both rows of the following commutative diagram are also cofiber sequences:
\[
\xymatrix{
	Sing_*(\mathbb{A}^2_k \setminus \{(0,0)\})  \ar[r] \ar[d]^{i_*} & Sing_*(\mathbb{A}^2_k) \ar[d]^{i_*} \ar[r] & Sing_*(\mathbb{A}^2_k/(\mathbb{A}^2_k \setminus \{(0,0)\}))  \ar[d]^{q_*} \\
	Sing_*(X(s) \setminus L) \ar[r] & Sing_*(X(s))  \ar[r] & Sing_*(X(s)/(X(s) \setminus L))
}
\]
Thus to prove that $i_*: Sing_*(\A^2_k) \to Sing_*(X(s))$ is a simplicial weak equivalence, by very weak five lemma \cite[Lemma 2.1]{df} it suffices to prove that 
$$i_*: Sing_*(\A^2_k\backslash\{(0,0)\}) \to Sing_*(X(s)\backslash L) \text{ and } $$
$$q_*: Sing_*(\mathbb{A}^2_k/(\mathbb{A}^2_k \setminus \{(0,0)\}) \to  Sing_*(X(s)/(X(s) \setminus L))$$
are simplicial weak equivalences.
% Nisnevich local weak equivalence suffices. 
In \cite[Section 3]{df}, it was shown that the left vertical map 
$$i: \mathbb{A}^2_k \setminus \{(0,0)\} \to X(s) \setminus L$$ 
is an $\mathbb{A}^1$-weak equivalence. So to prove the map 
$$i_*: Sing_*(\A^2_k\backslash\{(0,0)\}) \to Sing_*(X(s)\backslash L)$$
is a simplicial weak equivalence, it suffices to show that
% we only need to prove they are 
both the spaces $Sing_*(\A^2_k\backslash\{(0,0)\}) $ and  $Sing_*(X(s)\backslash L)$ are $\A^1$-local. 

The projection, 
$$SL_2 \equiv \{xw-yz=1 \subset \mathbb{A}^4_k\} \to \mathbb{A}^2_k \setminus \{0,0\} \text{ defined as } (x,y,z,w) \mapsto (x,y)$$
is an line bundle. Therefore, the induced map 
$$Sing_*(SL_2) \to Sing_*(\mathbb{A}^2_k \setminus \{0,0\})$$
is a simplicial weak equivalence (Proposition \ref{sing-etale-bundle}). Since, $Sing_*(SL_2)$ is $\mathbb{A}^1$-local \cite[Remark 3.3.8]{asokII}, so $Sing_*(\mathbb{A}^2_k \setminus \{0,0\})$ is also $\mathbb{A}^1$-local. 
By \cite[Proposition 3.1]{df}, there is a fourfold $W \in Sm/k$ such that $W$ admits Zariski locally trivial $\mathbb{A}^1$-bundles 
$$p_s: W \to X(s) \setminus L \text{ and } p_1: W \to X(1)\setminus L.$$
Therefore, the induced maps
$$(p_s)_*: Sing_*(W) \to Sing_*(X(s) \setminus L) \text{ and } (p_1)_*: Sing_*(W) \to Sing_*(X(1) \setminus L)$$
are Nisnevich local weak equivalences (Proposition \ref{sing-etale-bundle}).
 The projection map
$$p: X(1)\setminus L \to \mathbb{A}^2_k \setminus \{0,0\} \mathrm{\ defined\ as } (x,y,z,t) \mapsto (x,t) $$
is an line bundle. So again by Proposition \ref{sing-etale-bundle}, the induced map
$$p_*: Sing_*(X(1) \setminus L) \to Sing_*(\mathbb{A}^2_k \setminus \{0,0\})$$ 
is a Nisnevich local weak equivalence. Thus, $Sing_*(X(1) \setminus L)$ is $\A^1$-local and consequently \\
$Sing_*(X(s)\backslash L)$ is also $\mathbb{A}^1$-local. Thus the $\A^1$-weak equivalence $$i_*: Sing_*(\mathbb{A}^2_k \setminus \{0,0\}) \to Sing_*(X(s) \setminus L)$$ is a Nisnevich local weak equivalence.\\
Now we will prove that $q_*$ is a Nisnevich local weak equivalence.


Consider the projection
%$$\rho: X(s) \equiv \{f(x,y,z,t)=0 | f=  x^mz - y^r-t^s-x\} \to \mathbb{A}^3 \text{ given by } (x,y,z,t) \mapsto (y,z,t).$$
$$\rho: X(s) \equiv \{x^mz = y^r+t^s+x\} \to \mathbb{A}^3_k \equiv \Spec\ k[y,z,t]$$
$$ \text{ given by } (x,y,z,t) \mapsto (y,z,t).$$
Suppose, $f(x,y,z,t) = x^mz - y^r-t^s-x \in k[x,y,z,t]$ that defines $X(s)$. Firstly, we will find some open subset of $X(s)$, where $\rho$ is \'etale.
$$\nabla f = (mx^{m-1}z-1, -ry^{r-1}, x^m, -st^{s-1}).$$
Suppose, $\Tilde{x} = (x_0,y_0,z_0,t_0) \in X(s)$ is a $k$-point. So the tangent space $T_{\Tilde{x}}(X(s))$ is given by (as a vector subspace of $k^4$):
$$(mx_0^{m-1}z_0-1)x -ry_0^{r-1}y+ x_0^mz -st_0^{s-1}t = 0.$$
The induced map between the tangent spaces
$$\rho_*: T_{\Tilde{x}}(X(s)) \to k^3 \text{ is given by } (x, y,z,t) \mapsto (y,z,t).$$
The map $\rho_*$ is an isomorphism on the open subset $X(s) \backslash \{x^{m-1}z = \frac{1}{m}\}$ of $X(s)$. 

Thus $\rho|_{X(s) \backslash \{x^{m-1}z = \frac{1}{m}\}}: X(s) \backslash \{x^{m-1}z = \frac{1}{m}\} \to \mathbb{A}^3_k$ is \'etale.

Suppose, $V = X(s) \backslash (\{x^{m-1}z = \frac{1}{m}\} \cup \{x^{m-1}z = 1\})$. Then $V$ is an open subset of $X(s)$ and $L := X(s) \cap \{x = y=0\} \subset V$ and also $\rho|_V : V \to \mathbb{A}^3_k$ is an \'etale morphism. The inverse image $\rho|_V^{-1}(\{y=t=0\}) = L$. Let us denote the closed subscheme $\{y=t=0\}\subset \A^3$ also by $L$.
%So by Theorem \ref{purity sing}, the morphisms $\alpha_*$ and $g_*$ (as described in Lemma \ref{purity sing lemma}) are local weak equivalences for the pair $(V, L)$. So in particular, there is a local weak equivalence 
%$$Sing_*(V/(V \setminus L)) \sim Sing_*(Th(N_{V, L})).$$
%$$Sing_*(V/(V \setminus L)) \sim Sing_*(\mathbb{A}^3_k/(\mathbb{A}^3_k \setminus \{(0,z,0)\})).$$
Consider the following elementary distinguished squares in the Nisnevich topology (in fact the left square is in the Zariski topology, all morphisms are open immersions):
 \[
\xymatrix{
V \setminus L \ar[r] \ar[d]^j & V \ar[d]^j & V \setminus L \ar[r] \ar[d]^{\rho|_V} & V \ar[d]^{\rho|_V} \\
X(s) \setminus L \ar[r] & X(s) & \A^3_k \backslash L \ar[r] & \A^3_k
}
\]
%Thus the morphism of the pairs $(V, V \setminus L) \to (X, X \setminus L)$ induces the isomorphism of the Nisnevich sheaves
%$$V/(V \setminus L) \cong (X/(X \setminus L)).$$
%Applying the $Sing_*$ functor,
They induce the following isomorphisms of the Nisnevich sheaves
$$j: (V/(V \setminus L)) \cong (X(s)/(X(s) \setminus L)) \mathrm{\ and \ } (\rho|V)_*: (V/(V \setminus L)) \cong (\A^3_k/(\A^3_k\backslash L)).$$
Therefore, after applying the $Sing_*$ functor, we have the isomorphisms.
$$j_*: Sing_*(V/(V \setminus L)) \to Sing_*(X(s)/(X(s) \setminus L)) \mathrm{\ and \ }$$ $$(\rho|_V)_*: Sing_*(V/(V \setminus L)) \to Sing_*(\mathbb{A}^3_k/(\mathbb{A}^3_k \setminus L)).$$ 
Thus, the projection map $\rho: X(s) \to \mathbb{A}^3_k$ induces Nisnevich local weak equivalence 
$$\rho_*: Sing_*(X(s)/(X(s) \setminus L)) \to Sing_*(\mathbb{A}^3_k/(\mathbb{A}^3_k \setminus L)).$$
Therefore, to show that 
$$q_*: Sing_*(\mathbb{A}^2_k/(\mathbb{A}^2_k \setminus \{(0,0)\})) \to  Sing_*(X(s)/(X(s) \setminus L))$$ 
is a Nisnevich local weak equivalence, it is enough to show that the morphism
$$\rho \circ i: \mathbb{A}^2_k \to \mathbb{A}^3_k \text{ given by }(y,t) \mapsto (y,0,t)$$
induces Nisnevich local weak equivalence
$$Sing_*(\mathbb{A}^2_k/(\mathbb{A}^2_k \setminus \{(0,0)\})) \to Sing_*(\mathbb{A}^3_k/(\mathbb{A}^3_k \setminus L)).$$
Consider, the projection
$$\text{pr}: \mathbb{A}^3_k \to \mathbb{A}^2_k \text{ defined as }(y,z,t) \mapsto (y,t).$$
The morphisms $\text{pr}$ and $\text{pr}|_{\mathbb{A}^3_k \setminus L}: \mathbb{A}^3_k \setminus L \to \mathbb{A}^2_k \setminus \{(0,0)\}$ are line bundles. Therefore, by Proposition \ref{sing-etale-bundle}, the induced morphisms
$$\text{pr}_*: Sing_*(\A^3_k) \to Sing_*(\A^2_k) \text{ and }$$ 
$$(\text{pr}|_{\mathbb{A}^3_k \setminus L})_*: Sing_*(\mathbb{A}^3_k \setminus L) \to Sing_*(\mathbb{A}^2_k \setminus \{(0,0)\})$$
are Nisnevich local weak equivalences. The morphism $\rho \circ i$ is the zero section of $\text{pr}$, so the morphisms
$$Sing_*(\A^2_k) \to Sing_*(\A^3_k) \text{ and }Sing_*(\A^2_k \setminus \{(0,0)\}) \to Sing_*(\A^3_k \setminus L),$$
induced by $\rho \circ i$ are Nisnevich local weak equivalences. Therefore, by \cite[Lemma 2.11, Section 2.2]{mv}, the morphism 
$$Sing_*(\A^2_k)/(Sing_*(\A^2_k \setminus \{(0,0)\})) \to Sing_*(\A^3_k)/(Sing_*(\A^3_k \setminus L))$$
is a Nisnevich local weak equivalence. Finally using Lemma \ref{key lemma}, the morphism
$$Sing_*(\mathbb{A}^2_k/(\mathbb{A}^2_k \setminus \{(0,0)\})) \to Sing_*(\mathbb{A}^3_k/(\mathbb{A}^3_k \setminus L)),$$
induced by $\rho \circ i$, is a Nisnevich local weak equivalence. This completes the proof. Therefore, we conclude that $i_*: Sing_*(\A^2_k) \to Sing_*(X(s))$ is a simplicial weak equivalence, so $Sing_*(X(s))$ is simplicially contractible. Thus $Sing_*(X)$ is $\A^1$-local.



%To prove this, consider the following commutative diagram, where each row is a cofiber sequence (using Lemma \ref{key lemma}):
%\[\xymatrixcolsep{3pc}
%\xymatrix{{Sing_*(\A^2_k\backslash \{(0,0)\})}\ar[d]^{(\rho \circ i)_*} \ar[r] & {Sing_*(\A^2_k)} \ar[d]^{(\rho \circ i)_*} \ar[r] & {Sing_*(\A^2_k/\A^2_k \backslash\{(0,0)\})}\ar[d]\\
%{Sing_*(\A^3_k \backslash L)} \ar[r] & {Sing_*(\A^3_k)} \ar[r] & {Sing_*(\A^3_k/\A^3_k \backslash L)}}
%\]
%Here $(p\circ i)((y,t))=(y,0,t)$. Clearly $(p\circ i)_*$ is a local weak equivalence and $(p\circ i)|_{\A^2\backslash\{(0,0)\}}$ is also Nisnevich local weak equivalence as $Sing_*(\A^3 \backslash L) = Sing_*(\A^2 \backslash\{(0,0)\})\times Sing_*(L=A^1) $. Thus the morphism between their cofibers, induced by $p\circ i$ is also Nisnevich local weak equivalence; Which concludes the proof. 
%$$q: \mathbb{A}^3_k \to \mathbb{A}^2_k \text{ defined as } (y,z,t) \mapsto (y,t).$$
%The line bundle $q$ restricts to the line bundle
%$$q: \mathbb{A}^3_k \setminus \{(0,z,0)\} \to \mathbb{A}^2_k \setminus \{(0,0)\}$$
%Therefore, $q$ induces Nisnevich local weak equivalences
%$$Sing_*(\mathbb{A}^3_k) \to Sing_*(\mathbb{A}^2_k) \text{ and } Sing_*(\mathbb{A}^3_k \setminus \{(0,z,0)\}) \to Sing_*(\mathbb{A}^2_k \setminus \{(0,0)\})$$
% So $q$ induces local weak equivalence
%$$q_*:Sing_*(\mathbb{A}^3_k/(\mathbb{A}^3_k \setminus \{(0,z,0)\})) \xrightarrow{ \sim} Sing_*(\mathbb{A}^2/(\mathbb{A}^2 \setminus \{(0,0)\})).$$
%Since, $q \circ (p \circ i)$ is idenety, so $p \circ i$ induces local weak eqivalence
%$$Sing_*(\mathbb{A}^2/(\mathbb{A}^2 \setminus \{(0,0)\})) \xrightarrow{\sim} Sing_*(\mathbb{A}^3_k/(\mathbb{A}^3_k \setminus \{(0,z,0)\}))$$
%
%
%\
%
%
%\
%
%
%\
%
%
%
%
%The open immersion gives an étale morphism for the pair $(V,L) \to (X(s), L)$. So by Lemma \ref{sub lemma}, the morphisms $\alpha_*$ and $g_*$ are local weak equivalences for the pair $(X(s), L)$. Thus in particular,
%$$Sing_*(X(s)/(X(s) \setminus L)) \sim Sing_*(N_{X(s), L}/(N_{X(s),L} \setminus L)).$$
%But, $Sing_*(N_{X(s), L}/(N_{X(s),L} \setminus L)) = Sing_*(\mathbb{A}^3/(\mathbb{A}^3 \setminus \{x=z=0\}))$, where $\mathbb{A}^3= \Spec\ k[x,y,z]$.


%Now there is an étale morphism 
%$$\phi: X(s) \to \mathbb{A}^3 \text{ given by } (x,y,z,t) \mapsto (x,y,z).$$
%The morphism $\phi$ is étale, each component of $\phi$ is a projection, so it is smooth and $\phi$ is of relative dimension $0$. 
%$$\text{The inverse image } \phi^{-1}(\{x=y=0\}) = X(s) \cap \{x=y=t=0\} = L.$$
%So by Theorem \ref{purity sing}, there is a local weak equivalence 
%$$Sing_*(Th(N_{X(s), L})) \sim Sing_*(X(s)/(X(s) \setminus L)).$$
%Since $L \cong \mathbb{A}^1$, so the normal bundle $N_{X(s), L} \to L$ is the trivial rank $2$-bundle $\mathbb{A}^3 %\to \mathbb{A}^1$. So there is a local weak equivalence 
%%$$Sing_*(Th(N_{X(s), L})) \sim Sing_*(\mathbb{A}^3/(\mathbb{A}^3 \setminus \{x= y =0\})).$$
%The trivial line bundle $\mathbb{A}^3 \to \mathbb{A}^2$ given by $(x,y,z) \mapsto (x,z)$ induces line bundle 
%$$\mathbb{A}^3 \setminus \{x= z =0\} \to \mathbb{A}^2 \setminus \{(0,0)\}.$$ 
%So by Corollary \ref{sing line bundle}, there are local weak equivalences
%$$Sing_*(\mathbb{A}^3) \sim Sing_*(\mathbb{A}^2)$$ 
%and 
%$$ Sing_*(\mathbb{A}^3 \setminus \{x= z =0\}) \sim Sing_*(\mathbb{A}^2 \setminus \{(0,0)\}).$$
%Therefore, by \cite[Lemma 2.11, Section 2.2]{mv}, there is a local weak equivalence
%$$Sing_*(\mathbb{A}^3/(\mathbb{A}^3 \setminus \{x= z =0\})) \sim Sing_*(\mathbb{A}^2/(\mathbb{A}^2 \setminus \{0,0\})).$$
%So in the following commutative diagram
% \[
%\xymatrix{
%  Sing_*(X(s)/(X(s) \setminus L)) \ar[d]  & Sing_*(V/(V \setminus L))  \ar[l]  \ar[r]  & Sing_*(\mathbb{A}^3/(\mathbb{A}^3 \setminus \{x= z =0\}))  \ar[dll] \\
% Sing_*(\mathbb{A}^2/(\mathbb{A}^2 \setminus \{0,0\})) 
%}
%\]
%the left horizontal arrow is an isomorphism of simplicial sheaves and the right horizontal arrow is a local weak equivalence. The morphism 
%$$Sing_*(\mathbb{A}^3/(\mathbb{A}^3 \setminus \{x= z =0\})) \to Sing_*(\mathbb{A}^2/(\mathbb{A}^2 \setminus \{0,0\}))$$
%is also a local weak equivalence. Thus, 
%the left vertical map is a local weak equivalence.
%
%
%
%
%
%
%
%
%
%
%Consider the hyperplane $P$ in $X(s)$ Let $L$ be the line $X(s) \cap \{x=y=0\}$. Then $X(1) \setminus L \cong (\mathbb{A}^2 \setminus \{(0,0)\}) \times_k \mathbb{A}^1_k$, which is $\mathbb{A}^1$-weakly equivalent to $\mathbb{A}^2_k \setminus \{(0,0)\}$. 
%
%Following \cite[Section 3.2]{df}, there is a $\mathbb{G}_a$-action on $X(s)$ given by the locally nilpotent derivation 
%$$D: \Gamma(X(s), \mathcal{O}_{X(s)}) \to \Gamma(X(s), \mathcal{O}_{X(s)}) \text{ defined as }$$
%$$x \mapsto 0, \ y \mapsto x^m, \ z \mapsto ry^{r-1}, \ t \mapsto 0.$$
%The algebraic quotient is given by the projection
%$$\text{pr}: X(s) \to \mathbb{A}^2_k \text{ given by } (x,y,z,t) \mapsto (x,t).$$
%The fixed point set for this action is $L = X(s) \cap \{x= y =0\}$. The algebraic quotient restricts to the morphism 
%$$q: X(s) \setminus L \to \mathbb{A}^2 \setminus \{(0,0)\}.$$
%The restriction $q$ factors through an étale locally trivial $\mathbb{A}^1$-bundle $$\rho: X(s) \setminus L \to \mathcal{C},$$
%where $\mathcal{C}$ is a smooth algebraic space, followed by a morphism 
%$$\delta: \mathcal{C} \to \mathbb{A}^2_k \setminus \{(0, 0)\}$$ 
%\cite[Theorem 3.2]{df}. The algebraic space $\mathcal{C}$ along with the morphism $\delta$ depend only on $r$ and do not depend on $s$.
%%There is a $\mathbb{G}_a$-equivariant morphism given by the projection
%%$$\text{pr}: X(s) \setminus L \to \mathbb{A}^2 \setminus \{(0,0)\} \text{ given by } (x,y,z,t) \mapsto (x,t).$$
%%Then $\text{pr} \circ i: \mathbb{A}^2 \setminus \{0,0\} \to \mathbb{A}^2 \setminus \{0,0\}$ is given by $(y,t) \mapsto (-y^r-t^s, t)$.
%%$$\mathbb{G}_a \times X \to X \text{ defined as }$$ 
%%$$\lambda. (x,y,z,t) = (x, y+x^m. \lambda, z+ry^{r-1}. \lambda+ \frac{r(r-1)}{2!}x^m y^{r-2}. \lambda^2 %+\dots+\frac{r(r-1)\dots 1}{r!}x^{m(r-1)}. \lambda^r).$$
%
%Therefore, we have the following commutative diagram of cofiber sequences
%  \[
%\xymatrix{
%X(s) \setminus L  \ar[r] \ar[d]^q & X(s) \ar[d]^{\text{pr}} \ar[r] & X(s)/(X(s) \setminus L)  \ar[d] \\
%\mathbb{A}^2 \setminus \{(0,0)\} \ar[r] & \mathbb{A}^2 \ar[r] & \mathbb{A}^2/(\mathbb{A}^2 \setminus \{(0,0)\})
%}
%\]
%Since $Sing_*$ preserves cofibrations, applying the $Sing_*$ functor and by Lemma \ref{key lemma}, we have the following commutative diagram where each row is a cofiber sequence
% \begin{equation}
% \label{koras diagram}
%\xymatrix{
%Sing_*(X(s) \setminus L) \ar[r] \ar[d]^{q_*} & Sing_*(X(s)) \ar[d]^{\text{pr}_*} \ar[r] & Sing_*(X(s)/(X(s) \setminus L)) \ar[d] \\
%Sing_*(\mathbb{A}^2 \setminus \{(0,0)\}) \ar[r] & Sing_*(\mathbb{A}^2) \ar[r] & Sing_*(\mathbb{A}^2/(\mathbb{A}^2 \setminus \{(0,0)\})
%}
%\end{equation}
%%There is a smooth algebraic space $\mathcal{C}$ along with a morphism $\delta: \mathcal{C} \to \mathbb{A}^2 \setminus \{(0,0)\}$ such that for every $s \geq 1$, there is an étale locally trivial $\mathbb{A}^1$-bundle $\rho: X(s) \to \mathcal{C}$, so that $q = \delta \circ \rho$. 
%By Lemma \ref{sing etale bundle}, the induced morphism $\rho_*: Sing_*(X(s)) \to Sing_*(\mathcal{C})$ is a Nisnevich local weak equivalence. Taking $s=1$, the morphism 
%$$q: X(1) \setminus L \to \mathbb{A}^2 \setminus \{(0,0)\}$$ 
%is the trivial line bundle. So by Corollary \ref{sing line bundle}, the induced morphism $q_*: Sing_*(X(1) \setminus L) \to Sing_*(\mathbb{A}^2 \setminus \{(0,0)\})$ is a Nisnevich local weak equivalence. Thus 
%$$\delta_*: Sing_*(\mathcal{C}) \to Sing_*(\mathbb{A}^2 \setminus \{(0,0)\})$$ 
%is also a Nisnevich local weak equivalence. Therefore, for any $s \geq 1$, the morphism $$q_*: Sing_*(X(s) \setminus L) \to Sing_*(\mathbb{A}^2 \setminus \{(0,0)\})$$ 
%is a Nisnevich local weak equivalence. 
%Now we will show that the right vertical morphism
%$$Sing_*(X(s)/(X(s) \setminus L)) \to Sing_*(\mathbb{A}^2/(\mathbb{A}^2 \setminus \{(0,0)\}) $$
%is also a Nisnevich local weak equivalence. 
%
%
%
%
%
%
%
%
%
%
%Therefore, in diagram \ref{koras diagram}, the left most and the right most vertical maps are local weak equivalence. Thus by very weak five lemma (\cite[Lemma 2.1]{df}), $Sing_*(X(s))$ is simplicially contractible.
\end{proof}
Thus by using \cite[Lemma 4.1]{df}, we have the following corollary
\begin{corollary}
    Suppose, $X$ is the threefold contained in $\mathbb{A}^4_k = \Spec\ k[x,y,z,t]$ given by
    $$X \equiv x^mz = y^r+t^s+xq(x),$$
    where $m,r,s$ are integers, $m \geq 2$ and $r,s \geq 1$ are coprime and $q(x) \in k[x]$ is such that $q(0) \in k^*$. Then $Sing_*(X)$ is simplicially contractible.
\end{corollary}
%	\begin{corollary}
	%	$Sing_*(\A^n)$is contractible
%	\end{corollary}
	
	%\begin{proof}
	%	$p:\A^n \rightarrow \A^{n-1}$ is an $\A^1$-bundle hence $p_*:Sing_*(\A^n)\rightarrow Sing_*(\A^{n-1})$ is nisnevich local weak equivalence. Hence $Sing_*(\A^n)$ nisnevich local weak equivalent to $Sing_*(k)$ which is trivial. 
	%\end{proof}
	



   

%write the computations of horn filling for the case $n=2$. Suppose, we have given a morphism 

%$$\Lambda^2_l \to Sing_*(\mathbb{A}^1_k)(U), \text{ where } U \in Sm/k, l \in \{0, 1, 2\}$$
%\textbf{For} $\mathbf{l=0}$, we have given 
%$$\phi_1, \phi_2: k[T] \to \frac{A[x_0, x_1]}{(x_0+x_1 - 1)}$$ 
%such that $d^1 \circ \phi_2 = d^1 \circ \phi_1$.
%Suppose $\phi_1$ and $\phi_2$ are given by 
%$$T \xmapsto{\phi_1} \overline{f_1} \text{ and } T \xmapsto{\phi_2} \overline{f_2} \text{ respectively.}$$
%Define $\Phi: k[T] \to \frac{A[x_0, x_1, x_2]}{(x_0+x_1+x_2-1)}$ as $T \mapsto \overline{f_1{x_0+x_1,x_2}+f_2(x_0,x_1+x_2) - f_2(x_0+x_1, x_2)}$

  %\bibitem{asokII} Aravind Asok, Marc Hoyois and Matthias Wendt; \emph{Affine representability results in $\mathbb{A}^1$–homotopy theory, II: Principal bundles and homogeneous spaces}, Geometry Topology 22 (2018) 1181–1225.

   %  \bibitem{asokI} Aravind Asok, Marc Hoyois and Matthias Wendt; \emph{Affine representability results in $\mathbb{A}^1$–homotopy theory, I}, DUKE MATHEMATICAL JOURNAL, Vol. 166, No. 10.
		% cR.V. Gurjar a,∗,

\begin{remark}
Thus Theorem \ref{application purity koras} says that if $X$ is a Koras-Russell threefold of the first kind, then $Sing_*(X)(\mathcal{O})$ is a contractible  simplicial set, for any Henselian local scheme $\mathcal{O}$.
% in the Nisnevich local model structure. 
This in particular, implies that $Sing_*(X)$
%$Ex(Sing_*(X))$ 
is $\mathbb{A}^1$-local.
%or being $Ex(Sing_*(X))(U)$ is contractible simplicial set, for any $U \in Sm/k$; where
%%s the fibrant replacement functor for the Nisnevich local model strucure. 
If the Koras-Russell threefold of the first kind were a counterexample to the Zariski cancellation in characteristics zero, i.e. $X \times_k \mathbb{A}^1_\C \cong \mathbb{A}^4_\C$, then for any $U \in Sm/k$, $Sing_*(X)(U)$ has to be a Kan fibrant and contractible simplicial set, as it is a retract of $Sing_*(\A^4_\C)(U)$ which is Kan fibrant and contractible (Example \ref{kan fibrant An} and Lemma \ref{contractible An}). So a natural question is whether $X$ is $\mathbb{A}^1$-naive in the sense of \cite[Definition 2.1.1]{asokII} or equivalently, $Sing_*(X)$ satisfies the affine Nisnevich excision \cite[Section 2.1]{asokI}.
\end{remark}




%\end{proof}


	\section{$\A^1$-invariance of $\pi_0^{\A^1}$ of smooth affine surfaces} \label{a1 invariance section}
In this section we will prove that for any smooth affine complex surface $X$, the $\mathbb{A}^1$-connected component sheaf $\pi_0^{\mathbb{A}^1}(X)$ is $\mathbb{A}^1$-invariant. For this let us recall the definition of $\mathbb{A}^1$-connected component sheaf first. 
\begin{definition} \cite[Section 3.2]{mv}
Suppose, $\sX$ is a space over $k$. The $\A^1$-connected component of $\sX$, denoted by $\pi_0^{\A^1}(\sX)$, is a Nisnevich sheaf of sets on $Sm/k$ associated to the presheaf 
$$U \in Sm/k \mapsto Hom_{\sH(k)}(U, \sX).$$
The space $\sX$ is called $\A^1$-connected, if $\pi_0^{\A^1}(\sX)$ is isomorphic to $\Spec\ k$. 
\end{definition}
\begin{remark}
\begin{enumerate}

\item  For a space $\sX$ the canonical morphism $\sX \to \pi_0^{\A^1}(\sX)$ is an epimorphism \cite[Corollary 2.1.5]{am}.
\item The space $\sX$ is $\A^1$-connected, if and only if $\pi_0^{\A^1}(\sX)(\Spec\ F)$ is trivial, for every finitely generated, separable field extension $F$ over $k$ \cite[Lemma 3.3.6]{ictp}.
\end{enumerate}
\end{remark}
\begin{definition}($\A^1$-invariant, see also Definition \ref{a1 local definition}) \label{invariant definition}  \hfill
	\begin{enumerate}
		%\item A simplicial presheaf $\sX$ 
		%in $\Delta^{op}PSh(Sm/k)$ 
		%is called $\A^1$-local if for any simplicial presheaf $\sY$, 
		%in $\Delta^{op}PSh(Sm/k)$ 
		%the map
		%$$Hom_{\sH_s(k)}(\sY,\sX)\to Hom_{\sH_s(k)}(\sY\times\A^1,\sX),$$
		%induced by projection $\sY\times\A^1 \to \sY$, is a bijection of sets. \par
		\item	A presheaf of sets $\sF$ is called an $\A^1$-invariant if for any $U\in Sm/k$ the map
		$$\sF(U)\to \sF(U\times\A^1_k)$$ 
		induced by the projection $U\times\A^1_k \to U$ is a bijection of sets \cite[Definition 1.7]{mor}.
		\item A scheme $X$ over $k$ is called an $\A^1$-rigid, if $X$ is $\A^1$-invariant, as representable sheaf \cite[Example 2.4]{mv}. 
	\end{enumerate}
\end{definition}
\begin{remark} \cite[Lemma 2.16]{mazza} \label{mazza}
A presheaf of sets $\sF$  on $Sm/k$ is $\A^1$-invariant if and only if for every $U \in Sm/k$ the morphisms 
$$i_0^*, i_1^*: \sF(\A^1_k \times_k U) \to \sF(U),$$
induced by the $0$-section and $1$-section respectively, are equal.
\end{remark}
For a space $\sX$, the presheaf of sets $U \in Sm/k \mapsto Hom_{\sH(k)}(U, \sX)$ is $\A^1$-invariant, by the construction of $\sH(k)$. Morel conjectured that the sheaf $\pi_0^{\A^1}(\sX)$ is also $\A^1$-invariant \cite[Conjecture 1.12]{mor}. The Conjecture is false, in general. Ayoub constructed a space $\sX$ for which $\pi_0^{\A^1}(\sX)$ is not $\A^1$-invariant and this $\sX$ is not a representable sheaf \cite[Construction 8]{ayoub}.
However, $\pi_0^{\A^1}(\sX)$ is known to be $\A^1$-invariant for the following spaces $\sX$:
\begin{enumerate}
\item The space $\sX$ is $\A^1$-connected. 
\item The space $\sX$ is itself an $\A^1$-invariant Nisnevich sheaf or $\sX$ is an $\A^1$-rigid scheme (for example, abelian varieties, proper open subscheme of $\A^1_k$, a smooth projective curve of genus $g>0$) \cite[Example 2.1.10]{am}.
\item The space $\sX$ is a motivic $H$-group or $\sX$ is a homogeneous space for a motivic $H$-group \cite[Theorem 4.18]{u}.
\item The space $\sX$ is a smooth toric variety \cite[Lemma 4.2, Lemma 4.4]{w}.
\item The space $\sX$ is a smooth projective surface (over any field, if $\sX$ is a non-uniruled surface \cite[Corollary 3.15]{bhs} and over an algebraically closed field of characteristic zero, if $\sX$ is a ruled surface \cite[Theorem 1.2]{bs}).
\end{enumerate}

In this section, we will prove that if $X$ is an affine surface over an algebraically closed field of characteristic zero, then $\pi_0^{\A^1}(X)$ is an $\A^1$-invariant sheaf (Theorem \ref{connected-base-new}).

Let us first recall the construction and properties of the universal $\A^1$-invariant sheaf $\sL(\sF)$ from \cite{bhs}, associated to a Nisnevich sheaf of sets $\sF$ on $Sm/k$.
\begin{definition} \cite[Definition 2.9]{bhs}
%Suppose, $\sF$ is a Nisnevich sheaf of sets on $Sm/k$
Given a Nisnevich sheaf of sets $\sF$ on $Sm/k$, the sheaf $\sS(\sF)$ is defined to be the Nisnevich sheaf of sets associated to the presheaf $\sS^{\text{pre}}(\sF)$, which is defined as
$$\sS^{\text{pre}}(\sF)(U) = \sF(U)/\sim, \text{ for }U \in Sm/k,$$
where $\sF(U)/\sim$ the quotient corresponding to the equivalence relation ``$\sim$" on $\sF(U)$, generated by the naive $\A^1$-homotopies; briefly for $\alpha, \beta \in \sF(U)$, $\alpha \sim \beta$ if there are naive $\A^1$-homotopies $H_1, \dots, H_m \in \sF(U)$, satisfying $H_j\circ i_1 = H_{j+1}\circ i_0$, for every $j$, such that $H_1 \circ i_0 = \alpha$ and $H_m \circ i_1 = \beta$ (here $i_0, i_1: U \to \A^1_U$ are the $0$-section and the $1$-section respectively).

For an integer $n > 1$, the sheaf $\sS^n(\sF)$ is defined inductively
$$\sS^n(\sF) := \sS(\sS^{n-1}(\sF))$$
and $\sS^0(\sF) = \sF$. The quotient map $\sF(U) \to \sF(U)/\sim$ induces the canonical morphism
$$\eta: \sF \to \sS(\sF)$$
which is an epimorphism of Nisnevich sheaves. For every $n$, this $n$-th iteration $\sS^n$ defines a functor on the category of Nisnevich sheaves on $Sm/k$ and the canonical epimorphism $\sS^n(\sF) \to \sS^{n+1}(\sF)$ defines the natural transformation $\sS^n \to \sS^{n+1}$. The sheaf $\sL(\sF)$ is defined to be the filtered colimit with respect to the morphisms $\sS^n(\sF) \to \sS^{n+1}(\sF)$:
$$\mathcal{L}(\mathcal{F}) := \varinjlim_{n \geq 0}\mathcal{S}^n(\mathcal{F}).$$
There is an induced morphism $\sF \to \sL(\sF)$, which is an epimorphism.
\end{definition}
\begin{remark}
\begin{enumerate}
\item For $X \in Sm/k$, the sheaf $\sS(X)$ is the $\A^1$-chain connected component sheaf $\pi_0^{ch}(X)$, defined by Asok-Morel \cite[Definition 2.2.4]{am}; it is the Nisnevich sheaf associated to the presheaf $U \mapsto \pi_0(Sing_*(X)(U))$ \cite[Remark 2.10]{bhs}.
\item For a sheaf $\sF$, the canonical epimorphism $\sF \to \pi_0^{\A^1}(\sF)$ factors through the epimorphism $\sS(\sF) \to \pi_0^{\A^1}(\sF)$.
\item The sheaf $\sL(\sF)$ is $\A^1$-invariant \cite[Theorem 2.13]{bhs}.
\item For a sheaf $\sF$, both the sheaves $\pi_0^{\A^1}(\sF)$ and $\sL(\sF)$ satisfy the same universal property \cite[Remark 2.15]{bhs}: given a morphism $\sF \to \sG$ to an $\A^1$-invariant sheaf $\sG$, it factors uniquely through the canonical morphisms $\sF \to \pi_0^{\A^1}(\sF)$ and $\sF \to \sL(\sF)$. Since $\sL(\sF)$ is $\A^1$-invariant, this induces the morphism $\pi_0^{\A^1}(\sF) \to \sL(\sF)$, which is an epimorphism and  this morphism is an isomorphism if and only if $\pi_0^{\A^1}(\sF)$ is $\A^1$-invariant \cite[Corollary 2.18]{bhs}.
\item The canonical morphism $\pi_0^{\A^1}(\sF) \to \sL(\sF)$ induces the bijection over the sections $\Spec\ K$, for every finitely generated separable field extension $K/k$ (\cite[Corollary 2.2]{brs}, \cite[Corollary 2.15]{cb}).
\item If $X \in Sm/k$ is a proper, non-uniruled surface, then $\sS(X) \cong \sS^2(X)$ \cite[Theorem 3.14]{bhs}. We will prove (Theorem \ref{non-uniruled}) that if $X \in Sm/k$ is an affine, non $\A^1$-uniruled surface (Definition \ref{ruled-defn}), then also 
$$\sS(X) \cong \sS^2(X).$$
\end{enumerate}
\end{remark}
Let us recall the following two classes of affine varieties with dominant family of affine lines. Suppose, $X \in Sm/k$ is an affine variety, where $k$ is an algebraically closed field field of characteristic $0$.
\begin{definition}\label{ruled-defn}
\begin{enumerate}
%begin{definition}
\item $X$ is said to be \textbf{$\mathbb{A}^1$-uniruled or log-uniruled} if there is a dominant generically finite morphism $H: \mathbb{A}^1_k \times_k Y \to X$ for some $k$-variety $Y$. 
\item $X$ is said to be \textbf{$\mathbb{A}^1$-ruled} or affine-ruled if there is a Zariski open dense subset $U$ of $X$ such that $U$ is isomorphic to $\mathbb{A}^1_k \times_k Z$ for some $k$-variety $Z$ \cite[Section 2.1]{miyanishicrm}.
%\end{definition}
\end{enumerate}
\end{definition}
%\begin{remark}
%\begin{enumerate}
%\item An $\A^1$-ruled affine $k$-variety is also $\A^1$-uniruled.
%\item An $\A^1$-uniruled affine $k$-variety has negative logarithmic Kodaira dimension.
% (see \cite[??]{miyanishicrm}).
%\item 
%\end{enumerate}
%\end{remark}
In case of smooth affine surfaces, we have the following equivalences
%We can classify 
%$\A^1$-ruled or $\A^1$-uniruled surfaces are distinguished by the negativity
%according to their logarithmic Kodaira dimension (denoted by $\bar{k}(X)$). From 
(\cite[Chapter 2, Theorem 2.1.1 and Chapter 3, Lemma 1.3.1]{miyanishicrm} and \cite[Lemma 1.8]{russell}):
\begin{theorem} \label{all uniruled equivalences}
Let $X$ be a smooth affine surface over an algebraically closed field $k$ of characteristic zero.
%\%begin{definition}
%$X$ is said to be \textbf{$\mathbb{A}^1$-ruled} or affine-ruled if there is a Zariski open dense subset $U$ of $X$ such that $U$ is isomorphic to $\mathbb{A}^1_k \times_k Z$ for some $k$-variety $Z$ \cite[Section 2.1]{miyanishicrm}.
%\end{definition}
%Let $X$ be an affine surface over $\C$. w
%We can classify $\A^1$-ruled or $\A^1$-uniruled surfaces, according to their logarithmic Kodaira dimension (denoted by $\bar{k}(X)$). From \cite[Chapter 2, Theorem 2.1.1 and Chapter 3, Lemma 1.3.1]{miyanishicrm} and \cite[Lemma 1.8]{russell},
%, \cite[Chapter 2, Theorem 2.1.1 and Chapter 3, Lemma 1.3.1] and \cite[Chapter 3, Lemma 1.3.1], \cite[Lemma 1.8]{russell}{miyanishicrm}, 
Then 
%the we have 
the following are equivalent:
\begin{enumerate}
	\item $X$ is $\mathbb{A}^1$-uniruled.
	\item $X$ has negative logarithmic Kodaira dimension.
	\item $X$ is $\mathbb{A}^1$-ruled.
	\item $X$ admits an $\mathbb{A}^1$-fibration $\pi: X \to C$.
\end{enumerate}
\end{theorem}
We will prove the $\A^1$-invariance of $\pi_0^{\A^1}(X)$, for a smooth affine surface $X$ over an algebraically closed field $k$ of characteristic zero, dividing in two classes: one is non $\A^1$-uniruled affine surfaces and the other one is $\A^1$-uniruled affine surfaces. Equivalently in terms of the logarithmic Kodaira dimension (denoted by $\bar{k}(X)$): one class consists such $X$ having $\bar{k}(X)$ is non-negative 
%logarithmic Kodaira dimension, that is $\bar{k}(X) \in \{0,1,2\}$ 
and the other class consists such $X$ having $\bar{k}(X)$ is negative.
%surfaces having negative logarithmic Kodaira dimension, that is $\bar{k}(X) = - \infty$.
%	Surfaces can be classified by their logarithmic Kodaira dimension, denoted by $\kappa(X)$ for a surface $X$. 
	\subsection{$\A^1$-invariance of $\pi_0^{\A^1}(X)$, for non $\A^1$-uniruled, affine surface $X$:}
	
	Suppose, $X$ is a smooth affine non $\A^1$-uniruled surface over an algebraically closed field $k$ of characteristic $0$. In this subsection, we will prove that $\pi_0^{\A^1}(X)$ is $\A^1$-invariant, by proving $\sS(X) \cong \sS^2(X)$ (Theorem \ref{non-uniruled}, compare this with \cite[Theorem 3.14]{bhs} in case of proper, non-uniruled surfaces).

%\begin{fact}
Suppose, $\alpha: \A^1_k \to X$ is a morphism, where $X \in Sm/k$ is affine. Then the image of $\alpha$ (set-theoretic image) is closed in $X$. Indeed, if the image $\text{Im}(\alpha)$ is non-constant, then we extend $\alpha$ to $\Tilde{\alpha}: \mathbb{P}^1_k \to \bar{X}$ ($\bar{X}$ is a smooth compactification of $X$). The morphism $\Tilde{\alpha}$ is a projective morphism, since $\Tilde{\alpha}$ factors as
$$\mathbb{P}^1_k \xrightarrow{\text{graph of } \Tilde{\alpha}} \mathbb{P}^1_k \times_k \bar{X} \xrightarrow{\text{projection}} \overline{X}.$$
So $\Tilde{\alpha}$ is a proper morphism and hence $\text{Im}(\Tilde{\alpha})$ is closed in $\bar{X}$. Since $X$ is affine, so any morphism from $\mathbb{P}^1_k$ to $X$ is constant; thus the point of infinity of $\mathbb{P}^1_k$ maps to a $k$-point in $\bar{X} \setminus X$. Therefore, $\text{Im}(\alpha) = \text{Im}(\Tilde{\alpha}) \cap X$ is closed in $X$. 

 Before proving the main theorem, we need the following technical lemma, which we will use in the proof.

%\end{fact}
%and $X$ is not an $\mathbb{A}^1$-uniruled surface.

%By an affine line in $X$, we mean that there is a non-constant morphism $\gamma: \mathbb{A}^1_k \to X$. The image of an affine line $\gamma: \mathbb{A}^1_k \to X$ is closed in $X$. Indeed, if the image $Im(\gamma)$ is non-constant, then we extend $\gamma$ to $\Tilde{\gamma}: \mathbb{P}^1_k \to \overline{X}$ ($\overline{X}$ is a smooth compactification of $X$). The morphism $\Tilde{\gamma}$ is a projective morphism, since $\Tilde{\gamma}$ factors as
%$$\mathbb{P}^1_k \xrightarrow{\text{graph of } \Tilde{\gamma}} \mathbb{P}^1_k \times \overline{X} \xrightarrow{\text{projection}} \overline{X}.$$
%So $\Tilde{\gamma}$ is a proper morphism and hence $Im(\Tilde{\gamma})$ is closed. Since any morphism from $\mathbb{P}^1_k$ to $X$ is constant, so the point of infinity of $\mathbb{P}^1_k$ maps to a point in $\overline{X} \setminus X$. Therefore, $Im(\gamma) = Im(\Tilde{\gamma}) \cap X$ is closed in $X$.  

%\begin{lemma} \label{not contained1}
%Suppose, $\theta: U \to X$ is a morphism, where $U, X \in Sm/k$ irreducible. Assume that there is an $\alpha: \mathbb{A}^1_k \to X$ such that $Im(\theta) \subset Im(\alpha)$. Then $\theta = \alpha(0) \in \mathcal{S}(X)(U)$, where $\alpha(0)$ is considered as a morphism
%$$U \to Spec  \ k \xrightarrow{\alpha(0)} X$$.
%\end{lemma}
%\begin{proof}
%Since $Im(\theta)$ is contained in the image of $\mathbb{A}^1_k$, so there is a morphism $\Tilde{\theta}: U \to \mathbb{A}^1_k$ such that $\theta = \alpha \circ \Tilde{\theta}$.
%Consider the following morphisms:
%$$
%\begin{tikzcd}
%\mathcal{S}(U)
%\xrightarrow{\mathcal{S}(\Tilde{\theta})} 
 %\mathcal{S}(\mathbb{A}^1_k) \xrightarrow{\mathcal{S}(\alpha)}
 %\mathcal{S}(X)
%\end{tikzcd}
%$$
%Since $\mathcal{S}(\mathbb{A}^1_k)$ is isomorphic to the trivial sheaf $\Spec\ k$, so the morphism $\mathcal{S}(\alpha)$ is equal to the morphism $\alpha(0)$.
%Here the the morphisms from $\mathcal{S}(U)$ are $\mathcal{S}(\Tilde{\theta})$ and $\mathcal{S}(\Tilde{\gamma})$, induced by $\Tilde{\theta}$ and $\Tilde{\gamma}$ respectively. The rightmost morphism is $\mathcal{\alpha}$, induced by $\alpha$. 
%Here 
%\begin{align*}
%$$(\mathcal{S}(\alpha) \circ \mathcal{S}(\Tilde{\theta}))(Id_U) = \mathcal{S}(\theta)(Id_U) =  \theta$$
%\end{align*}
%Therefore, $\theta = \alpha(0) \in \mathcal{S}(X)(U)$.
%Similarly, 
%\begin{align*}
%\mathcal{S}(\alpha) \circ \mathcal{S}(\Tilde{\gamma})(Id_U) &= \mathcal{S}(\gamma)(Id_U) = \gamma                             
%\end{align*}
%The sheaf $\mathcal{S}(\mathbb{A}^1_k)$ is the trivial sheaf $\Spec\ k$, so $\gamma = \theta \in \mathcal{S}(X)(U)$.
%\end{proof}

%Before proving the main theorem
\begin{lemma} \label{not contained}
Suppose, $\theta: U \to X$ is a morphism, where $U, X \in Sm/k$ and $X$ is affine. Assume that there is an $\alpha: \mathbb{A}^1_k \to X$ such that 
%$$\mathrm{pr}$$
$\mathrm{Im}(\theta) \subset \mathrm{Im}(\alpha)$ (Notations \ref{notations}). Then $\theta = \alpha(0) \in \mathcal{S}(X)(U)$, where $\alpha(0)$ is considered as the morphism
$$U \to Spec  \ k \xrightarrow{\alpha(0)} X$$.
\end{lemma}
\begin{proof}
Since $\text{Im}(\theta)$ is contained in the image of $\mathbb{A}^1_k$, so there is a morphism $\Tilde{\theta}: U \to \mathbb{A}^1_k$ such that $\theta = \alpha \circ \Tilde{\theta}$. Indeed, if there is some irreducible component, say $U_0$ of $U$ such that its image is constant (a $k$-point), then $\theta|_{U_0}$ always factors through $\mathbb{A}^1_k$ (by taking a $k$-point from the fiber of $\alpha$ over that point), otherwise $\theta|_{U_0}: U_0 \to \text{Im}(\alpha)$ is dominant ($\text{Im}(\alpha)$ is a closed subscheme of $X$). 
%If $\theta|_{U_0}$ is dominant to $Im(\alpha)$, then 
In this case, $\theta|_{U_0}$ factors through the normalisation of $\text{Im}(\alpha)$, which is $\mathbb{A}^1_k$ itself. Therefore, $\theta$ always factors through $\mathbb{A}^1_k$. 
Consider the naive $\mathbb{A}^1$-homotopy $H: \mathbb{A}^1_U \to X$ defined as
$$(t,x) \mapsto \alpha(t\Tilde{\theta}(x)).$$
Then $H(0, x) = \alpha(0)$ and $H(1, x) = \theta(x)$.
Therefore, $\theta = \alpha(0) \in \mathcal{S}(X)(U)$.
%Similarly, 
%\begin{align*}
%\mathcal{S}(\alpha) \circ \mathcal{S}(\Tilde{\gamma})(Id_U) &= \mathcal{S}(\gamma)(Id_U) = \gamma                             
%\end{align*}
%The sheaf $\mathcal{S}(\mathbb{A}^1_k)$ is the trivial sheaf $\Spec\ k$, so $\gamma = \theta \in \mathcal{S}(X)(U)$.
\end{proof}
\begin{corollary} \label{not contained corr}
Suppose, $\theta, \gamma: U \to X$,
% are two morphisms
where $U, X \in Sm/k$ and $X$ is affine. Assume that there are $\alpha_1, \dots, \alpha_n : \mathbb{A}^1_k \to X$ such that $\alpha_i(1) = \alpha_{i+1}(0)$, for every $i$ and there are $p, q \leq n$ such that $\mathrm{Im}(\theta) \subset \mathrm{Im}(\alpha_p)$ and $\mathrm{Im}(\gamma) \subset \mathrm{Im}(\alpha_q)$. Then, $\theta = \gamma \in \mathcal{S}(X)(U)$.
\end{corollary}
\begin{proof}
Since $\text{Im}(\theta) \subset \text{Im}(\alpha_p) $ and $\text{Im}(\gamma) \subset \text{Im}(\alpha_q)$, so 
$$\theta = \alpha_p(0) \in \mathcal{S}(X)(U) \text{ and } \gamma = \alpha_q(0) \in \mathcal{S}(X)(U)$$
by Lemma \ref{not contained}. Therefore, $\theta =  \gamma \in \mathcal{S}(X)(U)$, since $\alpha_p$ and $\alpha_q$ are joined by a chain of $\mathbb{A}^1_k$'s.
\end{proof}
%\begin{proposition}
%The canonical epimorphism
%$$\mathcal{S}(X) \to \mathcal{S}^2(X)$$
%is an isomorphism.
%\end{proposition}
%\begin{remark}
%Suppose, $\alpha: \mathbb{A}^1_k \to X$ is a is a line in $X$. Then if the 
%\end{remark}
%\begin{lemma} \label{not contained}
%Suppose, $\theta, \gamma: U \to X$ are two morphisms, where $U, X \in Sm/k$. Assume that there is an $\alpha: \mathbb{A}^1_k \to X$ such that $Im(\gamma), Im(\theta) \subset Im(\alpha)$. Then $\theta = \gamma \in \mathcal{S}(X)(U)$.
%\end{lemma}
%\begin{proof}
%Since $Im(\gamma), Im(\theta)$ is contained in a line, so there are $\Tilde{\gamma}, \Tilde{\theta}: U \to \mathbb{A}^1_k$ such that $\gamma = \alpha \circ \Tilde{\gamma}$ and $\theta = \alpha \circ \Tilde{\theta}$.
%Consider the following commutative diagraam:
%$$
%\begin{tikzcd}
%\mathcal{S}(U)
%\ar[r, bend left] \ar[r, bend right]
%& \mathcal{S}(\mathbb{A}^1_k) \ar[r]
%& \mathcal{S}(X)
%\end{tikzcd}
%$$
%Here the the morphisms from $\mathcal{S}(U)$ are $\mathcal{S}(\Tilde{\theta})$ and $\mathcal{S}(\Tilde{\gamma})$, induced by $\Tilde{\theta}$ and $\Tilde{\gamma}$ respectively. The rightmost morphism is $\mathcal{\alpha}$, induced by $\alpha$. Here 
%\begin{align*}
%\mathcal{S}(\alpha) \circ \mathcal{S}(\Tilde{\theta})(Id_U) &= \mathcal{S}(\theta)(Id_U) =  \theta
%\end{align*}
%Similarly, 
%\begin{align*}
%\mathcal{S}(\alpha) \circ \mathcal{S}(\Tilde{\gamma})(Id_U) &= \mathcal{S}(\gamma)(Id_U) = \gamma                             
%\end{align*}
%The sheaf $\mathcal{S}(\mathbb{A}^1_k)$ is the trivial sheaf $\Spec\ k$, so $\gamma = \theta \in \mathcal{S}(X)(U)$.
%\end{proof}
%\begin{theorem} \label{S and S2}
%The canonical epimorphism
%$$\mathcal{S}(X) \to \mathcal{S}^2(X)$$
%is an isomorphism.
%\end{theorem}
%\begin{proof}
%We need to show that the canonical surjection
%$$\mathcal{S}(X)(U) \to \mathcal{S}^2(X)(U)$$
%is a bijection, for every  smooth Henselian local scheme $U \in Sm/k$.

%If possible assume that there are two morphisms $\gamma, \theta: U \to X$ such that $\theta \neq \gamma \in \mathcal{S}(X)(U)$ and $\theta = \gamma \in \mathcal{S}^2(X)(U)$. Thus there is a chain of naive $\mathbb{A}^1$-homotopies $H_1, \dots, H_m \in \mathcal{S}(X)(\mathbb{A}^1_U)$ such that $H_1(0) = \theta$ and $H_m(1) = \gamma$. 
%We can assume all $H_i$'s are non-constant. We can also assume there is a non-constant $H \in \mathcal{S}(X)(\mathbb{A}^1_U)$ such that $H(0) = \theta$ and $H(1) = \gamma$.
%Moreover, since $\gamma \neq \theta \in \mathcal{S}(X)(U)$ so by  Lemma \ref{not contained}, there is no chain of lines $\alpha_1, \dots, \alpha_n: \mathbb{A}^1_k \to X$ with $\alpha_i(1) = \alpha_{i+1}(0)$, for every $i \leq n$ such that $Im(\theta) \subset Im(\alpha_p)$ and $Im(\gamma) \subset \alpha_q$, for some $p, q \leq n$. 
%contains both $Im(\theta)$ and $Im(\gamma)$.

%Since the canonical morphism $\eta: X \to \mathcal{S}(X)$ is an epimorphism, so there is a Nisnevich covering $f: V \to \mathbb{A}^1_U$ and a morphism $\phi: V \to X$ which makes the following diagram commutative:
 %$$\begin{tikzcd}
 %V \ar[r, "\phi"] \ar[d, "f"] &X \ar[d, "\eta"] \\
%\mathbb{A}^1_{U} \ar[r,"H"] 
%&\mathcal{S}(X)
%\end{tikzcd}$$ 
%Moreover, since $H$ is non-constant, there exists a Nisnevich covering $V^{\prime} \to V \times_{\mathbb{A}^1_U} V$ such that 
%$$\phi \circ \text{pr}_1|_{V^\prime} \neq \phi \circ \text{pr}_2|_{V^\prime} : V^{\prime} \to X,$$
%where $\text{pr}_1, \text{pr}_2: V \times_{\mathbb{A}^1_U} V \to V$ are the projection maps.
%Thus there is a chain of non-constant naive $\mathbb{A}^1$-homotopies $G_1, \dots, G_l : \mathbb{A}^1_{V^{\prime}} \to X$ such that $G_1(0) = \phi \circ \text{pr}_1|_{V^\prime}$ and $G_l(1) = \phi \circ \text{pr}_2|_{V^\prime}$.
%Since any Nisnevich covering of a Henselian local scheme is trivial there is a lift $\theta^\prime: U \to V$ such that the following diagram:
%$$\begin{tikzcd}[column sep=50pt, row sep=50 pt]
%& V \ar[r, "\phi"] \ar[d, "f"] &X \ar[d] \\
 %U \ar[ur, "\theta^\prime"] \ar[r, "i_0"] \ar[urr, dotted, "H(0)", near end]
%& \mathbb{A}^1_{U} \ar[r,"H"] 
%&\mathcal{S}(X)
%\end{tikzcd}$$ 
%Here $i_0: U \to \mathbb{A}^1_U$ is the $0$-section. Here the left most triangle is commutative and the square is commutative. However, the other triangle is not commutative, i.e $\phi \circ \theta^\prime \neq H(0)$, in general. But as the square is commutative, $\phi \circ \theta^\prime = \Tilde{\theta}(say) = \theta \in \mathcal{S}(X)(U)$.
%Similarly considering the $1$-section $i_1: U \to \mathbb{A}^1_U$, there is a lift $\gamma^\prime: U \to V$ such that we have the following diagram:
%$$\begin{tikzcd}[column sep=50pt, row sep=50 pt]
%& V \ar[r, "\phi"] \ar[d, "f"] &X \ar[d] \\
 %U \ar[ur, "\gamma^\prime"] \ar[r, "i_1"] \ar[urr, dotted, "H(1)", near end]
%& \mathbb{A}^1_{U} \ar[r,"H"] 
%&\mathcal{S}(X)
%\end{tikzcd}$$ 
%Here the left most triangle is commutative and the square is commutative. However, the other triangle is not commutative, i.e $\phi \circ \gamma^\prime \neq H(1)$, in general. But as the square is commutative, $\phi \circ \gamma^\prime = \Tilde{\gamma}(say) = \gamma \in \mathcal{S}(X)(U)$. Since, $\theta \neq \gamma \in \mathcal{S}(X)(U)$, so $\Tilde{\theta} \neq \tilde{\gamma} \in \mathcal{S}(X)(U)$.

%The triangle is not commutative, in general however $\phi \circ \gamma^\prime = \Tilde{\gamma}(say) = \gamma \in \mathcal{S}(X)(U)$
%The image of $\phi$ contains $Im(\Tilde{\theta})$ and $Im(\Tilde{\gamma})$, $Im(\Tilde{\theta}) \neq Im(\Tilde{\gamma}) \in \mathcal{S}(X)(U)$.
%Suppose, $V  = \coprod_{i =1}^n V_i$, where $V_i$'s are the irreducible components of $V$, which are also the connected components of $V$. Since $U$ is irreducible, $Im(\Tilde{\theta}) \subset \overline{\phi(V_i)}$ and $Im(\Tilde{\gamma}) \subset \overline{\phi(V_j)}$, for some $i, j$. 
%If there are non-constant homotopies 
%We complete the proof considering the following cases.



%\textbf{Case 1:} Suppose, $Im(\Tilde{\gamma}), Im(\Tilde{\theta}) \subset \overline{\phi(V_i)}$, for every $i$. There is an irreducible component $V_0$ of $V^{\prime}$ such that
%$$\phi \circ \text{pr}_1|_{V_0} \neq \phi \circ \text{pr}_2|_{V_0}.$$
%We can assume all the naive $\mathbb{A}^1$-homotopies $G_1|_{\mathbb{A}^1_{V_0}}, \dots, G_l|_{\mathbb{A}^1_{V_0}}$ are non-constant.
%Suppose, $V_0$ maps to $V_r \times_{\mathbb{A}^1_k} V_s$, for some $r, s \leq n$. Since the projection maps and the morphism $V^{\prime} \to V \times_{\mathbb{A}^1_U}$ are \'etale maps, so $\overline{Im(G_1|_{\mathbb{A}^1_{V_0}}(0))} = \overline{\phi(V_r)}$ and $\overline{Im(G_l|_{\mathbb{A}^1_{V_0}}(1))} = \overline{\phi(V_s)}$. Since every $G_i|_{\mathbb{A}^1_{V_0}}$ is non-constant, so $Im(G_i|_{\mathbb{A}^1_{V_0}})$ contains a line. If for every $i$, $\overline{Im(G_i|_{\mathbb{A}^1_{V_0}})}$ is contained in the image of some $\beta: \mathbb{A}^1_k \to X$, then there is a chain of $\mathbb{A}^1_k$'s in $X$ that contains both $Im(\Tilde{\theta}), Im(\Tilde{\gamma})$, which would imply that $\Tilde{\theta} = \Tilde{\gamma} \in \mathcal{S}(X)(U)$, by Corollary \ref{not contained corr} and it is a contradiction. Therefore, there is some $j$ such that $G_j|_{\mathbb{A}^1_{V_0}}$ is non-constant and dominant. This is a contradiction, since $X$ is not an $\mathbb{A}^1$-uniruled surface.




%\textbf{Case 2:} Suppose, $Im(\Tilde{\theta}) \subset \overline{\phi(V_r)}$ and $Im(\Tilde{\gamma}) \subset \overline{\phi(V_s)}$, for some $r, s$. Also assume that $Im(\Tilde{\gamma}) \nsubseteq \overline{\phi(V_r)}$. There is an irreducible component $V_0$ of $V^\prime$ that maps to $V_r \times_{\mathbb{A}^1_U} V_s$. Since the projection maps and the morphism $V^{\prime} \to V \times_{\mathbb{A}^1_U}$ are \'etale maps, so $\overline{Im(\phi \circ \text{pr}_1|_{V_0})} = \overline{\phi(V_r)}$ and $\overline{Im(\phi \circ \text{pr}_2|_{V_0})} = \overline{\phi(V_s)}$. Since $Im(\Tilde{\gamma}) \nsubseteq \overline{\phi(V_r)}$, so
%$$\phi \circ \text{pr}_1|_{V_0} \neq \phi \circ \text{pr}_2|_{V_0}.$$
%So we can assume all the naive $\mathbb{A}^1$-homotopies $G_1|_{\mathbb{A}^1_{V_0}}, \dots, G_l|_{\mathbb{A}^1_{V_0}}$ are non-constant. Since every $G_i|_{\mathbb{A}^1_{V_0}}$ is non-constant, so $Im(G_i|_{\mathbb{A}^1_{V_0}})$ contains a line. If for every $i$, $\overline{Im(G_i|_{\mathbb{A}^1_{V_0}})}$ is contained in the image of some $\beta: \mathbb{A}^1_k \to X$, then there is a chain of $\mathbb{A}^1_k$'s in $X$ that contains both $Im(\Tilde{\theta}), Im(\Tilde{\gamma})$, which would imply that $\Tilde{\theta} = \Tilde{\gamma} \in \mathcal{S}(X)(U)$, by Corollary \ref{not contained corr} and it is a contradiction. Therefore, there is some $j$ such that $G_j|_{\mathbb{A}^1_{V_0}}$ is non-constant and dominant. This is a contradiction, since $X$ is not an $\mathbb{A}^1$-uniruled surface.




%Thus in both cases we arrive at a contradiction. Therefore, for every smooth henselian local scheme $U \in Sm/k$, the morphism
%$$\mathcal{S}(X)(U) \to \mathcal{S}^2(X)(U)$$
%is injective. Hence, the canonical epimorphism
%$$\mathcal{S}(X) \to \mathcal{S}^2(X)$$
%is an isomorphism.

%If for every $i$, $\overline{Im(G_i|_{\mathbb{A}^1_{V_0}})}$ is a line, then $Im(\Tilde{\theta}), Im(\Tilde{\gamma})$ are contained in a chain of $\mathbb{A}^1_k$'s.
%\end{proof}


%\begin{proposition} \label{ghost homotopy implies dominant}
%Let $X \in Sm/k$ be an 
%non $\mathbb{A}^1$-uniruled 
%affine surface and $U \in Sm/k$ be irreducible. 
%Then there is $W \in Sm/k$ irreducible al. 
%Suppose that there is a non-constant homotopy $H \in \mathcal{S}(X)(\mathbb{A}^1_U)$. Then there is $W \in Sm/k$ irreducible along with a dominant morphism $G: \mathbb{A}^1_W \to X$ such that $G(0,-) \neq G(1, -): W \to X$. 
%for some $W \in$ 
%\end{proposition}
%\begin{proof}
%The canonical morphism $\eta: X \to \mathcal{S}(X)$ is an epimorphism, so there is a Nisnevich covering $f: V \to \mathbb{A}^1_U$ and a morphism $\phi: V \to X$ which makes the following diagram commutative:
 %$$\begin{tikzcd}
 %V \ar[r, "\phi"] \ar[d, "f"] &X \ar[d, "\eta"] \\
%\mathbb{A}^1_U \ar[r,"H"] 
%&\mathcal{S}(X)
%\end{tikzcd}$$ 
%Since $H$ is non-constant, the morphisms $\phi \circ \text{pr}_1|_{V^\prime}$ and $\phi \circ \text{pr}_2|_{V^\prime}$ (here $\text{pr}_1, \text{pr}_2: V \times_{\mathbb{A}^1_U} V \to V$ are the projection maps) are not same in $X(V^\prime)$, for any Nisnevich covering $V^\prime \to V \times_{\mathbb{A}^1_U} V$. However from the commutative square, $\phi \circ \text{pr}_1 = \phi \circ \text{pr}_2 \in \mathcal{S}(X)(V \times_{\mathbb{A}^1_U} V)$.
%, but they , since $H$ is non-co but they. 
%Thus there is a Nisnevich covering $V^\prime \to V \times_{\mathbb{A}^1_U} V$ and a chain of non-constant naive $\mathbb{A}^1$-homotopies $G_1, \dots, G_l : \mathbb{A}^1_{V^{\prime}} \to X$ such that $G_1(0) = \phi \circ \text{pr}_1|_{V^\prime}$ and $G_l(1) = \phi \circ \text{pr}_2|_{V^\prime}$.


%Since $H$ is non-constant, The morphisms $\phi \circ \text{pr}_1$ and $\phi \circ \text{pr}_2$ (here $\text{pr}_1, \text{pr}_2: V \times_{\mathbb{A}^1_U} V \to V$ are the projection maps) are same in $\mathcal{S}(X)(V \times_{\mathbb{A}^1_U} V)$, but they are not same in $X(V^\prime)$, for any Nisnevich covering $V^\prime \to V \times_{\mathbb{A}^1_U} V$, since $H$ is non-co but they . Thus there is a Nisnevich covering $V^\prime \to V \times_{\mathbb{A}^1_U} V$ and a chain of non-constant naive $\mathbb{A}^1$-homotopies $G_1, \dots, G_l : \mathbb{A}^1_{V^{\prime}} \to X$ such that $G_1(0) = \phi \circ \text{pr}_1|_{V^\prime}$ and $G_l(1) = \phi \circ \text{pr}_2|_{V^\prime}$.
%Consider the following commutative diagram for the $0$-section of $H$:
%$$\begin{tikzcd}[column sep=50pt, row sep=50 pt]
%U^\prime \ar[r, "\theta^\prime"] \ar[d, "f^\prime"] &V \ar[r, "\phi"] \ar[d, "f"] &X \ar[d, "\eta"] \\
% U \ar[r, "i_0"]  
%& \mathbb{A}^1_{U} \ar[r,"H"] 
%&\mathcal{S}(X)
%\end{tikzcd}$$ 
%Here $i_0:U \to \mathbb{A}^1_U$ is the $0$-section and the left most square is a pullback square. The morphism $f^\prime: U^\prime \to U$ is a Nisnevich covering. From the commutativity of the diagram we have, 
%$$H(0)|_{U^\prime} = \phi \circ \theta^\prime \in \mathcal{S}(X)(U^\prime).$$
%Consider the following commutative diagram for the $1$-section of $H$:
%$$\begin{tikzcd}[column sep=50pt, row sep=50 pt]
%U^{\prime\prime} \ar[r, "\theta^{\prime\prime}"] \ar[d, "f^{\prime\prime}"] &V \ar[r, "\phi"] \ar[d, "f"] &X \ar[d, "\eta"] \\
% U \ar[r, "i_1"]  
%& \mathbb{A}^1_{U} \ar[r,"H"] 
%&\mathcal{S}(X)
%\end{tikzcd}$$ 
%Here $i_1:U \to \mathbb{A}^1_U$ is the $1$-section and the left most square is a pullback square. The morphism $f^{\prime \prime}: U^{\prime \prime }\to U$ is a Nisnevich covering. From the commutativity of the diagram we have, 
%$$H(1)|_{U^{\prime \prime}} = \phi \circ \theta^{\prime \prime} \in \mathcal{S}(X)(U^{\prime \prime}).$$
%Consider the following pullback diagram:
%$$\begin{tikzcd}
 %\Tilde{U} \ar[r, "\psi^{\prime}"] \ar[d, "\psi^{\prime\prime}"] &U^\prime \ar[d, "f^\prime"] \\
%U^{\prime\prime} \ar[r,"f^{\prime \prime}"] 
%&U
%\end{tikzcd}$$ 
%Here $\Tilde{U}$ is a common refinement of $U^\prime$ and $U^{\prime \prime}$. Thus we have
%$$H(0)|_{\Tilde{U}} = \phi \circ \theta^{\prime} \circ \psi^\prime \in \mathcal{S}(X)(\Tilde{U}) \text{ and } H(1)|_{\Tilde{U}} = \phi \circ \theta^{\prime \prime} \circ \psi^{\prime \prime} \in \mathcal{S}(X)(\Tilde{U}).$$
%Moreover since $H$ is non-constant, we have 
%$$\phi \circ \theta^{\prime} \circ \psi^\prime \neq  \phi \circ \theta^{\prime \prime} \circ \psi^{\prime \prime} \in \mathcal{S}(X)(\Tilde{U}).$$
%Thus there is an irreducible component (which is also a connected component) $W$ of $\Tilde{U}$ such that 
%$$(\phi \circ \theta^{\prime} \circ \psi^\prime)|_W \neq  (\phi \circ \theta^{\prime \prime} \circ \psi^{\prime \prime})|_W \in \mathcal{S}(X)(W).$$
%Suppose, $(\phi \circ \theta^{\prime} \circ \psi^\prime)|_W = \theta_1 \in X(W)$ and $(\phi \circ \theta^{\prime \prime} \circ \psi^{\prime \prime})|_W = \theta_2 \in X(W)$. The sections $\theta_1 \neq \theta_2 \in \mathcal{S}(X)(W)$.
%Suppose, $V  = \coprod_{i =1}^n V_i$, where $V_i$'s are the irreducible components of $V$, which are also the connected components of $V$. Since $W$ is irreducible, $Im(\theta_1) \subset \overline{\phi(V_i)}$ and $Im(\theta_2) \subset \overline{\phi(V_j)}$, for some $i, j$. 
%If there are non-constant homotopies 
%We complete the proof considering the following cases.



%\textbf{Case 1:} Suppose, $Im(\theta_1), Im(\theta_2) \subset \overline{\phi(V_p)}$, for every $p$. There is an irreducible component $V_0$ of $V^{\prime}$ such that
%$$\phi \circ \text{pr}_1|_{V_0} \neq \phi \circ \text{pr}_2|_{V_0}.$$
%We can assume that the naive $\mathbb{A}^1$-homotopy $G_1|_{\mathbb{A}^1_{V_0}}$ is non-constant.
%Suppose, $V_0$ maps to $V_r \times_{\mathbb{A}^1_U} V_s$, for some $r, s \leq n$. Since the projection maps and the morphism $V^{\prime} \to V \times_{\mathbb{A}^1_U} V$ are the \'etale maps, so $\overline{Im(G_1|_{\mathbb{A}^1_{V_0}}(0))} = \overline{\phi(V_r)}$. 
%and $\overline{Im(G_l|_{\mathbb{A}^1_{V_0}}(1))} = \overline{\phi(V_s)}$. 
%Since $G_1|_{\mathbb{A}^1_{V_0}}$ is non-constant, so $Im(G_1|_{\mathbb{A}^1_{V_0}})$ contains a line (say $L$). Since $X$ is affine, $L$ is closed in $X$. If $\overline{Im(G_1|_{\mathbb{A}^1_{V_0}})} = L$, then by our assumption $Im(\theta_1), Im(\theta_2) \subset L$, since $\overline{Im(G_1|_{\mathbb{A}^1_{V_0}}(0))} = \overline{\phi(V_r)}$. This implies that $\theta_1 = \theta_2 \in \mathcal{S}(X)(W)$, by Corollary \ref{not contained corr} and it is a contradiction. Therefore, the non-constant homotopy $G_1|_{\mathbb{A}^1_{V_0}}$ is dominant.
%  is contained in the image of some $\beta: \mathbb{A}^1_k \to X$, then there is a chain of $\mathbb{A}^1_k$'s in $X$ that contains both $Im(\Tilde{\theta}), Im(\Tilde{\gamma})$, which would imply that $\Tilde{\theta} = \Tilde{\gamma} \in \mathcal{S}(X)(U)$, by Corollary \ref{not contained corr} and it is a contradiction. Therefore, there is some $j$ such that $G_j|_{\mathbb{A}^1_{V_0}}$ is non-constant and dominant. This is a contradiction, since $X$ is not an $\mathbb{A}^1$-uniruled surface.




%\textbf{Case 2:} Suppose there is $r < n$ such that $Im(\theta_1), Im(\theta_2) \subset \overline{\phi(V_1)}, \dots, \overline{\phi(V_r)}$ (renumbering if necessary). Also assume that $Im(\theta_1)  \nsubseteq \overline{\phi(V_p)}$ and $Im(\theta_2)  \nsubseteq \overline{\phi(V_p)}$, for every $p> r$. There is an irreducible component $V_0$ of $V^\prime$ that maps to $V_r \times_{\mathbb{A}^1_U} V_n$. Since the projection maps and the morphism $V^{\prime} \to V \times_{\mathbb{A}^1_U} V$ are the \'etale maps, so $\overline{Im(\phi \circ \text{pr}_1|_{V_0})} = \overline{\phi(V_r)}$ and $\overline{Im(\phi \circ \text{pr}_2|_{V_0})} = \overline{\phi(V_n)}$. Since $Im(\theta_1), Im(\theta_2)  \nsubseteq \overline{\phi(V_n)}$, so
%$Im(\Tilde{\gamma}) \nsubseteq \overline{\phi(V_r)}$, so
%$$\phi \circ \text{pr}_1|_{V_0} \neq \phi \circ \text{pr}_2|_{V_0}.$$
%We can assume that the naive $\mathbb{A}^1$-homotopy $G_1|_{\mathbb{A}^1_{V_0}}$ is non-constant. So $Im(G_1|_{\mathbb{A}^1_{V_0}})$ contains a line (say $L$). Since $X$ is affine, $L$ is closed in $X$.  If $\overline{Im(G_1|_{\mathbb{A}^1_{V_0}})} = L$, then by our assumption $Im(\theta_1), Im(\theta_2) \subset L$, since $G_1|_{\mathbb{A}^1_{V_0}}(0) = \phi \circ \text{pr}_1|_{V_0}$. This implies that $\theta_1 = \theta_2 \in \mathcal{S}(X)(W)$, by Corollary \ref{not contained corr} and it is a contradiction. Therefore, the non-constant homotopy $G_1|_{\mathbb{A}^1_{V_0}}$ is dominant.




%\textbf{Case 3:} Suppose that $Im(\theta_1) \subset \overline{\phi(V_r)}$ and $Im(\theta_2) \subset \overline{\phi(V_s)}$, for some $r, s$. Also assume that $Im(\theta_2) \nsubseteq \overline{\phi(V_r)}$. There is an irreducible component $V_0$ of $V^\prime$ that maps to $V_r \times_{\mathbb{A}^1_U} V_s$. Since the projection maps and the morphism $V^{\prime} \to V \times_{\mathbb{A}^1_U} V$ are the \'etale maps, so $\overline{Im(\phi \circ \text{pr}_1|_{V_0})} = \overline{\phi(V_r)}$ and $\overline{Im(\phi \circ \text{pr}_2|_{V_0})} = \overline{\phi(V_s)}$. Since $Im(\theta_2) \nsubseteq \overline{\phi(V_r)}$, so
%$$\phi \circ \text{pr}_1|_{V_0} \neq \phi \circ \text{pr}_2|_{V_0}.$$
%So we can assume all the naive $\mathbb{A}^1$-homotopies $G_1|_{\mathbb{A}^1_{V_0}}, \dots, G_l|_{\mathbb{A}^1_{V_0}}$ are non-constant. Since every $G_i|_{\mathbb{A}^1_{V_0}}$ is non-constant, so $Im(G_i|_{\mathbb{A}^1_{V_0}})$ contains a line. Since $X$ is affine, a line is closed in $X$. If for every $i$, $\overline{Im(G_i|_{\mathbb{A}^1_{V_0}})}$ is a line,
% contained in the image of some $\beta: \mathbb{A}^1_k \to X$, 
%then there is a chain of $\mathbb{A}^1_k$'s in $X$ that contains both $Im(\theta_1), Im(\theta_2)$, which would imply that $\theta_1 = \theta_2 \in \mathcal{S}(X)(W)$, by Corollary \ref{not contained corr} and it is a contradiction. Therefore, there is some $j$ such that $\overline{Im(G_j|_{\mathbb{A}^1_{V_0}})}$ properly contains a line. Thus the non-constant homotopy $G_j|_{\mathbb{A}^1_{V_0}}$ is dominant. 
%This is a contradiction, since $X$ is not an $\mathbb{A}^1$-uniruled surface.




%\textbf{Case 4:}  Suppose that $Im(\theta_1) \subset \overline{\phi(V_r)}$ and $Im(\theta_2) \subset \overline{\phi(V_s)}$, for some $r, s$. Also assume that $Im(\theta_1) \nsubseteq \overline{\phi(V_s)}$. This is similar to Case 3. So proceeding as in Case 3, there is an irreducible component $V_0$ of $V^\prime$ and there is some $j$ such that the homotopy $G_j|_{\mathbb{A}^1_{V_0}}$ is non-constant and dominant.

%Since $H$ is non-constant, there is a Nisnevich covering $V^{\prime} \to V \times_{\mathbb{A}^1_k} V$ such that 
%$$\phi \circ \text{pr}_1|_{V^\prime} \neq \phi \circ \text{pr}_2|_{V^\prime} \in X(V^{\prime}),$$
%where $\text{pr}_1, \text{pr}_2: V \times_{\mathbb{A}^1_k} V \to V$ are the projection maps.
%So there is a chain of non-constant naive $\mathbb{A}^1$-homotopies $G_1, \dots, G_l : \mathbb{A}^1_{V^{\prime}} \to X$ such that $G_1(0) = \phi \circ \text{pr}_1|_{V^\prime}$ and $G_l(1) = \phi \circ \text{pr}_2|_{V^\prime}$.
%\end{proof}

%\begin{theorem}
%Suppose $X$ is an affine surface which is not an $\mathbb{A}^1$-uniruled. Then the canonical morphism $\mathcal{S}(X) \to \mathcal{S}^2(X)$ is an isomorphism.
%\end{theorem}
%\begin{proof}
%Suppose $U \in Sm/k$. The canonical morphism $X \to \mathcal{S}(X)$ is an epimorphism. We need to prove that the map
%$$\mathcal{S}(X)(U) \to \mathcal{S}^2(X)(U)$$
%is injective. Suppose $\alpha, \beta \in \mathcal{S}(X)(U)$ such that $\alpha = \beta \in \mathcal{S}^2(X)(U)$. If possible, assume that $\alpha \neq \beta \in \mathcal{S}(X)(U)$. Since $\alpha = \beta \in \mathcal{S}^2(X)(U)$, there is a Nisnevich covering $U^\prime \to U$ such that $\alpha|_{U^\prime} = \beta|_{U^\prime} \in \mathcal{S}^{\text{pre}}(\mathcal{S}(X))(U^\prime)$. Since $\alpha \neq \beta \in \mathcal{S}(X)(U)$, so $\alpha|_{U^\prime} \neq \beta|_{U^\prime} \in \mathcal{S}(X)(U^\prime)$. Thus there is an irreducible component $W$ of $U^\prime$ such that $\alpha|_W \neq \beta|_W \in \mathcal{S}(X)(W)$. Since $\alpha|_{U^\prime} = \beta|_{U^\prime} \in \mathcal{S}^{\text{pre}}(\mathcal{S}(X))(U^\prime)$, the sections $\alpha|_W = \beta|_W \in \mathcal{S}^{2, \text{pre}}(X)(W)$. Therefore there is a chain of non-constant homotopies $H_1, \dots, H_n \in \mathcal{S}(X)(\mathbb{A}^1_W)$ such that $H_1(0) = \alpha|_W$ and $H_n(1) = \beta|_W$. By Proposition \ref{ghost homotopy implies dominant}, there is $V \in Sm/k$ irreducible along with a non-constant homotopy $G: \mathbb{A}^1_V \to X$ which is a dominant morphism. This is a contradiction, since $X$ is not an $\mathbb{A}^1$-uniruled surface. Therefore, $\alpha = \beta \in \mathcal{S}(X)(U)$. Hence the map $\mathcal{S}(X)(U) \to \mathcal{S}^2(X)(U)$ is injective, for every $U \in Sm/k$. Hence the canonical morphism $\mathcal{S}(X) \to \mathcal{S}^2(X)$ is an isomorphism.
%\end{proof}


\begin{remark} (\cite[Remark 2.6]{choudhury2}) \label{ghost homotopy useful remark}
The morphism $\eta: X \to \sS(X)$ is an epimorphism. A homotopy $H \in \sS(X)(\A^1_U)$ (called an $\A^1$-ghost homotopy \cite[Definition 3.2]{bhs}) is given by a Nisnevich covering $f: V \to \A^1_U$ and a morphism $\phi: V \to X$ which makes the following diagram commutative:
\[
\xymatrix{V \ar[d]^f \ar[r]^{\phi} &X \ar[d]^{\eta}\\
\A^1_U \ar[r]^H &\sS(X)}
\]
Since, the sections
$$ \eta \circ \phi \circ \text{pr}_1 = \eta \circ \phi \circ \text{pr}_2 \in \sS(X)(V \times_{\A^1_U} V),$$
where $\text{pr}_1, \text{pr}_2: V \times_{\A^1_U} V \to V$ are the projections; so there is a Nisnevich covering $V^\prime \to  V \times_{\A^1_U} V$ along with a chain of homotopies
$$G_1, \dots, G_m : \A^1_{V^\prime} \to X \text{ such that } G_1(0)=\phi \circ \text{pr}_1|_{V^\prime} \text{ and }  G_m(1)=\phi \circ \text{pr}_2|_{V^\prime}$$
%$$G_{i}(1)=G_{i+1}(0), \text{for every }i \text{ and } G_1(0)=\phi \circ \text{pr}_1|_{V^\prime}, G_m(1)=\phi \circ \text{pr}_2|_{V^\prime}$$
(Notations \ref{notations}). If the sections 
$$\phi \circ \text{pr}_1|_W = \phi \circ \text{pr}_2|_W: W \to X$$
are equal for some Nisnevich covering $W \to V \times_{\A^1_U} V$, then (since $X$ is a Nisnevich sheaf) there is a morphism $\Tilde{\phi}: \A^1_U \to X$ that lifts $\phi$, that is $\Tilde{\phi} \circ f = \phi$. Thus, $\eta \circ \Tilde{\phi} = H$. But this implies that 
$$H(0) = H(1) \in \sS(X)(U)$$ 
(Notations \ref{notations}). Therefore, if we assume that $H \in \sS(X)(\A^1_U)$ is a non-constant homotopy (Definition \ref{non-constant-homotopy-defn}), then for every Nisnevich covering 
$$W \to  V \times_{\A^1_U} V,$$
 the morphisms
$$\phi \circ \text{pr}_1|_{W} \neq \phi \circ \text{pr}_2|_{W}: W \to X.$$
So if $H$ is a non-constant homotopy, we can assume that $G_i: \A^1_{V^\prime} \to X$ is also a non-constant homotopy, for every $i$.
% morphism $\phi: V \to X$
%where $G_i(0)=G \circ i_0$ and $G_i(1) = $
% $\alpha \in \sF(V)$ such that th
\end{remark}
We will quickly recall the following notations and terminologies from \cite{choudhury2}.
\begin{notation}\label{notations} \cite[Notations 2.8]{choudhury2} Let $\sF$ be a Nisnevich sheaf of sets in $Sm/k$, and $U\in Sm/k$.
\begin{enumerate}
\item Let $H \in \sF(\A^1_U)$ be a homotopy and $i_0,i_1: U \to \A^1_U$ be the 0-section and 1-section respectively. We denote the elements $H\circ i_0 \in \sF(U)$ by $H(0)$ and $H\circ i_1 \in \sF(U)$ by $H(1)$
% of $\sF(U)$ induced by the sections $i_0,i_1$ by $H(0):=H\circ i_0$ and $H(1):=H\circ i_1$ 
respectively.
\item A chain of homotopies $H_1,\dots,H_n \in \sF(\A^1_U)$ means a collection of homotopies $H_i\in\sF(\A^1_U)$ such that $H_i(1)=H_{i+1}(0),$ for every $i$.
\item For a morphism $f: X \to Y$ of schemes, by $\text{Im}(f)$ we mean the set theoretic image of $f$, that is the set $f(X)$, which is a subset of $Y$. If the set $\text{Im}(f)$ consists more than one element, then we say $f: X \to Y$ is a non-constant morphism. Note that this is different from non-constant homotopy $H \in \sF(\A^1_U)$, which means that $H(0) \neq H(1) \in \sF(U)$
\item For a scheme $X$ over $k$, by ``there is a line in $X$" (or $X$ contains a line) we mean there exists a non-constant morphism $\gamma: \A^1_k \to X$. For a subset $Y \subset X$, by ``$Y$ contains a line" means that there is a line in $X$, given by $\gamma: \A^1_k \to X$ such that $Im(\gamma) \subset Y$.

\end{enumerate}
\end{notation}
\begin{definition}\label{non-constant-homotopy-defn} (\cite[Definition 2.4]{choudhury2}, see also \cite[Example 2.5]{choudhury2})
	Suppose, 
	$\sF$ is a Nisnevich sheaf of sets on $Sm/k$ and 
	$X, U \in Sm/k$. A homotopy $H \in \sF(\A^1_U)$ is said to be a non-constant homotopy, if the sections
	$$H (0) \neq H (1) \in \sF(U),$$
	
\end{definition}
\begin{remark} \label{line remark}
If there is a non-constant homotopy $H: \A^1_U \to X$ (Definition \ref{non-constant-homotopy-defn}), for some $X, U \in Sm/k$, $U$ is irreducible, then there is a line in $X$ (Notations \ref{notations}). Indeed, since $H: \A^1_U \to X$ is a non-constant, homotopy so $H(0) \neq H(1): U \to X$. Thus there is $x \in U(k)$ such that 
$$H(0)(x) \neq H(1)(x).$$
So the composition of morphisms
$$\gamma: \A^1_k \xrightarrow{id, x} \A^1_U \xrightarrow{H} X$$
is a line in $X$ (Notations \ref{notations}).
\end{remark}
\begin{theorem} \label{non-uniruled}
Let $k$ be an algebraically closed field of characteristic zero. Suppose that $X \in Sm/k$ is an affine surface, which is not $\mathbb{A}^1$-uniruled. Then the canonical morphism $\mathcal{S}(X) \to \mathcal{S}^2(X)$ is an isomorphism.
\end{theorem}
\begin{proof}
Suppose $U \in Sm/k$. The canonical morphism $\mathcal{S}(X) \to \mathcal{S}^2(X)$ is an epimorphism of Nisnevich sheaves (Remark \ref{ghost homotopy useful remark}). We will prove that the map
$$\mathcal{S}(X)(U) \to \mathcal{S}^2(X)(U)$$
is injective. Suppose $\alpha, \beta \in \mathcal{S}(X)(U)$ such that $\alpha = \beta \in \mathcal{S}^2(X)(U)$. If possible, assume that $\alpha \neq \beta \in \mathcal{S}(X)(U)$. Since $\alpha = \beta \in \mathcal{S}^2(X)(U)$, there is a Nisnevich covering $U^\prime \to U$ such that $\alpha|_{U^\prime} = \beta|_{U^\prime} \in \mathcal{S}^{\text{pre}}(\mathcal{S}(X))(U^\prime)$. Since $\alpha \neq \beta \in \mathcal{S}(X)(U)$, so $\alpha|_{U^\prime} \neq \beta|_{U^\prime} \in \mathcal{S}(X)(U^\prime)$. 
Thus there is an irreducible component $W$ of $U^\prime$ such that $\alpha|_W \neq \beta|_W \in \mathcal{S}(X)(W)$. Since, 
$$\alpha|_{U^\prime} = \beta|_{U^\prime} \in \mathcal{S}^{\text{pre}}(\mathcal{S}(X))(U^\prime),$$
so the sections $\alpha|_W = \beta|_W \in \mathcal{S}^{\text{pre}}(\mathcal{S}(X))(W)$. Therefore there is a non-constant homotopy $H \in \mathcal{S}(X)(\mathbb{A}^1_W)$ (Definition \ref{non-constant-homotopy-defn}) such that $H(0) = \alpha|_W$. 
%Since the canonical morphism $\eta: X \to \mathcal{S}(X)$ is an epimorphism, so there is a Nisnevich covering $f: V \to \mathbb{A}^1_U$ and a morphism $\phi: V \to X$ which makes the following diagram commutative:
 %$$\begin{tikzcd}
 %V \ar[r, "\phi"] \ar[d, "f"] &X \ar[d, "\eta"] \\
%\mathbb{A}^1_U \ar[r,"H"] 
%&\mathcal{S}(X)
%\end{tikzcd}$$ 
%Consider the following commutative diagram for the $0$-section of $H$:
%$$\begin{tikzcd}[column sep=50pt, row sep=50 pt]
%W^\prime \ar[r, "\theta^\prime"] \ar[d, "f^\prime"] &V \ar[r, "\phi"] \ar[d, "f"] &X \ar[d, "\eta"] \\
 %W \ar[r, "i_0"]  
%& \mathbb{A}^1_{W} \ar[r,"H"] 
%&\mathcal{S}(X)
%\end{tikzcd}$$ 
%Here $i_0:W \to \mathbb{A}^1_W$ is the $0$-section and the left most square is a pullback square. The morphism $f^\prime: W^\prime %\to W$ is a Nisnevich covering. From the commutativity of the diagram we have, 
%$$H(0)|_{W^\prime} = \phi \circ \theta^\prime \in \mathcal{S}(X)(W^\prime).$$
%Consider the following commutative diagram for the $1$-section of $H$:
%$$\begin{tikzcd}[column sep=50pt, row sep=50 pt]
%W^{\prime\prime} \ar[r, "\theta^{\prime\prime}"] \ar[d, "f^{\prime\prime}"] &V \ar[r, "\phi"] \ar[d, "f"] &X \ar[d, "\eta"] \\
 %W \ar[r, "i_1"]  
%& \mathbb{A}^1_{W} \ar[r,"H"] 
%&\mathcal{S}(X)
%\end{tikzcd}$$ 
%Here $i_1:W \to \mathbb{A}^1_W$ is the $1$-section and the left most square is a pullback square. The morphism $f^{\prime \prime}: W^{\prime \prime }\to W$ is a Nisnevich covering. From the commutativity of the diagram we have, 
%$$H(1)|_{W^{\prime \prime}} = \phi \circ \theta^{\prime \prime} \in \mathcal{S}(X)(W^{\prime \prime}).$$
%Consider the following pullback diagram:
%$$\begin{tikzcd}
 %\Tilde{W} \ar[r, "\psi^{\prime}"] \ar[d, "\psi^{\prime\prime}"] &W^\prime \ar[d, "f^\prime"] \\
%W^{\prime\prime} \ar[r,"f^{\prime \prime}"] 
%&W
%\end{tikzcd}$$ 
%Here $\Tilde{W}$ is a common refinement of $W^\prime$ and $W^{\prime \prime}$. The morphism $f^\prime \circ \psi^\prime=f^{\prime \prime} \circ \psi^{\prime \prime}: \Tilde{W} \to W$ is a Nisnevich covering and we have
%$$H(0)|_{\Tilde{W}} = \phi \circ \theta^{\prime} \circ \psi^\prime \in \mathcal{S}(X)(\Tilde{W}) \text{ and } H(1)|_{\Tilde{W}} = \phi \circ \theta^{\prime \prime} \circ \psi^{\prime \prime} \in \mathcal{S}(X)(\Tilde{W}).$$
%Thus we are situation of Proposition \ref{ghost homotopy implies dominant 1}. 
In the next Proposition (Proposition \ref{ghost homotopy implies dominant 1}) we will prove that there is a dominant morphism $G: \mathbb{A}^1_Y \to X$, with $Y$ irreducible such that 
$$G(0, -) \neq G(1, -): Y \to X.$$ 
We can take here $dim(Y) = 1$. Thus we arrive at a contradiction, since $X$ is not an $\mathbb{A}^1$-uniruled surface. Hence the canonical map
$$\mathcal{S}(X) \to \mathcal{S}^2(X)$$
is an isomorphism.
Therefore, it remains to prove that if there is a non-constant homotopy $H \in \sS(X)(\A^1_U)$ (Definition \ref{non-constant-homotopy-defn}), then $X$ admits a dominant morphism $G: \A^1_W \to X$, which is also a non-constant homotopy (Definition \ref{non-constant-homotopy-defn}), for some $W \in Sm/k$ irreducible.
%and $H_n(1) = \beta|_W$. Consider one such non-constant homotopy, say $H_1 \in \mathcal{S}(X)(\mathbb{A}^1_W)$. 

%We will construct a non-constant homotopy $H \in \mathcal{S}(X)(\math)$
 %By Proposition \ref{ghost homotopy implies dominant}, there is $V \in Sm/k$ irreducible along with a non-constant homotopy $G: \mathbb{A}^1_V \to X$ which is a dominant morphism. This is a contradiction, since $X$ is not an $\mathbb{A}^1$-uniruled surface. Therefore, $\alpha = \beta \in \mathcal{S}(X)(U)$. Hence the map $\mathcal{S}(X)(U) \to \mathcal{S}^2(X)(U)$ is injective, for every $U \in Sm/k$. Hence the canonical morphism $\mathcal{S}(X) \to \mathcal{S}^2(X)$ is an isomorphism.
\end{proof}
The next proposition is similar to \cite[Proposition 2.9]{choudhury2}. 

%Let $X \in Sm/k$ be an 
%non $\mathbb{A}^1$-uniruled 
%affine surface and $U \in Sm/k$ be irreducible. 
%Then there is $W \in Sm/k$ irreducible al. 
%Suppose that there is a non-constant homotopy $H \in \mathcal{S}(X)(\mathbb{A}^1_U)$. The canonical morphism $\eta: X \to \mathcal{S}(X)$ is an epimorphism, so there is a Nisnevich covering $f: V \to \mathbb{A}^1_U$ and a morphism $\phi: V \to X$ which makes the following diagram commutative:
%\[
%\xymatrix{V \ar[d]^f \ar[r]^{\phi} &X \ar[d]^{\eta}\\
%\A^1_U \ar[r]^H &\sS(X)}
%\]
%Suppose that the $0$-section and $1$-section of $H$ lift to $V$ upto a Nisnevich covering. This means that there is a Nisnevich covering $\Tilde{U} \to U$ and there are morphisms $\theta, \psi: \Tilde{U} \to V$ such that the following diagrams are commutative:
 
%\[\xymatrixcolsep{3pc}
%\xymatrix{\tilde{U}\ar[d] \ar[r]^{\theta} &V \ar[d]^f \ar[r]^{\phi} &X \ar[d]^{\eta} &\tilde{U}\ar[d] \ar[r]^{\psi} &V \ar[d]^f \ar[r]^{\phi} &X \ar[d]^{\eta}\\
%U \ar[r]^{i_0} &\A^1_U \ar[r]^H &\sS(X) &U \ar[r]^{i_1} &\A^1_U \ar[r]^H &\sS(X)} 
%\]

 
%Here $i_0, i_1:U \to \mathbb{A}^1_U$ are the $0$-section and $1$-section respectively. If $H$ is any non-constant homotopy on $\mathbb{A}^1_U$, then we can reduce to this situation. By taking the pullback of $f$ with respect to the $0$-section and the $1$-section, there are Nisnevich covering $U^\prime \to U$ and $U^{\prime \prime} \to U$ such that $H(0)|_{U^\prime}$ and $H(1)|_{U^{\prime \prime}}$ lift to $V$. Now taking a common refinement $\Tilde{U}$ of $U^\prime$ and $U^{\prime \prime}$, both $H(0)|_{\Tilde{U}}$ and $H(1)_{\Tilde{U}}$ lift to $V$.
\begin{proposition} \label{ghost homotopy implies dominant 1}
Let $X \in Sm/k$ be an 
%non $\mathbb{A}^1$-uniruled 
affine surface and $U \in Sm/k$ be irreducible. 
%Then there is $W \in Sm/k$ irreducible al. 
Suppose that $H \in \mathcal{S}(X)(\mathbb{A}^1_U)$ is a non-constant homotopy (Definition \ref{non-constant-homotopy-defn}).
% as above i.e. there is a Nisnevich covering $\Tilde{U} \to U$ such that both $H(0)|_{\Tilde{U}}$ and $H(1)|_{\Tilde{U}}$ lift to $V$. 
Then there is $W \in Sm/k$ irreducible, along with a dominant morphism $G: \mathbb{A}^1_W \to X$ such that $G(0,-) \neq G(1, -): W \to X$ (i.e $G$ is also a non-constant homotopy in the sense of Definition \ref{non-constant-homotopy-defn}). 
%for some $W \in$ 
\end{proposition}
\begin{proof}
The homotopy $H \in \sS(X)(\A^1_U)$ is non-constant. Following Remark \ref{ghost homotopy useful remark}, 
%the morphisms $\phi \circ \text{pr}_1|_{V^\prime}$ and $\phi \circ \text{pr}_2|_{V^\prime}$
% (here $\text{pr}_1, \text{pr}_2: V \times_{\mathbb{A}^1_U} V \to V$ are the projection maps) 
%are not same in $X(V^\prime)$, for any Nisnevich covering $V^\prime \to V \times_{\mathbb{A}^1_U} V$. However since $H$ lifts to $X$ over $V$, the sections $\phi \circ \text{pr}_1 = \phi \circ \text{pr}_2 \in \mathcal{S}(X)(V \times_{\mathbb{A}^1_U} V)$.
%, but they , since $H$ is non-co but they. 
%Thus 
there is a Nisnevich covering $V^\prime \to V \times_{\mathbb{A}^1_U} V$ and a chain of non-constant naive $\mathbb{A}^1$-homotopies (Notations \ref{notations}) $G_1, \dots, G_m : \mathbb{A}^1_{V^{\prime}} \to X$ such that 
$$G_1(0) = \phi \circ \text{pr}_1|_{V^\prime} \text{ and } G_m(1) = \phi \circ \text{pr}_2|_{V^\prime}.$$
We claim that $H(0)$ and $H(1)$ (Notations \ref{notations}) lift to $V$, over some Nisnevich covering $\Tilde{U} \to U$. This means that there is a Nisnevich covering $\Tilde{U} \to U$ and there are morphisms $\theta, \psi: \Tilde{U} \to V$ such that the following diagrams are commutative:
 
\[\xymatrixcolsep{3pc}
\xymatrix{\tilde{U}\ar[d] \ar[r]^{\theta} &V \ar[d]^f \ar[r]^{\phi} &X \ar[d]^{\eta} &\tilde{U}\ar[d] \ar[r]^{\psi} &V \ar[d]^f \ar[r]^{\phi} &X \ar[d]^{\eta}\\
U \ar[r]^{i_0} &\A^1_U \ar[r]^H &\sS(X) &U \ar[r]^{i_1} &\A^1_U \ar[r]^H &\sS(X)} 
\]

 
%Here $i_0, i_1:U \to \mathbb{A}^1_U$ are the $0$-section and $1$-section respectively. If $H$ is any non-constant homotopy on $\mathbb{A}^1_U$, then we can reduce to this situation. 
Indeed, by taking the pullback of $f$ with respect to $i_0, i_1: U \to \A^1_U$,
% the $0$-section and the $1$-section, 
there are Nisnevich covering $U^\prime \to U$ and $U^{\prime \prime} \to U$ such that $H(0)|_{U^\prime}$ and $H(1)|_{U^{\prime \prime}}$ lift to $V$. Now taking a common refinement $\Tilde{U}$ of $U^\prime$ and $U^{\prime \prime}$, both the sections $H(0)|_{\Tilde{U}}$ and $H(1)_{\Tilde{U}}$ lift to $V$.
%\begin{proposition} \label{ghost homotopy implies dominant 1}


%the $0$-section and $1$-section of $H$ lift to $V$ over a Nisnevich covering $\Tilde{U} \to U$. 
So,
$$H(0)|_{\Tilde{U}} =\eta \circ \phi \circ \theta \in \mathcal{S}(X)(\Tilde{U}) \text{ and } H(1)|_{\Tilde{U}} = \eta \circ \phi \circ \psi \in \mathcal{S}(X)(\Tilde{U}).$$
Since $H$ is a non-constant homotopy, so 
$$\eta \circ \phi \circ \theta \neq \eta \circ \phi \circ \psi \in \mathcal{S}(X)(\Tilde{U}).$$
Thus there is an irreducible component (which is also a connected component) $W$ of $\Tilde{U}$ such that 
$$(\eta \circ \phi \circ \theta)|_W \neq (\eta \circ \phi \circ \psi)|_W \in \mathcal{S}(X)(W).$$
Suppose, $(\phi \circ \theta)|_W = \theta_1 \in X(W)$ and $(\phi \circ \psi)|_W = \theta_2 \in X(W)$. The sections 
$$\eta \circ \theta_1 \neq \eta \circ \theta_2 \in \mathcal{S}(X)(W).$$
So by Corollary \ref{not contained corr}, there exists no chain of affine lines $\gamma_1, \dots, \gamma_n: \A^1_k \to X$ (Notations \ref{notations}) in $X$ such that $\text{Im}(\theta_1) \subset \text{Im}(\gamma_i)$ and $\text{Im}(\theta_2) \subset \text{Im}(\gamma_j)$, for some $i,j$.

%Suppose, $(\phi \circ \theta^{\prime} \circ \psi^\prime)|_W = \theta_1 \in X(W)$ and $(\phi \circ \theta^{\prime \prime} \circ \psi^{\prime \prime})|_W = \theta_2 \in X(W)$. The sections $\theta_1 \neq \theta_2 \in \mathcal{S}(X)(W)$.
Suppose, $V  = \coprod_{i =1}^n V_i$, where $V_i$'s are the irreducible components of $V$, which are also the connected components of $V$. Since $W$ is irreducible, $\text{Im}(\theta_1) \subset \overline{\phi(V_i)}$ and $\text{Im}(\theta_2) \subset \overline{\phi(V_j)}$, for some $i, j$ (Notations \ref{notations}(3)). 
%If there are non-constant homotopies 
We complete the proof considering the following cases.

\

\textbf{Case 1:} \textbf{Suppose, $\text{Im}(\theta_1), \text{Im}(\theta_2) \subseteq \overline{\phi(V_p)}$, for every $p$.}

  Since $H$ is a non-constant homotopy (Definition \ref{non-constant-homotopy-defn}), there is an irreducible component $V_0$ of $V^{\prime}$ such that (Remark \ref{ghost homotopy useful remark})
$$\phi \circ \text{pr}_1|_{V_0} \neq \phi \circ \text{pr}_2|_{V_0}: V_0 \to X.$$
We can assume that the naive $\mathbb{A}^1$-homotopy $G_1|_{\mathbb{A}^1_{V_0}}$ is non-constant.
Suppose, $V_0$ maps to $V_r \times_{\mathbb{A}^1_U} V_s$, for some $r, s \leq n$. Since the projection maps and the morphism $V^{\prime} \to V \times_{\mathbb{A}^1_U} V$ are the \'etale maps, so $\overline{\text{Im}(G_1|_{\mathbb{A}^1_{V_0}}(0))} = \overline{\phi(V_r)}$. 
%and $\overline{Im(G_l|_{\mathbb{A}^1_{V_0}}(1))} = \overline{\phi(V_s)}$. 
Since $G_1|_{\mathbb{A}^1_{V_0}}$ is a non-constant homotopy (Notations \ref{notations}), so $\text{Im}(G_1|_{\mathbb{A}^1_{V_0}})$ contains a line, say $\gamma: \A^1_k \to X$ (Remark \ref{line remark}). Since $X$ is affine, so $\text{Im}(\gamma)$ is closed in $X$. If $\overline{\text{Im}(G_1|_{\mathbb{A}^1_{V_0}})} = \text{Im}(\gamma)$, then by our assumption in this case, $\text{Im}(\theta_1), \text{Im}(\theta_2) \subset \text{Im}(\gamma)$, since $\overline{\text{Im}(G_1|_{\mathbb{A}^1_{V_0}}(0))} = \overline{\phi(V_r)}$. This implies that $\theta_1 = \theta_2 \in \mathcal{S}(X)(W)$, by Corollary \ref{not contained corr} and it is a contradiction. Therefore, the non-constant homotopy $G_1|_{\mathbb{A}^1_{V_0}}$ is dominant.
%  is contained in the image of some $\beta: \mathbb{A}^1_k \to X$, then there is a chain of $\mathbb{A}^1_k$'s in $X$ that contains both $Im(\Tilde{\theta}), Im(\Tilde{\gamma})$, which would imply that $\Tilde{\theta} = \Tilde{\gamma} \in \mathcal{S}(X)(U)$, by Corollary \ref{not contained corr} and it is a contradiction. Therefore, there is some $j$ such that $G_j|_{\mathbb{A}^1_{V_0}}$ is non-constant and dominant. This is a contradiction, since $X$ is not an $\mathbb{A}^1$-uniruled surface.

\


\textbf{Case 2:} \textbf{Suppose there is some $p$ such that $\text{Im}(\theta_1) \nsubseteq \overline{\phi(V_p)}$.}

Case 2 consists two subcases:

%. Also assume that there is $\mathbf{r < n}$ such that 
%$$\text{Im}(\theta_1), \text{Im}(\theta_2) \subset \overline{\phi(V_1)}, \dots, \overline{\phi(V_r)}$$ 
%(renumbering if necessary). Also assume that 
%$$\text{Im}(\theta_1)  \nsubseteq \overline{\phi(V_p)} \text{ and } \text{Im}(\theta_2)  \nsubseteq \overline{\phi(V_p)}, \text{ for every } p > r.$$
%, for every $p> r$.
%} 
     \textbf{Subcase 2(a):} \textbf{Renumbering if necessary, assume that there is some $\mathbf{r < n}$ such that 
$$\text{Im}(\theta_1), \text{Im}(\theta_2) \subseteq \overline{\phi(V_1)}, \dots, \overline{\phi(V_r)}$$ 
and $\text{Im}(\theta_1)  \nsubseteq \overline{\phi(V_p)}$, for some $p >r$.
}

There is an irreducible component $V_0$ of $V^\prime$ that maps to $V_r \times_{\mathbb{A}^1_U} V_p$. Since the projection maps and the morphism $V^{\prime} \to V \times_{\mathbb{A}^1_U} V$ are the \'etale maps, so 
$$\overline{\text{Im}(\phi \circ \text{pr}_1|_{V_0})} = \overline{\phi(V_r)} \text{ and } \overline{\text{Im}(\phi \circ \text{pr}_2|_{V_0})} = \overline{\phi(V_p)}.$$
Since $\text{Im}(\theta_1) \nsubseteq \overline{\phi(V_p)}$, so
%$Im(\Tilde{\gamma}) \nsubseteq \overline{\phi(V_r)}$, so
$$\phi \circ \text{pr}_1|_{V_0} \neq \phi \circ \text{pr}_2|_{V_0}: V_0 \to X.$$
Thus we can assume that the naive $\mathbb{A}^1$-homotopy $G_1|_{\mathbb{A}^1_{V_0}}$ is a non-constant homotopy. So $\text{Im}(G_1|_{\mathbb{A}^1_{V_0}})$ contains a line $\gamma: \A^1_k \to X$ (Remark \ref{line remark}). Since $X$ is affine, so $\text{Im}(\gamma)$ is closed in $X$.  If $\overline{\text{Im}(G_1|_{\mathbb{A}^1_{V_0}})} = \text{Im}(\gamma)$, then by our assumption $\text{Im}(\theta_1), \text{Im}(\theta_2) \subset \text{Im}(\gamma)$, since $G_1|_{\mathbb{A}^1_{V_0}}(0) = \phi \circ \text{pr}_1|_{V_0}$. This implies that $\theta_1 = \theta_2 \in \mathcal{S}(X)(W)$, by Corollary \ref{not contained corr} and it is a contradiction. Therefore, the non-constant homotopy $G_1|_{\mathbb{A}^1_{V_0}}$ is dominant.

\
\textbf{Subcase 2(b):} \textbf{There exist $r,s$ such that 
$$\text{Im}(\theta_1) \subseteq \overline{\phi(V_r)}, \ \text{Im}(\theta_2) \subseteq \overline{\phi(V_s)} \text{ and } \text{Im}(\theta_2) \nsubseteq \overline{\phi(V_r)}.$$
}


%\textbf{Case 3:} Suppose that $Im(\theta_1) \subset \overline{\phi(V_r)}$ and $Im(\theta_2) \subset \overline{\phi(V_s)}$, for some $r, s$. Also assume that $Im(\theta_2) \nsubseteq \overline{\phi(V_r)}$. 
There is an irreducible component $V_0$ of $V^\prime$ that maps to $V_r \times_{\mathbb{A}^1_U} V_s$. Since the projection maps and the morphism $V^{\prime} \to V \times_{\mathbb{A}^1_U} V$ are the \'etale maps, so 
$$\overline{\text{Im}(\phi \circ \text{pr}_1|_{V_0})} = \overline{\phi(V_r)} \text{ and } \overline{\text{Im}(\phi \circ \text{pr}_2|_{V_0})} = \overline{\phi(V_s)}.$$
Since $\text{Im}(\theta_2) \nsubseteq \overline{\phi(V_r)}$, so
$$\phi \circ \text{pr}_1|_{V_0} \neq \phi \circ \text{pr}_2|_{V_0}.$$
We can assume all the naive $\mathbb{A}^1$-homotopies $G_1|_{\mathbb{A}^1_{V_0}}, \dots, G_m|_{\mathbb{A}^1_{V_0}}$ are non-constant (Definition \ref{non-constant-homotopy-defn}). Since every $G_i|_{\mathbb{A}^1_{V_0}}$ is non-constant, so $\text{Im}(G_i|_{\mathbb{A}^1_{V_0}})$ contains a line, say $\gamma_i: \A^1_k \to X$ (Remark \ref{line remark}). Since $X$ is affine, $\text{Im}(\gamma)$ is closed in $X$. If for every $i$, $\overline{\text{Im}(G_i|_{\mathbb{A}^1_{V_0}})} = \text{Im}(\gamma_i)$, 
%is a line,
% contained in the image of some $\beta: \mathbb{A}^1_k \to X$, 
then there is a chain of affine lines $\gamma_1, \dots, \gamma_m: \A^1_k \to X$ such that (Notations \ref{notations})
$$\text{Im}(\theta_1) \subset \text{Im}(\gamma_1) \text{ and } \text{Im}(\theta_2) \subset \text{Im}(\gamma_m),$$
% that contains both $Im(\theta_1), Im(\theta_2)$, 
since
$$G_1|_{\mathbb{A}^1_{V_0}}(0) = \phi \circ \text{pr}_1|_{V_0} \text{ and } G_m|_{\mathbb{A}^1_{V_0}}(1) = \phi \circ \text{pr}_2|_{V_0}$$ 
and this would imply that $\theta_1 = \theta_2 \in \mathcal{S}(X)(W)$, by Corollary \ref{not contained corr}. It is a contradiction. Therefore, there is some $j$ such that $\overline{\text{Im}(G_j|_{\mathbb{A}^1_{V_0}})}$ properly contains a line. Thus the non-constant homotopy $G_j|_{\mathbb{A}^1_{V_0}}$ is dominant. 
%This is a contradiction, since $X$ is not an $\mathbb{A}^1$-uniruled surface.

\

\textbf{Case 3:} \textbf{Suppose there is some $q$ such that $\text{Im}(\theta_2) \nsubseteq \overline{\phi(V_q)}$.}

Case 3 is similar to Case 2. Hence, in this Case also, there exists a non-constant homotopy $G: \A^1_Y \to X$ (Definition \ref{non-constant-homotopy-defn}), for some $Y \in Sm/k$ irreducible such that $G$ is a dominant morphism.
%, where $y \in Sm/k$ irre %\ref{\textbf{Case 2:} \textbf{Suppose there is some $p$ such that $\text{Im}(\theta_1) \nsubseteq \overline{\phi(V_p)}$.}})
%\textbf{Case 4:}  Suppose that $Im(\theta_1) \subset \overline{\phi(V_r)}$ and $Im(\theta_2) \subset \overline{\phi(V_s)}$, for some $r, s$. Also assume that $Im(\theta_1) \nsubseteq \overline{\phi(V_s)}$. This is similar to Case 3. So proceeding as in Case 3, there is an irreducible component $V_0$ of $V^\prime$ and there is some $j$ such that the homotopy $G_j|_{\mathbb{A}^1_{V_0}}$ is non-constant and dominant.

Therefore, the Proposition is proved.
%Since $H$ is non-constant, there is a Nisnevich covering $V^{\prime} \to V \times_{\mathbb{A}^1_k} V$ such that 
%$$\phi \circ \text{pr}_1|_{V^\prime} \neq \phi \circ \text{pr}_2|_{V^\prime} \in X(V^{\prime}),$$
%where $\text{pr}_1, \text{pr}_2: V \times_{\mathbb{A}^1_k} V \to V$ are the projection maps.
%So there is a chain of non-constant naive $\mathbb{A}^1$-homotopies $G_1, \dots, G_l : \mathbb{A}^1_{V^{\prime}} \to X$ such that $G_1(0) = \phi \circ \text{pr}_1|_{V^\prime}$ and $G_l(1) = \phi \circ \text{pr}_2|_{V^\prime}$.
\end{proof}
%\begin{definition}
%An affine $k$-variety $X$ is said to be \textbf{$\mathbb{A}^1$-uniruled or log-uniruled} if there is a dominant generically finite morphism $H: \mathbb{A}^1_k \times_k Y \to X$ for some $k$-variety $Y$. 
%\end{definition}
%\begin{theorem} \label{non-uniruled}
%Let $k$ be an algebraically closed field of characteristic zero. Suppose that $X \in Sm/k$ is an affine surface, which is not an $\mathbb{A}^1$-uniruled. Then the canonical morphism $\mathcal{S}(X) \to \mathcal{S}^2(X)$ is an isomorphism.
%\end{theorem}
%\begin{proof}
%Suppose $U \in Sm/k$. The canonical morphism $\mathcal{S}(X) \to \mathcal{S}^2(X)$ is an epimorphism of Nisnevich sheaves (Remark \ref{}). We will prove that the map
%$$\mathcal{S}(X)(U) \to \mathcal{S}^2(X)(U)$$
%is injective. Suppose $\alpha, \beta \in \mathcal{S}(X)(U)$ such that $\alpha = \beta \in \mathcal{S}^2(X)(U)$. If possible, assume that $\alpha \neq \beta \in \mathcal{S}(X)(U)$. Since $\alpha = \beta \in \mathcal{S}^2(X)(U)$, there is a Nisnevich covering $U^\prime \to U$ such that $\alpha|_{U^\prime} = \beta|_{U^\prime} \in \mathcal{S}^{\text{pre}}(\mathcal{S}(X))(U^\prime)$. Since $\alpha \neq \beta \in \mathcal{S}(X)(U)$, so $\alpha|_{U^\prime} \neq \beta|_{U^\prime} \in \mathcal{S}(X)(U^\prime)$. 
%Thus there is an irreducible component $W$ of $U^\prime$ such that $\alpha|_W \neq \beta|_W \in \mathcal{S}(X)(W)$. Since, 
%$$\alpha|_{U^\prime} = \beta|_{U^\prime} \in \mathcal{S}^{\text{pre}}(\mathcal{S}(X))(U^\prime),$$
%so the sections $\alpha|_W = \beta|_W \in \mathcal{S}^{\text{pre}}(\mathcal{S}(X))(W)$. Therefore there is a non-constant homotopy $H \in \mathcal{S}(X)(\mathbb{A}^1_W)$ such that $H(0) = \alpha|_W$. 
%Since the canonical morphism $\eta: X \to \mathcal{S}(X)$ is an epimorphism, so there is a Nisnevich covering $f: V \to \mathbb{A}^1_U$ and a morphism $\phi: V \to X$ which makes the following diagram commutative:
 %$$\begin{tikzcd}
 %V \ar[r, "\phi"] \ar[d, "f"] &X \ar[d, "\eta"] \\
%\mathbb{A}^1_U \ar[r,"H"] 
%&\mathcal{S}(X)
%\end{tikzcd}$$ 
%Consider the following commutative diagram for the $0$-section of $H$:
%$$\begin{tikzcd}[column sep=50pt, row sep=50 pt]
%W^\prime \ar[r, "\theta^\prime"] \ar[d, "f^\prime"] &V \ar[r, "\phi"] \ar[d, "f"] &X \ar[d, "\eta"] \\
 %W \ar[r, "i_0"]  
%& \mathbb{A}^1_{W} \ar[r,"H"] 
%&\mathcal{S}(X)
%\end{tikzcd}$$ 
%Here $i_0:W \to \mathbb{A}^1_W$ is the $0$-section and the left most square is a pullback square. The morphism $f^\prime: W^\prime %\to W$ is a Nisnevich covering. From the commutativity of the diagram we have, 
%$$H(0)|_{W^\prime} = \phi \circ \theta^\prime \in \mathcal{S}(X)(W^\prime).$$
%Consider the following commutative diagram for the $1$-section of $H$:
%$$\begin{tikzcd}[column sep=50pt, row sep=50 pt]
%W^{\prime\prime} \ar[r, "\theta^{\prime\prime}"] \ar[d, "f^{\prime\prime}"] &V \ar[r, "\phi"] \ar[d, "f"] &X \ar[d, "\eta"] \\
 %W \ar[r, "i_1"]  
%& \mathbb{A}^1_{W} \ar[r,"H"] 
%&\mathcal{S}(X)
%\end{tikzcd}$$ 
%Here $i_1:W \to \mathbb{A}^1_W$ is the $1$-section and the left most square is a pullback square. The morphism $f^{\prime \prime}: W^{\prime \prime }\to W$ is a Nisnevich covering. From the commutativity of the diagram we have, 
%$$H(1)|_{W^{\prime \prime}} = \phi \circ \theta^{\prime \prime} \in \mathcal{S}(X)(W^{\prime \prime}).$$
%Consider the following pullback diagram:
%$$\begin{tikzcd}
 %\Tilde{W} \ar[r, "\psi^{\prime}"] \ar[d, "\psi^{\prime\prime}"] &W^\prime \ar[d, "f^\prime"] \\
%W^{\prime\prime} \ar[r,"f^{\prime \prime}"] 
%&W
%\end{tikzcd}$$ 
%Here $\Tilde{W}$ is a common refinement of $W^\prime$ and $W^{\prime \prime}$. The morphism $f^\prime \circ \psi^\prime=f^{\prime \prime} \circ \psi^{\prime \prime}: \Tilde{W} \to W$ is a Nisnevich covering and we have
%$$H(0)|_{\Tilde{W}} = \phi \circ \theta^{\prime} \circ \psi^\prime \in \mathcal{S}(X)(\Tilde{W}) \text{ and } H(1)|_{\Tilde{W}} = \phi \circ \theta^{\prime \prime} \circ \psi^{\prime \prime} \in \mathcal{S}(X)(\Tilde{W}).$$
%Thus we are situation of Proposition \ref{ghost homotopy implies dominant 1}. 
%By Proposition \ref{ghost homotopy implies dominant 1}, there is a dominant morphism $G: \mathbb{A}^1_Y \to X$, with $Y$ irreducible such that $G(0, -) \neq G(1, -): Y \to X$. We can take $dim(Y) = 1$ here. Thus we arrive at a contradiction, since $X$ is not an $\mathbb{A}^1$-uniruled surface. Hence the canonical map
%$$\mathcal{S}(X) \to \mathcal{S}^2(X)$$
%is an isomorphism.
%and $H_n(1) = \beta|_W$. Consider one such non-constant homotopy, say $H_1 \in \mathcal{S}(X)(\mathbb{A}^1_W)$. 

%We will construct a non-constant homotopy $H \in \mathcal{S}(X)(\math)$
 %By Proposition \ref{ghost homotopy implies dominant}, there is $V \in Sm/k$ irreducible along with a non-constant homotopy $G: \mathbb{A}^1_V \to X$ which is a dominant morphism. This is a contradiction, since $X$ is not an $\mathbb{A}^1$-uniruled surface. Therefore, $\alpha = \beta \in \mathcal{S}(X)(U)$. Hence the map $\mathcal{S}(X)(U) \to \mathcal{S}^2(X)(U)$ is injective, for every $U \in Sm/k$. Hence the canonical morphism $\mathcal{S}(X) \to \mathcal{S}^2(X)$ is an isomorphism.
%\end{proof}

\begin{remark}
If $X$ is not an $\mathbb{A}^1$-uniruled surface, then in each cases $\overline{\text{Im}(G_i)}$ has dimension $1$. Since the image of an affine line is closed in $X$, so $\overline{\text{Im}(G_i)}$ is the image of an affine line. Thus $\theta_1 = \theta_2 \in \mathcal{S}(X)(W)$.
\end{remark}



%\begin{theorem}
%The canonical surjection
%$$\mathcal{S}(X)(\Spec\ k) \to \mathcal{S}^2(X)(\Spec\ k)$$
%is a bijection.
%\end{theorem}
%\begin{proof}
%If possible assume that, there are $k$-points $x, y$ with $x \neq y \in \mathcal{S}(X)(\Spec\ k)$, but $x = y \in \mathcal{S}^2(X)(\Spec\ k)$. Thus there is a chain of non-constant naive $\mathbb{A}^1$-homotopies $H_1, \dots, H_m \in \mathcal{S}(X)(\mathbb{A}^1_k)$ such that $H_1(0) = x$ and $H_m(1) = y$. We can assume all $H_i$'s are non-constant. We can also assume there is non-constant $H \in \mathcal{S}(X)(\mathbb{A}^1_k)$ such that $H(0) = x$ and $H(1) = y$.

%Since the canonical morphism $\eta: X \to \mathcal{S}(X)$ is an epimorphism, so there is a Nisnevich covering $f: V \to \mathbb{A}^1_k$ and a morphism $\phi: V \to X$ which makes the following diagram commutative:
 %$$\begin{tikzcd}
 %V \ar[r, "\phi"] \ar[d, "f"] &X \ar[d, "\eta"] \\
%\mathbb{A}^1_{k} \ar[r,"H"] 
%&\mathcal{S}(X)
%\end{tikzcd}$$ 
%Since $H$ is non-constant, there is a Nisnevich covering $V^{\prime} \to V \times_{\mathbb{A}^1_k} V$ such that 
%$$\phi \circ \text{pr}_1|_{V^\prime} \neq \phi \circ \text{pr}_2|_{V^\prime} \in X(V^{\prime}),$$
%where $\text{pr}_1, \text{pr}_2: V \times_{\mathbb{A}^1_k} V \to V$ are the projection maps.
%So there is a chain of non-constant naive $\mathbb{A}^1$-homotopies $G_1, \dots, G_l : \mathbb{A}^1_{V^{\prime}} \to X$ such that $G_1(0) = \phi \circ \text{pr}_1|_{V^\prime}$ and $G_l(1) = \phi \circ \text{pr}_2|_{V^\prime}$.
%There are $a, b \in V(k)$ such that $f(a) = 0$ and $f(b) = 1$. Since $x \neq y \in \mathcal{S}(X)(\Spec\ k)$, the image of $\phi$ contains two $k$-points $\phi(a), \phi(b)$ 
%(which we again denote by $x$ and $y$) in $X$ 
%such that $\phi(a) \neq \phi(b) \in \mathcal{S}(X)(\Spec\ k)$. 

%Suppose, $V  = \coprod_{i =1}^n V_i$, where $V_i$'s are the irreducible components of $V$, which are also the connected components of $V$. We complete the proof considering the following cases.

%\

%\textbf{Case 1:} Suppose, $\phi(a), \phi(b) \in \overline{\phi(V_i)}$, for every $i$. There is an irreducible component $V_0$ of $V^{\prime}$ such that
%$$\phi \circ \text{pr}_1|_{V_0} \neq \phi \circ \text{pr}_1|_{V_0}.$$
%We can assume that all the naive $\mathbb{A}^1$-homotopies $G_1|_{\mathbb{A}^1_{V_0}}, \dots, G_l|_{\mathbb{A}^1_{V_0}}$ are non-constant. Suppose, $V_0$ maps to $V_r \times_{\mathbb{A}^1_k} V_s$, for some $r, s \leq n$. Since the projection maps and the morphism $V^{\prime} \to V \times_{\mathbb{A}^1_U}$ are \'etale maps, so $\overline{Im(G_1|_{\mathbb{A}^1_{V_0}}(0))} = \overline{\phi(V_r)}$ and $\overline{Im(G_l|_{\mathbb{A}^1_{V_0}}(1))} = \overline{\phi(V_s)}$. Since every $G_i|_{\mathbb{A}^1_{V_0}}$ is non-constant, so $Im(G_i|_{\mathbb{A}^1_{V_0}})$ contains a line. If for every $i$, $\overline{Im(G_i|_{\mathbb{A}^1_{V_0}})}$ is contained in the image of some $\beta: \mathbb{A}^1_k \to X$, then then there is a chain of $\mathbb{A}^1_k$'s in $X$ that contains both $\phi(a)$ and $\phi(b)$, which would imply that $\phi(a) = \phi(b) \in \mathcal{S}(X)(\Spec\ k)$ and it is a contradiction. Therefore, there an some $j$ such that $G_j|_{\mathbb{A}^1_{V_0}}$ is non-constant and dominant. This is a contradiction, since $X$ is not an $\mathbb{A}^1$-uniruled surface.


% There is some $i$ such that there are $x^{\prime} \in  \overline{Im(G_i|_{\mathbb{A}^1_{V_0}}(0))}(k)$ and $y^{\prime} \in \overline{Im(G_i|_{\mathbb{A}^1_{V_0}}(1))}(k)$ and $x^\prime \neq y^\prime \in \mathcal{S}(X)(\Spec\ k)$. Indeed, if for every $j$, $z \in \overline{Im(G_j|_{\mathbb{A}^1_{V_0}}(0))}(k)$, $w \in \overline{Im(G_j|_{\mathbb{A}^1_{V_0}}(1))}(k)$, $z= w \in \mathcal{S}(X)(\Spec\ k)$. This implies that $\forall z \in \overline{Im(G_1|_{\mathbb{A}^1_{V_0}}(0))}(k)$, $w \in \overline{Im(G_l|_{\mathbb{A}^1_{V_0}}(1))}(k)$, we  must have $z = w \in \mathcal{S}(X)(\Spec\ k)$. Therefore, $x = y \in \mathcal{S}(X)(\Spec\ k)$, since $\overline{Im(G_1|_{\mathbb{A}^1_{V_0}}(0))} = \overline{\phi(V_r)}$ and $\overline{Im(G_l|_{\mathbb{A}^1_{V_0}}(1))} = \overline{\phi(V_s)}$. This is a contradiction and so our claim is true. The image of $G_i|_{\mathbb{A}^1_{V_0}}$ contains a line, since $G_i|_{\mathbb{A}^1_{V_0}}$ is non-constant. Both $x^{\prime}, y^{\prime}$ can not be in that line, since $x^{\prime} \neq y^{\prime} \in \mathcal{S}(X)(\Spec\ k)$. Therefore, $G_i|_{\mathbb{A}^1_{V_0}}$ is dominant, which is a contradiction.

%\

%\textbf{Case 2:} Suppose, $x \in \overline{\phi(V_i)}$, for every $i$ and there is $r$ such that $y \notin \overline{\phi(V_1)}, \dots, \overline{\phi(V_r)}$ and $y \in \overline{\phi(V_{r+1})}, \dots, \overline{\phi(V_n)}$. There is an irreducible component $V_0$ that maps to $V_1 \times_{\mathbb{A}^1_k} V_{r+1}$. Since, $\overline{Im(\phi \circ \text{pr}_1|_{V_0})} = \overline{\phi(V_1)}$ and $\overline{\phi \circ \text{pr}_2|_{V_0}} = \overline{\phi(V_{r+1})}$ and $y \notin \overline{\phi(V_{1})}$, 
%$$\phi \circ \text{pr}_1|_{V_0} \neq \phi \circ \text{pr}_2|_{V_0}.$$
%We can assume all the naive $\mathbb{A}^1$-homotopies $G_1|_{\mathbb{A}^1_{V_0}}, \dots, G_l|_{\mathbb{A}^1_{V_0}}$ are non-constant. There is some $i$ such that there are $x^{\prime} \in  \overline{Im(G_i|_{\mathbb{A}^1_{V_0}}(0))}(k)$ and $y^{\prime} \in \overline{Im(G_i|_{\mathbb{A}^1_{V_0}}(1))}(k)$ and $x^\prime \neq y^\prime \in \mathcal{S}(X)(\Spec\ k)$.


%\


%\textbf{Case 2:} Suppose that $\phi(a) \in \overline{\phi(V_r)}$ and $\phi(b) \in \overline{\phi(V_s)}$, for some $r, s$. Also assume that $\phi(b) \notin \overline{\phi(V_r)}$. There is an irreducible component $V_0$ that maps to $V_r \times_{\mathbb{A}^1_k} V_{s}$. Since the projection maps and the morphism $V^{\prime} \to V \times_{\mathbb{A}^1_U}$ are \'etale maps, so $\overline{Im(\phi \circ \text{pr}_1|_{V_0})} = \overline{\phi(V_r)}$ and $\overline{Im(\phi \circ \text{pr}_2|_{V_0})} = \overline{\phi(V_s)}$. Since $\phi(b) \notin \overline{\phi(V_r)}$, so
%$$\phi \circ \text{pr}_1|_{V_0} \neq \phi \circ \text{pr}_2|_{V_0}.$$
%So we can assume all the naive $\mathbb{A}^1$-homotopies $G_1|_{\mathbb{A}^1_{V_0}}, \dots, G_l|_{\mathbb{A}^1_{V_0}}$ are non-constant. Since every $G_i|_{\mathbb{A}^1_{V_0}}$ is non-constant, so $Im(G_i|_{\mathbb{A}^1_{V_0}})$ contains a line. If for every $i$, $\overline{Im(G_i|_{\mathbb{A}^1_{V_0}})}$ is contained in the image of some $\beta: \mathbb{A}^1_k \to X$, then there is a chain of $\mathbb{A}^1_k$'s in $X$ that contains both $\phi(a)$ and $\phi(b)$, 
%lie in a chain of $\mathbb{A}^1_k$'s, 
%which would imply that $\phi(a) = \phi(b) \in \mathcal{S}(X)(\Spec\ k)$ and it is a contradiction. Therefore, there an some $j$ such that $G_j|_{\mathbb{A}^1_{V_0}}$ is non-constant and dominant. This is a contradiction, since $X$ is not an $\mathbb{A}^1$-uniruled surface.

%\

%Thus in both the cases we arrive at a contradiction. Therefore, the canonical epimorphism
%$$\mathcal{S}(X) \to \mathcal{S}^2(X)$$
%is injective over the section $\Spec\ k$ and hence the morphism 
%$$\mathcal{S}(X)(\Spec\ k) \to \mathcal{S}^2(X)(\Spec\ k)$$
%is a bijection. 
%Since, $\overline{Im(\phi \circ \text{pr}_1|_{V_0})} = \overline{\phi(V_r)}$ and $\overline{\phi \circ \text{pr}_2|_{V_0}} = \overline{\phi(V_{s})}$ and $y \notin \overline{\phi(V_{r})}$, 
%$$\phi \circ \text{pr}_1|_{V_0} \neq \phi \circ \text{pr}_2|_{V_0}.$$
%We can assume all the naive $\mathbb{A}^1$-homotopies $G_1|_{\mathbb{A}^1_{V_0}}, \dots, G_l|_{\mathbb{A}^1_{V_0}}$ are non-constant. There is some $i$ such that there are $x^{\prime} \in  \overline{Im(G_i|_{\mathbb{A}^1_{V_0}}(0))}(k)$ and $y^{\prime} \in \overline{Im(G_i|_{\mathbb{A}^1_{V_0}}(1))}(k)$ and $x^\prime \neq y^\prime \in \mathcal{S}(X)(\Spec\ k)$.
%\end{proof}

\begin{corollary} \label{termination}
Suppose, $X \in Sm/k$ is an affine surface, which is not an $\A^1$-uniruled surface. Then the canonical epimorphism 
$$\mathcal{S}(X) \to \mathcal{L}(X)$$
is an isomorphism.
\end{corollary}
\begin{proof}
Since the canonical morphism $\mathcal{S}(X) \to \mathcal{S}^2(X)$ is an isomorphism by Theorem \ref{non-uniruled}, so all the canonical morphisms $\mathcal{S}^n(X) \to \mathcal{S}^{n+1}(X)$ is an isomorphism, for every $n$. 
%So the sequence $\{\mathcal{S}^n(X)\}_n$ terminates. 
Since the sheaf $\mathcal{L}(X) =  \underset{n}{\varinjlim} \; \mathcal{S}^n(X)$, 
%is a colimit of $\mathcal{S}^n(X)$'s, 
so the canonical morphism
$$\mathcal{S}(X) \to \mathcal{L}(X)$$
is an isomorphism. 
\end{proof}
\begin{corollary} \label{invariant non-uniruled}
Let $k$ be an algebraically closed field of characteristic $0$. Suppose that $X \in Sm/k$ is an affine and not an $\mathbb{A}^1$-uniruled surface. Then $\pi_0^{\mathbb{A}^1}(X)$ is $\mathbb{A}^1$-invariant.
\end{corollary}
\begin{proof}
The canonical epimorphism $\mathcal{S}(X) \to \pi_0^{\mathbb{A}^1}(X)$ gives an epimorphism $\Phi: \mathcal{L}(X) \to \pi_0^{\mathbb{A}^1}(X)$, by Corollary \ref{termination}. There is a canonical epimorphism $\Psi: \pi_0^{\mathbb{A}^1}(X) \to \mathcal{L}(X)$. Since $\mathcal{L}(X)$ is the universal $\mathbb{A}^1$-invariant sheaf, so $\Psi \circ \Phi$ is identity. Hence $\Phi$ is an isomorpism.
%Both the morphisms $\Psi \circ \Phi$ and $\Phi \circ \Psi$ are the identities, due to the universal properties of $\mathcal{L}(X)$ and $\pi_0^{\mathbb{A}^1}(X)$ respectively. 
Therefore,
$$\pi_0^{\mathbb{A}^1}(X) \cong \mathcal{L}(X) \cong \mathcal{S}(X)$$
and consequently, $\pi_0^{\mathbb{A}^1}(X)$ is $\mathbb{A}^1$-invariant.
\end{proof}
\begin{corollary}
Let $X \in Sm/k$ be an affine surface, where $k$ is an algebraically closed field of characteristic zero. Suppose that $X$ has non-negative logarithmic Kodaira dimension. Then $\pi_0^{\mathbb{A}^1}(X)$ is $\mathbb{A}^1$-invariant. For example, if $X$ is the Ramanujan surface \cite[Section 3]{ra} or $X$ is a tom Dieck-Petrie surface, then $\pi_0^{\mathbb{A}^1}(X)$ is $\mathbb{A}^1$-invariant.
\end{corollary}
\begin{proof}
By \cite[Lemma 1.8]{russell} an $\mathbb{A}^1$-uniruled affine surface has negative logarithmic Kodaira dimension. So the corollary follows from Corollary \ref{invariant non-uniruled}.
\end{proof}
\begin{remark}
Corollary \ref{invariant non-uniruled} gives an affirmative answer to \cite[Conjecture 2.2.8]{am}, in case of affine, non $\A^1$-uniruled, smooth surfaces $X$. However, in general the Conjecture is known to be false \cite[Section 4]{bhs}. Also, note that if $Sing_*(X)$ is $\mathbb{A}^1$-local, then the epimorphism 
$$\mathcal{S}(X) \to \pi_0^{\mathbb{A}^1}(X)$$
is an isomorphism \cite[Remark 2.2.9]{am}. Thus it is natural to ask for an affine non $\A^1$-uniruled, smooth surface $X$, whether $Sing_*(X)$ is $\mathbb{A}^1$-local.
\end{remark}
	
%	\subsection{Surfaces with negative logarithmic Kodaira dimension}
\subsection{$\A^1$-invariance of $\pi_0^{\A^1}(X)$, for $\A^1$-uniruled, affine surface $X$:}
\


Let $X$ be a smooth affine surface over an algebraically closed field $k$ of characteristic zero. In this subsection, we will prove that if $X$ is an $\A^1$-uniruled surface (Definition \ref{ruled-defn}), then $\pi_0^{\A^1}(X)$ is $\A^1$-invariant (Theorem \ref{connected-base-new}). Recall that, $X$ is $\A^1$-uniruled if and only if $X$ admits an $\A^1$-fibration $\pi:X \to C$, onto a smooth curve $C$ (Theorem \ref{all uniruled equivalences}) and the $\A^1$-fibration $\pi$ factors as $\pi= \alpha \circ \theta$ (Theorem \ref{factor-A1-fibration}), where $\theta: X \to \sC$ is an \'etale locally trivial $\A^1$-bundle over an algebraic space $\sC$ and $\alpha: \sC \to C$ is the structure morphism which is surjective, quasi-finite and birational. We have proved that the morphism $\theta$ is an $\A^1$-weak equivalence (Corollary \ref{factorisation A1 equivalence}). The $\A^1$-fibration $\pi: X \to C$ can be divided into two following classes -

\begin{enumerate}
	\item $C$ is an $\A^1$-rigid curve
	\item $C$ is not $\A^1$-rigid, hence $\A^1$-connected curve.
\end{enumerate} 
 We will prove the $\A^1$-invariance of $\pi_0^{\A^1}(X)$ by proving $\pi_0^{\A^1}(\sC)$ is $\A^1$-invariant in each case.
%for an $\A^1$-uniruled affine surface $X$, we will divide the $\A^1$-uniruled surfaces $X$ in two classes, according to the base curve $C$ of the $\A^1$-fibration $\pi: X \to C$ (by Theorem \ref{all uniruled equivalences}): one class consists such $X$ that admits an $\A^1$-fibration onto an $\A^1$-rigid smooth curve $C$ (Definition \ref{invariant definition}(2)) and the other
%%, for which $C$ to be an $\A^1$-rigid curve and the other 
%class consists such $X$ that admits an $\A^1$-fibration $\pi: X \to C$ such that $C$ is not an $\A^1$-rigid curve.
%% the base curve $C$ not to be $\A^1$-rigid. 
%From the classification of projective curves based on their genus, a smooth curve $C$ is not $\A^1$-rigid if and only if 
%%$C$ to be $\A^1$-connected and this implies 
%$C$ is isomorphic to the affine line $\A^1_k$ or the projective line $\P^1_k$. 

First we consider the case of $\A^1$-fibration $\pi: X\to C$ onto an $\A^1$-rigid smooth curve $C$. The following theorem says that if the base curve $C$ is $\mathbb{A}^1$-rigid, then any morphism $\mathbb{A}^1_k \to \sC$ is constant, which means the morphism factors through $\Spec k$. 

		\begin{theorem}\label{rigid-base}
			Let $X$ be a smooth affine surface, which admits an $\mathbb{A}^1$-fibration $\pi: X \to C$ onto an  
%$C$ is a smooth 
$\A^1$-rigid, smooth curve $C$ and
%$Cand  $\pi: X \rightarrow C \mathrm{\ an\ }\A^1$-fibration,  
$\sC$ be the algebraic space 
appearing in the factorisation of $\pi$, by Theorem \ref{factor-A1-fibration}. Then any morphism $H: \A^1_k \to \sC$, factors through $\Spec k$.
% $\mathrm{\ is}$ also $\A^1$-rigid.
		\end{theorem}
		
		\begin{proof}
			Suppose, the $\A^1$-fibration $\pi: X \to C$ factorises as (Theorem \ref{factor-A1-fibration})
 %Let 
$$X \xrightarrow{\theta} \mathcal{C}\xrightarrow{\alpha} C,$$ 
%is the factorisation as in \cite{doub}, 
where $\sC$ is an algebraic space along with the morphism $\alpha$, which is surjective, quasi-finite, birational morphism to $C$ and $\theta$ is an étale $\A^1$-bundle. 
%We need to show that for a hensel local ring $\sO$, any morphism $H: \mathbb{A}^1_{\mathcal{O}} \to \mathcal{C}$ always factors through $\Spec\ \mathcal{O}$. We will prove it in two steps:
%\begin{enumerate}
%\item Firstly we will prove that any morphism $H: \mathbb{A}^1_L \to \mathcal{C}$ 
%is constant, that is $H$ 
%factors through $\Spec\ L$, for every finitely generated separable extension $L/k$.
%\item Using Part (1), we will prove that any morphism $H: \mathbb{A}^1_\mathcal{O} \to \mathcal{C}$ factors through $\Spec\ \mathcal{O}$, for every Henselian local ring $\mathcal{O}$.
%\end{enumerate}

%\textbf{Proof of Step 1:} 
Suppose, $H: \A^1_k \rightarrow \sC$ is a morphism. Since $C$ is $\mathbb{A}^1$-rigid, so the morphism $\alpha \circ H: \mathbb{A}^1_k \to C$  
%as $\A^1_k \xrightarrow{\beta}\sC \xrightarrow{\alpha}C$ is a constant map it 
factors as
$$\A^1_k \xrightarrow{\pr} \Spec\ k  \xrightarrow{x} C,$$
we have the following commutative diagram 

\[
\xymatrix{
\A^1_k \ar[r]^H \ar[d]^{\pr}& \sC \ar[d]^\alpha\\
\Spec\ k \ar[r]^x & C
}\]


              Suppose, $x: \Spec\ k \to C$ maps to a $k$-point of $C$ (as it cannot map to the generic point of $C$). Since $\alpha$ is quasi-finite, so the fiber $\alpha^{-1}(x)$ consists only finitely many $k$-points. If Im$(H)\subset \sC$ is a subscheme, then $H$ is constant, since the topological space $\overline{\text{Im}(H)}$ is irreducible and connected . Otherwise Im$(H)$ must maps to a non-reduced point $\mathfrak{c}$ of $\sC$ (a point of $\sC$ over which the fiber is non-reduced) with $\kappa(\mathfrak{c})=k[t]/(t^m)$ for some $m>1$ . Then we have an étale covering $\phi: U \to \sC$ which have an unique $k$-point $y\mapsto \mathfrak{c}$ and $H$ factors through $U$. Since the only $k$-algebra map between $\kappa(\mathfrak{c})=k[t/(t^m)]\to \kappa(y)=k$ is $t\mapsto 0$ and $H$ factors through $U$, it must be constant and factors through $\Spec k$. 
 %So $\text{Im}(H)\subset \alpha^{-1}(x)$. Since $\alpha$ is quasi-finite, so $\alpha^{-1}(x)$ consists only finitely many $k$-points of $\mathcal{C}$. But the topological space $\overline{\text{Im}(H)}$ is irreducible, therefore, $H$ is constant. Hence $H$ factors through $\Spec\ L$.
%Now take $\tilde{x}$ be any $k$-rational point of $X$ then $\A^1_k \xrightarrow{\tilde{x}} X$ is a lift of $\beta$. Thus $\beta$ is constant.

%It is enough to show any morphism $\A^1_k\rightarrow\sC$ and $\A^1_\sO \rightarrow \sC$ factors through $spec(k)$ and $\sO$ respectively, for any separably generated field $k$ and any hensel local ring $\sO$.\\
			

%First let $\beta: \A^1_k \rightarrow \sC$ any morphism as $\A^1_k \xrightarrow{\beta}\sC \xrightarrow{\alpha}C$ is a constant map it factors through $\A^1 \xrightarrow{\pr} spec(k) \xrightarrow{x} C$, hence $\im(\beta)\subset \alpha^{-1}(x)$. Now take $\tilde{x}$ be any $k$-rational point of $X$ then $\A^1_k \xrightarrow{\tilde{x}} X$ is a lift of $\beta$. Thus $\beta$ is constant.\\
			
%\textbf{Proof of Step 2:} Suppose,  $\beta: \A^1_\sO \rightarrow \sC$ is a morphism. We have the following pullback diagram:
		%
%			\[\xymatrixcolsep{4pc}
		%	\xymatrix{{\mathbb{A}^1_{\mathcal{O}}\times_{\mathcal{C}}X \cong \mathbb{A}^1_{\mathcal{O}}%\times_k\mathbb{A}^1_k} \ar[d]^{\beta_* \rho} \ar[r] & X \ar[d]^{\rho}\\
	%			{\mathbb{A}^1_{\mathcal{O}}} \ar[ru]|{\tilde{\beta}} \ar[r]_\beta & {\mathcal{C}}}
	%		\]
			
	%		As $\rho$ is an étale locally trivial $\A^1$-bundle, so $\beta^*\rho$ is also an étale locally trivial $\A^1$-bundle. Since $\Aut(\A^1) = \G_a\rtimes \G_m$, so an étale locally trivial $\A^1$-bundle over $\mathbb{A}^1_\mathcal{O}$ is trivial (by the same argument, as in the proof of Theorem \ref{sing-etale-bundle}). Thus $\beta_*\rho$ is isomorphic to the trivial $\mathbb{A}^1$-bundle $\mathbb{A}^1_{\mathcal{O}} \times_k \mathbb{A}^1_k \to \mathbb{A}^1_{\mathcal{O}}$. So $\beta_*\rho$ admits a section.
%Consider the composition of mophisms $\tilde{\beta}$ defined as
% $$\tilde{\beta}: \A^1_{\sO}\xrightarrow{s_0}\A^1_{\sO}\times_k \A^1_k \cong \A^1_{\sO}\times_{\sC} X \rightarrow X,$$ where $s_0$ is the zero section. If $\im(s_0)$ is inside a single fiber of the map $\A^1_{\sO}\times_k \A^1_k \cong \A^1_{\sO}\times_{\sC} X \rightarrow X$ then $\tilde{\beta}$ is constant hence $\beta$ is also constant. Otherwise we see $(\tilde{\beta} \circ i_0(\sO\rightarrow \A^1_{\sO}) \circ \pr(\A^1_{\sO}\rightarrow \sO))|_{\A^1_{k(\eta)}}=\tilde{\beta}|_{\A^1_{k(\eta)}}$ where $\eta$ is the generic point of $\sO$. Now for any $x\in \A^1_{\sO}$ we have $y \in \A^1_{\sO}$ a $k(\eta)$ point of $\A^1_{\sO}$ such that $x\in \overline{\{y\}}$, but $\tilde{\beta}$ maps every $k(\eta)$ points and generic point of $\A^1_{\sO}$ to a $k(\eta)$ point of $X$ , say $\eta'$. Thus $\im(\tilde{\beta})\subset \overline{\{\eta'\}}=\im(\tilde{\beta}\circ \pr(\A^1_{\sO}\rightarrow \sO) \circ i_0(\sO\rightarrow\A^1_{\sO}))$. Hence $\tilde{\beta}$ factors through $spec(\sO)$, subsequently $\beta$ also factor through it. 
			
			
		\end{proof}

Before proving the $\A^1$-invariance of $\pi_0^{\A^1}(X)$,
%To do that 
we need the
%the case of $\A^1$-connected curve, 
 this following technical result, which says that the algebraic space $\sC$ in Theorem \ref{factor-A1-fibration} is $\A^1$-invariant if and only if any morphism $\mathbb{A}^1_k \to \sC$ is constant. 
 %By ``a non-constant morphism $f: X \to \sC$", where $X$ is a smooth scheme over $k$ and $\sC$ is a smooth algebraic space over $k$, we mean that the induced morphism between the associated topological spaces is non-constant.
 %and by $f$ is dominant we mean that.
%		\begin{lemma}\label{A1 in algebraic space}
%			Suppose, $\mathcal{C}$ is a finite type, smooth algebraic space (not necessarily separated) over $k$ and $U \in Sm/k$ ($k$ is an algebraically closed field). Suppose that there are two morphisms $f,g: U \to \mathcal{C}$ such that for every $u \in U(k)$, the compositions of the morphisms
%			$$f(u): \Spec\ k \xrightarrow{u} U \xrightarrow{f} \mathcal{C}$$
%			and
%			$$g(u): \Spec\ k \xrightarrow{u} U \xrightarrow{g} \mathcal{C}$$
%			are equal. Then, $f=g \in \mathcal{C}(U)$.
%		\end{lemma}
%		\begin{proof}
%			We will follow the proof in \cite[Lemma 2.18]{haine}. Consider the equiliser of $f$ and $g$ which is given by the fiber product in the category of algebraic spaces:
%		
%			
%			\[\xymatrixcolsep{3pc}
%			\xymatrix{ E \ar[d]^\phi \ar[r] & \sC \ar[d]^\Delta \\
%				U \ar[r]^{(f,g)\ \ } & {\sC\times\sC}}
%			\]
%			Here $\Delta$ is the diagonal morphism. 
%%the diagonal morphism 
%Since $\mathcal{C}$ is an algebraic space, $\Delta$ is representable by schemes \cite[Theorem 4.2.1]{alpher} and the %diagonal morphism 
%$\Delta$ is a monomorphism \cite[Theorem 5.5.10]{alpher}. Therefore, the algebraic space $E$ is a smooth scheme of finite type over $k$ (not necessarily separated) and the morphism
%			$$\phi: E \to U$$
%			is a monomorphism (in the category of schemes). Moreover from the hypothesis, we have the morphism $\phi$ restricts to a surjective maps between the $k$-points
%			$$\phi: E(k) \to U(k).$$
%			This implies, $\phi$ is dominant; so $dim(E) \geq dim(U)$.
%			Now, we claim that $\phi$ is an immersion that is, $\phi$ can be factored as a closed immersion, followed by an open immersion. 
%			Firstly, by \cite[Tag 05VH]{stack} $\phi$ is an unramified morphism. Therefore, by \cite[Theorem 1.2]{rydh}, $\phi$ can be factored as
%			$$E \xrightarrow{i} E_\phi \xrightarrow{e} U,$$
%			where $E_\phi$ is an algebraic space, $i$ is a closed immersion and $e$ is an \'etale morphism. So, $dim(E) \leq dim(E_\phi)$ and $dim(E_\phi) = dim(U)$. Hence, 
%			$$dim(E) = dim(E_\phi) = dim(U).$$
%			Therefore, the closed immersion $i$ must be surjective and hence is an isomorphism. So the algebraic space $E_\phi$ is an scheme and the \'etale morphism $e$ is an open immersion \cite[Tag 025G]{stack}. 
%			
%			Since $\phi = e \circ i$, so $\phi(E)\subset e(E_\phi)$. But $\phi$ is surjective on the closed points, so the closed subset $U \setminus e(E_\phi)$ does not contain any closed point. Hence the open immersion $e$ must be surjective. So $e$ is an isomorphism. 
%%Similarly, the closed immersion must also be surjective, so $i$ is also an isomorphism. 
%Therefore, $\phi$ is an isomorphism.
%			Hence, $f = g$.
%			
%		\end{proof}
%Therefore, using \cite[Lemma 2.16]{mazza}, we have the following corollary:
%\begin{corollary} \label{not a1 rigid}
%Suppose, $\mathcal{C}$ is a finite type, smooth algebraic space over an algebraically closed field $k$ (not necessarily seperated). Then $\mathcal{C}$ is $\mathbb{A}^1$-invariant if and only if every morphism $\gamma: \mathbb{A}^1_k \to \mathcal{C}$ is constant.
%\end{corollary}
%\begin{proof}
%One direction is from the definition of $\mathbb{A}^1$-invariant sheaf. For the other direction, assume that every morphism $\mathbb{A}^1_k \to \mathcal{C}$ is constant. If possible, $\mathcal{C}$ is not $\mathbb{A}^1$-invariant Then by \cite[Lemma 2.16]{mazza}, there is $H \in \mathcal{C}(\mathbb{A}^1_U)$, for some $U \in Sm/k$ such that $H \circ i_0 \neq H \circ i_1$, where $i_0, i_1: U \to \mathbb{A}^1_U$ are the $0$-section and the $1$-section respectively. Since, $H \circ i_0 \neq H \circ i_1$, so by Lemma \ref{A1 in algebraic space}, there is $u \in U(k)$ such that the following two compositions of the morphisms are not equal:
%$$\Spec\ k \xrightarrow{u} U \xrightarrow{H \circ i_0} X \text{ and } \Spec\ k \xrightarrow{u} U \xrightarrow{H \circ i_1} X.$$
%Now, consider the morphism $\gamma: \mathbb{A}^1_k \to \mathcal{C}$ defined as the following compositions of the morphisms:
%$$\mathbb{A}^1_k \times_k \Spec\ k \xrightarrow{(\text{identity}, u)} \mathbb{A}^1_k \times_k U \xrightarrow{H} \mathcal{C}.$$
%Since, $H \circ i_0 \neq H \circ i_1$, so the morphism $\gamma$ is non-constant. It is a contradiction. So the algebraic space $\mathcal{C}$ is $\mathbb{A}^1$-invariant.
%\end{proof}
%\begin{remark}
%Lemma \ref{A1 in algebraic space} can be thought of as a version of the similar well known result in case of schemes: for two finite type schemes $X,Y$ over an algebraically closed field $k$ with $X$ is reduced, if $f,g: X \to Y$ are two morphisms such that for every $x \in X(k)$, the compositions of morphisms
%$$\Spec\ k \xrightarrow{x} X \xrightarrow{f} Y \text{ and } \Spec\ k \xrightarrow{x} X \xrightarrow{g} Y$$
%are equal, then $f=g$.
% where $k$ is
%\end{remark}
%We need the following technical lemma:
%\begin{definition}
%$X$ is said to be \textbf{$\mathbb{A}^1$-ruled} or affine-ruled if there is a Zariski open dense subset $U$ of $X$ such that $U$ is isomorphic to $\mathbb{A}^1_k \times_k Z$ for some $k$-variety $Z$ \cite[Section 2.1]{miyanishicrm}.
%\end{definition}

%Let $X$ be an affine surface over $\C$. we can classify them using logarithmic kodaira dimension. From \cite[Lemma 1.8]{russell}, \cite[Chapter 2, Theorem 2.1.1]{miyanishicrm} and \cite[Chapter 3, Lemma 1.3.1]{miyanishicrm}, we have the following equivalences:
%\begin{enumerate}
%	\item $X$ is $\mathbb{A}^1$-uniruled.
%	\item $X$ has negative logarithmic Kodaira dimension.
%	\item $X$ is $\mathbb{A}^1$-ruled.
%	\item $X$ admits an $\mathbb{A}^1$-fibration $\pi: X \to C$.
%\end{enumerate}

%	Let X be a smooth affine surface with $\bar{\kappa}(X)=-\infty$, then it admits an $\A^1$-fibration $\pi: X \to C$ onto a smooth curve $C$. The curve $C$ is either $\A^1$-rigid, otherwise $C$ is $\A^1$-connected. Here we will prove the $\A^1$-invariance of $\pi_0^{\A^1}(X)$ in two steps, according to the base curve $C$ is $\mathbb{A}^1$-rigid or $\mathbb{A}^1$-connected. %considering  for two types of curves, the surface admitting the $\A^1$-fibration to.

\begin{lemma}\label{not-a1-rigid}
	Let $X$ be a smooth affine surface, which admits an $\mathbb{A}^1$-fibration $\pi: X \to C$ onto a  
	%$C$ is a smooth 
	smooth curve $C$ and
	%$Cand  $\pi: X \rightarrow C \mathrm{\ an\ }\A^1$-fibration,  
	$\sC$ is the algebraic space 
	%appeared 
%as in the factorisation of $\pi$ 
	in Theorem \ref{factor-A1-fibration}. If $\sC$ is not $\A^1$-invariant, then there exists a non-constant morphism $\gamma:\A^1_k \to \sC$.
\end{lemma}

\begin{proof}
The $\A^1$-fibration $\pi: X \to C$
%	Let $X \to C$ be an $\A^1$-fibration 
from a smooth affine surface $X$ to a smooth curve $C$ factors through an \'etale locally trivial $\A^1$-bundle $\theta: X \to \sC$
% and let $\sC$ be the 
over a smooth algebraic space $\sC$ (Theorem \ref{factor-A1-fibration}). Assume that $\sC$ is not $\A^1$-invariant.
Then by \cite[Lemma 2.16]{mazza}, there is $H\in \sC(\A^1_U)$ for some $U \in Sm/k$ such that 
$H(0) \neq H(1) \in \sC(U)$ (Notations \ref{notations}). Since $\sC$ is an \'etale sheaf on $Sm/k$, we can moreover assume $U$ to be affine.
%, where 
%$s_0,s_1:U\to\A^1_U$
%$$H(0)= H \circ i_0 \text{ and } H(1)=H \circ i_1; \ i_0, i_1: U \to \A^1_U$$ 
%are the $0$-section and $1$-section
% zero and one section 
%respectively. 
The pullback $\theta^*H: \A^1_U\times_\sC X \to \A^1_U$ is an \'etale $\A^1$-bundle over $\A^1_U$.
An étale $\A^1$-bundle over $\A^1_U$ corresponds to a cohomology class in $H^1_{\text{\'et}}(\A^1_U,\Aut(\A^1_k))$ and 
$$H^1_{\text{\'et}}(\A^1_U,\Aut(\A^1_k)) \cong H^1_{\text{Zar}}(\A^1_U,\Aut(\A^1_k)) \cong H^1_{\text{Zar}}(U,\Aut(\A^1_k)),$$ 
so we can choose a Zariski affine trivialisation of $\{U_i\}_i$ of $U$ such that $\theta^*H|_{\A^1_{U_i}}$ is trivial, for every $i$. Thus there is some $i,\ H|_{\A^1_{U_i}}(0) \neq H|_{\A^1_{U_i}}(1) \in \sC(U_i)$. Therefore, 
%without loss of generality 
we can assume that there is $H \in \sC(\A^1_U)$, where $U \in Sm/k$ is irreducible, affine and $H$ satisfies $H(0) \neq H(1) \in \sC(U)$ and also the pullback $\theta^*H: \A^1_U\times_\sC X \to \A^1_U$ is a trivial $\A^1$-bundle.
%connected.
Consider the following diagram:	
  
  \begin{equation}\label{main-lemma-diagram}
  	\xymatrix{
  		H^{-1}(V)\times_\sC X \ar@{^{(}->}[r] \ar[d]^{\theta^*H}&\A^1_U\times_\sC X \ar[r]^-{H'} \ar[d]_{\theta^*H}& X \ar[d]^\theta & V\times_\sC X \ar@{_{(}->}[l] \ar[d]_\theta\\
  		H^{-1}(V) \ar@{^{(}->}[r] \ar@/^/[u]^{s_0}&\A^1_U \ar@/_/[u]_{s_0} \ar[r]^H&\sC & V \ar@{_{(}->}[l] \ar@/_/[u]_{s_0}
  	}
  \end{equation}
  
%  Here the pullback bundle $\theta^*H: \A^1_U\times_\sC X \to \A^1_U$ is trivial $\A^1$-bundle. 
Since the $\A^1$-fibration $\pi$ is generically trivial, there is an open subspace $V\subset\sC$
% be the open subspace
, which is also a scheme,
% part of the algebraic space 
over which the bundle is trivial, that is $\theta: V \times_\sC X \to V$ is a trivial $\A^1$-bundle. We denote all the zero sections of corresponding bundles by $s_0$. So the following diagram is a pullback diagram

  \begin{equation}\label{main-lemma-diagram-pulback}
  \xymatrixcolsep{7pc}
  \xymatrix{
  H^{-1}(V)\times_\sC X \cong H^{-1}(V)\times \A^1_k \ar[d]^{\theta^*H} \ar[r]^-{(H,id)}& V\times \A^1_k \cong V\times_\sC X \ar[d]_{\theta} \\
  H^{-1}(V)\ar[r]^H \ar@/^/[u]^{s_0} & V \ar@/_/[u]_{s_0}}
  \end{equation}
%Define, $\widetilde{H}$ to be the composition of morphisms
%$$ \A^1_U\xrightarrow{s_0}\A^1_U\times_\sC X \xrightarrow{H'} X.$$
%Then, $\theta \circ \widetilde{H} =H$.
%Since the sections $H(0) \neq H(1) \in \sC(U)$,
% and $\theta \circ \widetilde{H} = H$, 
%so, 
%$$\widetilde{H}(0) \neq \widetilde{H}(1) \in X(U).$$


%  is a pullback diagram, t
Therefore by commutativity of diagram (\ref{main-lemma-diagram}) we have 
$$(H^\prime \circ s_0)(H^{-1}(V)) = s_0(H(H^{-1}(V))).$$
Since, $H$ is a non-constant homotopy (that is $H(0) \neq H(1) \in \sC(U)$) and the algebraic space $\sC$ is of dimension $1$, so the restriction
$$H: H^{-1}(V) \to V$$
is dominant. So,
$$\overline{(H'\circ s_0)(H^{-1}(V))}= \overline{s_0(V)} \subseteq V\times_\sC X$$

%$$(H'|_{H^{-1}(V)}\circ s_0)(H^{-1}(V))=s_0(V).$$
%$$(H^\prime \circ s_0)(H^{-1}(V)) = s_0(V).$$
 Thus, 
%dim$\overline{(H'|_{H^{-1}(V)}\circ s_0)}= \dim\overline{(s_0(V))}=1$ in $V\times_\sC X$.
dim$\overline{(H'\circ s_0)(H^{-1}(V))}= \dim\overline{(s_0(V))}=1$, as closures in $V\times_\sC X$.

 Let $\widetilde{H}$ be the composition of the morphisms
$$ \A^1_U\xrightarrow{s_0}\A^1_U\times_\sC X \xrightarrow{H'} X.$$
Then, $\widetilde{H}$ is a lift of $H$, that is $\theta \circ \widetilde{H} =H$.


 \begin{figure}
	
	\includegraphics[width=0.5\textwidth]{Surface-2}
	\caption{Non-constant map from $\A^1_k$}
\end{figure}
\graphicspath{{./Surface-2.jpg}}
Since, $\overline{(H'\circ s_0)(H^{-1}(V))}$ as closure in $V\times_\sC X$ is of dimension $1$ and $V\times_\sC X$, $H^{-1}(V)$ are open subschemes of $X$, $\A^1_U$
% \times_\sC X$
respectively, so
% from commutativity of diagram (\ref{main-lemma-diagram}) we have 
$$\dim\overline{\widetilde{H}(\A^1_U)}= 1, \text{ as closure in } X.$$ 
%as $\dim\overline{(H'|_{H^{-1}(V)}\circ s_0)}=1$ and $H^{-1}(V)\times_\sC X$ is open dense in $\A^1_U \times_\sC X$.\par 
Since the sections $H(0) \neq H(1) \in \sC(U)$ and $\widetilde{H}$ is a lift of $H$,
%$\theta \circ \widetilde{H} = H$, 
so, 
$$\widetilde{H}(0) \neq \widetilde{H}(1) \in X(U).$$

%Since the sectins $H(0) \neq H(1) \in \sC(U)$ and $\theta \circ \widetilde{H} = H$, so $\widetilde{H}(0) \neq \widetilde{H}(1) \in X(U)$. 
Thus there exists $x \in U(k)$ such that $\widetilde{H}(0)(x)\neq \widetilde{H}(1)(x)$. Now define $\widetilde{\gamma}: \A^1_k \to X$ as
$$\widetilde{\gamma}: \A^1_k \times \Spec(k) \xrightarrow{(id\times x)} \A^1_k \times U \xrightarrow{\widetilde{H}} X.$$

%  Now we have a lift $\widetilde{H}:\A^1_U \to X$ of $H:\A^1_U \to \sC$ such that Im($\widetilde{H}(\A^1_U)$) does not lie inside a single fiber and $t_0^*(\widetilde{H})\neq t_1^*(\widetilde{H})\in X(U)$. Since $\widetilde{H}\circ t_0,\ \widetilde{H}\circ t_1 :U\to X$ are not same t
%Therefore there exists $x \in U(k)$ such that $\widetilde{H}(0)(x)\neq \widetilde{H}(1)(x)$. Now define 
%$$\widetilde{\gamma}: \A^1_k \times \Spec(k) \xrightarrow{(id\times x)} \A^1_k \times U \xrightarrow{\widetilde{H}} X.$$

 %Clearly, the images $\widetilde{\gamma}(\A^1_k))\subset \overline{(\widetilde{H}(\A^1_U))}$,
 Since $X$ is affine, so the image $\widetilde{\gamma}(\A^1_k)$ is closed in $X$ and since $\widetilde{\gamma}(0) \neq \widetilde{\gamma}(1)$, so $\dim \widetilde{\gamma}(\A^1_k) = 1$. Now $\widetilde{\gamma}(\mathbb{A}^1_k) \subseteq \widetilde{H}(\mathbb{A}^1_U)$ and since, $\dim\overline{\widetilde{H}(\A^1_U)}= 1$, so
 $$\overline{\widetilde{H}(\A^1_U)} =   \widetilde{\gamma}(\mathbb{A}^1_k) = \widetilde{H}(\mathbb{A}^1_U).$$
 Since, $H(0) \neq H(1) \in \sC(U)$, so the image $\widetilde{H}(\mathbb{A}^1_U)$ 
 %Thus the image $\widetilde{H}(\mathbb{A}^1_U)$ 
 can not be contained in a single fiber over $\theta$.
% $$\$$
% by$\overline{\widetilde{H}(\A^1_U)}$ has dimension $1$. Thus 
%$$\widetilde{\gamma}(\A^1_k)= \overline{\widetilde{H}(\A^1_U)}$$
% and 
Therefore, $\gamma:=\theta \circ \widetilde{\gamma}:\A^1_k \to \sC$ is a non-constant morphism.
 
\end{proof}

Therefore, by Theorem \ref{rigid-base} and Lemma \ref{not-a1-rigid}, the algebraic space $\sC$ is also $\A^1$-invariant in this case.
\begin{corollary} \label{a1 invariance of algebraic space}
Let $X$ be a smooth affine surface, which admits an $\mathbb{A}^1$-fibration $\pi: X \to C$ onto an  
%$C$ is a smooth 
$\A^1$-rigid, smooth curve $C$ and
%$Cand  $\pi: X \rightarrow C \mathrm{\ an\ }\A^1$-fibration,  
$\sC$ be the algebraic space 
%appeared in the factorisation of $\pi$ 
in Theorem \ref{factor-A1-fibration}. Then $\sC$ is also $\A^1$-invariant.
\end{corollary}
The $\A^1$-invariant sheaves are fibrant objects in the $\A^1$-model structure. So, if $\sC$ is $\A^1$-invariant, then $\pi_0^{\A^1}(\sC) \cong \sC$. Thus by Corollary \ref{factorisation A1 equivalence} we have the following:
\begin{corollary} \label{rigid case a1 invariant}
Let $X$ be a smooth affine surface, which admits an $\mathbb{A}^1$-fibration $\pi: X \to C$ onto an  
%$C$ is a smooth 
$\A^1$-rigid, smooth curve $C$. Then $\pi_0^{\A^1}(X) \cong \sC$, where $\sC$ is the algebraic space in Theorem \ref{factor-A1-fibration}. Thus, in particular $\pi_0^{\A^1}(X)$ is $\A^1$-invariant.
\end{corollary}
Since two $\A^1$-invariant sheaves are isomorphic if and only if they are $\A^1$-weakly equivalent, so by
%firstly 
Corollary \ref{factorisation A1 equivalence} and Corollary \ref{a1 invariance of algebraic space}, we can describe the $\A^1$-homotopy types of $X$, if $C$ is an $\A^1$-rigid curve.
% can conclude
%have the following:
\begin{corollary} \label{homotopy types base rigid}
Suppose, $X_i$'s are two $\A^1$-uniruled, smooth, affine surfaces that admit $\A^1$-fibrations $\pi_i: X_i \to C_i$
% and $\pi_2: X_2 \to C_2$ respectively 
onto $\A^1$-rigid smooth curves $C_i$'s respectively, for $i=1,2$. Suppose that $\pi_i$ factors through $\theta_i: X \to \sC_i$, over smooth algebraic spaces $\sC_i$ respectively (as appeared in Theorem \ref{factor-A1-fibration}), for $i=1,2$.
% and $C_2$. Also assume that the the base curves $C_1$ and $C_2$ are $\A^1$-rigid.
%algebraic spaces $\sC_1$ and $\sC_2$ in the factorisation of $\pi_1$ and $\pi_2$ (as in Theorem \ref{}) respectively are $\A^1$-invariant. 
Then $X_1$ and $X_2$ are $\A^1$-weakly equivalent if and only if $\sC_1$ and $\sC_2$ are isomorphic.
\end{corollary}
\begin{remark}
More generally, the conclusion of Corollary \ref{homotopy types base rigid} is true, if we only assume that the algebraic spaces $\sC_1, \sC_2$ are $\A^1$-invariant. Note that by Corollary \ref{not-a1-rigid}, if the base curve $C$ of the $\A^1$-fibration $\pi:X \to C$ is $\A^1$-rigid, then the algebraic space $\sC$ is also $\A^1$-invariant. But there exists $X$ with an $\A^1$-fibration $X \to \A^1_k$ such that the algebraic space $\sC$ (by Theorem \ref{factor-A1-fibration}) is $\A^1$-invariant (Example \ref{rigid example}).
%for $\A^1_k$-connected base curves that is not true, $\sC$ can be $\A^1$-rigid (Example \ref{rigid example}).
%there are two affine surfaces $X_i$ 
%we assume that $X$ 
%admitting the $\A^1$-fibrations $\pi_i: X_i \to C_i$, $i=1,2$ such that the algebraic space $\sC_i$ (as  in Theorem \ref{factor-A1-fibration}) is $\A^1$-invariant. Then the $\A^1$-homotopy type of
\end{remark}
		\begin{example}
			\begin{enumerate}
				\item Let $X$ be the surface
                $$X\equiv\{(x-1)^2z=y(y-1)+(x-1)\}\backslash \{z=y(y-1)-1\}.$$
                $X$ is smooth by the Jacobian criterion and let $\pi=\pr_x$ be the $\A^1$-fibration to $\G_m=\A^1\backslash\{0\}$, an $\A^1$-rigid curve. In this case, the algebraic space $\sC$ is actually a scheme and it is $\G_m$ with double $1$, let us denote it by $\tilde{\G}_m$.\par $\tilde{\G}_m$ is $\A^1$-rigid as for any morphism $f:\A^1\to\tilde{\G}_m$ the composition morphism $$\A^1\xrightarrow{f}\tilde{\G}_m\xrightarrow{q}\G_m$$gives a constant map, thus $f$ must be constant. Here $q$ sends both of the $1$ to same point $1$ in $\tilde{\G}_m$ and identity otherwise. 
				\item Let $X\equiv\{(x-1)^2z=y^2+(x-1)\}\backslash \{z=y^2\} \xrightarrow{\pr_x}\G_m$ be the $\A^1$-fibration. It factors as $X\xrightarrow{\rho}\sC\xrightarrow{p}\G_m$. Then similarly any $f:\A^1\to \sC$ the composition $p \circ f$ is constant. Since $p$ is surjective $f$ must be constant.
			\end{enumerate}
		\end{example}
		
		
\begin{definition} \label{multi-section definition}
Suppose, $\pi: X \to C$ is an $\mathbb{A}^1$-fibration. A morphism $\gamma: C \to X$ is called a horizontal section (or a multi-section) of $\pi$, if the morphism $\pi \circ \gamma: C \to C$ is surjective. The multi-section $\gamma$ is said to be a section of $\pi$, if the morphism $\pi \circ \gamma: C \to C$ is the identity morphism on $C$.
\end{definition}	
\begin{remark} \label{multi-section remark}
\begin{enumerate}
\item A multi-section $\gamma: C \to X$ of an $\A^1$-fibration $\pi: X \to C$, intersects every fiber $\pi^{-1}(x)$ , for all $x \in C$. However, a multi-section does not always intersect every irreducible components of a degenerate fiber (Example \ref{multi-section example}(3)).
\item Suppose, the base of the $\mathbb{A}^1$-fibration $\pi$ is $\mathbb{A}^1_k$.  Then in Definition \ref{multi-section definition}, the surjectivity of $\pi \circ \gamma$ is equivalent to $\pi \circ \gamma$ is non-constant, since the complement of finitely many points (at least one) in $\mathbb{A}^1_k$ is $\mathbb{A}^1$-rigid. If there exists a multi-section $\gamma$ of $\pi: X \to \A^1_k$, then the algebraic space $\mathcal{C}$ in the factorisation of $\pi$ in Theorem \ref{factor-A1-fibration} admits a non-constant morphism $\theta \circ \gamma: \A^1_k \to \sC$, so $\sC$ is not $\mathbb{A}^1$-rigid. In Corollary \ref{multi-section exists}, we will prove the converse.
 \end{enumerate}
\end{remark}
\begin{example} \label{multi-section example}
\begin{enumerate}
\item A multi-section can be thought of a section of an $\mathbb{A}^1$-fibration, in a generalized sense. An (algebraic) line bundle has always a section (zero section). Since $\A^1$-fibration is generically trivial, so an $\A^1$-fibration $\pi:X \to C$ has always a rational section, that is there is a rational map $\gamma: C \dashrightarrow X$ such that $\pi \circ \gamma$ is identity (as rational maps). But an $\A^1$-fibration, or even an $\mathbb{A}^1$-bundle does not always admit a section  (\cite[Example 2]{doub}, see also Example \ref{multi-section example}(4)). 
\item Consider, the $\mathbb{A}^1$-fibration $\pi: X \to \mathbb{A}^1_{\mathbb{C}}$ given by the projection to the $x$-axis, where 
$$X \subset \mathbb{A}^3_\mathbb{C} \text{ is given by } x^nz = P(y) - x, \ P(y) \in \mathbb{C}[y] $$
\text{ is a non-constant polynomial} and $n \geq 2$.
Then the map 
$$\gamma: \mathbb{A}^1_{\mathbb{C}} \to X \text{ defined as } t \mapsto (P(t), t, 0)$$
is a multi-section of $\pi$ and moreover $\text{Im}(\gamma)$ intersects every irreducible component of the fiber $\pi^{-1}(0)$ (which is the only degenerate fiber of $\pi$). Thus both $X$ and the algebraic space $\mathcal{C}$, appeared in the factorisation of $\pi$ (by Theorem \ref{factor-A1-fibration}) are $\mathbb{A}^1$-chain connected. 
\item 	A multi-section of an $\A^1$-fibration $\pi: X \to \A^1_\C$ does not always intersect every irreducible components of the degenerate fibers of $\pi$. Consider the following example:

Let $X$ be the surface given by
$$X \equiv\{f(x,y,z)=x^2(x-1)^2z-y(y-1)^2+x(x-1)=0\} \subseteq \mathbb{A}^3_\mathbb{C} $$
and consider the $\A^1$-fibration
$X \xrightarrow{\pr_x}\A^1_\C$ projection to the $x$-axis.
%is an $\A^1$-fibration. 
The surface $X$ is smooth, since  
\begin{align*}
    \nabla(f)&=(2x(x-1)(2x-1)z+2x-1,3y^2-4y+1,x^2(x-1)^2) \\
    &\neq (0,0,0);
\end{align*}
%$$\neq (0,0,0)$$; 
as when $x=0,1;\ \text{ then } 2x(x-1)(2x-1)z+2x-1=-1,1$ respectively, otherwise $x^2(x-1)^2\neq 0$. The degenerate fibers are over $x=0,1$, given by $y(y-1)^2=0$. So both the fibers $\pi^{-1}(0)$ and $\pi^{-1}(1)$ have two irreducible components: one irreducible component is reduced, given by $y=0$ and other irreducible component is non-reduced of multiplicity $2$, given by $(y-1)^2=0$. The morphism
$$\gamma(t):\A^1_\C \to X \text{ defined as } t \mapsto (t,t(t-1),t^2-t-2)$$ 
is a multi-section of $\pi$.
%at $x(t)=t=0,1$ we have $y(t)=t(t-1)=0$ 
Observe that $\text{Im}(\gamma)$ intersects only the reduced components of the degenerate fibers (that is the component given by $y=0$), since $t(t-1) =0$, if $t = 0 \text{ or }1$. Thus the algebraic space $\sC$ (by Theorem \ref{factor-A1-fibration}) is not $\A^1$-invariant. So, $X$ is $\A^1$-connected (from the proof of Theorem \ref{connected-base-new}).
%making $X,\ \A^1$-connected (from Theorem \ref{connected-base-new}). w
We also know that if $\alpha:\A^1_\C \to X, \ \alpha(t)=(x(t), y(t),z(t))$ is a multi-section , then it cannot intersect the only the non-reduced components of the degenerate fibers, since in that case $x(t),\ (x(t)-1) $ both have to be a square polynomial. \par
	However, we do not know whether $X$ admits a multi-section which intersects every irreducible components of the degenerate fibers. An affirmative answer would imply that $X$ (and also the algebraic space $\sC$) is $\A^1$-chain connected.

\item Consider, a particular example of (2) of an $\mathbb{A}^1$-fibration $\pi: X \to \mathbb{A}^1_k$, the projection to the $x$-coordinate, where $X$
% \subset \mathbb{A}^3_k$ 
is given by $x^2z = y^2-x$. Then $\pi$ does not admit a section (indeed, if $\gamma \equiv (t, y(t), z(t))$ is a section of $\pi$, then we would have $t^2$ divides $(y(t)^2 - t)$ in $k[t]$, which is not possible), but the morphism $\gamma: \mathbb{A}^1_k \to X$ given by $t \mapsto (t^2,t,0)$ is a multi-section of $\pi$. It is natural to ask when an $\mathbb{A}^1$-fibration admits a multi-section. We will also see an example of an $\mathbb{A}^1$-fibration (Example \ref{rigid example}) over the base $\mathbb{A}^1_k$, which it does not admit even a multi-section.
\end{enumerate}
 
\end{example}

We have the following corollary to Lemma \ref{not-a1-rigid} on the existence of a multi-section of an $\mathbb{A}^1$-fibration over $\mathbb{A}^1_k$.
\begin{corollary} \label{multi-section exists}
Suppose, $\pi:X \to \mathbb{A}^1_k$ is an $\mathbb{A}^1$-fibration from an affine surface $X$ and $X \xrightarrow{\theta} \sC \xrightarrow{\alpha} \A^1_k$ is the factorisation of $\pi$ as in Theorem \ref{factor-A1-fibration}, where $\theta$ is an \'etale $\mathbb{A}^1$-bundle. Then the following are equivalent:
\begin{enumerate}
    \item $\pi$ admits a multi-section.
    \item The algebraic space $\mathcal{C}$ is not $\mathbb{A}^1$-invariant.
    \item There exists a non-constant morphism $\gamma: \A^1_k \to \sC$.
\end{enumerate}
%, if and only if 
\end{corollary}
\begin{proof}
The equivalence of (2) and (3) follows from Lemma \ref{not-a1-rigid} and (1) implies (2), by
%If $\pi$ admits a multi-section, then $\sC$ is not $\A^1$-invariant (
Remark \ref{multi-section remark}). Suppose, $\sC$ is not $\A^1$-invariant, by Lemma \ref{not-a1-rigid} there is a non-constant morphism $\gamma: \mathbb{A}^1_k \to \mathcal{C}$. Consider the pullback diagram
\[\xymatrixcolsep{3pc}
			\xymatrix{ \mathbb{A}^1_k \times_\mathcal{C} X \ar[d]_{\theta^*\gamma} \ar[r]_-{\text{pr}} & X \ar[d]^\theta \\
				\mathbb{A}^1_k \ar[r]^{\gamma } &\sC}
			\]
The morphism $\theta^*\gamma: \mathbb{A}^1_k \times_\mathcal{C} X \to \mathbb{A}^1_k$ is an \'etale $\mathbb{A}^1$-bundle. By the same way in the proof of Theorem \ref{sing-etale-bundle}, we have $H^1_{\text{\'et}}(\mathbb{A}^1_k, \text{Aff}_1) = \{0\}$, so any \'etale $\mathbb{A}^1$-bundle over $\mathbb{A}^1_k$ is trivial. So $\theta^*\gamma$ admits a section, that is there is a morphism $s: \mathbb{A}^1_k \to \mathbb{A}^1_k \times_\mathcal{C} X$ such that $\theta^*\gamma \circ s$ is the identity morphism. Suppose, 
$$\Tilde{\gamma}= \text{pr} \circ s: \mathbb{A}^1_k \to X.$$
Then $\Tilde{\gamma}$ is a lift of $\gamma$, that is $\theta \circ \Tilde{\gamma}  = \gamma$, the lower triangle in the diagram commutes:
\[\xymatrixcolsep{3pc}
			\xymatrix{ \mathbb{A}^1_k \times_\mathcal{C} X \ar[d]^{\theta^*\gamma} \ar[r]_-{\text{pr}} & X \ar[d]^\theta \\
				\mathbb{A}^1_k   \ar@/^/[u]^s \ar[r]_{\gamma} \ar[ur]_{\Tilde{\gamma}} & \sC}
			\]
Since, $\gamma$ is non-constant, so the image of $\Tilde{\gamma}$ is not contained in a fiber of $\theta$. Since $\mathbb{A}^1_k \setminus \{\text{finitely many points}\}$ is $\mathbb{A}^1$-rigid, so the morphism $\pi \circ \Tilde{\gamma}$ is surjective. Thus, $\Tilde{\gamma}$ is a multi-section of $\pi$.
% intersects
%Thus the \'etale $\mathbb{A}^1$-bundle $\theta^*\gamma: \mathbb{A}^1_k \times_\mathcal{C} X \to \mathbb{A}^1_k$ is isomorphic to the trivial bundle  over $\mathbb{A}^1_k$ is isomorphic 
\end{proof}
Suppose, if every fiber of the $\mathbb{A}^1$-fibration $\pi: X \to \mathbb{A}^1_k$ is irreducible.
%and the algebraic space  
Then from the construction of $\sC$ (Remark \ref{comments of structure theorem}),
a point of $\mathcal{C}$ corresponds to a point of the base $\mathbb{A}^1_k$, which agains corresponds to a fiber of $\pi$. Thus, if moreover $\mathcal{C}$ is not $\mathbb{A}^1$-invariant, then the non-constant morphism $\gamma: \mathbb{A}^1_k \to \mathcal{C}$ (Lemma \ref{not-a1-rigid}) is surjective on the points and the multi-section $\Tilde{\gamma}: \mathbb{A}^1_k \to X$ (Corollary \ref{multi-section exists}) intersects every fiber of $\pi$. Therefore, in this case both $\mathcal{C}$ and $X$ are $\mathbb{A}^1$-chain connected. We have the following corollary:
\begin{corollary}\label{irred-fibers}
			Suppose, $\pi:X \to \mathbb{A}^1_k$ is an $\mathbb{A}^1$-fibration from a smooth affine surface $X$ such that every fiber of $\pi$ is irreducible.  
%$X \xrightarrow{\theta} \sC \xrightarrow{p} \A^1$ be the factorisation, where $\theta$ is an étale bundle. 
Let $\sC$ be the algebraic space in Theorem \ref{factor-A1-fibration}. If $\sC$ is not $\A^1$-invariant, 
%(therefore, $\pi$ can have only finitely many non-reduced, irreducible fibers).
			Then both $\sC$ and $X$ are $\A^1$-chain connected. 
            %($\mathbf{P}^1 ??$)
%In particular $X$ is $\A^1$-connected
\end{corollary}
\begin{remark}
\begin{enumerate}
   \item In Example \ref{rigid example}, we will see that there exists an $\A^1$-fibration $\pi: X \to \A^1_k$ such that every fiber of $\pi$ is irreducible and the algebraic space $\sC$ (by Theorem \ref{factor-A1-fibration}) is $\A^1$-invariant. So $\pi$ does not admit a multi-section and $X$ is also not $\A^1$-connected, here $\pi_0^{\A^1}(X) \cong \sC$.
 \item   Suppose, $\pi: X \to \P^1_k$ is an $\A^1$-fibration from an affine surface $X$ such that every fiber of $\pi$ is irreducible. If moreover, the algebraic space $\sC$ (by Theorem \ref{factor-A1-fibration}) is not $\A^1$-invariant, then by Lemma \ref{not-a1-rigid}, there is a non-constant morphism $\gamma: \A^1_k \to \sC$, which can miss at most one point of $\sC$ (since $\sC$ is the algebraic space over $\P^1$) and consequently, the morphism $\Tilde{\gamma}: \A^1_k \to X$ (from Lemma \ref{not-a1-rigid}) may not intersect at most one fiber of $\pi$. Thus, we don't know whether $X$ or $\sC$ is $\A^1$-chain connected in this case. But $\sC$ (hence $X$ is also) is $\A^1$-connected (Theorem \ref{connected-base-new}).
 \end{enumerate}
\end{remark}
\begin{example}
Consider again Example \ref{multi-section example}(2) with $P(y) = (y-\alpha)^m$, for some $\alpha \in k$  fixed \cite[Example 5]{doub}
 $$X \equiv x^nz=(y-\alpha)^m-x \text{ and } \pi: X \to \A^1_k \text{ defined as }(x,y,z)\mapsto x.$$
Then, every fiber of $\pi$ is irreducible and $\pi$ admits a multi-section. So the algebraic space $\sC$ in Theorem \ref{factor-A1-fibration} is not $\A^1$-invariant. Thus, both $X$ and $\sC$ are $\A^1$-chain connected.
\end{example}

If every fiber of the $\A^1$-fibration $\pi; X \to C$, where $C \cong \A^1_k \text{ or } \P^1_k$, is reduced, then the algebraic space $\sC$ (in Theorem \ref{factor-A1-fibration}) is a scheme (possibly non-separated), 
%and $\sC$ is 
isomorphic to $C$ with multiple points. Therefore, $\sC$ is $\A^1$-chain connected and for each copies of $\A^1$ in $\sC$ would lift to an $\A^1$ in $X$ (since $\theta: X \to \sC$ is \'etale locally trivial $\A^1$-bundle, by Theorem \ref{factor-A1-fibration}). Thus, $X$ is also $\A^1$-chain connected (see Section \ref{section a1 homotopy types} for the $\A^1$-homotopy types of such surfaces). Hence we have the following corollary:
\begin{corollary}
    Suppose, $\pi:X \to C$, where $C \cong \A^1_k \text{ or } \P^1_k$, is an $\mathbb{A}^1$-fibration from a smooth affine surface $X$ such that every fiber of $\pi$ is reduced and $\sC$ is the algebraic space in Theorem \ref{factor-A1-fibration}. Then both $\sC$ and $X$ are $\A^1$-chain connected.
\end{corollary}
%(\textbf{reduced + A1, P1})
More generally, 
%if the base curve $C$ of the $\A^1$-fibration $\pi: X \to C$ is not $\A^1$-rigid, then 
we have the following result:
 %it follows that the morphism
		%\begin{lemma}\label{irred-fibers}
		%	Suppose, $\pi:X \to \mathbb{A}^1$ is an $\mathbb{A}^1$-fibration and $X \xrightarrow{\theta} \sC \xrightarrow{p} \A^1$ be the factorisation, where $\theta$ is an étale bundle. Moreover assume that $\sC$ is not $\A^1$-rigid and every fiber of $\pi$ is irreducible (therefore, $\pi$ can have only finitely many non-reduced, irreducible fibers).
		%	Then $\sC$ is ${A}^1$-chain connected. In particular $X$ is $\A^1$-connected
		%\end{lemma}
	%	\begin{proof}
%			First we assume every fiber of $\pi$ is irreducible.
		%	Suppose, $\pi$ is factored as 
		%	$$X \xrightarrow{\theta} \mathcal{C} \xrightarrow{p} \mathbb{A}^1.$$
%			If $\mathcal{C}$ is $\mathbb{A}^1$-rigid, then 
%			$$\pi_0^{\mathbb{A}^1}(X) \cong \pi_0^{\mathbb{A}^1}(\mathcal{C}) \cong \mathcal{C}$$
%			is $\mathbb{A}^1$-invariant.
		%	Since $\sC$ is not $\A^1$-rigid, by \cite[Lemma 2.16]{mazza} there is a morphism $H: \mathbb{A}^1_U \to \mathcal{C}$, for some $U \in Sm/k$ such that 
		%	$$H(0,-) \neq H(1,-0):U \to \mathcal{C}.$$
		%	By Lemma \ref{A1 in algebraic space}, there is $u \in U(k)$ such that
		%	$$H(0,u) \neq H(1,u).$$
		%	Consider, the composition of morphisms
		%	$$\gamma: \mathbb{A}^1 \xrightarrow{(Id,u)} \mathbb{A}^1_U \xrightarrow{H} \mathcal{C}$$
	%		is a non-constant morphism. Since $dim(\mathcal{C}) = 1$, so $\gamma$ must be dominant.
			
		%	We claim that $\gamma$ is surjective. Consider, the composition of morphisms
		%	$$\mathbb{A}^1 \xrightarrow{\gamma} \mathcal{C} \xrightarrow{p} \mathbb{A}^1,$$
	%		where $p$ appears in the factorisation of the $\mathbb{A}^1$-fibration $\pi$. The morphism $p \circ \gamma: \mathbb{A}^1 \to \mathbb{A}^1$ is surjective.
			
		%	Suppose, if possible there is $\alpha \in \mathcal{C}(k)$ such that $\alpha \notin Im(\gamma)$. Since, $p \circ \gamma$ is surjective, so there is $t \in \mathbb{A}^1(k)$ such that $p(\gamma(t)) = p(\alpha) = \beta \text{ (say)}$. But the fiber $\pi^{-1}(\beta)$ is irreducible. So,
		%	$$\gamma(t), \alpha \in \theta(\pi^{-1}(\beta)).$$
		%	But $\theta(\pi^{-1}(\beta))$ is just a $k$-point in $\mathcal{C}$. So, $\alpha = \gamma(t)$ and hence $\gamma$ is surjective. So, $\mathcal{C}$ is $\mathbb{A}^1$-connected. Thus $X$ is $\mathbb{A}^1$-connected.
	%	\end{proof}



%\begin{theorem} \label{connected-base}
%			Suppose $\pi: X \rightarrow C$ is an $\A^1$-fibration onto a smooth, rational curve $C$. Then $\pi_0^{\A^1}(X)$ is $\A^1$-invariant.
%		\end{theorem}
%		
%		\begin{proof}
%	%		Following \cite[Theorem 9]{doub}, t
%Following the notations in Remark \ref{comments of structure theorem}, the $\mathbb{A}^1$-fibration $\pi:X \to C$ factorises as:
%		$$X \xrightarrow{\theta} \mathcal{C} \xrightarrow{p} \check{C} \xrightarrow{\check{p}}C,$$ 
%		where,
%		\begin{enumerate}
%		\item 
%%The morphism $p \circ \theta$ is an $\mathbb{A}^1$-fibration, in which every fiber is irreducible. 
%If $\pi$ has reducible fiber, then $\check{C}$ is a scheme obtained by gluing finitely many copies of $C$; so in particular $\check{C}$ is a non-separated scheme, in this case. The points of $\check{C}$ correspond to an irreducible component of a fiber of $\pi$ and every fiber of $p \circ \theta$ is irreducible.
%		\item The morphism $\theta: X \to \mathcal{C}$ is an étale locally trivial $\mathbb{A}^1$-bundle. The points of $\mathcal{C}$ correspond to a point of $\check{C}$, but if $\pi$ has a non-reduced fiber, then $\mathcal{C}$ is an algebraic space of dimension $1$, which is not a scheme.
%		\end{enumerate}
%% (if $\pi$ has a reducible fiber and then )
%%be the factorisation of the $\A^1$-fibration, where $\check{C}$ is a (possibly \textit{non-separated}) scheme. 
%By Corollary \ref{factorisation A1 equivalence}, the morphism $\theta: X \to \mathcal{C}$ is an $\mathbb{A}^1$-weak equivalence; so 
%$$\pi_0^{\mathbb{A}^1}(X) \xrightarrow{\cong} \pi_0^{\mathbb{A}^1}(\mathcal{C}).$$
%Thus if $\sC$ is $\A^1$-rigid, then 
%$$\pi_0^{\mathbb{A}^1}(X) \cong \pi_0^{\A^1}(\sC) \cong \sC \text{ is } \A^1 \text{-invariant}.$$ 
%%so without loss of generality 
%So for the remaining case, we assume that $\sC$ is not $\A^1$-rigid. 
%%Then by 
%%imilar argument as in 
%%Corollary \ref{not a1 rigid}, there is a non-constant morphism $\gamma: \A^1_k \to \sC$. 
%In this case, we will prove that $\mathcal{C}$ is $\mathbb{A}^1$-connected. This would imply that $X$ is $\mathbb{A}^1$-connected.
%
%To prove that $\mathcal{C}$ is $\mathbb{A}^1$-connected, by \cite[Lemma 3.3.6]{ictp} we need to prove that $\pi_0^{\mathbb{A}^1}(\mathcal{C})(\Spec\ F)$ is trivial, for every finitely generated field extension $F/k$. The algebraic space $\mathcal{C}$ of dimension $1$ has only two types of points: the generic point, say $\eta$ having residue field $k(t)$ (since, the morphism $\alpha: \sC \to C$ in Remark \ref{comments of structure theorem} is birational and $C$ is a rational curve) and the $k$-points. Since we have assumed that $\sC$ is not $\A^1$-invariant, so by Corollary \ref{not a1 rigid}, 
%%, so in this case there is 
%$\sC$ admits a non-constant morphism $\gamma: \mathbb{A}^1_k \to \mathcal{C}$. Since $\sC$ is of dimension $1$, so at the level of topological spaces the image of $\gamma$ misses only finitely many $k$-points. Thus, there is an algebraic space $\sC^\prime$, which is open in $\sC$ such that $\gamma$ factors through $\sC^\prime$ 
%$$\A^1_k \xrightarrow{\gamma} \sC^\prime \xrightarrow{\text{open immersion}} \sC$$
%and the map $\gamma: \A^1_k \to \sC^\prime$ is surjective (at the level of topological spaces). We will use the same letter $\gamma: \A^1_k \to \sC$ or $\gamma: \A^1_k \to \sC^\prime$ interchangeably.
%
%% (thus $\gamma$ is dominant), 
%So it is enough to prove that $\pi_0^{\mathbb{A}^1}(\mathcal{C})(\Spec\ k)$ is trivial. Indeed, because of  the non-constant morphism $\gamma: \A^1_k \to \sC$, the points $\eta = x \in \sS(\sC)(\Spec\ k(t))$ and $x=y \in \sS(\sC)(\Spec\ k)$, for any two $k$-points $x,y$ in the image of $\gamma$ in $\sC$. Thus, for any points $y_1, y_2 \in \sC(F)$ such that $y_1, y_2 \in \sC^\prime(F)$, we have that $y_1=y_2 \in \pi_0^{\A^1}(\sC)(\Spec\ F)$. 
%
%We will prove that actually 
%$\mathcal{S}^3(\mathcal{C})(\Spec\ k)$ is trivial. 
%
%Let $\alpha , \beta :\Spec(k)\to \sC$ be two $k$-points of $\sC$. If both $\alpha, \beta \in \sC^\prime(k)$, then $\alpha=\beta \in \sS(\sC)(k)$. Let us assume that both $\alpha$ and $\beta$ are not in $\sC^\prime(k)$.
%%$\alpha \neq \beta \in \mathcal{S}(\mathcal{C})(\Spec\ k)$. 
%In this case, we will prove that $\alpha = \beta \in \sS^3(\sC)(k)$. That is, we will construct an $H \in \mathcal{S}^2(\mathcal{C})(\mathbb{A}^1_k)$ such that $H(0) = \alpha \in \sS^2(\sC)(k)$ and $H(1) = \beta \in \sS^2(\sC)(k)$.
%
%
%%$\alpha$ and $\beta$ are $\mathbb{A}^1$-ghost homotopic (Remark \ref{}). That is, we will construct an $H \in \mathcal{S}(\mathcal{C})(\mathbb{A}^1_k)$ such that $H(0) = \alpha$ and $H(1) = \beta$.
%
%
%% so they are same in $\pi_0^{\A^1}(\sC)$. 
%Since, $\pi$ can only has finitely many degenerate fibers and any two $k$-points in $\sC^\prime$ are naively $\mathbb{A}^1$-homotopic, so 
%%and the image of $\gamma$ can only Thus without loss of generality 
%we can assume that $\alpha \in \sC^\prime(k)$ with
%$$\alpha: \Spec(k)\xrightarrow{i_0} \A^1_k \xrightarrow{\gamma} \sC$$ 
%and the fiber of $\pi$ over $\alpha$ is irreducible and reduced, that is
%$$\pi^{-1}(\check{p} \circ p (\alpha)) \cong \A^1_k.$$ 
%%Since any two $k$-points in $\text{Im}(\gama)$ can be homotopped to a point; over which the fiber is irreducible and reduced via the $\A^1$-chain homotopy $\gamma: \A^1 \to \sC$.\par
%
%%			Following 
%%the notations in Remark \ref{comments of structure theorem},By construction of $\sC$ in \cite[Thm 9]{doub} i
%Suppose, $p(\beta)=\check{\beta}$ and the fiber of $p \circ \theta$ over $\check{\beta}$ is non-reduced. Then, following \cite[Theorem 9]{doub}, there is a germ of smooth curve $p_{\beta}:\tilde{C}_{\beta} \hookrightarrow X$ ($\tilde{C}_\beta$ is a rational curve, since $X$ is rational) intersecting the fiber $(p \circ \theta)^{-1}(\check{\beta})$ transeversally. The composition $\theta_{\beta}:=p \circ \theta \circ p_{\beta}:\tilde{C}_{\beta}\rightarrow\check{C}$ is a quasi-finite morphism onto its image and there is a unique point $\tilde{\beta} \in \tilde{C}_\beta$ such that $\tilde{\beta} \mapsto \check{\beta}$. We can assume that the image $\theta_\beta(\tilde{C}_\beta)$ is an affine open neighbourhood of $\check{\beta}$ and the fiber of $p \circ \theta$ over every point in $\theta_\beta(\tilde{C}_\beta) \setminus \check{\beta}$ is reduced. There can be finitely many points in $\check{C}$ that are outside $\text{Im}(p \circ \gamma)$. So by further shrinking $\tilde{C}_\beta$, we can assume that
%$$\theta_\beta(\tilde{C}_\beta) \setminus \{\check{\beta}\} \subset \text{Im}(p \circ \gamma)$$
%Then taking quotient $\Tilde{C_{\beta}}$ by the \'etale equivalence relation generated by $q_\beta$
%$$x,y \in \tilde{C}_\beta \setminus \{\tilde{\beta}\}, \  x\sim y \mathrm{\ if\ } \theta_{\beta}(x)=\theta_{\beta}(y),$$  
%we get the quotient map $q_{\beta}:\Tilde{C_{\beta}}\to (\Tilde{C_{\beta}}/\sim) =: \sC_{\beta} \subset \sC$ such that $q_\beta(\tilde{\beta}) = \beta$.
%%containing the unique point $\Tilde{\beta} \mapsto \beta$
%%$\Tilde{\beta} \mapsto \check{\beta}$ 
%and over every other point of $\sC_{\beta}$, the fiber of $\pi$ is reduced and irreducible. Also, since 
%$$\theta_\beta \text{ satisfies } \theta_\beta(\tilde{C}_\beta) \setminus \{\check{\beta}\} \subset \text{Im}(p \circ \gamma),$$ 
%so, $q_\beta$ also satisfies
%$$q_\beta(\tilde{C}_\beta \setminus \{\tilde{\beta}\}) \subset \sC^\prime,$$
%in the level of topological spaces.
%This defines a $\A^1$-ghost homotopy as follows $-$
%			
%			\begin{equation}
%				\label{main-ghost-homotopy}
%%		
%\xymatrixcolsep{5pc}
%\xymatrix{{V \times_{\mathbb{A}^1_k} V} \ar@<-.5ex>[r]_<<<<<<<<<<{\pr_2} \ar@<.5ex>[r]^<<<<<<<<<<{\pr_1} &{V:=\tilde{C}_{\beta}\sqcup (\A^1_k\backslash\{0\}}) \ar[dd] \ar[r]^<<<<<<<<<<{H=q_\beta \sqcup \alpha}&\sC \\
%	\\
%	{\mathrm{Spec}(k)} \ar@<-.5ex>[r] \ar@<.5ex>[r] \ar@<-.5ex>[uur]_{i_1} \ar@<.5ex>[uur]^{\tilde{\beta}} &{\A^1_k}}
%			\end{equation}
%		
%			Firstly, note that $X$ is a rational surface (Remark \ref{properties a1 fibration}). So every tangent direction and therefore normal direction of the fiber can be given by germ of a rational curve. Thus we can choose $\tilde{C}_{\beta}$ to be a rational curve and $\phi: \tilde{C}_{\beta}\cong \A^1_k \backslash\{\mathrm{finite\ set\ not\ containing\ }0\}:\ \tilde{\beta}\mapsto 0$. So
%$$V:=\tilde{C}_{\beta} \sqcup (\A^1_k \backslash \{0\}) $$
%is a Nisnevich cover of $\A^1_k$ in the diagram \ref{main-ghost-homotopy} above is actually a Zariski cover given by their corresponding open immersions. Suppose, $H : V \to \mathcal{C}$ is given by 
%$$H|_{\tilde{C}_\beta} = q_{\beta}:\Tilde{C_{\beta}}\to (\Tilde{C_{\beta}}/\sim) =: \sC_{\beta} \subset \sC $$
%and $H|_{\mathbb{A}^1_k \setminus \{0\}}$ is the constant morphism to $\alpha$. The Zariski covering $V \to \mathbb{A}^1_k$ has the lifts $\Tilde{\beta} \in \Tilde{C}_\beta \subset V$ and $1 \in \mathbb{A}^1_k \setminus \{0\} \subset V$ of $0,1 \in \mathbb{A}^1_k$ respectively. And $H$ maps $\Tilde{\beta}$ to $\beta$ and $1$ to $\alpha$.
%Now,
%$$V \times_{\mathbb{A}^1_k} V = \Tilde{C}_\beta \sqcup (\mathbb{A}^1_k \setminus\{0\}) \sqcup (\Tilde{C}_\beta \times_{\mathbb{A}^1_k} \mathbb{A}^1_k \setminus\{0\}) \sqcup (\mathbb{A}^1_k \setminus\{0\} \times_{\mathbb{A}^1_k} \Tilde{C}_\beta)$$
%and there are two maps
%$$\text{pr}_1, \text{pr}_2: V \times_{\mathbb{A}^1_k} V \to V.$$
%To give an $\mathbb{A}^1$-ghost homotopy between $\alpha$ and $\beta$ (that would be $H \in \mathcal{S}^2(\mathcal{C})(\mathbb{A}^1_k)$) we need to give an $\mathbb{A}^1$-ghost homotopy in $\sS(\sC)$ between 
%$$H \circ \text{pr}_1 \text{ and } H \circ \text{pr}_2, \text{ upto some Nisnevich covering of }V \times_{\mathbb{A}^1_k} V .$$
%Now, the morphisms $\text{pr}_1$ and $\text{pr}_2$ are same to the connected components $\Tilde{C}_\beta$ and $\mathbb{A}^1_k \setminus\{0\}$. The other two connected components are both isomorphic to $\tilde{C}_{\beta}\backslash\{\tilde{\beta}\}$. But on this component, $H \circ \text{pr}_1$ factors as
%$$\tilde{C}_{\beta}\backslash\{\tilde{\beta}\} \hookrightarrow \Tilde{C}_\beta \xrightarrow{q_\beta} \mathcal{C}$$
%and $H \circ \text{pr}_2$ factors as
%$$\tilde{C}_{\beta}\backslash\{\tilde{\beta}\} \xrightarrow{\phi \circ \text{ inclusion}} \mathbb{A}^1_k \setminus \{0\} \xrightarrow{\alpha} \mathcal{C}.$$
%Since, $q_\beta(\tilde{C}_\beta \setminus \{\tilde{\beta}\}) \subset \sC^\prime$, so both $H \circ \text{pr}_1|_{\tilde{C}_{\beta}\backslash\{\tilde{\beta}\}}$ and $H \circ \text{pr}_2|_{\tilde{C}_{\beta}\backslash\{\tilde{\beta}\}}$ factor through $\sC^\prime$.
%
%%and similarly for $H \circ \text{pr}_2$. Therefore, we are yet to give 
%Now we will produce an $\mathbb{A}^1$-ghost homotopy in $\sS(\sC)$ between
%%$$H \circ \text{pr}_1|_{\tilde{C}_{\beta}\backslash\{\tilde{\beta}\}} \text{ and }H \circ \text{pr}_2|_{\tilde{C}_{\beta}\backslash\{\tilde{\beta}\}}.$$
%$$H \circ \text{pr}_1|_{W} \text{ and }H \circ \text{pr}_2|_{W},$$
%for some Nisnevich covering $W \to \tilde{C}_{\beta}\backslash\{\tilde{\beta}\}$. 
%%This can be done because $q_\beta(\tilde{C}_\beta \setminus \{\tilde{\beta}\}) \subset \sC^\prime$ and $\sS^2(\sC^\prime) \cong \Spec\ k$ \cite[]{}.
%The non-constant morphism $\gamma: \mathbb{A}^1_k \to \mathcal{C}^\prime$ gives surjective morphism on field valued points, that is for every $F/k$ finitely generated field extension, the map
%$$\gamma: \A^1_k(\Spec\ F) \to \sC^\prime(\Spec\ F)$$
%is surjective. Consider the following commutative diagram of sets (the field valued sections of Nisnevich sheaves):
%%. Consider the commutative diagram in the category of étale sheaves in $Sm/k$
% \[
%\xymatrix{
%\mathbb{A}^1_k(\Spec\ F) \ar[r]^\gamma \ar[d] & \mathcal{C}^\prime(\Spec\ F) \ar[d] \\
%\mathcal{S}(\mathbb{A}^1_k)(\Spec\ F) \ar[r]_{\mathcal{S}(\gamma)} & \mathcal{S}(\mathcal{C}^\prime)(\Spec\ F)
%}
%\]
%The both vertical maps and the upper horizontal map are surjective. Therefore, the map
%$$\mathcal{S}(\mathbb{A}^1_k)(\Spec\ F) \to \mathcal{S}(\mathcal{C}^\prime)(\Spec\ F)$$
%is also an surjective. Since, $\mathcal{S}(\mathbb{A}^1_k) \cong \Spec\ k$, so
%$\mathcal{S}(\mathcal{C}^\prime)(\Spec\ F)$ is trivial, for every finitely generated field extension $F/k$. Therefore, by \cite[Theorem 3.2]{naive connectedness sawant}, $\sS^2(\sC^\prime) \cong \Spec\ k$.
%
%Now $H \circ \text{pr}_1|_{\tilde{C}_{\beta}\backslash\{\tilde{\beta}\}}, H \circ \text{pr}_2|_{\tilde{C}_{\beta}\backslash\{\tilde{\beta}\}} \in \sS^2(\sC^\prime)(\tilde{C}_{\beta}\backslash\{\tilde{\beta}\})$. Since, $\sS^2(\sC^\prime)$ is the trivial sheaf, so there is a Nisnevich covering $W \to \tilde{C}_{\beta}\backslash\{\tilde{\beta}\}$ and there is a chain of homotopies 
%$$G = (G_0,\dots,G_r):\A^1_W\to  \sS(\sC^\prime)$$
%%(since $\mathcal{C}$ is of dimension $1$ and $\gamma$ is non-constant, so $\text{Im}(\gamma)$ can be considered as an open subspace in $\mathcal{C}$) 
%such that 
%$$G_0(0, -) = H \circ \text{pr}_1|_{W} \text{ and } G_r(1, -) = H \circ \text{pr}_2|_{W}.$$
%Since, $\sC^\prime$ is an open subspace of $\sC$, so this $G$ is a homotopy in $\sS(\sC)$
%$$G = (G_0,\dots,G_r):\A^1_W\to  \sS(\sC^\prime) \to \sS(\sC)$$
%%(since $\mathcal{C}$ is of dimension $1$ and $\gamma$ is non-constant, so $\text{Im}(\gamma)$ can be considered as an open subspace in $\mathcal{C}$) 
%such that 
%$$G_0(0, -) = H \circ \text{pr}_1|_{W} \in \sS(\sC)(W) \text{ and } G_r(1, -) = H \circ \text{pr}_2|_{W} \in \sS(\sC)(W).$$
%%\cong \Spec\ k.$$
%Therefore, this defines $\A^1$-ghost homotopy in $\sS^2(\sC)$ between $\alpha$ and $\beta$. 
%Hence, $\alpha = \beta \in \sS^3(\sC)(k)$. 
%%Therefore, $\sS^3(\sC)(\Spec\ k)$ is trivial. This completes the proof.
%
%%Thus any two sections in $\mathcal{C}^\prime(W)$ are $\mathbb{A}^1$-chain homotopic, after some Nisnevich covering.
%
%
%
%%This is immediate, because by shrinking $\tilde{C}_{\beta}$ if necessary; we can assume that
%%$$\text{Im}(q_\beta) \subseteq \text{Im}(\gamma)$$
%%and therefore, both $H \circ \text{pr}_1|_{\tilde{C}_{\beta}\backslash\{\tilde{\beta}\}}$ and $H \circ \text{pr}_2|_{\tilde{C}_{\beta}\backslash\{\tilde{\beta}\}}$ factor through $\text{Im}(\gamma)$. So there is a Nisnevich covering $W \to \tilde{C}_{\beta}\backslash\{\tilde{\beta}\}$ and there is a chain of $\mathbb{A}^1$-homotopies 
%%$$G = (G_0,\dots,G_r):\A^1_W\to \text{Im}(\gamma) \to \sC$$
%%(since $\mathcal{C}$ is of dimension $1$ and $\gamma$ is non-constant, so $\text{Im}(\gamma)$ can be considered as an open subspace in $\mathcal{C}$) such that 
%%$$G_0(0, -) = H \circ \text{pr}_1|_{W} \text{ and } G_r(1, -) = H \circ \text{pr}_2|_{W}.$$
%%where, $W \to \tilde{C}_{\beta}\backslash\{\tilde{\beta}\}$ is.
%%and
%%$$G_r(1,-) = H \circ \text{pr}_2|_{\tilde{C}_{\beta}\backslash\{\tilde{\beta}\}}.$$
%
%%Indeed, denoting $\text{Im}(\gamma) = \mathcal{C}^\prime$, we have non-constant morphism $\gamma: \mathbb{A}^1_k \to \mathcal{C}^\prime$. Consider the commutative diagram in the category of étale sheaves in $Sm/k$
%% \[
%%\xymatrix{
%%\mathbb{A}^1_k \ar[r]^\gamma \ar[d] & \mathcal{C}^\prime \ar[d] \\
%%\mathcal{S}(\mathbb{A}^1_k) \ar[r]_{\mathcal{S}(\gamma)} & \mathcal{S}(\mathcal{C}^\prime)
%%}
%%\]
%%The both vertical maps and the upper horizontal map are epimorphisms. Therefore, the map
%%$$\mathcal{S}(\mathbb{A}^1_k) \to \mathcal{S}(\mathcal{C}^\prime)$$
%%is also an epimorphism. Since, $\mathcal{S}(\mathbb{A}^1_k) \cong \Spec\ k$, so
%%$$\mathcal{S}(\mathcal{C}^\prime) \cong \Spec\ k.$$
%%Thus any two sections in $\mathcal{C}^\prime(W)$ are $\mathbb{A}^1$-chain homotopic, after some Nisnevich covering.
%
%%For example, any morphism $\phi: W \to \mathcal{C}$ such that $\text{Im}(\phi) \subseteq \text{Im}(\gamma)$ is homotopic to the constant morphism $\psi: W \xrightarrow{\alpha = \gamma(0)} \mathcal{C}$. The homotopy is given by
%
%%\
%
%
%% secondly, $W$ is a connected component of $V\times_{\A^1}V$ thus a nisnevich cover; $f_0$ is inclusion and $f_1$ is $\phi$ composed with inclusion. Now $q_{\beta}\circ f_0(W)\subset \im(\gamma)$ as $\im(\gamma)$ is dense and $\beta \notin (q_{\beta}\circ f_0)(W)$ and we can shrink $\tilde{C}_{\beta}$ if necessary; and $\alpha:\A^1_k\to \sC$ is constant map to $\alpha \in \im(\gamma)$ so $\alpha \circ f_1$ is also constant. Im($\gamma$) is $\A^1$-chain conneced so we have a $\A^1$-chain homotopy $$H = (h_0,\dots,h_r):\A^1_W\to \im(\gamma)\subset \sC$$ connecting $f_0\circ q_{\beta}\mathrm{\ and \ } f_1\circ \alpha$. This is a $\A^1$-ghost homotopy between $\alpha\mathrm{\ and \ }\beta $, thus $\alpha=\beta\in \sS^2(\sC)$. 
%
%
%If the fiber over $\beta$ is reduced, then the fiber of $p \circ \theta$ over $p(\beta)$ is irreducible and reduced, hence
%$$(p \circ \theta)^{-1}(p(\beta)) \cong \theta^{-1}(\beta) \cong \mathbb{A}^1_k.$$
%So around $\beta$, the morphism $\theta$ admits a local section and thus the rational curve $\Tilde{C}_\beta$ passing through $\Tilde{\beta}$ intersects a fiber $\theta$ at most one point. Thus from the construction of $\mathcal{C}$, the rational curve $\Tilde{C}_\beta$ is embedded as an open subscheme of $\mathcal{C}$.
%%  from the construction of $$
%
%Then it is a scheme point of $\sC$ and in that case $\sC_{\beta}$ is actually a scheme. Then the same construction as before would prove $\alpha = \beta \in \sS^3(\sC)(k)$. Therefore, $\sS^3(\sC)(\Spec\ k)$ is trivial. 
%%This completes the proof.
%%give an $\A^1$-ghost homotopy between $\alpha$ and $\beta$. So, $\sS^2(\sC)(\Spec\ k)$ is trivial.
%% but that does not change anything in the constructed $\A^1$-ghost homotopy. 
%Hence, $\sC$ is $\A^1$-connected and so, $X$ is also $\A^1$-connected.
%			\end{proof}
			
			\begin{theorem}\label{connected-base-new}
            Let $X$ be a smooth affine surface, which admits an $\A^1$-fibration
%Suppose 
$\pi: X \rightarrow C$ is an onto a smooth, rational curve $C$. Then $\pi_0^{\A^1}(X)$ is $\A^1$-invariant.
			\end{theorem}
		\begin{proof}
			Following the notations in Theorem \ref{factor-A1-fibration} and Construction \ref{construction of algebraic space}
            %remark \ref{comments of structure theorem}
            , the $\mathbb{A}^1$-fibration $\pi:X \to C$ factorises as:
			$$X \xrightarrow{\theta} \mathcal{C} \xrightarrow{p} \check{C} \xrightarrow{\check{p}}C.$$ 
			In this factorisation $\theta:X\to\sC$ is a \'etale $\A^1$-bundle over the algebraic space $\sC$, which is not a scheme if $\pi$ has a non-reduced fiber. Each point of $\check{C}$ corresponds to an irreducible component of a fiber of $\pi$. 
            %So $\check{C}$ is a non-separated curve, if $\pi$ Therefore, if every fiber of $\pi$ is irreducible, then $\check{C} \xrightarrow{\check{p}} C$ is an isomorphism 
            If every fiber of $\pi$ is reduced, then $\sC\xrightarrow{p}\check{C}$ is an isomorphism.
            %and $\check{C}$ is a (possibly non-separated) scheme.
			By Corollary \ref{factorisation A1 equivalence}, the morphism $\theta: X \to \mathcal{C}$ is an $\mathbb{A}^1$-weak equivalence; therefore 
			$$\pi_0^{\mathbb{A}^1}(X) \xrightarrow{\cong} \pi_0^{\mathbb{A}^1}(\mathcal{C}).$$
			Thus if $\sC$ is $\A^1$-invariant, then 
			$$\pi_0^{\mathbb{A}^1}(X) \cong \pi_0^{\A^1}(\sC) \cong \sC \text{ is also } \A^1 \text{-invariant}.$$ 
			For the remaining case, assume that $\sC$ is not $\A^1$-invariant. In this case, we will prove that $\mathcal{C}$ is $\mathbb{A}^1$-connected. This would imply that $X$ is also $\mathbb{A}^1$-connected.

To prove that $\mathcal{C}$ is $\mathbb{A}^1$-connected, by \cite[Lemma 3.3.6]{ictp} we need to prove that $\pi_0^{\mathbb{A}^1}(\mathcal{C})(Spec \ F)$ is trivial, for every finitely generated field extension $F/k$. The algebraic space $\mathcal{C}$ of dimension $1$ has only two types of points: the generic point, say $\eta$ having residue field $k(t)$ (the morphism $\alpha = \check{p} \circ p: \sC \to C$ is birational and $C$ is a rational curve) and the $k$-points. Since we have assumed that $\sC$ is not $\A^1$-invariant, so by Lemma \ref{not-a1-rigid}, 
%, so in this case there is 
there is a non-constant morphism $\gamma: \mathbb{A}^1_k \to \mathcal{C}$. Since $\sC$ is of dimension $1$, so $\text{Im}(\gamma)$ contains the generic point $\eta$. Thus, for any $x \in \sC(k)$ such that $x \in \text{Im}(\gamma)$, we have 
$$\eta = x \in \sS(\sC)(k(t)).$$
%Hence, $\sS(\sC)(k(t))$ is trivial. 
Hence to prove $\sC$ is $\A^1$-connected, we only need to show $\pi_0^{\A^1}(\sC)(k)$ is trivial. We will actually prove that $\sS^2(\sC)(k)$ is trivial, that is any two points 
$\alpha,\beta\in\sC(k)$ we will prove that $\alpha =\beta\in\sS^2(\sC)(k)$. If both the points $\alpha, \beta \in \text{Im}(\gamma)$, then $\alpha = \beta \in \sS(\sC)(k)$. So without loss of generality, we can assume that $\alpha \in \mathrm{Im}(\gamma)$ and $\beta \notin \mathrm{Im}(\gamma)$. Suppose, $p(\beta) = \check{\beta} \in \check{C}$.

Since $\gamma$ is non-constant and $\sC$ is of dimension $1$, so $\text{Im}(\gamma)$ misses only finitely many $k$-points of $\sC$.


%so at the level of topological spaces the image of $\gamma$ misses only finitely many $k$-points. Thus, there is an algebraic space $\sC^\prime$, which is open in $\sC$ such that $\gamma$ factors through $\sC^\prime$ 
%$$\A^1_k \xrightarrow{\gamma} \sC^\prime \xrightarrow{\text{open immersion}} \sC$$
%and the map $\gamma: \A^1_k \to \sC^\prime$ is surjective (at the level of topological spaces). We will use the same letter $\gamma: \A^1_k \to \sC$ or $\gamma: \A^1_k \to \sC^\prime$ interchangeably.

% (thus $\gamma$ is dominant), 
%So it is enough to prove that $\pi_0^{\mathbb{A}^1}(\mathcal{C})(Spec \ k)$ is trivial. Indeed, because of  the non-constant morphism $\gamma: \A^1_k \to \sC$, the points $\eta = x \in \sS(\sC)(Spec \ k(t))$ and $x=y \in \sS(\sC)(Spec \ k)$, for any two $k$-points $x,y$ in the image of $\gamma$ in $\sC$. Thus, for any points $y_1, y_2 \in \sC(F)$ such that $y_1, y_2 \in \sC^\prime(F)$, we have that $y_1=y_2 \in \pi_0^{\A^1}(\sC)(Spec \ F)$. 

%We will prove that actually 
%$\mathcal{S}^3(\mathcal{C})(Spec \ k)$ is trivial. 

%Let $\alpha , \beta :\Spec(k)\to \sC$ be two $k$-points of $\sC$. If both $\alpha, \beta \in \sC^\prime(k)$, then %$\alpha=\beta \in \sS(\sC)(k)$. Let us assume that both $\alpha$ and $\beta$ are not in $\sC^\prime(k)$.
%$\alpha \neq \beta \in \mathcal{S}(\mathcal{C})(Spec \ k)$. 
%In this case, we will prove that $\alpha = \beta \in \sS^3(\sC)(k)$. That is, we will construct an $H \in \mathcal{S}^2(\mathcal{C})(\mathbb{A}^1_k)$ such that $H(0) = \alpha \in \sS^2(\sC)(k)$ and $H(1) = \beta \in \sS^2(\sC)(k)$.

    
            
            
  %          In this case, we will prove that $\sC$ is $\A^1$-connected. By Lemma \ref{not-a1-rigid}, there is a non-constant morphism $\gamma: \A^1_k \to \sC$. Since $\sC$ is an algberaic space of dimension $1$, so $\gamma$ is a dominant morphism. Let $\eta$ be the generic point of $\sC$. Since $\gamma$ is dominant, therefore $\eta \in \mathrm{Im}(\gamma)$. $\eta$ is a $k(t)$ point of $\sC$, and for any any closed point $x$ of $\sC$, $\eta=x\in \sS(\sC)(k(t))$, as $k$ is an algebraically closed field every other point of $\sC$ is a $k$ point. Hence, to prove $\sC$ is $\A^1$-connected we only need to show $\pi_0^{\A^1}(\sC)(k)$ is trivial. We will actually proof $\sS^2(\sC)$ is trivial, by showing for any two closed point $\alpha,\beta\in\sC(k)$ we have $\alpha =\beta\in\sS^2(\sC)(k)$. Without loss of generality let us assume $\alpha \in \mathrm{Im}(\gamma)$ and $\beta \notin \mathrm{Im}(\gamma)$.\par
			Following construction of the algebraic space $\sC$ as in \cite[Theorem 9]{doub}, 
           % for any point $\check{\beta}\in\check{C}$ 
            there is a germ of smooth curve $\tilde{C}_\beta\hookrightarrow X$ intersecting the fiber $(p \circ \theta)^{-1}(\check{\beta})$
           % over $\check{\beta}$ 
            transversally. Since $X$ is a rational variety, therefore we can choose $\tilde{C}_\beta$ to be a rational curve \cite[Proposition 3.4]{bhs1}. 
            %We choose an open immersion $j:\tilde{C}_\beta \hookrightarrow \A^1_k, $ with there is $\tilde{\beta} \in \tilde{C}_\beta$ such that $j(\tilde{\beta}) = 0$.
             %$\tilde{\beta}\mapsto 0$.
            %There is some open subset $C_\beta$ given by the open immersion $j: C_\beta \hookrightarrow \A^1_k$ and let $\phi: C_\beta \to \tilde{C}_\beta \hookrightarrow X$ be the rational curve  Let $\beta$ be the unique point of $\sC$ such that, $p(\beta)=\check{\beta}$. 
            %Let $q= \theta\circ\phi: \tilde{C}_\beta \to \sC$. 
            By shrinking $\tilde{C}_\beta$ if necessary, we can assume the following:
			\begin{enumerate}
			\item There exist a unique point $\tilde{\beta}\in \tilde{C}_\beta$ such that $\theta(\tilde{\beta}) = \beta.$ 
			\item The morphism $\phi:= \tilde{C}_\beta \xrightarrow{\theta} \sC \xrightarrow{p} \check{C}$
            %$\phi:\tilde{C}_\beta\to X \to \check{C}$ 
            is quasi finite onto its image.
			\item Im$(\phi)$ is an affine, open neighbourhood of $\check{\beta}$ in $\check{C}$.
			\item For a point $y \in \text{Im}(\phi)$, if $y \neq \check{\beta}$, then the fiber $(p \circ \theta)^{-1}(y)$ is reduced and irreducible.
          %  $\check{\beta}$ is the only point of Im$(\phi)$ with degenerate fiber.  
			\item $U:=\theta(\tilde{C}_\beta) \setminus \{\beta\} \subseteq \text{Im}(\gamma)$, since $\text{Im}(\gamma)$ misses only finitely many $k$-points of $\sC$.
%            Im$(q)\backslash \beta\subset \mathrm{Im}(\gamma)$
            \item The \'etale $\A^1$-bundle $\theta: \theta^{-1}(U) \to U$ is the trivial $\A^1$-bundle over $U$, since the $\A^1$-fibration $\pi$ is generically trivial.
			\end{enumerate}
			
 	
 		
 %	Let $U:= q^{-1}(\mathrm{Im}(q)\backslash \beta)$(??). Clearly the étale $\A^1$-bundle is trivial over $U$.
 %Before giving the $\A^1$-ghost homotopy between $\alpha$ and $\beta$, first observe
 By (5), $\theta(\tilde{C}_\beta \setminus \{\tilde{\beta}\}) \subseteq \text{Im}(\gamma)$. We claim that there is a morphism 
 $$\tilde{\theta}: \tilde{C}_\beta \setminus \{\tilde{\beta}\} \to \A^1_k$$
 such that the restriction $\theta: \tilde{C}_\beta \setminus \{\tilde{\beta}\} \to \sC$ factors through $\tilde{\theta}$, that is
 $$\theta|_{\tilde{C}_\beta \setminus \{\tilde{\beta}\}} = \gamma \circ \tilde{\theta}.$$
% , therefore 
 %$$\theta|_{\tilde{C}_\beta \setminus \{\tilde{\beta}\}}: \tilde{C}_\beta \setminus \{\tilde{\beta}\} \to \sC$$ 
% is dominant. So there is a morphism $\tilde{\theta}: \tilde{C}_\beta \setminus \{\tilde{\beta}\} \to \A^1_k$ such that  the The restriction $\theta: \tilde{C}_\beta \setminus \{\tilde{\beta}\} \to \sC$ factors through $\$
Indeed following the proof of Lemma \ref{not-a1-rigid}, $\gamma$ is constructed as follows
$$\gamma= \theta \circ \tilde{\gamma}, \  \A^1_k\xrightarrow{\tilde{\gamma}} X \xrightarrow{\theta} \sC ,$$
such that $\tilde{\gamma}|_{\gamma^{-1}(U)}$ is the following composition of morphisms
$$\gamma^{-1}(U)\xrightarrow{s_0} \gamma^{-1}(U)\times_\sC X \to U \times_\sC X\xhookrightarrow{\text{open}} X,$$
where $s_0$ is the zero section of the morphism $\gamma^{-1}(U)\times_\sC X \to \gamma^{-1}(U)$ (denote it by $\theta_\gamma$), which is the trivial $\A^1$-bundle.
Consider the following commutative diagram. 
 	\[
 	\xymatrix{
 		\gamma^{-1}(U)\times_\sC X \ar[d]^{{\theta}_\gamma} \ar[r] & U\times_\sC X \ar[d]^{{\theta}_U} \ar[r] & X \ar[d]^{\theta} &  \\
 		\gamma^{-1}(U) \ar@/^/[u]^{s_0} \ar@{->>}[r] & U \ar@/^/[u]^{s_0} \ar@{^{(}->}[r] & \sC & \A^1_k \ar[ul]_{\tilde{\gamma}} \ar[l]^\gamma  }     
 	\]
    Here both the morphisms $\theta_\gamma:\gamma^{-1}(U)\times_\sC X \to \gamma^{-1}(U)$ and $\theta_U: U \times_\sC X \to U$ are the trivial $\A^1$-bundles and their corresponding zero sections are denoted by $s_0$ (by shrinking $\tilde{C}_\beta$, we can assume $\gamma^{-1}(U) \to U$ is surjective).
    
 From the diagram, we have Im$(\tilde{q}:U\xrightarrow{s_0}U\times_\sC X \to X)\subset \mathrm{Im}(\tilde{\gamma}) $.
 %and $q=\rho\circ\tilde{q}$. 
 Since $X$ is affine, so Im$(\tilde{\gamma})$ is closed in $X$ and the normalisation of $\text{Im}(\tilde{\gamma})$ is $\A^1_k$. Therefore $\tilde{q}$	factors through the normalisation map $\A^1_k \to \text{Im}(\tilde{\gamma})$. Thus, there exists a morphism $q^\prime: U \to \A^1_k$ such that $\tilde{q}: U \to X$ is the composition of the morphisms
 $$U\xrightarrow{q'}\A^1_k \xrightarrow{\text{normalisation}}\mathrm{Im}(\tilde{\gamma}) \hookrightarrow X.$$
 Consequently, 
 $$\theta|_{\tilde{C}_\beta \setminus \{\tilde{\beta}\}} = \gamma \circ \tilde{\theta}, \text{ where } \tilde{\theta}:=q^\prime \circ \theta|_{\tilde{C}_\beta \setminus \{\tilde{\beta}\}}.$$
 %, where $n$ is the normalisation map. 
 As $\tilde{C}_\beta$ is a rational curve, we choose an open immersion $j:\tilde{C}_\beta \hookrightarrow \A^1_k, $ such that $\tilde{\beta}\mapsto 0$. Let us denote the canonical immersion $\A^1_k\backslash\{0\}\hookrightarrow\A^1_k$ by $i$. 
 %Now we will give an $\A^1$-ghost homotopy between $\alpha$ and $\beta$ as follows.

 
 %Then $U \subset \tilde{C}_\beta \times_{\A^1_k} (\A^1_k\backslash \{0\})$. 
 Now the $\A^1$-ghost homotopy between $\alpha$ and $\beta$ is as follows:
 \begin{equation}
 	\label{new-main-ghost-homotopy}
 	%		
 	\xymatrixcolsep{5pc}
 	\xymatrix{V \times_{\mathbb{A}^1_k} V \ar@<-.5ex>[r]_-{\pr_2} \ar@<.5ex>[r]^-{\pr_1} &{V:=\tilde{C}_{\beta}\sqcup (\A^1_k\backslash\{0\}}) \ar[dd]^{(j,i)} \ar[r]^-{H=\theta|_{\tilde{C}_\beta } \sqcup \alpha}&\sC \\
 		\\
 		{\mathrm{Spec}(k)} \ar@<-.5ex>[r] \ar@<.5ex>[r] \ar@<-.5ex>[uur]_{i_1} \ar@<.5ex>[uur]^{\tilde{\beta}} &{\A^1_k}}
 \end{equation} 


In the diagram \ref{new-main-ghost-homotopy}, $(j,i): V \to \A^1_k$ is 
is a Nisnevich cover of $\A^1_k$, where
$$V:=\tilde{C}_{\beta} \sqcup (\A^1_k \backslash \{0\}),  $$
along with the open immersions $j$ and $i$ (so it is actually a Zariski covering).
%in the diagram \ref{main-ghost-homotopy} above is actually a Zariski cover 
%given by their corresponding open immersions $j:$. 
Suppose, $H : V \to \mathcal{C}$ is given by the restriction $\theta: \tilde{C}_\beta \to \sC$ and
%$$q_{\beta}:\Tilde{C_{\beta}}\to (\Tilde{C_{\beta}}/\sim) =: \sC_{\beta} \subset \sC $$
and $H|_{\mathbb{A}^1_k \setminus \{0\}}$ is the constant morphism to $\alpha$. The Zariski covering $V \to \mathbb{A}^1_k$ has the lifts $\Tilde{\beta} \in \Tilde{C}_\beta$ and $1 \in \mathbb{A}^1_k \setminus \{0\}$ of $0$ and $1$ in $\mathbb{A}^1_k$ respectively. And $H$ maps $\Tilde{\beta}$ to $\beta$ and $1$ to $\alpha$.
Now,
$$V \times_{\mathbb{A}^1_k} V = \Tilde{C}_\beta \sqcup (\mathbb{A}^1_k \setminus\{0\}) \sqcup (\Tilde{C}_\beta \times_{\mathbb{A}^1_k} \mathbb{A}^1_k \setminus\{0\}) \sqcup (\mathbb{A}^1_k \setminus\{0\} \times_{\mathbb{A}^1_k} \Tilde{C}_\beta)$$
and there are two morphisms
$$\text{pr}_1, \text{pr}_2: V \times_{\mathbb{A}^1_k} V \to V.$$
To give an $\mathbb{A}^1$-ghost homotopy between $\alpha$ and $\beta$ (that would be $H$) we need to give a naive $\mathbb{A}^1$-homotopy between 
$$H \circ \text{pr}_1 \text{ and } H \circ \text{pr}_2.$$
Now, the morphisms $\text{pr}_1$ and $\text{pr}_2$ are same to the connected components $\Tilde{C}_\beta$ and $\mathbb{A}^1_k \setminus\{0\}$. The other two connected components are both isomorphic to $\tilde{C}_{\beta}\backslash\{\tilde{\beta}\}$. But $H \circ \text{pr}_1$ factors as
$$\tilde{C}_{\beta}\backslash\{\tilde{\beta}\} \hookrightarrow \Tilde{C}_\beta \xrightarrow{\theta} \mathcal{C}$$
and $H \circ \text{pr}_2$ factors as
$$\tilde{C}_{\beta}\backslash\{\tilde{\beta}\} \xrightarrow{j} \mathbb{A}^1_k \setminus \{0\} \xrightarrow{\alpha} \mathcal{C}.$$
But we have proved that $\theta|_{\tilde{C}_\beta \setminus \{\tilde{\beta}\}}$ factors as
$$\tilde{C}_\beta \setminus \{\tilde{\beta}\} \xrightarrow{\tilde{\theta}} \A^1_k \xrightarrow{\gamma} \sC.$$
So there the $\mathbb{A}^1$-homotopy
$$G:\A^1_{\tilde{C}_\beta \setminus \{\tilde{\beta}\}} \to \sC$$
defined as $G = t (H \circ \text{pr}_1) + (1-t) (H \circ \text{pr}_2)$. Thus $H$ defines an $\A^1$-ghost homotopy between $\alpha$ and $\beta$. So, $\alpha = \beta \in \sS^2(\sC)(k)$.


%and similarly for $H \circ \text{pr}_2$. Therefore, we are yet to give a naive $\mathbb{A}^1$-homotopy between
%$$H \circ \text{pr}_1|_{\tilde{C}_{\beta}\backslash\{\tilde{\beta}\}} \text{ and }H \circ \text{pr}_2|_{\tilde{C}_{\beta}\backslash\{\tilde{\beta}\}}.$$
%$$H \circ \text{pr}_1|_{W} \text{ and }H \circ \text{pr}_2|_{W},$$
%for some Nisnevich covering $W \to \tilde{C}_{\beta}\backslash\{\tilde{\beta}\}$.
%This is immediate, because by shrinking $\tilde{C}_{\beta}$ if necessary; we can assume that
%$$\text{Im}(q_\beta) \subseteq \text{Im}(\gamma)$$
%and therefore, both $H \circ \text{pr}_1|_{\tilde{C}_{\beta}\backslash\{\tilde{\beta}\}}$ and $H \circ \text{pr}_2|_{\tilde{C}_{\beta}\backslash\{\tilde{\beta}\}}$ factor through $\text{Im}(\gamma)$. So there is a Nisnevich covering $W \to \tilde{C}_{\beta}\backslash\{\tilde{\beta}\}$ and there is a chain of $\mathbb{A}^1$-homotopies 
%$$G = (G_0,\dots,G_r):\A^1_W\to \text{Im}(\gamma) \to \sC$$




 %\  


 %\   


 
 % where $\alpha: \A^1_k\backslash \{0\} \to \sC$ denotes the constant map to the $k$-point $\alpha$. As $H\circ \pr_1 =q:U \to \sC$ factors through $\A^1_k$ and $H\circ \pr_2=\alpha:U\to \sC$ is constant, thus exists a lift $\alpha'$ to $\A^1_k$, we have a $\A^1$-chain homotopy $G:\A^1_U \to \A^1$ with $G(0)=q'$ and $G(1)= \alpha'$. Thus, the composition $\rho \circ n \circ G$ we have the required $\A^1$-ghost homotopy between $ H\circ \tilde{\beta} = \beta\mathrm{\ and\ } i_1\circ\alpha =\alpha$. 
 Hence $\sS^2(\sC)(k )$ is trivial. Since $\eta = x\in \sS(\sC)(k)$, for some $k$-point $x \in \text{Im}(\gamma)$, so for every finitely generated field extension $F/k$, $\sS^2(\sC)(F)$ is trivial and thus, $\pi_0^{\A^1}(\sC)(F)$ is trivial. Hence $\sC$ is $\A^1$-connected, by \cite[Lemma 3.3.6]{ictp}. Therefore, by Corollary \ref{factorisation A1 equivalence}, $X$ is also $\A^1$-connected. 
 \end{proof}
 	\begin{remark}
 	    Thus if $X$ is a rational, affine surface admitting an $\A^1$-fibration $\pi: X \to C$ 
       % is an $\A^1$-fibration from a smooth affine surface 
        onto a smooth curve $C$ (the curve $C$ must be a rational curve) and $\sC$ is the algebraic space (by Theorem \ref{factor-A1-fibration}), then by Corollary \ref{a1 invariance of algebraic space} and Theorem \ref{connected-base-new} we have the following:
        \begin{enumerate}
            \item If $C$ is a $\A^1$-rigid curve, then $\pi_0^{\A^1}(X) \cong \sC$.
            \item If $C \cong \A^1_k \text{ or } \P^1_k$, then 
            \begin{enumerate}
                \item either $\sC$ is $\A^1$-invariant, in that case $\pi_0^{\A^1}(X) \cong \sC$ (Example \ref{rigid example}).
                \item or $\sC$ is not $\A^1$-invariant. In this case $\sS^2(\sC)(F)$ is trivial, for every finitely generated separable field extension $F/k$. So both $\sC$ and $X$ are $\A^1$-connected.
            \end{enumerate}
        \end{enumerate}
 	\end{remark}
\begin{corollary} \label{negative kodaira invariant}
Let $X$ be a smooth, affine surface over an algebraically closed field of characteristic zero. Suppose, $X$ is $\A^1$-uniruled. Then $\pi_0^{\A^1}(X)$ is $\A^1$-invariant.
\end{corollary}
Therefore by Corollary \ref{rigid case a1 invariant} and Corollary \ref{negative kodaira invariant}, we have the $\A^1$-invariance of the $\pi_0^{\A^1}(X)$:
\begin{corollary}\label{main corollary}
Let $X$ be a smooth, affine surface over an algebraically closed field of characteristic zero. 
%Suppose, $X$ is $\A^1$-uniruled. 
Then $\pi_0^{\A^1}(X)$ is $\A^1$-invariant.
\end{corollary}

\begin{corollary} \label{naive spaces}
Suppose, $X \in Sm/k$ is an affine surface, which is $\A^1$-uniruled, admitting an $\A^1$-fibration $\pi: X \to C$. Also, assume that the algebraic space $\sC$ in Theorem \ref{factor-A1-fibration} is $\A^1$-invariant (for example, if $C$ is an $\A^1$-rigid curve, by Corollary \ref{a1 invariance of algebraic space}). Then $X$ is $\A^1$-naive \cite[Definition 2.1.1]{asokII}, in particular $Sing_*(X)$ is $\A^1$-local.
\end{corollary}
			\begin{example} \label{rigid example}
 There are $\mathbb{A}^1$-fibrations $X \to \mathbb{A}^1$ such that $\sC$ is $\mathbb{A}^1$-rigid. Consider the smooth affine surface $X\equiv\{x^2(x-1)^2z-y^2+x(x-1)=0\}$ and $\pr_x:X\to \A^1$ is the fibration. For any morphism $\gamma: \A^1 \to X$, $\pr_x\circ\gamma$ is constant. Clearly if $\gamma(t)=(x(t),y(t),z(t))$ then 
	\begin{align*}
		x(t)^2(x(t)-1)^2z(t)-y(t)^2+x(t)(x(t)-1)&=0\\
		x(t)(x(t)-1)[x(t)(x(t)-1)z(t)+1]&=y(t)^2
	\end{align*}
	Thus any root $t_0$ of $x(t)\mathrm{\ or\ }(x(t)-1)$ must be a root of $y(t)^2$. So the roots of $x(t)$ and $(x(t)-1)$ must have multiplicity of multiple of $2$ since $x(t_0)(x(t_0)-1)z(t_0)+1=1$ for any root $t_0$ of $x(t)\mathrm{\ or\ }(x(t)-1)$. Thus both of them are square polynomial. Let $x(t)=u(t)^2,\ x(t)-1=v(t)^2$, then $$1=u(t)^2-v(t)^2=(u(t)-v(t))(u(t)+v(t)).$$ So both $(u(t)+v(t))$ and $(u(t)-v(t))$ are constant then both $u(t),v(t)$ are constant, so $\pr_x\circ\gamma(t)=x(t)=u(t)^2$ is constant.\par Now, let $\sC$ is the algebraic space in the factorisation $X\xrightarrow{\theta}\sC\xrightarrow{p}\A^1$. Then by Lemma \ref{irred-fibers} any non-constant map $\gamma:\A^1_k\to \sC$ is surjective. Also as $X\to \sC$ is a étale bundle any non-constant map $\gamma: \A^1_k \to \sC$ lifts to $X$ (Corollary \ref{not-a1-rigid}) Which must lie in a single fiber. Thus there does not exist any non-constant map $\gamma:\A^1_k\to \sC$ so $\sC$ is $\A^1$-rigid.  
%Let $X \subseteq \mathbb{A}^3_k$ be the surface defined as
% $$X\equiv\{f(x,y,z) = x^2z-y^2(y-1)^2-xy(y-1) -2x(y-1) = 0\} \subseteq \mathbb{A}^3_k.$$ 
%Then 
%$$\nabla f = (2xz+y(y-1)-2(y-1), -2y(y-1)^2-2y^2(y-1)+x(y-1)+xy-2x, x^2)$$
%The surface is not smooth at $(0,1,0)$.
%%amma(t) \neq 0\}, \text{ then } \gamma_3(t) = \frac{\gamma_2(t)^2(\gamma_2(t)-1) - \gamma_1(t)}{}$$
%
%
%Clearly the morphism $\gamma:\A^1 \to X \mathrm{\ by\ }t\mapsto (t,1,0)$ is non-constant. It also intesects every fiber, but does not intersects the irreducible component $F_0$ of the fiber over $0$. However $\gamma$ descends to a nonconstant map to $\sC$. Hence, $\sC$ is $\A^1$-connected and so is $X$.   


	
			\end{example}

%\begin{remark}
%Thus, if the algebraic space $\mathcal{C}$ in the factorisation of the $\mathbb{A}^1$-fibration $\pi: X \to C$ is %$\mathbb{A}^1$-invariant (by Theorem \ref{factor-A1-fibration}), then the canonical map
%$$\mathcal{C} \to Sing_*(\mathcal{C})$$
%is an isomorphism. Thus $\mathcal{C}$ is $\mathbb{A}^1$-local or $\mathcal{C}$ is $\mathbb{A}^1$-fibrant. Thus the %simplicial weak equivalence
%$$Sing_*(X) \to Sing_*(\mathcal{C})$$
%gives there is a simplicial weak equivalence
%$$Sing_*(X) \to \mathcal{C}.$$
%On the other hand, the $\A^1$-weak equivalence $X \to \sC$
%%$Sing_*(X) \to Sing_*(\mathcal{C})$
%induces a simplicial weak equivalence
%$$Ex^{\mathbb{A}^1}(X) \to Ex^{\mathbb{A}^1}(\mathcal{C}).$$
%Since, $\mathcal{C}$ is $\mathbb{A}^1$-rigid, so $Ex^{\mathbb{A}^1}(\mathcal{C})$ is simplicially weakly %equivalent to $\mathcal{C}$. So, combining all these, we have the morphism
%$$Sing_*(X) \to Ex^{\mathbb{A}^1}(X)$$
%is a simplicial weak equivalence. T
%So in particular, in these cases $X$ is $\mathbb{A}^1$-naive. For example, if $X$ admits an $\mathbb{A}^1$-fibration to an $\mathbb{A}^1$-rigid curve (also if $X$ is the surface as in Example \ref{rigid example}), then $X$ is $\mathbb{A}^1$-naive in the sense of \cite[Definition 2.1.1]{asokII} or equivalently, $Sing_*(X)$ satisfies the affine Nisnevich excision (\cite[Section 2.1]{asokI}).

%\end{remark}
			\begin{remark}
 From Theorem \ref{connected-base-new} we see that if the algebraic space $\sC$ appearing in Theorem \ref{factor-A1-fibration} is not $\A^1$-invariant, then it is $\A^1$-connected by proving $\sS^2(\sC)(k)$ is trivial. %(equivalently, $X$ is $\A^1$-chain connected). 
 However we do not know whether in those cases $\sS(\sC)(k)$ is trivial (equivalently, $X$ is $\A^1$-chain connected) or not. Note that, for such an example of an $\A^1$-fibration $\pi:X \to C$ with $\sS(\sC)(k)$ non-trivial (and $\sS^2(\sC)(k)$ is trivial), the cannonical epimorphism $\sS(X)\to \pi_0^{\A^1}(X)$ is not an isomorphism, thus $Sing_*(X)$ is not $\A^1$-local (see \cite[Conjecture 2.2.8]{am}).
\end{remark} 
\section{$\mathbb{A}^1$-homotopy type of $\mathbb{A}^1$-ruled surfaces} \label{section a1 homotopy types}
In the previous section we have seen that an $\A^1$-fibration $\pi : X \to C$ where $X$ is a smooth affine surface and $C$ is a smooth curve factors as $X \xrightarrow{\theta}\sC\xrightarrow{p} C$ where $\theta$ is an étale $\A^1$-bundle, $\sC$ is a smooth algebraic space and $p$ is birational and quasi finite. We have also deduced that $\theta$ is an $\A^1$-weak equivalence.
 In this section, we describe the $\A^1$-homotopy types of an $\A^1$-uniruled surfaces in the following cases: %either the algebraic space $\sC$ is $\A^1$-invariant (in particular, 
the base curve $C$ of the $\A^1$-fibration $\pi: X \to C$ is $\A^1$-rigid
%, by Corollary \ref{}) 
or every fiber of the $\A^1$-fibration $\pi: X \to C$ is reduced (in this case, the algebraic space $\sC$ is actually a scheme).
%base curve $C$ of the $\A^1$-fibration is $$ 

%Suppose, $X$ admits an $\A^1$-fibration $\pi:X \to C$ onto a smooth curve $C$. Assume that the algebraic space $\sC$ in the factorisation of $\pi$ in Theorem \ref{} is $\A^1$-invariant. 

% can conclude
\subsection{The base curve $C$ is an $\mathbb{A}^1$-rigid curve}
Suppose, $X$ admits an $\mathbb{A}^1$-fibration $\pi: X \to C$ to an $\mathbb{A}^1$-rigid curve $C$. Then by Theorem \ref{factor-A1-fibration}, $\pi$ factors as
$$X \xrightarrow{\rho} \mathcal{C} \xrightarrow{p} C,$$
where $\rho$ is an étale locally trivial $\mathbb{A}^1$-bundle. By Corollary \ref{factorisation A1 equivalence}, $\rho$ is an $\mathbb{A}^1$-weak equivalence. Since $C$ is $\mathbb{A}^1$-rigid, so the algebraic space $\mathcal{C}$ is an $\mathbb{A}^1$-rigid Nisnevich sheaf (an algebraic space is an étale sheaf, so is a Nisnevich sheaf). Therefore, the $\mathbb{A}^1$-homotopy types of $\mathcal{C}$ is equivalent to the isomorphism types of $\mathcal{C}$ (Corollary \ref{homotopy types base rigid}). Thus the class of $\mathbb{A}^1$-homotopy types of a surfaces admitting an $\mathbb{A}^1$-fibration to an $\mathbb{A}^1$-rigid curve consists all the smooth one dimensional algebraic spaces (schemes or non-separated) that can be constructed in the factorisation. In particular this class consists all $\mathbb{A}^1$-rigid curves $C$ with multiple points (in the case when every fiber of $\pi$ is reduced), alongwith the non-separated algebraic spaces obtained from $C$, if there is a non-redued fiber.
%If the base of the $\A^1$-fibration is $\A^1$-rigid $X$ admits an $\A^1$

		\subsection{The base curve $C$ is $\mathbb{A}^1$-connected}

In this subsection, we will find the $\mathbb{A}^1$-homotopy types of the smooth $\mathbb{A}^1$-ruled surfaces $X$ that admit $\mathbb{A}^1$-fibration $\pi: X \to C$ to an $\mathbb{A}^1$-connected curve with  every fiber of $\pi$ is reduced. 
%We will also see that if a fiber of $\pi: X \to \mathbb{A}^1$ has an irreducible component of multiplicity $m$ (say), then there is a $m$-torsion element in $Pic(X)$.
% This in particular implies that if $X$ is a smooth affine surface that admits an $\mathbb{A}^1$-fibration $\pi: X \to \mathbb{A}^1$, 
  %accor $X$ acording to their $\mathbb{A}^1$-homotopy types, in the case if .

%\subsection{every fiber of $\pi$ is reduced}

		%If all of the fibers of an $\A^1$-fibration is reduced then the algebraic space $\sC$ in the appearing in the factorisation (of Theorem \ref{factor-A1-fibration}) $X\xrightarrow{\theta}\sC\xrightarrow{p}C$ is actually a scheme. The scheme will be non-separated in case of the $\A^1$-fibration is not an $\A^1$-bundle. Also since the surface have the same $\A^1$-homotopy type as those schemes $\sC$ in the factorisation, it is enough to know their $\A^1$-homotopy types. Since $\mathbb{P}^1$ is $\mathbb{A}^1$-weakly equivalent to $\mathbb{A}^1$ with double origins, so to describe the $\mathbb{A}^1$-homotopy types of $\mathbb{A}^1$-ruled surface $X$, it suffices to consider the two cases: one is $X$ admits an $\mathbb{A}^1$-fibration $\pi: X \to C$ with $C$ is an $\mathbb{A}^1$-rigid curve and the other is $X$ admits an $\mathbb{A}^1$-fibration $\pi: X \to \mathbb{A}^1$.



Since $C$ is $\mathbb{A}^1$-connected, therefore $C$ is isomorphic to $\mathbb{A}^1$ or $\P^1$.
Let us recall the purity isomorphism in $\mathcal{H}_\bullet(k)$ by Morel and Voevodsky. 
\begin{theorem} \cite[Theorem 2.23, Section 3.2]{mv} \label{homotopy purity}
	Let $i: Z \to X$ be a closed immersion of schemes in $Sm/k$ and $N_{X,Z} \to Z$ is the normal bundle over $Z$ of rank $n=\text{codim}_Z(X)$. Then there is an $\mathbb{A}^1$-weak equivalence of pointed sheaves
	$$X/(X \setminus i(Z)) \sim Th(N_{X, Z}),$$
	where for a vector bundle $\rho: V \to X$, the Thom space of $\rho$ is defined as
	$$Th(\rho):= V/(V \setminus s(X)),$$
	here $s: X \to V$ is the zero section of $\rho$.
\end{theorem}

		\begin{example}
\begin{enumerate}
	\item Let $X=\tilde{\A}^1$ denotes $\A^1$ with \textit{double} origins. It is the homotopy pushout of the following diagram
	
	
	\[
	\xymatrix{{\G_m} \ar@{^{(}->}[d] \ar[r] & {\A^1\cong*} \ar[d] \\
		{\A^1} \ar[r] & {\tilde{\A}^1\cong\A^1/\G_m\cong\P^1.}}
	\]
	As $\A^1/\G_m$ is the homotopy pushout since $\G_m\hookrightarrow\A^1$ is cofibration. Then $\A^1/\G^m$ is the Thom space of the $1$-dimetional vector bundle over a point. Hence, by Theorem \ref{homotopy purity} it is isomorphic to $\P^1$. 
		\item Let $X=\A^1_{0;n}$ denotes $\A^1$ with $n$ origins, formed by gluing $n$ copies of $\A^1$ along $\A^1\backslash\{0\}$ via identity. Also denote the origins by $0_1,\dots,0_n$. 
%		Take $U_{n-1} = X\backslash\{0_n\}= \A^1_{0,(n-1)} \mathrm{\ and},\ V=X\backslash\{0_1,\dots,0_{n-2}\}=\A^1_{0;2}$. $U\cup V=X \mathrm{\ and\ }U\cap V=\A^1$ is $\A^1$-contractible. Then by $\A^1$-Van-Kampen theorem \cite[Thm 5.1.1]{am} $$\pi_1^{\A^1}(X)\cong\pi_1^{\A^1}(U_{n-1})\star^{\A^1}\pi_1^{\A^1}(V).$$
%	So by induction we have $$\pi_1^{\A^1}(X)\cong\pi_1^{\A^1}(\P^1)\star\cdots\star\pi_1^{\A^1}(\P^1)\cong\pi_1^{\A^1}((\P^1)^{\vee (n-1)})$$ product of $(n-1)$ copies of $\pi_1^{\A^1}(\P^1)$.\par
	By the homotopy pushout diagram we see, 
	
	\[
	\xymatrix{{\G_m} \ar@{^{(}->}[d] \ar[r] & {\A^1\cong*} \ar[d] \\
		{\A^1_{(n-1);0}} \ar[r] & {\A^1_{n;0}\cong (\A^1_{(n-1);0}/\G_m). }
	}
	\]
	Let $Z=\{0_1,\dots,0_{(n-1)}\}\hookrightarrow\A^1_{(n-1);0}$ be the smooth closed embedding. Then by homotopy purity (Theorem \ref{homotopy purity}) we see, $$(\A^1_{(n-1);0}/\G_m)=Th(N_{Z,X})=\Sigma_T^1(Z)_+=\P^1\wedge(Z)_+=(\P^1)^{\vee (n-1)}$$ as $Z$ is the disjoint union of $(n-1)$ points, thus the normal bundle is trivial. Hence, the Thom space is wedge of $(n-1)$-many copies of $\P^1$. The calculations are done as it was mentioned in \cite[Proposition 5.3.4.3]{as}\\
	
	The Picard group of $\A^1_{0;n}$ is $\Pic(\A^1_{0;n})=\Z^{(n-1)}$, direct sum of $(n-1)$ copies of $\Z$. Indeed, using Weil divisors we see, $nc = \Div((x-c)^n)=0$ in $\Pic(\A^1_{0;n})$ if $c$ is not one of the origins. Let $q:\A^1_{0;n}\to \A^1$ by $0_i\mapsto 0.\ c\mapsto c$, then $\Div(q)=\Sigma_{1}^{n}(0_i)=0$ in $\Pic(\A^1_{0;n})$. Let $f:\A^1_{0;n}=\tilde{\A}^1 \to\A^1$ be any map, then $f(0_i)=f(0_j)$ for all $i,j$.  Let $U=\Spec(k(x)),\ T=\Spec(k[x]_{(x)})$ 
	
	
	\[\xymatrixcolsep{4pc}
	\xymatrix{{U=\Spec(k(x))} \ar@{^{(}->}[d] \ar[r]^<<<<<<<<<\eta & {\tilde{\A}^1} \ar[r]^f & {\A^1} \\
		{T=\Spec(k[x]_{(x)})} \ar@<-1.3px>[ur]_{s_j} \ar@<1.3px>[ur]^{s_i}
	}
	\]
	where $\eta$ maps to the generic point of $\tilde{\A}^1$. Then the two extension $s_i,\ s_j$ mapping the closed point to $0_i\mathrm{\ and \ }0_j$ respectively must be a unique map after composing with $f$. Thus $f(0_i)=f(0_j)$ for all $i,j$. Thus, we have $$\Pic(\A^1_{0;n})=\Z[0_1,\dots0_n]/\sum_{1}^{n}(0_i)\cong \Z^{(n-1)}.$$ Hence $\A^1_{0;n}\mathrm{\ are\ }\A^1_{0;m}$ is not $\A^1$-homotopic if $m\neq n$.
	
	
	
	\item Let $X=\A^1_{0,1;2,2}\ i.e.,\ \A^1$ with \textit{double} origin and \textit{double} 1. 
	\[
	\xymatrix{{\G_m} \ar@{^{(}->}[d] \ar[r]^{t\mapsto(t+1)} & {\A^1\cong*} \ar[d] \\
		{\A^1_{2;0}} \ar[r] & {\A^1_{0,1;2,2}\cong\A^1_{0;2}/\G_m}
	}
	\]
	By the above homotopy pushout diagram we see that it is $\A^1$-homotopic to $\A^1_{0;2}/\G_m$, which is also $\A^1$-homotpic to $\A^1_{0;3}\ i.e.,\ \A^1$ with $3$ origins.\\
	
	However they are not isomorphic as schemes. By similar argument as in the last example we have for any map $f:X=\A^1_{0,1;2,2}\to \A^1$ $f(0_i)=f(1_j)$ for all $1\leq i,j\leq 2$. Now, for any $g:X\to \A^1_{(0,3)}$ we have the composition $$X\xrightarrow{g}\A^1_{(0,3)}\xrightarrow{q}\A^1.$$ Here $q$ maps all of the origins to $0$ of $\A^1$ and identity elsewhere. Thus $g(0_i),g(1_j)\subset \{0_1,0_2,0_3\}$ since $q\circ g(0_i)=q\circ g(1_j)$ but $q(x)=q(y)$ if $x=y \mathrm{\ or\ } x,y \in \{0_1,0_2,0_3\}.$ Hence, $g$ cannot be injective.
	
	 
	
	\item From the previous case for $n\geq2$, $X=\P^1_{0;(n-1)}$ i.e., $\P^1$ with $(n-1)$ origins is isomorphic to $\A^1_{0;n}$. The case $n=2$  is done in the first example. Using induction it follows from the same homotpy purity diagram 
	
	\[
	\xymatrix{{\G_m} \ar@{^{(}->}[d] \ar[rr]&& {\A^1\cong*} \ar[d] \\
		{\P^1_{0;(n-2)}\cong\A^1_{0;(n-1)}} \ar[rr] && {\P^1_{0;(n-1)}\cong\A^1_{0;(n-1)}/\G_m\cong\A^1_{0;n}.}
	}
	\]
\end{enumerate}
		\end{example}
			
				
				
			
			
		Now we generalise this computations in the following lemma.
		\begin{lemma}
			Suppose, $C_1$ be the curve affine line with $m$ origins and Suppose there are $r$ many $k$-points $a_1, \dots, a_r \in \mathbb{A}^1$ along with the integers $n_1, \dots, n_r$ such that $\sum_{1}^{r}(n_i-1)=m-1$. $C_2$ be the affine line with the points $a_i,\ n_i$ many times denoted as,
			$C_2:= \mathbb{A}^1(a_1, n_1; \dots; a_r, n_r)$ 
			Then $C_1$ and $C_2$ are $\A^1$-weakly equivalent.
			
		\end{lemma}
		\begin{proof}
			 Let $C_2= \mathbb{A}^1(a_1, n_1; \dots, a_r, n_r)$. Without loss of generality we can assume $a_1=0,\ a_2=1$. By induction, also assume $\mathbb{A}^1(a_1, n_1; \dots, a_r, n_r-1) \cong \A^1_{0;1+(\sum_{0}^{(r-1)}(n_i-1))+(n_r-2)}$ the affine line with $(1+(\sum_{0}^{(r-1)}(n_i-1))+(n_r-2))$ many origins. From the pushout diagram we have $C_2$ is the homotopy pushout of the following diagram
			
			
			\[
			\xymatrix{{\G_m} \ar@{^{(}->}[d] \ar[rr]^{t\mapsto(a_r+t)}&& {\A^1 \cong *}\ar[d] \\
				{\A^1_{0;1+ (\sum_0^{r-1}(n_i-1))+(n_r-2)}\cong \mathbb{A}^1(a_1, n_1; \dots, a_r, n_r-1)} \ar[rr]&& X.
			}
			\]
			Thus $C_2\cong C_1$. Also by similar argument we can conclude that any morphism between $C_2\to C_1$ cannot be injective, thus they are not isomorphic.
		\end{proof}
%		\begin{corollary}
%			For a smooth affine surface $X$, admitting a fibration to an $\A^1$-connected smooth curve with reduced fibers, i.e. $\A^1\mathrm{\ or\ }\P^1$. $X$ is $\A^1$-homotopic to $\A^1_{0,n}$, i.e. $\A^1$ with $n$ origins for some $n$.
%		\end{corollary}
		\begin{remark}
			We see that for $\A^1$-connected smooth curve, their $\A^1$-homotopic type is not same as their isomorphic type, however for $\A^1$-rigid curves it is not the case. If $C$ is an $\A^1$-rigid curve (\textit{not-necessarily separated}) then $\pi_0^{\A^1}(C)=C$ as a sheaf, thus in this case isomorphic type is same as $\A^1$-homotpy type. 
		\end{remark}

\section{Appendix} \label{detailed horn formula}

\begin{point}[Boundary and degeneracy maps of $Sing_*(\sX)(U)$]\label{maps of sing}~
	
Let $\sX$ be a space and $U$ be a scheme over $k$. Here we describe the boundary and degeneracy maps of the simplicial set $Sing_*(\sX)(U)$. They are denoted by $d_i$ and $s_i$ respectively. 
$$d_i: Sing_*(\mathcal{X})(U)_n \to Sing_*(\mathcal{X})(U)_{n-1}, \ 0 \leq i \leq n$$
are induced by the maps $d^i: \Delta^{n-1}_a \to \Delta^n_a$, which are given by the morphism between $k$-algebras (we again denote it by $d^i$)
$$d^i:  \frac{k[x_0, x_1,\dots, x_n]}{(\sum_{i=0}^n x_i - 1)} \to  \frac{k[x_0, x_1,\dots, x_{n-1}]}{(\sum_{i=0}^{n-1} x_i - 1)}$$
defined as 
\[
\begin{cases}
	\overline{x_j} \mapsto \overline{x_j} & j < i \\
	\overline{x_j} \mapsto \overline{0} & j = i \\
	\overline{x_j} \mapsto \overline{x_{j-1}} & j > i
\end{cases}
\]
(here, for $f \in k[x_0,\dots, x_n]$ $\bar{f}$ denotes the class in $\frac{k[x_0, x_1,\dots, x_n]}{(\sum_{i=0}^n x_i - 1)}$). 
Similarly, here the degeneracy maps 
$$s_i: Sing_*(\mathcal{X})(U)_n \to Sing_*(\mathcal{X})(U)_{n+1}, \ 0 \leq i \leq n$$
are induced by the maps $s^i: \Delta^{n+1}_a \to \Delta^n_a$, 
which are given by the morphism between $k$-algebras (we again denote it by $s^i$)
$$s^i:  \frac{k[x_0, x_1,\dots, x_n]}{(\sum_{i=0}^n x_i - 1)} \to  \frac{k[x_0, x_1,\dots, x_{n+1}]}{(\sum_{i=0}^{n+1} x_i - 1)}$$
defined as 
\[
\begin{cases}
	\overline{x_j} \mapsto \overline{x_j} & j < i \\
	\overline{x_i} \mapsto \overline{x_i + x_{i+1}} & j = i \\
	\overline{x_j} \mapsto \overline{x_{j+1}} & j > i
\end{cases}
\]
\end{point}
%\begin{remark} \label{boundary and horn filling}
%    \begin{enumerate}
%        \item The standard cosimplicial object in $Sm/k$ is $\Delta^n_a$, which is isomorphic to $\A^n_k$ (Definition \ref{sing functor}). Let us fix an isomorphism
%        $$\Delta^n_a \to \A^n_k (:= \Spec k[x_0,\dots,x_{n-1}]) \text{ defined as } $$
%        $$x_i \mapsto x_i, \text{if }i < n \text{ and }x_n \mapsto 1-\sum_{i=0}^{n-1}x_i.$$
%     \item   The algebraic $l$-th horn ($0 \leq l \leq n$, denoted by $\Lambda^n_{l,a}$) is the (reducible) closed subvariety of $\Delta^n_a$ is generated by the images of 
%     $$d_i: \Delta^{n-1}_a \to \Delta^n_a, \text{ for }0 \leq i \leq n, \ i \neq l.$$ 
%     So, 
%        $$\Lambda^n_{l,a} \cong \Spec \ \frac{k[x_0,\dots, x_n]}{(x_0 \dots \hat{x_l} \dots x_n(\sum_{i=0}^n x_i -1))} \subseteq \Spec \ \frac{k[x_0, \dots x_n]}{(\sum_{i=0}^n x_i - 1)}.$$
%     \item Suppose, $X$ and $U$ are $k$-schemes. A morphism $\phi: \Lambda^n_{l} \to Sing_*(X)(U)$ is given by the morphisms
%     $$f_0, \dots, \hat{f_l},\dots, f_n: \Delta^{n-1}_a \times_k U \to X $$
%    $ \text{ such that } d_i(f_j) = d_{j-1}(f_i), \text{ for }i < j \text{ and }i,j \neq l$. This means a morphism $\phi: \Lambda^n_{l} \to Sing_*(X)(U)$ corresponds to a morphism (denote it by $\phi$ again)
%    $$\phi: \Spec \ \frac{A[x_0,\dots,x_n]}{(x_0 \dots \hat{x_l} \dots x_n(\sum_{i=0}^n x_i -1))} \to X,$$
%where $A = \sO(U)$. Thus, if $X \cong \A^m_k (= \Spec \ k[T_1,\dots,T_m])$, then $\phi$ can be extened to a morphism 
%$$\Tilde{\phi}: \Spec \ \frac{A[x_0, \dots,x_n]}{(\sum_{i=0}^n x_i -1)} \to X;$$
%since any morphism $k[T_1, \dots,T_m] \to B/I$ ($B$ is a $k$-algebra and $I$ is an ideal of $B$) can be lifted to $k[T_1, \dots,T_m] \to B$. So, $Sing_*(\A^m_k)(U)$ is a Kan fibrant (Definition \ref{kan fibrant definition}) simplicial set.
% \item The algebraic boundary (denoted by $\partial \Delta^n_a$) is the (reducible) closed subvariety of $\Delta^n_a$ is generated by the images of $d_i: \Delta^{n-1}_a \to \Delta^n_a$, for $0 \leq i \leq n$. So, 
%        $$\partial \Delta^n_a \cong \Spec \ \frac{k[x_0,\dots,x_n]}{(x_0 \dots x_n(\sum_{i=0}^n x_i - 1))} \subseteq \Spec \ \frac{k[x_0, \dots x_n]}{(\sum_{i=0}^n x_i - 1)}$$
%        or following the isomorphism in (1)
%        $$\partial \Delta^n_a \cong \Spec \frac{k[x_0, \dots, x_{n-1}]}{(x_0\dots x_{n-1}(\sum_{i=0}^{n-1} x_i-1))} \subseteq \Spec \ k[x_0, \dots, x_{n-1}].$$
%        \item Suppose, $X$ and $U$ are $k$-schemes. A morphism $\phi: \partial \Delta^n \to Sing_*(X)(U)$ is given by the morphisms
%     $$f_0,\dots, f_n: \Delta^{n-1}_a \times_k U \to X \text{ such that } d_i(f_j) = d_{j-1}(f_i), \text{ for }i < j.$$
%   This means a morphism $\phi: \partial \Delta^n \to Sing_*(X)(U)$ corresponds to a morphism (denote it by $\phi$ again)
%    $$\phi: \Spec \ \frac{A[x_0,\dots,x_n]}{(x_0 \dots x_n(\sum_{i=0}^n x_i -1))} \to X,$$
%where $A = \sO(U)$. Thus, if $X \cong \A^m_k$, then $\phi$ can be extened to a morphism 
%$$\Tilde{\phi}: \Spec \ \frac{A[x_0, \dots,x_n]}{(\sum_{i=0}^n x_i -1)} \to X.$$
%The lifting property of $Sing_*(\A^m_k)(U)$ with respect to the boundary inclusions $\partial\Delta^n \to \Delta^n$ is equivalent to the fact that the simplicial set $Sing_*(\A^m_k)(U)$ is a Kan fibrant and contractible (see also Lemma \ref{contractible An}).
%%simplicial set.
%%Since the inclusions of simplicial sets $\partial\Delta^n \to \Delta^n$ is 
%     \end{enumerate}
%\end{remark}
To give the horn filling formula for $Sing_*(\A^m_k)(U)$ for $U \in Sm/k$ first we recall the horn filling formula for simplicial groups.

\begin{lemma}\cite[Lemma 3.1]{curtis-simplicial}\label{horn-formula}
Let $G$ be a simplicial group. $(f_1,\dots,f_{l-1},f_{l+1},f_n) \in G_{n-1} \mathrm{\ such\ that\ }d_i(f_j)=d_{j-1}(f_i)\ i<j\neq k;\ $ be a horn. Then take -
	\begin{align*}
		g_0 &= s_0(f_0)& \\
		g_i &= g_{i-1}\cdot s_i(d_i(g_{i-1}))^{-1}\cdot s_i(f_i)\  &\mathrm{ for}\ 0<i<l\\
		g_i &= g_{i+1} \cdot s_{i-1}(d_i(g_{i+1}))^{-1}\cdot s_i(f_i)\  & \mathrm{for}\ l<i<n\\
		g_n &= g_{l-1}\cdot s_{n-1}(d_n(g_{l-1}))^{-1}\cdot s_{n-1}(f_n).&
	\end{align*}
	Where $\cdot$ is the group operation. When $k=0$ we set $g_{-1}=e_G$ the identity element. Then $g_l\in G_n$ satisfies the property $d_j(g_{l+1})=f_j$ for $0\leq j \leq n$.
\end{lemma}

Clearly for $U\in Sm/k\ Sing_*(\A^m_K)(U)$ is a simplicial group. The group operation $+$ of $Sing_n(\A^m_k)(U)$ is defined point-wise and using group structure of $\A^m_k$. Let $f,g \in Sing_n(\A^m_k)(U)=Hom(\Delta^n_a\times U,\A^m_k)$, then $(f+g) \in Hom (\Delta^n_a\times U,\A^m_k)$ defined as \[ (f+g)(x_0,\dots,x_n)= f(x_0,\dots,x_n)+g(x_0,\dots,x_n) \] where the second + is the group operation of the variety $\A^m_k$. It has $\mathrm{0}\in Hom(\Delta^n_a\times U,\A^m_k)$ constant map to $(0,\dots,0)$ as identity element and the inverse of $f$ is $-f$ which is defined pointwise.  Similarly the boundary and degeneracy maps are group homomorphism. Let $f,g \in Sing_n(\A^m_k)(U)=Hom(\Delta^n_a\times U,\A^m_k)$ then 
\begin{flalign*}
	&d_i(f+g)(x_0,\dots,x_{n-1})&\\
	&= (f+g)(x_0,\dots,x_{i-1},0,x_{i},\dots,x_{n-1})&\\
	&= f(x_0,\dots,x_{i-1},0,x_{i},\dots,x_{n-1})+ g(x_0,\dots,x_{i-1},0,x_{i},\dots,x_{n-1})&\\
	&=d_i(f)+d_i(g)&	
\end{flalign*} 


thus $d_i:Hom(\Delta^n_a \times U,\A^m_k)\to Hom(\Delta^{n-1}_a \times U,\A^m_k)$ is a group homomorphism. Likewise,
 \begin{align*}
 	&s_i(f+g)(x_0,\dots,x_{n+1})\\
 	&= (f+g)(x_0,\dots,x_{i-1},x_{i}+x_{i+1},\dots,x_{n+1})\\
 	&= f(x_0,\dots,x_{i-1},x_{i}+x_{i+1},\dots,x_{n+1})+ g(x_0,\dots,x_{i-1},x_{i}+x_{i+1},\dots,x_{n+1})\\
 	&=s_i(f)+s_i(g)	
 \end{align*} 
 thus $s_i:Hom(\Delta^n_a \times U,\A^m_k)\to Hom(\Delta^{n+1}_a \times U,\A^m_k)$ is a group homomorphism.
 
\begin{point}[Formula of horn filling of $Sing_*(\mathbb{A}^m_k)(U),\ U\in Sm/k$]\label{full hornfill An}~\\
	Let $f_0,\dots,f_{l-1},f_{l+1},f_n \in Sing_{n-1}(\A^m_k)(U)$ satisfying the compatibility condition $d_i(f_j)=d_{j-1}(f_i)$ for $i<j\neq l$ be a horn. Using Lemma \ref{horn-formula} and the group structure of $Sing_*(\A^m_k)(U)$ we define inductively -
	\begin{align*}
		g_0 &= s_0(f_0)& \\
		g_i &= g_{i-1}- s_i(d_i(g_{i-1})) + s_i(f_i)\  &\mathrm{ for}\ 0<i<l\\
		g_i &= g_{i+1} - s_{i-1}(d_i(g_{i+1})) + s_i(f_i)\  & \mathrm{for}\ l<i<n\\
		g_n &= g_{l-1} - s_{n-1}(d_n(g_{l-1})) + s_{n-1}(f_n).&
	\end{align*}
Similarly for $l=0$ we take $g_{-1}=0\in Hom(\Delta^n_a \times U,\A^m)$, in that case $g_n=s_{n-1}(f_n)$. Then $g_{l+1}\in Sing_n(\A^m_k)(U)$ satisfies $d_j(g_{l+1})=f_j \mathrm{\ for\ } 0\leq j\leq n$. For algebraic description of the horn filling see \cite[7.1.1]{Biman-Thesis}.
	
	
\end{point}
%	\begin{align*}
%		Sing_*(\mathbb{A}^m_k)(U)_n & = Hom_{Sm/k}(\Delta^n_a \times_k U, \mathbb{A}^m_k) \\
%		& \cong Hom_{k-alg}(k[T_1,\dots, T_m], \frac{A[x_0, x_1,\dots, x_n]}{(\sum_{i=0}^n x_i - 1)}),
%	\end{align*}
%	where $A$ is the ring of regular functions $\mathcal{O}(U)$. For $n \geq 1$ and $0 \leq l \leq n$, suppose we are given 
%	$$\phi_0,\dots,\phi_{l-1}, \phi_{l+1},\dots, \phi_n: k[T_1,\dots, T_m] \to \frac{A[x_0, x_1,\dots, x_{n-1}]}{(\sum_{i=0}^{n-1} x_i - 1)}$$
%	such that $d^i \circ \phi_j = d^{j-1} \circ \phi_i$, if $i < j \text{ and } \ i,j \neq l$, here we again denote the map $Id_A \otimes d^i$ by $d^i$ ($Id_A$ is the identity map on $A$):
%	$$Id_A \otimes d^i: A \otimes_k \frac{k[x_0, x_1,\dots, x_{n-1}]}{(\sum_{i=0}^{n-1} x_i - 1)}  \to A \otimes_k \frac{k[x_0, x_1,\dots, x_{n-2}]}{(\sum_{i=0}^{n-2} x_i - 1)}$$
%	and denote the map $Id_A \otimes s^i$ by $s^i$ again:
%	$$Id_A \otimes s^i: A \otimes_k \frac{k[x_0, x_1,\dots, x_{n-1}]}{(\sum_{i=0}^{n-1} x_i - 1)}  \to A \otimes_k \frac{k[x_0, x_1,\dots, x_{n}]}{(\sum_{i=0}^{n} x_i - 1)}$$
%	Since $T_p$'s are independent variables, suppose that each $\phi_i$ is given by 
%	$$T_p \xmapsto{\phi_i} \overline{f_i^{(p)}}, \ f_i^{(p)} \in A[x_0, x_1,\dots, x_{n-1}], \ 1 \leq p \leq m, \ 0 \leq i \leq n, i \neq l$$
%	and $\overline{f_i^{(p)}}$ is the class of $f_i^{(p)}$ in $\frac{A[x_0, x_1,\dots, x_{n-1}]}{(\sum_{i=0}^{n-1} x_i - 1)}$. We need to find some 
%	$$\Phi: k[T_1,\dots, T_m] \to \frac{A[x_0, x_1,\dots, x_n]}{(\sum_{i=0}^n x_i - 1)}$$ 
%	such that $d^i \circ \Phi = \phi_i, \text{ for all } \ i \neq l$. 
%	
%	If such $\Phi$ exists, then $\Phi$ is determined by the image of $T_i$.
%	The inclusion of $k$-algebras 
%	$$j_i: k[T] \xrightarrow{T \mapsto T_i } k[T_1, \dots, T_m]$$
%	gives a $l$-th horn $\Lambda^n_l$ in $Sing_*(\mathbb{A}^1_k)(U)$ given by $(n-1)$-simplices $\phi_0 \circ j_i, \dots, \phi_{l-1}\circ j_i, \phi_{l+1} \circ j_i, \dots, \phi_n \circ j_i$ of $Sing_*(\mathbb{A}^1_k)(U)$. Therefore it is enough to give the horn filling formula for $m = 1$.
%	
%	
%	
%	
%	 Suppose that $\Phi$ is given by $T_j \mapsto \overline{f_j}$. Then we would have for every $j \neq l$, $d^i(\overline{f_j}) = \overline{f_j^{(i)}}$. Therefore it is enough to prove the proposition for $m = 1$. \par
%	Using \cite[Theorem 3.1]{nlab}, we write the formulas of horn filling of $Sing_*(\mathbb{A}^m_k)(U)$. 
%	Suppose we are given 
%	$$\phi_0,\dots,\phi_{l-1}, \phi_{l+1},\dots, \phi_n: k[T] \to \frac{A[x_0, x_1,\dots, x_{n-1}]}{(\sum_{i=0}^{n-1} x_i - 1)}$$
%	satisfying the compatibility conditions. Assume that each $\phi_i$ is given by $T \mapsto \bar{f_i}$. 
%	Define two $k$-algebra morphisms
%	$$T^j, S^j: \frac{A[x_0, x_1,\dots, x_n]}{(\sum_{i=0}^n x_i - 1)} \to \frac{A[x_0, x_1,\dots, x_n]}{(\sum_{i=0}^n x_i - 1)}$$ 
%	as $T^j = s^j \circ d^j$ and $S^j = s^j \circ d^{j+1}$, i.e.
%	$$\overline{f(x_0,\dots, x_n)} \xmapsto{T^j} \overline{f(x_0,\dots, x_{j-1},0,x_j+x_{j+1},x_{j+2},\dots,x_n)}$$
%	$$\overline{f(x_0,\dots,x_n)} \xmapsto{S^j} \overline{f(x_0,\dots,x_{j-1},x_j+x_{j+1},0,x_{j+2},\dots,x_n)}$$
%	Consider three cases. \\
%	\textbf{Case 1:} Suppose, $l=0$. We are given the $(n-1)$-simplices as $\phi_1,\dots, \phi_n$ and each $\phi_i$ is given by 
%	$$T_p \mapsto \overline{f_i^{(p)}}, \ p= 1, \dots, m.$$
%	Define $\Phi:  k[T_1, \dots, T_m] \to \frac{A[x_0, x_1,\dots, x_{n}]}{(\sum_{i=0}^{n} x_i - 1)}$ as 
%\begin{equation*}
%	\begin{aligned}
%		T_p \mapsto & \overline{\mathlarger{\mathlarger{\sum}_{i=0}^{n-1} f_{i+1}^{(p)}(x_0,\dots,x_{i-1},x_i+x_{i+1},x_{i+2},\dots,x_n)}} \\
%		+ &\mathlarger{\mathlarger{\sum}_{r=0}^{n-2} \ \mathlarger{\sum}_{\substack{t \neq 0 \\ 0 \leq i_1 <\dots<i_t \leq r}} (-1)^t S^{i_1} \circ S^{i_2} \circ\dots\circ S^{i_t}}\\[-20pt]
%		& \hspace{5cm}(\overline{f_{r+2}^{(p)}(x_0,\dots,x_r,x_{r+1}+x_{r+2},x_{r+3},\dots, x_n)})
%	\end{aligned}
%\end{equation*}	
%
%	Thus $\Phi$ is defined as
%	$$T_p \longmapsto \sum_{i=0}^{n-1} s^i \overline{f_{i+1}^{(p)}} + \sum_{r =0}^{n-2} \ \sum_{\substack{t \neq 0 \\ 0 \leq i_1 < \dots < i_t \leq r }}(-1)^t S^{i_1} \circ \dots S^{i_t}s^{r+1}\overline{f_{r+2}^{(p)}},$$
%	for $p = 1, \dots, m$.
%	\
%	
%	\textbf{Case 2:} Suppose, $l=n$. We are given the $(n-1)$-simplices $\phi_0, \dots, \phi_{n-1}$ and $\phi_i$ and each $\phi_i$ is defined as 
%	$$T_p \mapsto \overline{f_i^{(p)}}, \ p= 1, \dots, m.$$
%	Define $\Phi:  k[T_1, \dots, T_m] \to \frac{A[x_0, x_1,\dots, x_{n}]}{(\sum_{i=0}^{n} x_i - 1)}$ as 
%	\begin{equation*}
%		\begin{aligned}
%			T_p \longmapsto & \overline{\mathlarger{\mathlarger{\sum}_{i=0}^{n-1} f_{i}^{(p)}(x_0,\dots,x_{i-1},x_i+x_{i+1},x_{i+2},\dots,x_n)}} \\ 
%			+ &\mathlarger{\mathlarger{\sum}_{r=1}^{n-1} \ \mathlarger{\sum}_{\substack{t \neq 0 \\ r \leq i_t <\dots<i_1 \leq n-1}} (-1)^t T^{i_1} \circ T^{i_2} \circ \dots \circ T^{i_t}}\\[-20pt]
%		~	&\hspace{5cm}(\overline{f_{r-1}^{(p)}(x_0,\dots,x_{r-1}+x_r,x_{r+1},\dots, x_n)})
%		\end{aligned}
%	\end{equation*}
%
%	Thus $\Phi$ is defined as
%	$$T \longmapsto \sum_{i=0}^{n-1} s^i \overline{f_{i}^{(p)}} + \sum_{r =1}^{n-1} \ \sum_{\substack{t \neq 0 \\ r \leq i_t < \dots < i_1 \leq n-1 }}(-1)^t T^{i_1} \circ \dots T^{i_t}s^{r-1}\overline{f_{r-1}^{(p)}},$$
%	for $p = 1, \dots, m$.
%	
%	
%	\textbf{Case 3:} Suppose, $ 0 < l < n$. We are given $\phi_0, \dots, \phi_{l-1}, \phi_{l+1}, \dots, \phi_n$ and each $\phi_i$ is defined as 
%	$$T_p \mapsto \overline{f_i^{(p)}}, \ p= 1, \dots, m.$$
%	Define $\Phi:  k[T_1, \dots, T_m] \to \frac{A[x_0, x_1,\dots, x_{n}]}{(\sum_{i=0}^{n} x_i - 1)}$ as 
%	\begin{equation*}
%		\begin{aligned}
%			T_p \longmapsto & \overline{\mathlarger{\mathlarger{\sum}_{i=0}^{l-1} f_{i}^{(p)}(x_0,\dots,x_{i-1},x_i+x_{i+1},x_{i+2},\dots,x_n)}} \\
%			+ &  \overline{\mathlarger{\mathlarger{\sum}_{i=0}^{n-l-1} f_{n-i}^{(p)}(x_0,\dots,x_{n-i-1}+x_{n-i},\dots,x_n)}} \\
%			+ &\mathlarger{\mathlarger{\sum}_{s=1}^{l-1} \ \mathlarger{ \sum}\limits_{\substack{(t,v) \neq (0,0) \\ l \leq i_1 <\dots<i_t \leq n-1 \\ s \leq j_v <\dots< j_1 \leq l-1 }} (-1)^{t+v} S^{i_1} \circ\dots S^{i_t} \circ T^{j_1} \circ \dots \circ T^{j_v}}\\[-25pt]
%		~	& \hspace{5cm}{ (\overline{f_{s-1}^{(p)}(x_0,\dots,x_{s-1}+x_{s},\dots, x_n)})}\\[15pt]
%		~	+ & \mathlarger{\mathlarger{\sum}_{\substack{r \neq 0 \\ l \leq i_1 <\dots<i_r \leq n-1 }} (-1)^r S^{i_1} \circ\dots\circ S^{i_r} (\overline{f_{l-1}^{(p)}(x_0,\dots,x_{l-1}+x_{l},\dots,x_n)})} \\
%			+ & \mathlarger{\mathlarger{\sum}_{t=1}^{n-1-l} \ \mathlarger{\sum}_{\substack{v \neq 0 \\ l \leq i_1 <\dots<i_v \leq n-1-t}} (-1)^v S^{i_1} \circ S^{i_2} \circ\dots \circ S^{i_v}}\\[-20pt]
%		~	& \hspace{4.5cm} (\overline{f_{n-t+1}^{(p)}(x_0,\dots,x_{n-t}+x_{n-t+1},\dots, x_n)}),
%		\end{aligned}
%	\end{equation*}
%	\begin{align*}
%		T_p \longmapsto & \overline{\mathlarger{\mathlarger{\sum}_{i=0}^{l-1} f_{i}^{(p)}(x_0,\dots,x_{i-1},x_i+x_{i+1},x_{i+2},\dots,x_n)}} \\
%		+ &  \overline{\mathlarger{\mathlarger{\sum}_{i=0}^{n-l-1} f_{n-i}^{(p)}(x_0,\dots,x_{n-i-1}+x_{n-i},\dots,x_n)}} \\
%		+ &\mathlarger{\mathlarger{\sum}_{s=1}^{l-1} \ \mathlarger{ \sum}\limits_{\substack{(t,v) \neq (0,0) \\ l \leq i_1 <\dots<i_t \leq n-1 \\ s \leq j_v <\dots< j_1 \leq l-1 }} (-1)^{t+v} S^{i_1} \circ\dots S^{i_t} \circ T^{j_1} \circ}\\& \hspace{5cm}\vspace{5cm}{\dots \circ T^{j_v}(\overline{f_{s-1}^{(p)}(x_0,\dots,x_{s-1}+x_{s},\dots, x_n)})}\\
%		+ & \mathlarger{\mathlarger{\sum}_{\substack{r \neq 0 \\ l \leq i_1 <\dots<i_r \leq n-1 }} (-1)^r S^{i_1} \circ\dots\circ S^{i_r} \overline{f_{l-1}^{(p)}(x_0,\dots,x_{l-1}+x_{l},\dots,x_n)}} \\
%		+ & \mathlarger{\mathlarger{\sum}_{t=1}^{n-1-l} \ \mathlarger{\sum}_{\substack{v \neq 0 \\ l \leq i_1 <\dots<i_v \leq n-1-t}} (-1)^v S^{i_1} \circ S^{i_2} \circ \dots S^{i_v}(\overline{f_{n-t+1}^{(p)}(x_0,\dots,x_{n-t}+x_{n-t+1},\dots, x_n)})},
%	\end{align*}
%	for $p = 1, \dots m$. This completes the formulas of horn filling of $Sing_*(\mathbb{A}^m_k)(U)$.
%\end{point}

\

\



\


\


%\begin{lemma}
%    Suppose, $\theta: X \to \sC$ is an \'etale locally trivial $\A^1$-bundle. Also assume that given a finitely generated field extension $L/k$ and a morphism $\gamma: \A^1_L \to \sC$, the morphism $\gamma$ factors through $\Spec\ L$. Then any morphism $H: \A^1_\sO \to \sC$ factors through $\sO$, for every Henselian local scheme $\sO$. Thus, $\sC$ is $\A^1$-invariant.
%\end{lemma}
%\begin{proof}
%    Consider the commutative diagram
%    \[\xymatrixcolsep{5pc}
%		\xymatrix{\mathbb{A}^1_\mathcal{O}\times_k \mathbb{A}^1_k \ar[rd]_{p_1} \ar[r]^{\cong}_{\Phi = (Id,f)} &\mathbb{A}^1_\mathcal{O}\times_\mathcal{C}X \ar[d]^{\pr_1} \ar[r]^<<<<{\pr_2} &X \ar[d]^{\theta}\\
%		 &\A^1_{\sO} \ar[r]^H &\sC}
%		\]
%        The isomorphism $\Phi$ is given by $\Phi \equiv (Id, f)$, $Id: \A^1_\sO \to \A^1_\sO$ is the identity morphism and $f: \mathbb{A}^1_\mathcal{O}\times_k \mathbb{A}^1_k \to X$ is a morphism such that 
%        $$\theta \circ f =H \circ p_1, \ p_1: \mathbb{A}^1_\mathcal{O}\times_k \mathbb{A}^1_k \to \A^1_\sO$$
%       % and
%        $$\text{ that is } H(t,x)=\theta(f(t,x,s)) \text{ and }\Phi((t,x),s)= ((t,x), f(t,x,s))$$
%        That means fixing $x_0 \in \sO$, $f(t,x_0,s)$ lies in fiber of $\theta$ over $H(0,x_0)$, for every $t,s \in \A^1_k$.
%        First observe that if $s_0 \neq s_1 \in \A^1_k$, then 
%        $$f(t,x,s_0) \neq f(t,x,s_1), \text{ since } \Phi \text{ is an isomorphism}.$$
%        Fix $x_0 \in \sO,t_0 \in \A^1_k, s_0 \in \A^1_k$.
%        We claim that  the map 
%        
%        $\A^1_{k(x_0)} \to X$ given by, $s \mapsto f(t_0,x_0,s)$ surjects onto the fiber of $\theta$ over $H(0,x_0)$ (this map factors through the fiber).
%
%        It is enough to any morphism (over $k$) $\A^1_L \to A^1_F \setminus \{P\}$ must be constant, where $L, F$ are fields over $k$ and $L/F$ is extension. In the level of rings
%        the map is given by ($h \in F[T]$ is some irreducible polynomial determining $P$)
%        $$F[T]_h \to L[T]$$
%        so $h$ must maps to an element of $L$, so this claim is true.
%
%        Thus $s \mapsto f(t_0,x_0,s)$ surjects onto the fiber, fixing $x_0$ and $t_0$.
%        Therefore, $t \mapsto f(t,x_0,s_0)$, fixing $x_0, s_0$ is constant, since $\Phi$ is isomorphism. That means, $f$ is independent of variable $t$.
%
%        So, if we take a section $\phi: \A^1_\sO \to \mathbb{A}^1_\mathcal{O}\times_k \mathbb{A}^1_k$ given by $(t,x) \mapsto (t,x,0)$, then the lift of $H$ to $X$ that is $\Tilde{H}: \A^1_\sO \to X$ factors through $\sO$.
%\end{proof}
%\begin{lemma}
%    Suppose, $\theta: X \to \sC$ is an \'etale locally trivial $\A^1$-bundle. Also assume that given a finitely generated field extension $L/k$ and a morphism $\gamma: \A^1_L \to \sC$, the morphism $\gamma$ factors through $\Spec\ L$. Then any morphism $H: \A^1_\sO \to \sC$ factors through $\sO$, for every Henselian local scheme $\sO$. Thus, $\sC$ is $\A^1$-invariant.
%\end{lemma}
%\begin{proof}
%    Consider the commutative diagram
%    \[\xymatrixcolsep{5pc}
%		\xymatrix{\mathbb{A}^1_\mathcal{O}\times_k \mathbb{A}^1_k \ar[rd]_{p_1} \ar[r]^{\cong}_{\Phi = (Id,f)} &\mathbb{A}^1_\mathcal{O}\times_\mathcal{C}X \ar[d]^{\pr_1} \ar[r]^<<<<{\pr_2} &X \ar[d]^{\theta}\\
%		 &\A^1_{\sO} \ar[r]^H &\sC}
%		\]
%        The isomorphism $\Phi$ is given by $\Phi \equiv (Id, f)$, $Id: \A^1_\sO \to \A^1_\sO$ is the identity morphism and $f: \mathbb{A}^1_\mathcal{O}\times_k \mathbb{A}^1_k \to X$ is a morphism such that 
%        $$\theta \circ f =H \circ p_1, \ p_1: \mathbb{A}^1_\mathcal{O}\times_k \mathbb{A}^1_k \to \A^1_\sO$$
%       % and
%        $$\text{ that is } H(t,x)=\theta(f(t,x,s)) \text{ and }\Phi((t,x),s)= ((t,x), f(t,x,s))$$
%        That means fixing $x_0: \Spec\ k(x_0) \to \sO$, the map
%        $$f_{x_0}:\A^1_{k(x_0)} \times_k \A^1_k \to X \text{ as  the composition}$$
%        $$\A^1_{k(x_0)} \times_k \A^1_k \xrightarrow{(Id, x_0, Id)} \mathbb{A}^1_\mathcal{O}\times_k \mathbb{A}^1_k \xrightarrow{f} X$$
%        $$ \theta \circ f_{x_0}=H \circ p_1 \circ (Id, x_0, Id)$$
%        $$H \circ p_1 \circ (Id, x_0, Id) \circ i_0: \A^1_{k(x_0)} \to \sC$$
%        where $i_0: \A^1_{k(x_0)} \to \A^1_{k(x_0)} \times_k \A^1_k$
%  $$p_1 \circ (Id, x_0, Id) \circ i_0 : \A^1_{k(x_0)} \to \A^1_{\sO} \text{ as } t \mapsto (t, x_0)$$
%  factors through $\Spec\ k(x_0) \xrightarrow{H(0, x_0)} \sC$. So $\text{Im}(f_{x_0})$ lies in the fiber of $\theta$ over $H(0,x_0)$.
%
%  Thus, $f_{x_0}$ factors through $\Spec\ k(x_0) \times_{\sC, H(0, x_0),f} X$. That is $f_{x_0}$ equals 
%  $$\A^1_{k(x_0)} \times_k \A^1_k \to \Spec\ k(x_0) \times_{\sC, H(0, x_0),f} X \to X.$$
%  Now $\Spec\ k(x_0) \times_{\sC, H(0, x_0),f} X \cong \A^1_{k(x_0)}$, since it is fiber of $\theta$ over $H(0,x_0)$.
%  First observe that since $\Phi$ is an isomorphism, so
%  $$\Spec\ k \xrightarrow{(t_0, s_0)} \A^1_{k(x_0)} \times_k \A^1_k \xrightarrow{f_{x_0}} X$$
%  is not equal to
%  $$\Spec\ k \xrightarrow{(t_0, s_1)} \A^1_{k(x_0)} \times_k \A^1_k \xrightarrow{f_{x_0}} X$$
%  their images are distinct, for fixing $s_0, t_0, s_1 \in \A^1_k(k)$ and $s_0 \neq s_1$.
%
%  Now we claim that fixing $t_0 \in \A^1_k(k)$, the map
%  $$\A^1_{k(x_0)} \xrightarrow{(t_0, Id)} \A^1_{k(x_0)} \times_k \A^1_k \xrightarrow{f_{x_0}} X$$ 
%  surjects onto the fiber, that is
%  $$\A^1_{k(x_0)} \xrightarrow{(t_0, Id)} \A^1_{k(x_0)} \times_k \A^1_k \to \Spec\ k(x_0) \times_{\sC, H(0, x_0),f} X (\cong \A^1_{k(x_0)})$$
%  is non-constant, hence surjective.
%
%  This implies the map
%  $$\A^1_{k(x_0)} \xrightarrow{(Id, s_0)} \A^1_{k(x_0)} \times_k \A^1_k \to \Spec\ k(x_0) \times_{\sC, H(0, x_0),f} X \to X$$
%  is constant, fixing $s_0 \in \A^1_k(k)$.
%
%  So, if we take a section $\phi: \A^1_\sO \to \mathbb{A}^1_\mathcal{O}\times_k \mathbb{A}^1_k$ given by $(t,x) \mapsto (t,x,0)$, we claim that
%  $$\Tilde{H}:= \text{pr}_2 \circ \Phi \circ \phi: \A^1_{\sO} \to X$$ 
%  factors through $\sO$
%  $$\theta \circ \text{pr}_2 \circ \Phi \circ \phi = H$$
%We will show that $\Tilde{H} = H \circ j_0 \circ p$, where
%$$p: \A^1_\sO \to \sO, \ j_0: \sO \to \A^1_\sO.$$
%From prevous, fixing $\Spec\ k(x_0) \to \sO$, $\Tilde{H}|_{\A^1_{k(x_0)}}$ factors through $\Spec\ k(x_0)$. 
%
%So the maps $\Tilde{H}$ and $H \circ j_0 \circ p$ are pointwise same, thereforefore they are same; we can follow the equiliser argument and morphism between scheme s isa a locally closed immersion.
%  
%  
%        \  
%        
%        \
%        
%        $f(t,x_0,s)$  lies in fiber of $\theta$ over $H(0,x_0)$, for every $t,s \in \A^1_k$.
%        First observe that if $s_0 \neq s_1 \in \A^1_k$, then 
%        $$f(t,x,s_0) \neq f(t,x,s_1), \text{ since } \Phi \text{ is an isomorphism}.$$
%        Fix $x_0 \in \sO,t_0 \in \A^1_k, s_0 \in \A^1_k$.
%        We claim that  the map 
%        
%        $\A^1_{k(x_0)} \to X$ given by, $s \mapsto f(t_0,x_0,s)$ surjects onto the fiber of $\theta$ over $H(0,x_0)$ (this map factors through the fiber).
%
%        It is enough to any morphism (over $k$) $\A^1_L \to A^1_F \setminus \{P\}$ must be constant, where $L, F$ are fields over $k$ and $L/F$ is extension. In the level of rings
%        the map is given by ($h \in F[T]$ is some irreducible polynomial determining $P$)
%        $$F[T]_h \to L[T]$$
%        so $h$ must maps to an element of $L$, so this claim is true.
%
%        Thus $s \mapsto f(t_0,x_0,s)$ surjects onto the fiber, fixing $x_0$ and $t_0$.
%        Therefore, $t \mapsto f(t,x_0,s_0)$, fixing $x_0, s_0$ is constant, since $\Phi$ is isomorphism. That means, $f$ is independent of variable $t$.
%
%        So, if we take a section $\phi: \A^1_\sO \to \mathbb{A}^1_\mathcal{O}\times_k \mathbb{A}^1_k$ given by $(t,x) \mapsto (t,x,0)$, then the lift of $H$ to $X$ that is $\Tilde{H}: \A^1_\sO \to X$ factors through $\sO$.
%\end{proof}

	
	





		
\bibliographystyle{alpha}
	\begin{thebibliography}{alpha}
		\bibitem{asa} T. Asanuma; \emph{Polynomial fiber rings of algebras over noetherian rings}, Invent. math. 87, 101-127 (1987).
		
		\bibitem{asok homology} A. Asok; \emph{Birational invariants and $\mathbb{A}^1$-connectedness}, J. Reine Angew. Math. 681 (2012) 39-64.
		
		\bibitem{asokdoran} A. Asok, B. Doran; \emph{$\mathbb{A}^1$-homotopy groups, excision, and solvable quotients}, Advances in Mathematics, 221(2009) Pages - 1144-1190.
		
		
		\bibitem{adf} A. Asok, B. Doran, J. Fasel;  \emph{Smooth models of motivic spheres and the clutching construction}, Int. Math. Res. Not. 2017 (6): 1890-1925. 
		
		%\bibitem{russell} P. Russell, \emph{On Affine-Ruled Rational Surfaces}, Mathematische Annalen 255 (1981): 287-302.
		\bibitem{am} A. Asok and F. Morel; \emph{Smooth varieties up to $\mathbb{A}^{1}$-homotopy and algebraic h-cobordisms,} Adv. Math. 227(5) (2011) 1990–2058.
		
		\bibitem{as} A. Asok, P. A. \O stv\ae r; \emph{$\mathbb{A}^1$-homotopy theory and contractible varieties: a survey}, Homotopy Theory and Arithmetic Geometry – Motivic and Diophantine Aspects. Lecture Notes in Mathematics, vol 2292. Springer, Cham. https://doi.org/10.1007/978-3-030-78977-05.
		
		
		\bibitem{bhs} C. Balwe, A. Hogadi, A. Sawant; \emph{$\mathbb{A}^1$-connected components of schemes}, Adv Math, Volume 282, 2016.
		
		\bibitem{bhs1} C. Balwe, A. Hogadi, A. Sawant; \emph{Geometric criteria for $\mathbb{A}^1$-connectedness and applications to norm varieties}, Journal of Algebraic Geometry, 32 (2023), no. 4, 677-696.
		\bibitem{bhs2} C. Balwe, A. Hogadi, A. Sawant; \emph{Strong $\mathbb{A}^1$-invariance of $\mathbb{A}^1$-connected components of reductive algebraic groups}, Journal of Topology, 16 (2023), no. 2, 634-649.
		\bibitem{brs}C. Balwe, B. Rani, A. Sawant;\emph{ Remarks on iterations of the $\mathbb{A}^1$-chain connected components construction}, Annals of K-Theory, 7 (2022), no. 2, 385-394.
		\bibitem{bs} C. Balwe, A. Sawant; \emph{$\mathbb{A}^1$-connected components of ruled surfaces}, Geometry and Topology, 26 (2022), 321-376.
		\bibitem{u} U. Choudhury; \emph{Connectivity of motivic $H$-spaces}. Algebr. Geom. Topol., Volume 14, Number 1(2014), 37-55.
\bibitem{balwe reductive} C. Balwe, A. Sawant; \emph{$\A^1$-connectedness in reductive algebraic groups}, Transactions of the American Mathematical Society.
		\bibitem{cb} U. Choudhury, B. Roy; \emph{$\mathbb{A}^1$-connected components and characterisation of $\mathbb{A}^2$}, Journal für die reine und angewandte Mathematik (Crelle's Journal), Volume 2024, Issue 807, Page - 55-80. 
		
		\bibitem{df} A.Dubouloz, J.Fasel; \emph{Families of $\mathbb{A}^1$-contractible affine threefolds}, Algebraic Geometry 5 (1) (2018) 1–14, doi:10.14231/AG-2018-001.
		
		\bibitem{dpo} A. Dubouloz, S. Pauli, P. A. \O stv\ae r; \emph{$\mathbb{A}^1$-contractibility of affine modifications}, International Journal of Mathematics, Vol. 30, No. 14 (2019) 1950069.
		
		\bibitem{fr} G. Freudenberg; \emph{Algebraic Theory of Locally Nilpotent Derivations}, Encyclopedia of Mathematical Sciences, Vol.136, Invariant Theory and Algebraic Transformation Groups, VII (Springer-Verlag, 2017), second edition.
		
		\bibitem{gurjar miyanishi} R. V. Gurjar, M. Miyanishi, \emph{Affine lines on logarithmic $\mathbb{Q}$-homology planes}, Math. Ann. 294:3 (1992), 463–482.
		
		\bibitem{hart} R. Hartshorne; \emph{Algebraic Geometry}, Springer - Verlag, New York, 1982.
		
		\bibitem{iitaka} S. Iitakaa; \emph{Algebraic Geometry: An Introduction to Birational Geometry of Algebraic Varieties}, Graduate Texts in Mathemati, Vol. 76, Springer-Verlag, new York, 1982.
		
		\bibitem{kaliman} S. Kaliman, L. Makar-Limanov; \emph{On the Russell-Koras contractible threefolds}, J. Algebraic Geom., 6 no. 2 (1997), 247–268.
		
		\bibitem{kamb} T. Kambayashi; \emph{On the Absence of Nontrivial Separable Forms of the Affine Plane }, Journal Of Algebra 35, (1975), 449-456.
		
		\bibitem{koizumi} Junnosuke Koizumi; \emph{Zeroth $\mathbb{A}^1$-homology of smooth proper varieties}, https://doi.org/10.48550/arXiv.2101.04951.

\bibitem{ayoub} Joseph Ayoub; \emph{Counterexamples to F. Morel's conjecture on $\pi_0^{\mathbb{A}^1}$}, Comptes Rendus Mathématique, Volume 361 (2023), 1087\textendash1090.
		
		%\bibitem{lang} J. Lang, \emph{An example related to the affine theorem of Castelnuovo}, Michigan Math. J.
		
		\bibitem{may} J. P. May; \emph{Simplicial Objects in Algebraic Topology}, The University of Chicago Press, Chicago and London.
		
		\bibitem{MWV} C. Mazza, V. Voevodsky, C. Weibel; \emph{Lecture Notes on Motivic Cohomology}, Clay Mathematics Monographs Volume: 2; 2006; 216 pp.
		
		\bibitem{miyanishi} M. Miyanishi; \emph{Non-complete algebtraic surfaces}, Lecture Notes in Mathematics 857, Springer Berlin, Heidelberg.
		
		
		%\bibitem{miyaruss} M. Miyanishi, P. Russell, \emph{Purely Inseparable Coverings of Exponent one of the Affine Plane}, Journal of Pure and Applied Algebra 28 (1983) Page: 279-317, North-Holland. 
		
		\bibitem{mor} F. Morel; \emph{$\mathbb{A}^1$-algebraic topology over a field}, Lecture Notes in Mathematics, Springer Verlag, 2012, Volume 2052
		ISBN: 978-3-642-29513-3.
		
		\bibitem{mv} F. Morel, V. Voevodsky; \emph{$\mathbb{A}^1$-homotopy theory of schemes}, Publications Math\'ematiques de l'IHES 90 (1990)  p.~45-143.
		
		%\bibitem{murthy} M. P. Murthy, R. G. Swan; \emph{Vector Bundles over Affine Surfaces}, Inventiones math. 36, 125 - 165 (1976).
		
		\bibitem{ra} C. P. Ramanujam, \emph{A topological characterization of the affine plane as an algebraic
			variety}, Ann. Math. 94 (1971) 69-88.
		
		%\bibitem{russell} P. Russell, \emph{On Affine-Ruled Rational Surfaces}, Mathematische Annalen 255 (1981): 287-302.
		
		%\bibitem{russell2} P. Russell, \emph{Forms of the affine line and its additive group}, Pacific Journal of Mathematics, Vol. 32, No. 2, 1970.
		
		\bibitem{sathaye} A. Sathaye; \emph{Polynomial Rings in Two Variables Over a D.V.R.: A Criterion}, Invent. Math. 74, 159-168 (1983).  
		
		\bibitem{serre} J. P. Serre; \emph{Faisceaux alg\'ebriques coh\'erents}. Ann. of Math. 61, 197-278 (1955). 
		
		\bibitem{stack} The Stacks project authors; \emph{The stacks project, https://stacks.math.columbia.edu}, 2022.

\bibitem{ictp} F. Morel; \emph{An introduction to A1-homotopy theory}, in Contemporary Developments in Algebraic K-theory. ICTP Lecture Notes, vol. XV (electronic) (Abdus Salam International Central Theoretical Physics, Trieste, 2004), pp. 357–441.
		
		\bibitem{sawant} A. Sawant; \emph{$\mathbb{A}^1$-connected components of schemes}, Ph. D Thesis, submitted to Tata Institute of Fundamental Research, Mumbai.
		
		\bibitem{yuri} Y. Shimizu; \emph{Relative $\mathbb{A}^1$-homology and its Applications}, Homology, Homotopy and Applications, Vol. 24(1), 2022, Pages - 129-141.
		
		%\bibitem{voev} V. Voevodsky; \emph{On motivic cohomology with $\mathbb{Z}/l$-cofficients}, Ann. of Math. (2) 174 (2011), no 1, 401-438.
		
		\bibitem{voev2} V. Voevodsky; \emph{Motivic cohomology with $\mathbb{Z}/2$-cofficients}, Publ. Math. Inst. Hautes \'Etudes Sci. No. 98 (2003), 59-104.  
		
		\bibitem{miyanishifibration} M. Miyanishi; \emph{Singularities of normal affine surfaces containing cylinderlike open sets}, J. Algebra 68 (1981), no. 2, 268–275.
		
		\bibitem{ouali} M’hammed El Kahoui, Mustapha Ouali; \emph{Fixed point free locally nilpotent derivations of $\mathbb{A}^2$-fibrations}, Journal of Algebra 372 (2012) 480–487.
		\bibitem{masuda} Rajendra Vasant Gurjar, Kayo Masuda, Masayoshi Miyanishi; \emph{Affine Space Fibrations}, De Gruyter Studies in Mathematics, Volume 79.
		\bibitem{wright}  T. Kambayashi and D. Wright; \emph{Flat families of affine lines are affine-line bundles};
		Illinois J. Math. 29(4) (1985) 672–681.
		\bibitem{miyanishicrm}  M. Miyanishi; \emph{Open algebraic surfaces}, CRM Monograph Series 12, Amer. Math. Soc., Providence, RI, 2001. 
		\bibitem{russell} P. Russell; \emph{On Affine-Ruled Rational Surfaces}, Mathematische Annalen 255 (1981), pages - 287-302.
		%\bibitem{asok} A. Asok and J. Fasel; \emph{Vector bundles on algebraic varieties}, ICM 2022 Proceedings, arXiv:2111.03107.
		\bibitem{kambayashi} T. Kambayashi, M. Miyanishi; \emph{On flat fibrations by the affine line}, Illinois J. Math. 22 (1978) 662–671.
		\bibitem{threefolds} R.V. Gurjar, K. Masuda, M. Miyanishi; \emph{$\mathbb{A}^1$-fibrations on affine threefolds}, Journal of Pure and Applied Algebra, Volume 216, Issue 2, February 2012, Pages 296-313.
		%\bibitem{stack} Authors, The Stacks Project: The Stacks project, https://stacks.math.columbia.edu, 2023.
		\bibitem{stack} The Stacks project authors; \emph{The stacks project, https://stacks.math.columbia.edu}, 2024.
		\bibitem{morel ictp} F. Morel; \emph{An introduction to $\mathbb{A}^1$-homotopy theory}, Contemporary developments in
		algebraic K-theory, 357–441 (electronic), ICTP Lect. Notes, XV, Abdus Salam Int.
		Cent. Theoret. Phys., Trieste, 2004.
		
		\bibitem{doub} A.Dubouloz; \emph{$\mathbb{A}^1$-cylinders over smooth $\mathbb{A}^1$-fibered affine surfaces}.
		
		\bibitem{alpher} Jarod Alpher; \emph{Stacks and Moduli}, updated Jan 2026.
		
		\bibitem{mazza} C.Mazza, V. Voevodsky, C. Weibel; \emph{Lecture Notes on motivic cohomology}, Clay Mathematics Monographs, Volume 2.
		
		\bibitem{JardineSimplicial} P.G. Goerss and J.F. Jardine, \emph{Simplicial homotopy theory}, Progress in Mathematics, 174, Birkh\"auser, Basel, 1999; MR1711612
		
		\bibitem{rydh} D. Rydh; \emph{The canonical embedding of an unramified morphism in an étale morphism}, Math. Z. (2011) 268:707–723, DOI: 10.1007/s00209-010-0691-8.

        \bibitem{milne} James S. Milne; \emph{`Etale Cohomology}.

        \bibitem{nlab} The nLab authors; \emph{simplicial group}, \url{https://ncatlab.org/nlab/show/simplicial+group}.

        \bibitem{hovey} Mark Hovey; \emph{Model Categories}.
        
        \bibitem{choudhury2} Utsav Choudhury and Biman Roy; \emph{$\mathbb{A}^1$-homotopy type of $\mathbb{A}^2 \setminus \{(0, 0)\}$}. Advances in Mathematics, 489:Paper No.110806, 31, 2026.

         \bibitem{dugger} Daniel Dugger; \emph{Universal Homotopy Theories}, Advances in Mathematics, Pages: 144-176, Volume 164, Published in 2001.
  \bibitem{asokII} Aravind Asok, Marc Hoyois and Matthias Wendt; \emph{Affine representability results in $\mathbb{A}^1$–homotopy theory, II: Principal bundles and homogeneous spaces}, Geometry Topology 22 (2018) 1181–1225.

     \bibitem{asokI} Aravind Asok, Marc Hoyois and Matthias Wendt; \emph{Affine representability results in $\mathbb{A}^1$–homotopy theory, I}, DUKE MATHEMATICAL JOURNAL, Vol. 166, No. 10.
\bibitem{haine} Magnus Carlson, Peter J. Haine and Sebastian Wolf; \emph{Reconstruction of schemes from their \'etale topoi}, available at arXiv:2407.19920v1.

\bibitem{koras} M. Koras and P. Russell; \emph{Contractible threefolds and $\mathbb{C}^*$-actions on $\mathbb{C}^3$}, J. Algebraic Geom. 6 (1997), 4, 671–695.
\bibitem{hko} M. Hoyois, A. Krishna, P. A. Østvær; \emph{$\A^1$-contractibility of Koras-Russell threefolds}, Algebraic Geometry 3 (2016), 407-423.

\bibitem{w} M. Wendt; \emph{On the $\mathbb{A}^1$-fundamental groups of smooth toric varieties}, Adv. Math. 223 (2010) 352\textendash378.

\bibitem{steenrod} Norman Steenrod; \emph{The Topology of fiber Bundles}; Princeton Mathematical Series.

\bibitem{serre thesis} Jean-Pierre Serre,; \emph{Homologie singulirre des espaces fibrrs}, Applications, Ann. of Math. 54 (1951), 425 - 505. 

\bibitem{brauer group} Jean-Louis Colliot-Thélène , Alexei N. Skorobogatov; \emph{The Brauer–Grothendieck Group}, Springer Cham.

\bibitem{krishnakumar} Adrien Dubouloz, Krishna Kumar Madhavan Vijayalakshmi, Paul Arne Østvær; \emph{Relative $\A^1$-Contractibility of Smooth Schemes}, available at https://doi.org/10.48550/arXiv.2605.07638.

\bibitem{naive connectedness sawant} A. Sawant; \emph{Naive vs. genuine $\A^1$-connectedness}, available at  arXiv:1703.05935v1.

\bibitem{jardine} R. Jardine, P. G. Goerss; \emph{Simplicial Homotopy Theory}.
\bibitem{curtis-simplicial} E.~B. Curtis, \emph{Simplicial homotopy theory}, Advances in Math. {\bf 6} (1971), 107--209 (1971); MR0279808
\bibitem{Biman-Thesis} B. Roy; \emph{$\A^1$ homotopy types of $\A^2$ and $\A^2\backslash(0,0)$.}, [PhD thesis, Indian Statistical Institute] Retrived from \url{http://hdl.handle.net/10263/7485}
%\bibitem{makarlimanov} S. Kaliman and L. Makar-Limanov; \emph{On the Russell-Koras contractible threefolds}, J. Algebraic Geom. 6 (1997), 2, 247–268.
		%, K. Masuda b
		% M. Miyanishi cA
		%1
		
		
		
		%\bibitem{as} A. Asok, P. A. \O stv\ae r; \emph{$\mathbb{A}^1$-homotopy theory and contractible varieties: a survey}, Homotopy Theory and Arithmetic Geometry – Motivic and Diophantine Aspects. Lecture Notes in Mathematics, vol 2292. Springer, Cham. https://doi.org/10.1007/978-3-030-78977-05.
		
		
		%\bibitem{asan} T. Asanuma; \emph{Purely inseparable $k$-forms of affine algebraic curves}, in Contemporary Mathematics, Affine Algebraic Geometry 369, American Mathematical Society, Providence, RI, 2005, 31-46.
		
		%\bibitem{ayoub} J. Ayoub; \emph{Une version relative de la conjecture des p\'eriodes de Kontsevich-Zagier}, Annals of Mathematics 181 (2015), 905–992.
		
		%\bibitem{ayoub2} J. Ayoub; \emph{L’alg\'ebre de Hopf et le groupe de Galois motiviques d’un corps de caract\'eristique nulle, I}, J. Reine Angew. Math. 693 (2014), pp. 1–149.
		
		%\bibitem{ayoub3} J. Ayoub; \emph{L’alg\'ebre de Hopf et le groupe de Galois motiviques d’un corps de caract\'eristique nulle, II}, J. Reine Angew. Math. 693 (2014), pp. 151–226.
		
		
		
		
		
		
		%\bibitem{bad} L. B\v{a}descu: \emph{Algebraic Surfaces}, Translated by Vladimir Ma\c{s}ek; Springer Edition, 2001.
		
		
		%U.Choudhury, B. Roy; \emph{Characterisation of the Affine Plane using $\mathbb{A}^1$-Homotopy Theory}, preprint, arXiv:2112.13089.
		
		%\bibitem{d} D. Dugger; \emph{Universal Homotopy Theories}, Advances in Mathematics, Volume 164, Issue 1, 1 December 2001, Pages 144-176
		
		%\bibitem{dhi} D. Dugger, S. Hollander and D. C. Isaksen; \emph{Hypercovers and simpilcial presheaves}, Mathematical Proceedings of the Cambridge Philosophical Society, 136, pp 9-51 doi:10.1017/S0305004103007175. 
		
		%\bibitem{di} D.Dugger, D. Isaksen, \emph{The motivic Adams spectral sequence}, Geom. Topol 14 (2010), no. 2, 967-1014. 
		
		%\bibitem{dutta} A. K. Dutta, N. Gupta, A. Lahiri; \emph{On Separable $\mathbb{A}^2$ and $\mathbb{A}^3$-Forms}, Nagoya Math. J., 239 (2020), 346-354, DOI 10.1017/nmj.2018.45.
		
		
		
		%\bibitem{gr} C. Greither; \emph{Forms of the affine line and their genus}, J. Pure Appl. Algebra 39 (1986), 105-118.
		
		%\bibitem{gurmiya} R. V. Gurjar, M. Miyanishi; \emph{Affine surfaces with $\bar{k} \leq 1$}, ``Algebraic geometry and Commutative algebra in honour of Masayoshi Nagata", Kinokuniya, Tokyo, 1987, pp. 99-124.
		
		%\bibitem{gurmiya2} R. V. Gurjar, M. Miyanishi; \emph{Affine lines on logarithmic Q-homology planes}, Math. Ann. 294, 463-482 (1992).
		
		%\bibitem{gurpa} R. V. Gurjar, A. J. Parameswaran; \emph{Affine lines on Q-homology planes}, J. Math. Kyoto Univ. (JMKYAZ) 35-1 (1995) 63 - 77.
		
		
		
		%\bibitem{hump} J. E. Humphreys; \emph{Linear Algebraic Groups}, GTM 21, Springer-Verlag.
		
		%\bibitem{iit} S., Iitaka; \emph{On logarithmic Kodaira dimension of algebraic varieties}, Complex Analysis and Algebraic Geometry, W.L. Baily, Jr. and T. Shioda ed., Iwanami Shoten, Tokyo, Tokyo, 1977, 175-189.
		
		%\bibitem{isak} D. C. Isaksen, \emph{Stable Stems}, Mem. Am. Math. Soc. 262, 1269 (2019).
		
		%\bibitem{isakwangxu} D.C. Isaksen, G. Wang, Z. Xu, \emph{More Stable Stems}, arXiv:2001.04511v2 (17 June 2020).
		
		%\bibitem{isakwangxu2} D.C. Isaksen, G. Wang, Z. Xu, \emph{Classical and $\mathbb{C}$-motivic Adams charts (2020)}, s.wayne.edu/isaksen/adams-charts (October 2022).
		
		%\bibitem{tkamb} T. Kambayashi; \emph{on Fujita's strong cancellation theorem for the affine plane}, J. Fac. Sci. Univ. Tokyo, 27 (1980), 535-548.
		
		%\bibitem{kamb} T. Kambayashi; \emph{On the Absence of Nontrivial Separable Forms of the Affine Plane }, Journal Of Algebra 35, (1975), 449-456.
		
		%\bibitem{koch} S. Kochman; \emph{Stable Homotopy Groups of Spheres}, Springer Lecture Notes 1423, 1990.
		
		%\bibitem{kollar} J. Koll\'ar; \emph{Rational curves on algebraic varieties}, in: Ergeb. Math. Grenzgeb. (3) (Results in Mathematics and
		%Related Areas. 3rd Series. A Series of Modern Surveys in Mathematics), vol. 32, Springer-Verlag, Berlin, 1996.
		
		%\bibitem{miyasugie} M. Miyanishi and T. Sugie; \emph{Homology planes with quotient singularities}, J. Math. Kyoto Univ.
		
		%\bibitem{miyat} M. Miyanishi and S. Tsunoda; \emph{Absence of the affine lines on the homology planes of general type}, J. Math. Kyoto Univ. (JMKYAZ) 32-3 (1992) 443 - 450.
		
		
		
		%\bibitem{rachet} R. Achet; \emph{Picard group of the forms of the affine line and of the additive group}, Journal of Pure and Applied Algebra, 2017, 221 (11), pp.2838-2860., 10.1016, hal-01377384v2.
		
		
		
		
		
		
		
		
		%\bibitem{milne} James S. Milne; \emph{Lectures on étale cohomology}.
		
		
		
		
		
		
	\end{thebibliography}
    \end{document}
